\input{preamble}

% OK, start here.
%
\begin{document}

\title{Modules sur les sites}


\maketitle

\phantomsection
\label{section-phantom}

\tableofcontents

\section{Introduction}
\label{section-introduction}

\noindent
Dans ce chapitre, nous développons les notions de base relatives aux faisceaux
de modules sur les topos annelés ou les sites annelés. Nous établissons d'abord
quelques faits fondamentaux sur les faisceaux abéliens, puis nous introduisons
les sites annelés et les topos annelés. Nous étudions ensuite certaines notions
très élémentaires concernant les (pré)faisceaux de $\mathcal{O}$-modules,
parallèlement au traitement des (pré)faisceaux de $\mathcal{O}$-modules dans le
chapitre consacré aux faisceaux sur les espaces. Cela fait, nous reprenons une
grande partie de la discussion du chapitre sur les faisceaux de modules (voir
Faisceaux de modules, Section \ref{modules-section-introduction}). Les
références de base sont \cite{FAC}, \cite{EGA} et \cite{SGA4}.






\section{Préfaisceaux abéliens}
\label{section-abelian-pre-sheaves}

\noindent
Soit $\mathcal{C}$ une catégorie. Les préfaisceaux abéliens ont été introduits
dans Sites, Sections \ref{sites-section-presheaves} et
\ref{sites-section-sheaves}, puis étudiés plus avant dans Sites, Section
\ref{sites-section-sheaves-algebraic-structures}. Nous suivrons la convention
de cette dernière référence : un préfaisceau abélien est un préfaisceau
d'ensembles muni, sur chacun de ses ensembles de sections, d'une loi d'addition
compatible avec les applications de restriction. Rappelons que la catégorie
des préfaisceaux abéliens sur $\mathcal{C}$ est notée
$\textit{PAb}(\mathcal{C})$.

\medskip\noindent
La catégorie $\textit{PAb}(\mathcal{C})$ est abélienne au sens d'Homologie,
Définition \ref{homology-definition-abelian-category}. Étant donnée une
application de préfaisceaux $\varphi : \mathcal{G}_1 \to \mathcal{G}_2$, le
noyau de $\varphi$ est le préfaisceau abélien
$U \mapsto \Ker(\mathcal{G}_1(U) \to \mathcal{G}_2(U))$, et le conoyau de
$\varphi$ est
le préfaisceau $U \mapsto
\Coker(\mathcal{G}_1(U) \to \mathcal{G}_2(U))$. Puisque la catégorie des
groupes abéliens est abélienne, il s'ensuit que $\Coim = \Im$, cette égalité
étant vraie au-dessus de chaque $U$. Une suite de préfaisceaux abéliens
$$
\mathcal{G}_1 \longrightarrow
\mathcal{G}_2 \longrightarrow
\mathcal{G}_3
$$
est exacte si et seulement si
$\mathcal{G}_1(U) \to \mathcal{G}_2(U) \to \mathcal{G}_3(U)$
est une suite exacte de groupes abéliens pour tout $U \in \Ob(\mathcal{C})$.
Nous laissons les vérifications au lecteur.

\begin{lemma}
\label{lemma-limits-colimits-abelian-presheaves}
Soit $\mathcal{C}$ une catégorie.
\begin{enumerate}
\item Toutes les limites et toutes les colimites existent dans
$\textit{PAb}(\mathcal{C})$.
\item Toutes les limites et toutes les colimites commutent à la prise des
sections au-dessus des objets de $\mathcal{C}$.
\end{enumerate}
\end{lemma}

\begin{proof}
Soit $\mathcal{I} \to \textit{PAb}(\mathcal{C})$, $i \mapsto \mathcal{F}_i$,
un diagramme. Nous pouvons simplement définir des préfaisceaux abéliens $L$ et
$C$ par les règles
$$
L : U \longmapsto \lim_i \mathcal{F}_i(U)
$$
and
$$
C : U \longmapsto \colim_i \mathcal{F}_i(U).
$$
Il existe manifestement des applications de préfaisceaux abéliens
$L \to \mathcal{F}_i$ et $\mathcal{F}_i \to C$, obtenues à partir des
applications correspondantes sur les groupes de sections au-dessus de chaque
$U$. On vérifie immédiatement que $L$ et $C$, munis de ces applications, sont
respectivement la limite et la colimite du diagramme dans
$\textit{PAb}(\mathcal{C})$. Cela démontre (1) et (2). Nous omettons les
détails.
\end{proof}


\section{Faisceaux abéliens}
\label{section-abelian-sheaves}

\noindent
Soit $\mathcal{C}$ un site. La catégorie des faisceaux abéliens sur
$\mathcal{C}$ est notée $\textit{Ab}(\mathcal{C})$. C'est la sous-catégorie
pleine de $\textit{PAb}(\mathcal{C})$ formée des préfaisceaux abéliens dont le
préfaisceau d'ensembles sous-jacent est un faisceau. Les propriétés
($\alpha$) -- ($\zeta$) de Sites, Section
\ref{sites-section-sheaves-algebraic-structures}, sont satisfaites ; voir
Sites, Proposition
\ref{sites-proposition-functoriality-algebraic-structures-topoi}. En
particulier, le foncteur d'inclusion
$\textit{Ab}(\mathcal{C}) \to \textit{PAb}(\mathcal{C})$ admet un adjoint à
gauche, à savoir le foncteur de faisceautisation
$\mathcal{G} \mapsto \mathcal{G}^\#$.

\medskip\noindent
Nous suggérons au lecteur de démontrer le lemme sur une feuille de brouillon
plutôt que d'en lire la démonstration.

\begin{lemma}
\label{lemma-abelian-abelian}
Soit $\mathcal{C}$ un site. Soit $\varphi : \mathcal{F} \to \mathcal{G}$
un morphisme de faisceaux abéliens sur $\mathcal{C}$.
\begin{enumerate}
\item La catégorie $\textit{Ab}(\mathcal{C})$ est une catégorie abélienne.
\item Le noyau $\Ker(\varphi)$ de $\varphi$ coïncide avec le noyau de
$\varphi$ comme morphisme de préfaisceaux.
\item Le morphisme $\varphi$ est injectif (Homologie, Définition
\ref{homology-definition-injective-surjective}) si et seulement si $\varphi$ est
injectif comme application de préfaisceaux (Sites, Définition
\ref{sites-definition-presheaves-injective-surjective}), si et seulement si
$\varphi$ est injectif comme application de faisceaux (Sites, Définition
\ref{sites-definition-sheaves-injective-surjective}).
\item Le conoyau $\Coker(\varphi)$ de $\varphi$ est le faisceau associé au
conoyau de $\varphi$ comme morphisme de préfaisceaux.
\item Le morphisme $\varphi$ est surjectif (Homologie, Définition
\ref{homology-definition-injective-surjective}) si et seulement si $\varphi$ est
surjectif comme application de faisceaux (Sites, Définition
\ref{sites-definition-sheaves-injective-surjective}).
\item Un complexe de faisceaux abéliens
$$
\mathcal{F} \to \mathcal{G} \to \mathcal{H}
$$
est exact en $\mathcal{G}$ si et seulement si, pour tout
$U \in \Ob(\mathcal{C})$ et tout $s \in \mathcal{G}(U)$ dont l'image dans
$\mathcal{H}(U)$ est nulle, il existe un recouvrement
$\{U_i \to U\}_{i \in I}$ dans $\mathcal{C}$ tel que chaque $s|_{U_i}$
appartienne à l'image de $\mathcal{F}(U_i) \to \mathcal{G}(U_i)$.
\end{enumerate}
\end{lemma}

\begin{proof}
Nous affirmons qu'Homologie, Lemme
\ref{homology-lemma-adjoint-get-abelian}, s'applique aux catégories
$\mathcal{A} = \textit{Ab}(\mathcal{C})$ et
$\mathcal{B} = \textit{PAb}(\mathcal{C})$, ainsi qu'aux foncteurs
$a : \mathcal{A} \to \mathcal{B}$ (l'inclusion) et
$b : \mathcal{B} \to \mathcal{A}$ (la faisceautisation). Vérifions les
hypothèses d'Homologie, Lemme \ref{homology-lemma-adjoint-get-abelian}.
L'hypothèse (1) dit que $\mathcal{A}$ et $\mathcal{B}$ sont des
catégories additives, que $a$ et $b$ sont des foncteurs additifs, et que $a$
est adjoint à droite de $b$. Les deux premières assertions sont claires, et
l'adjonction est Sites, Section
\ref{sites-section-sheaves-algebraic-structures} ($\epsilon$). L'hypothèse
(2) dit que $\textit{PAb}(\mathcal{C})$ est abélienne, ce que nous avons vu à
la Section \ref{section-abelian-pre-sheaves}, et que la faisceautisation est
exacte à gauche, ce qui est Sites, Section
\ref{sites-section-sheaves-algebraic-structures} ($\zeta$). Enfin, la dernière
hypothèse est $ba \cong \text{id}_\mathcal{A}$, ce qui est Sites, Section
\ref{sites-section-sheaves-algebraic-structures} ($\delta$). Ainsi Homologie,
Lemme \ref{homology-lemma-adjoint-get-abelian}, s'applique et nous concluons
que $\textit{Ab}(\mathcal{C})$ est abélienne.

\medskip\noindent
La démonstration d'Homologie, Lemme
\ref{homology-lemma-adjoint-get-abelian}, montre que $\Ker(\varphi)$ et
$\Coker(\varphi)$ sont les faisceaux associés respectivement au noyau et au
conoyau de $\varphi$ comme morphisme de préfaisceaux abéliens. Cela démontre
(4). Puisque le noyau est un égalisateur (donc une limite) et que la
faisceautisation commute aux limites finies, nous en déduisons (2).

\medskip\noindent
L'assertion (2) implique (3). L'assertion (4) implique (5) d'après notre
description de la faisceautisation. La caractérisation de l'exactitude en (6)
résulte de (2), de (5), et du fait que la suite est exacte si et seulement si
$\Im(\mathcal{F} \to \mathcal{G}) =
\Ker(\mathcal{G} \to \mathcal{H})$.
\end{proof}

\noindent
On peut encore formuler la partie (6) du lemme en disant qu'une suite de
faisceaux abéliens
$$
\mathcal{F}_1 \longrightarrow
\mathcal{F}_2 \longrightarrow
\mathcal{F}_3
$$
est exacte si et seulement si le faisceau associé à
$U \mapsto \Im(\mathcal{F}_1(U) \to \mathcal{F}_2(U))$
est égal au noyau de $\mathcal{F}_2 \to \mathcal{F}_3$.

\begin{lemma}
\label{lemma-limits-colimits-abelian-sheaves}
Soit $\mathcal{C}$ un site.
\begin{enumerate}
\item Toutes les limites et toutes les colimites existent dans
$\textit{Ab}(\mathcal{C})$.
\item Les limites coïncident avec les limites correspondantes de préfaisceaux
abéliens sur $\mathcal{C}$ (autrement dit, elles commutent à la prise des
sections au-dessus des objets de $\mathcal{C}$).
\item Les sommes directes finies coïncident avec les sommes directes finies
correspondantes dans la catégorie des préfaisceaux abéliens sur $\mathcal{C}$.
\item Une colimite est le faisceau associé à la colimite correspondante dans la
catégorie des préfaisceaux abéliens.
\item Les colimites filtrantes sont exactes.
\end{enumerate}
\end{lemma}

\begin{proof}
D'après le Lemme \ref{lemma-limits-colimits-abelian-presheaves}, les limites et
les colimites de préfaisceaux abéliens existent et s'obtiennent en prenant les
limites et les colimites au niveau des sections au-dessus des objets.

\medskip\noindent
Soit $\mathcal{I} \to \textit{Ab}(\mathcal{C})$,
$i \mapsto \mathcal{F}_i$, un diagramme. Soit $\lim_i \mathcal{F}_i$ sa
limite comme préfaisceau abélien. D'après Sites, Lemme
\ref{sites-lemma-limit-sheaf}, c'est un faisceau abélien. On voit alors sans
difficulté que $\lim_i \mathcal{F}_i$ est la limite du diagramme dans
$\textit{Ab}(\mathcal{C})$. Cela démontre l'existence des limites et (2).

\medskip\noindent
D'après Catégories, Lemme \ref{categories-lemma-adjoint-exact}, et puisque la
faisceautisation est adjointe à gauche au foncteur d'inclusion, nous voyons que
$\colim_i \mathcal{F}$ existe et est le faisceau associé à la colimite dans
$\textit{PAb}(\mathcal{C})$. Cela démontre l'existence des colimites et (4).

\medskip\noindent
Dans toute catégorie abélienne, les sommes directes finies coïncident avec les
produits finis. Ainsi (3) résulte de (2).

\medskip\noindent
Preuve de (5). L'assertion signifie qu'étant donné un système
$0 \to \mathcal{F}_i \to \mathcal{G}_i \to \mathcal{H}_i \to 0$
de suites exactes de faisceaux abéliens indexé par un ensemble filtrant $I$, la
suite
$0 \to \colim \mathcal{F}_i \to \colim \mathcal{G}_i \to
\colim \mathcal{H}_i \to 0$ est également exacte. Un argument formel utilisant
Homologie, Lemme \ref{homology-lemma-check-exactness}, et la définition des
colimites montre que la suite
$\colim \mathcal{F}_i \to \colim \mathcal{G}_i \to \colim \mathcal{H}_i \to 0$
est exacte. Remarquons que
$\colim \mathcal{F}_i \to \colim \mathcal{G}_i$ est le faisceau associé à
l'application entre les colimites de préfaisceaux ; celle-ci est injective
puisque chacune des applications
$\mathcal{F}_i \to \mathcal{G}_i$ l'est. La faisceautisation étant exacte, on
conclut.
\end{proof}







\section{Préfaisceaux abéliens libres}
\label{section-free-abelian-presheaf}

\noindent
Afin de préparer la notation de la définition suivante, convenons de noter le
groupe abélien libre sur un ensemble $S$ par\footnote{Dans d'autres chapitres,
la notation $\mathbf{Z}[S]$ désigne parfois l'anneau de polynômes sur
$\mathbf{Z}$ en les éléments de $S$.}
$\mathbf{Z}[S] = \bigoplus_{s \in S} \mathbf{Z}$. Il est caractérisé par la
propriété
$$
\Mor_{\textit{Ab}}(\mathbf{Z}[S], A)
=
\Mor_{\textit{Sets}}(S, A)
$$
Autrement dit, la construction $S \mapsto \mathbf{Z}[S]$ est adjointe à
gauche au foncteur d'oubli $\textit{Ab} \to \textit{Sets}$.

\begin{definition}
\label{definition-free-abelian-presheaf-on}
Soit $\mathcal{C}$ une catégorie. Soit $\mathcal{G}$ un préfaisceau
d'ensembles. Le {\it préfaisceau abélien libre}
$\mathbf{Z}_\mathcal{G}$ sur $\mathcal{G}$ est le préfaisceau abélien défini
par la règle
$$
U \longmapsto \mathbf{Z}[\mathcal{G}(U)].
$$
Dans le cas particulier $\mathcal{G} = h_X$ d'un préfaisceau représentable
associé à un objet $X$ de $\mathcal{C}$, nous utilisons la notation
$\mathbf{Z}_X = \mathbf{Z}_{h_X}$. Autrement dit,
$$
\mathbf{Z}_X(U) = \mathbf{Z}[\Mor_\mathcal{C}(U, X)].
$$
\end{definition}

\noindent
Cette construction est manifestement fonctorielle en le préfaisceau
$\mathcal{G}$. Elle est en fait adjointe au foncteur d'oubli
$\textit{PAb}(\mathcal{C}) \to \textit{PSh}(\mathcal{C})$. Voici l'énoncé
précis.

\begin{lemma}
\label{lemma-obvious-adjointness}
Soit $\mathcal{C}$ une catégorie. Soient $\mathcal{G}$ et $\mathcal{F}$ des
préfaisceaux d'ensembles. Soit $\mathcal{A}$ un préfaisceau abélien. Soit $U$
un objet de $\mathcal{C}$. Alors nous avons
\begin{align*}
\Mor_{\textit{PSh}(\mathcal{C})}(h_U, \mathcal{F})
& =
\mathcal{F}(U), \\
\Mor_{\textit{PAb}(\mathcal{C})}(\mathbf{Z}_\mathcal{G}, \mathcal{A})
& =
\Mor_{\textit{PSh}(\mathcal{C})}(\mathcal{G}, \mathcal{A}), \\
\Mor_{\textit{PAb}(\mathcal{C})}(\mathbf{Z}_U, \mathcal{A})
& =
\mathcal{A}(U).
\end{align*}
Toutes ces égalités sont fonctorielles.
\end{lemma}

\begin{proof}
Omis.
\end{proof}

\begin{lemma}
\label{lemma-coproduct-sum-free-abelian-presheaf}
Soit $\mathcal{C}$ une catégorie. Soit $I$ un ensemble. Pour chaque
$i \in I$, soit $\mathcal{G}_i$ un préfaisceau d'ensembles. Alors
$$
\mathbf{Z}_{\coprod_i \mathcal{G}_i}
=
\bigoplus\nolimits_{i \in I} \mathbf{Z}_{\mathcal{G}_i}
$$
dans $\textit{PAb}(\mathcal{C})$.
\end{lemma}

\begin{proof}
Omis.
\end{proof}



\section{Faisceaux abéliens libres}
\label{section-free-abelian-sheaf}

\noindent
Voici la notion de faisceau abélien libre sur un faisceau d'ensembles.

\begin{definition}
\label{definition-free-abelian-sheaf-on}
Soit $\mathcal{C}$ un site. Soit $\mathcal{G}$ un préfaisceau d'ensembles. Le
{\it faisceau abélien libre} $\mathbf{Z}_\mathcal{G}^\#$ sur $\mathcal{G}$
est le faisceau abélien $\mathbf{Z}_\mathcal{G}^\#$ associé au préfaisceau
abélien libre sur $\mathcal{G}$. Dans le cas particulier
$\mathcal{G} = h_X$ d'un préfaisceau représentable associé à un objet $X$ de
$\mathcal{C}$, nous utilisons la notation $\mathbf{Z}_X^\#$.
\end{definition}

\noindent
Cette construction est manifestement fonctorielle en le préfaisceau
$\mathcal{G}$. Elle fournit en fait un adjoint au foncteur d'oubli
$\textit{Ab}(\mathcal{C}) \to \Sh(\mathcal{C})$. Voici l'énoncé précis.

\begin{lemma}
\label{lemma-obvious-adjointness-sheaves}
Soit $\mathcal{C}$ un site. Soient $\mathcal{G}$ et $\mathcal{F}$ des
faisceaux d'ensembles. Soit $\mathcal{A}$ un faisceau abélien. Soit $U$ un
objet de $\mathcal{C}$. Alors nous avons
\begin{align*}
\Mor_{\Sh(\mathcal{C})}(h_U^\#, \mathcal{F})
& =
\mathcal{F}(U), \\
\Mor_{\textit{Ab}(\mathcal{C})}(\mathbf{Z}_\mathcal{G}^\#,
\mathcal{A})
& =
\Mor_{\Sh(\mathcal{C})}(\mathcal{G}, \mathcal{A}), \\
\Mor_{\textit{Ab}(\mathcal{C})}(\mathbf{Z}_U^\#, \mathcal{A})
& =
\mathcal{A}(U).
\end{align*}
Toutes ces égalités sont fonctorielles.
\end{lemma}

\begin{proof}
Omis.
\end{proof}

\begin{lemma}
\label{lemma-may-sheafify-before-abelianize}
Soit $\mathcal{C}$ un site. Soit $\mathcal{G}$ un préfaisceau d'ensembles.
Alors $\mathbf{Z}_\mathcal{G}^\# = (\mathbf{Z}_{\mathcal{G}^\#})^\#$.
\end{lemma}

\begin{proof}
Omis.
\end{proof}








\section{Sites annelés}
\label{section-ringed-sites}

\noindent
Dans ce chapitre, nous travaillons principalement avec des faisceaux de modules
sur un site annelé. Il nous faut donc définir cette notion.

\begin{definition}
\label{definition-ringed-site}
Sites annelés.
\begin{enumerate}
\item Un {\it site annelé} est un couple $(\mathcal{C}, \mathcal{O})$, où
$\mathcal{C}$ est un site et $\mathcal{O}$ un faisceau d'anneaux sur
$\mathcal{C}$. Le faisceau $\mathcal{O}$ est appelé le {\it faisceau
structural} du site annelé.
\item Soient $(\mathcal{C}, \mathcal{O})$ et
$(\mathcal{C}', \mathcal{O}')$ des sites annelés. Un {\it morphisme de sites
annelés}
$$
(f, f^\sharp) :
(\mathcal{C}, \mathcal{O})
\longrightarrow
(\mathcal{C}', \mathcal{O}')
$$
est la donnée d'un morphisme de sites $f : \mathcal{C} \to \mathcal{C}'$ (voir
Sites, Définition \ref{sites-definition-morphism-sites}) et d'une application
de faisceaux d'anneaux $f^\sharp : f^{-1}\mathcal{O}' \to \mathcal{O}$, ce qui,
par adjonction, revient à donner une application de faisceaux d'anneaux
$f^\sharp : \mathcal{O}' \to f_*\mathcal{O}$.
\item Soient
$(f, f^\sharp) :
(\mathcal{C}_1, \mathcal{O}_1) \to (\mathcal{C}_2, \mathcal{O}_2)$ and
$(g, g^\sharp) :
(\mathcal{C}_2, \mathcal{O}_2) \to (\mathcal{C}_3, \mathcal{O}_3)$
des morphismes de sites annelés. Nous définissons alors la {\it composition des
morphismes de sites annelés} par la règle
$$
(g, g^\sharp) \circ (f, f^\sharp) = (g \circ f, f^\sharp \circ g^\sharp).
$$
Nous utilisons ici la composition des morphismes de sites définie dans Sites,
Définition \ref{sites-definition-composition-morphisms-sites}, et
$f^\sharp \circ g^\sharp$ désigne le morphisme de faisceaux d'anneaux
$$
\mathcal{O}_3 \xrightarrow{g^\sharp} g_*\mathcal{O}_2
\xrightarrow{g_*f^\sharp} g_*f_*\mathcal{O}_1 = (g \circ f)_*\mathcal{O}_1
$$
\end{enumerate}
\end{definition}






\section{Topos annelés}
\label{section-ringed-topoi}

\noindent
Un topos annelé est simplement un site annelé, à ceci près que la notion de
morphisme de topos annelés diffère de celle de morphisme de sites annelés.

\begin{definition}
\label{definition-ringed-topos}
Topos annelés.
\begin{enumerate}
\item Un {\it topos annelé} est un couple
$(\Sh(\mathcal{C}), \mathcal{O})$
où $\mathcal{C}$ est un site et $\mathcal{O}$ un faisceau d'anneaux sur
$\mathcal{C}$. Le faisceau $\mathcal{O}$ est appelé le {\it faisceau
structural} du topos annelé.
\item Soient $(\Sh(\mathcal{C}), \mathcal{O})$ et
$(\Sh(\mathcal{C}'), \mathcal{O}')$ des topos annelés. Un {\it morphisme de
topos annelés}
$$
(f, f^\sharp) :
(\Sh(\mathcal{C}), \mathcal{O})
\longrightarrow
(\Sh(\mathcal{C}'), \mathcal{O}')
$$
est la donnée d'un morphisme de topos
$f : \Sh(\mathcal{C}) \to \Sh(\mathcal{C}')$ (voir Sites, Définition
\ref{sites-definition-topos}) et d'une application de faisceaux d'anneaux
$f^\sharp : f^{-1}\mathcal{O}' \to \mathcal{O}$, ce qui, par adjonction,
revient à donner une application de faisceaux d'anneaux
$f^\sharp : \mathcal{O}' \to f_*\mathcal{O}$.
\item Soient
$(f, f^\sharp) :
(\Sh(\mathcal{C}_1), \mathcal{O}_1)
\to (\Sh(\mathcal{C}_2), \mathcal{O}_2)$ and
$(g, g^\sharp) :
(\Sh(\mathcal{C}_2), \mathcal{O}_2) \to
(\Sh(\mathcal{C}_3), \mathcal{O}_3)$
des morphismes de topos annelés. Nous définissons alors la {\it composition des
morphismes de topos annelés} par la règle
$$
(g, g^\sharp) \circ (f, f^\sharp) = (g \circ f, f^\sharp \circ g^\sharp).
$$
Nous utilisons ici la composition des morphismes de topos définie dans Sites,
Définition \ref{sites-definition-topos}, et $f^\sharp \circ g^\sharp$ désigne
le morphisme de faisceaux d'anneaux
$$
\mathcal{O}_3 \xrightarrow{g^\sharp} g_*\mathcal{O}_2
\xrightarrow{g_*f^\sharp} g_*f_*\mathcal{O}_1 = (g \circ f)_*\mathcal{O}_1
$$
\end{enumerate}
\end{definition}

\noindent
Tout morphisme de topos annelés est la composée d'une équivalence de topos
annelés et d'un morphisme de topos annelés associé à un morphisme de sites
annelés. Voici l'énoncé précis.

\begin{lemma}
\label{lemma-morphism-ringed-topoi-comes-from-morphism-ringed-sites}
Soit $(f, f^\sharp) :
(\Sh(\mathcal{C}), \mathcal{O}_\mathcal{C})
\to (\Sh(\mathcal{D}), \mathcal{O}_\mathcal{D})$
un morphisme de topos annelés. Il existe une factorisation
$$
\xymatrix{
(\Sh(\mathcal{C}), \mathcal{O}_\mathcal{C})
\ar[rr]_{(f, f^\sharp)}
\ar[d]_{(g, g^\sharp)}
& &
(\Sh(\mathcal{D}), \mathcal{O}_\mathcal{D}) \ar[d]^{(e, e^\sharp)}
\\
(\Sh(\mathcal{C}'), \mathcal{O}_{\mathcal{C}'})
\ar[rr]^{(h, h^\sharp)} & &
(\Sh(\mathcal{D}'), \mathcal{O}_{\mathcal{D}'})
}
$$
où
\begin{enumerate}
\item $g : \Sh(\mathcal{C}) \to \Sh(\mathcal{C}')$ est une équivalence de
topos induite par un foncteur cocontinu spécial
$\mathcal{C} \to \mathcal{C}'$ (voir Sites, Définition
\ref{sites-definition-special-cocontinuous-functor}),
\item $e : \Sh(\mathcal{D}) \to \Sh(\mathcal{D}')$ est une équivalence de
topos induite par un foncteur cocontinu spécial
$\mathcal{D} \to \mathcal{D}'$ (voir Sites, Définition
\ref{sites-definition-special-cocontinuous-functor}),
\item $\mathcal{O}_{\mathcal{C}'} = g_*\mathcal{O}_\mathcal{C}$
et $g^\sharp$ est l'application évidente,
\item $\mathcal{O}_{\mathcal{D}'} = e_*\mathcal{O}_\mathcal{D}$
et $e^\sharp$ est l'application évidente,
\item les sites $\mathcal{C}'$ et $\mathcal{D}'$ possèdent des objets finaux
et des produits fibrés (autrement dit, toutes les limites finies),
\item $h$ est un morphisme de sites induit par un foncteur continu
$u : \mathcal{D}' \to \mathcal{C}'$ qui commute à toutes les limites finies
(autrement dit, qui satisfait les hypothèses de Sites, Proposition
\ref{sites-proposition-get-morphism}), et
\item étant données des familles quelconques de faisceaux $\mathcal{F}_i$
(resp.\ $\mathcal{G}_j$) sur $\mathcal{C}$ (resp.\ $\mathcal{D}$), nous pouvons
supposer que chacun est un faisceau représentable sur $\mathcal{C}'$
(resp.\ $\mathcal{D}'$).
\end{enumerate}
De plus, si $(f, f^\sharp)$ est une équivalence de topos annelés, nous pouvons
choisir le diagramme de sorte que
$\mathcal{C}' = \mathcal{D}'$,
$\mathcal{O}_{\mathcal{C}'} = \mathcal{O}_{\mathcal{D}'}$
et que $(h, h^\sharp)$ soit l'identité.
\end{lemma}

\begin{proof}
Cela résulte de Sites, Lemme
\ref{sites-lemma-morphism-topoi-comes-from-morphism-sites}, et de Sites,
Remarques \ref{sites-remark-morphism-topoi-comes-from-morphism-sites} et
\ref{sites-remark-equivalence-topoi-comes-from-morphism-sites}. Il suffit de
transporter simultanément les faisceaux d'anneaux. Nous omettons certains
détails.
\end{proof}







\section{2-morphismes de topos annelés}
\label{section-2-category}

\noindent
Cette brève section concerne la notion de $2$-morphisme de topos annelés.

\begin{definition}
\label{definition-2-morphism-ringed-topoi}
Soient
$f, g :
(\Sh(\mathcal{C}), \mathcal{O}_\mathcal{C})
\to
(\Sh(\mathcal{D}), \mathcal{O}_\mathcal{D})$
deux morphismes de topos annelés. Un {\it 2-morphisme de $f$ vers $g$} est
la donnée d'une transformation de foncteurs $t : f_* \to g_*$ telle que
$$
\xymatrix{
& \mathcal{O}_\mathcal{D}
\ar[ld]_{f^\sharp}
\ar[rd]^{g^\sharp} \\
f_*\mathcal{O}_\mathcal{C} \ar[rr]^t & &
g_*\mathcal{O}_\mathcal{C}
}
$$
soit commutatif.
\end{definition}

\noindent
Nous représentons parfois $t$ graphiquement de la manière suivante :
$$
\xymatrix{
(\Sh(\mathcal{C}), \mathcal{O}_\mathcal{C})
\rrtwocell^f_g{t}
&
&
(\Sh(\mathcal{D}), \mathcal{O}_\mathcal{D})
}
$$
Comme dans
Sites, Section \ref{sites-section-2-category}
donner un 2-morphisme $t : f_* \to g_*$ équivaut à donner
$t : g^{-1} \to f^{-1}$ (habituellement noté par le même symbole), de sorte
que le diagramme
$$
\xymatrix{
f^{-1}\mathcal{O}_\mathcal{D}
\ar[rd]_{f^\sharp}  & &
g^{-1}\mathcal{O}_\mathcal{D} \ar[ll]^t \ar[ld]^{g^\sharp} \\
& \mathcal{O}_\mathcal{C}
}
$$
soit commutatif. Comme dans
Sites, Section \ref{sites-section-2-category}
les axiomes d'une 2-catégorie stricte sont satisfaits pour les compositions
horizontale et verticale définies comme dans la référence citée.










\section{Préfaisceaux de modules}
\label{section-presheaves-modules}

\noindent
Soit $\mathcal{C}$ une catégorie. Soit $\mathcal{O}$ un préfaisceau d'anneaux
sur $\mathcal{C}$. Nous n'avons pas encore défini les préfaisceaux de
$\mathcal{O}$-modules ; faisons-le maintenant.

\begin{definition}
\label{definition-presheaf-modules}
Soit $\mathcal{C}$ une catégorie et soit $\mathcal{O}$ un préfaisceau
d'anneaux sur $\mathcal{C}$.
\begin{enumerate}
\item Un {\it préfaisceau de $\mathcal{O}$-modules} est la donnée d'un
préfaisceau abélien $\mathcal{F}$ et d'une application de préfaisceaux
d'ensembles
$$
\mathcal{O} \times \mathcal{F} \longrightarrow \mathcal{F}
$$
tel que, pour tout objet $U$ de $\mathcal{C}$, l'application
$\mathcal{O}(U) \times \mathcal{F}(U) \to \mathcal{F}(U)$
définisse une structure de $\mathcal{O}(U)$-module sur le groupe abélien
$\mathcal{F}(U)$.
\item Un {\it morphisme $\varphi : \mathcal{F} \to \mathcal{G}$ de
préfaisceaux de $\mathcal{O}$-modules} est un morphisme de préfaisceaux
abéliens $\varphi : \mathcal{F} \to \mathcal{G}$ tel que le diagramme
$$
\xymatrix{
\mathcal{O} \times \mathcal{F} \ar[r] \ar[d]_{\text{id} \times \varphi} &
\mathcal{F} \ar[d]^{\varphi} \\
\mathcal{O} \times \mathcal{G} \ar[r] &
\mathcal{G}
}
$$
soit commutatif.
\item L'ensemble des morphismes de $\mathcal{O}$-modules ci-dessus est noté
$\Hom_\mathcal{O}(\mathcal{F}, \mathcal{G})$.
\item La catégorie des préfaisceaux de $\mathcal{O}$-modules est notée
$\textit{PMod}(\mathcal{O})$.
\end{enumerate}
\end{definition}

\noindent
Supposons que $\mathcal{O}_1 \to \mathcal{O}_2$ soit un morphisme de
préfaisceaux d'anneaux sur la catégorie $\mathcal{C}$. Si $\mathcal{F}$ est un
préfaisceau de $\mathcal{O}_2$-modules, nous pouvons alors considérer
$\mathcal{F}$ comme un préfaisceau de $\mathcal{O}_1$-modules au moyen de la composée
$$
\mathcal{O}_1 \times \mathcal{F}
\to
\mathcal{O}_2 \times \mathcal{F}
\to
\mathcal{F}.
$$
Nous notons parfois ce préfaisceau
$\mathcal{F}_{\mathcal{O}_1}$ pour signaler la restriction des scalaires. Nous
appelons ce préfaisceau la {\it restriction de $\mathcal{F}$}. Nous obtenons le
foncteur de restriction des scalaires
$$
\textit{PMod}(\mathcal{O}_2)
\longrightarrow
\textit{PMod}(\mathcal{O}_1)
$$

\medskip\noindent
Réciproquement, étant donné un préfaisceau de $\mathcal{O}_1$-modules
$\mathcal{G}$, nous pouvons construire un préfaisceau de
$\mathcal{O}_2$-modules
$\mathcal{O}_2 \otimes_{p, \mathcal{O}_1} \mathcal{G}$
par la règle
$$
U \longmapsto
\left(\mathcal{O}_2 \otimes_{p, \mathcal{O}_1} \mathcal{G}\right)(U)
=
\mathcal{O}_2(U) \otimes_{\mathcal{O}_1(U)} \mathcal{G}(U)
$$
pour $U \in \Ob(\mathcal{C})$, avec les applications de restriction évidentes.
L'indice $p$ signifie « préfaisceau » et non « point ». Ce préfaisceau est
appelé le préfaisceau produit tensoriel. Nous obtenons le foncteur
d'{\it extension des scalaires}
$$
\textit{PMod}(\mathcal{O}_1)
\longrightarrow
\textit{PMod}(\mathcal{O}_2)
$$

\begin{lemma}
\label{lemma-adjointness-tensor-restrict-presheaves}
Avec $\mathcal{C}$, $\mathcal{O}_1 \to \mathcal{O}_2$, $\mathcal{F}$ et
$\mathcal{G}$ comme ci-dessus, il existe une bijection canonique
$$
\Hom_{\mathcal{O}_1}(\mathcal{G}, \mathcal{F}_{\mathcal{O}_1})
=
\Hom_{\mathcal{O}_2}(
\mathcal{O}_2 \otimes_{p, \mathcal{O}_1} \mathcal{G},
\mathcal{F}
)
$$
Autrement dit, les foncteurs de restriction et d'extension des scalaires
définis ci-dessus sont adjoints l'un de l'autre.
\end{lemma}

\begin{proof}
Cela résulte du fait que, pour un homomorphisme d'anneaux $A \to B$, les
foncteurs de restriction et d'extension des scalaires sont adjoints l'un de
l'autre.
\end{proof}


\section{Faisceaux de modules}
\label{section-sheaves-modules}

\begin{definition}
\label{definition-sheaf-modules}
Soit $\mathcal{C}$ un site. Soit $\mathcal{O}$ un faisceau d'anneaux sur
$\mathcal{C}$.
\begin{enumerate}
\item Un {\it faisceau de $\mathcal{O}$-modules} est un préfaisceau de
$\mathcal{O}$-modules $\mathcal{F}$, voir Définition
\ref{definition-presheaf-modules}, tel que le préfaisceau de groupes abéliens
sous-jacent à $\mathcal{F}$ soit un faisceau.
\item Un {\it morphisme de faisceaux de $\mathcal{O}$-modules} est un
morphisme de préfaisceaux de $\mathcal{O}$-modules.
\item Étant donnés des faisceaux de $\mathcal{O}$-modules
$\mathcal{F}$ et $\mathcal{G}$, nous notons
$\Hom_\mathcal{O}(\mathcal{F}, \mathcal{G})$ l'ensemble de leurs morphismes de
faisceaux de $\mathcal{O}$-modules.
\item La catégorie des faisceaux de $\mathcal{O}$-modules est notée
$\textit{Mod}(\mathcal{O})$.
\end{enumerate}
\end{definition}

\noindent
Cette définition garde un certain sens lorsque $\mathcal{O}$ n'est qu'un
préfaisceau d'anneaux, bien que nous ne connaissions aucun exemple où cela soit
utile. Nous éviterons donc l'expression « faisceaux de $\mathcal{O}$-modules »
lorsque $\mathcal{O}$ n'est pas un faisceau d'anneaux.



\section{Faisceautisation des préfaisceaux de modules}
\label{section-sheafification-presheaves-modules}

\begin{lemma}
\label{lemma-sheafification-presheaf-modules}
Soit $\mathcal{C}$ un site. Soit $\mathcal{O}$ un préfaisceau d'anneaux sur
$\mathcal{C}$ et soit $\mathcal{F}$ un préfaisceau de $\mathcal{O}$-modules.
Soit $\mathcal{O}^\#$ le faisceau associé à $\mathcal{O}$ comme préfaisceau
d'anneaux, voir Sites, Section
\ref{sites-section-sheaves-algebraic-structures}. Soit $\mathcal{F}^\#$ le
faisceau associé à $\mathcal{F}$ comme préfaisceau de groupes abéliens. Il
existe une unique application de faisceaux d'ensembles
$$
\mathcal{O}^\# \times \mathcal{F}^\#
\longrightarrow
\mathcal{F}^\#
$$
qui rende commutatif le diagramme
$$
\xymatrix{
\mathcal{O} \times \mathcal{F} \ar[r] \ar[d] &
\mathcal{F} \ar[d] \\
\mathcal{O}^\# \times \mathcal{F}^\# \ar[r] &
\mathcal{F}^\#
}
$$
et qui munisse $\mathcal{F}^\#$ d'une structure de faisceau de
$\mathcal{O}^\#$-modules. De plus, si $\mathcal{G}$ est un faisceau de
$\mathcal{O}^\#$-modules, tout morphisme de préfaisceaux de
$\mathcal{O}$-modules $\mathcal{F} \to \mathcal{G}$ (où $\mathcal{G}$ est
considéré par restriction comme un $\mathcal{O}$-module) se factorise de façon
unique sous la forme $\mathcal{F} \to \mathcal{F}^\# \to \mathcal{G}$, où
$\mathcal{F}^\# \to \mathcal{G}$ est un morphisme de
$\mathcal{O}^\#$-modules.
\end{lemma}

\begin{proof}
Omis.
\end{proof}

\noindent
Cela signifie précisément que le foncteur
$i : \textit{Mod}(\mathcal{O}^\#) \to \textit{PMod}(\mathcal{O})$
(qui combine la restriction des scalaires et l'inclusion des faisceaux dans
les préfaisceaux) et le foncteur de faisceautisation du lemme
${}^\# : \textit{PMod}(\mathcal{O}) \to \textit{Mod}(\mathcal{O}^\#)$
sont adjoints. En formule,
$$
\Mor_{\textit{PMod}(\mathcal{O})}(\mathcal{F}, i\mathcal{G})
=
\Mor_{\textit{Mod}(\mathcal{O}^\#)}(\mathcal{F}^\#, \mathcal{G})
$$
Un cas important se présente lorsque $\mathcal{O}$ est déjà un faisceau
d'anneaux. La formule devient alors
$$
\Mor_{\textit{PMod}(\mathcal{O})}(\mathcal{F}, i\mathcal{G})
=
\Mor_{\textit{Mod}(\mathcal{O})}(\mathcal{F}^\#, \mathcal{G})
$$
puisque $\mathcal{O} = \mathcal{O}^\#$ dans ce cas.

\begin{lemma}
\label{lemma-sheafification-exact}
Soit $\mathcal{C}$ un site. Soit $\mathcal{O}$ un préfaisceau d'anneaux sur
$\mathcal{C}$. Le foncteur de faisceautisation
$$
\textit{PMod}(\mathcal{O}) \longrightarrow \textit{Mod}(\mathcal{O}^\#), \quad
\mathcal{F} \longmapsto \mathcal{F}^\#
$$
est exact.
\end{lemma}

\begin{proof}
Cela est vrai parce que la faisceautisation
$\textit{PAb}(\mathcal{C}) \to \textit{Ab}(\mathcal{C})$ est exacte. Voir la
discussion de la Section \ref{section-abelian-sheaves}.
\end{proof}

\noindent
Soit $\mathcal{C}$ un site. Soit
$\mathcal{O}_1 \to \mathcal{O}_2$ un morphisme de faisceaux d'anneaux sur
$\mathcal{C}$. À la Section \ref{section-presheaves-modules}, nous avons défini
les foncteurs de restriction et d'extension des scalaires correspondants sur
les préfaisceaux de modules.

\medskip\noindent
Si $\mathcal{F}$ est un faisceau de $\mathcal{O}_2$-modules, la restriction
$\mathcal{F}_{\mathcal{O}_1}$ de $\mathcal{F}$ est manifestement un faisceau de
$\mathcal{O}_1$-modules. Nous obtenons le foncteur de restriction des scalaires
$$
\textit{Mod}(\mathcal{O}_2)
\longrightarrow
\textit{Mod}(\mathcal{O}_1)
$$

\medskip\noindent
Réciproquement, étant donné un faisceau de $\mathcal{O}_1$-modules
$\mathcal{G}$, le préfaisceau de $\mathcal{O}_2$-modules
$\mathcal{O}_2 \otimes_{p, \mathcal{O}_1} \mathcal{G}$
n'est en général pas un faisceau. Nous définissons donc le {\it faisceau produit
tensoriel}
$\mathcal{O}_2 \otimes_{\mathcal{O}_1} \mathcal{G}$
par la formule
$$
\mathcal{O}_2 \otimes_{\mathcal{O}_1} \mathcal{G}
=
(\mathcal{O}_2 \otimes_{p, \mathcal{O}_1} \mathcal{G})^\#
$$
comme le faisceau associé à la construction précédente sur les préfaisceaux.
Nous obtenons le foncteur d'{\it extension des scalaires}
$$
\textit{Mod}(\mathcal{O}_1)
\longrightarrow
\textit{Mod}(\mathcal{O}_2)
$$

\begin{lemma}
\label{lemma-adjointness-tensor-restrict}
Avec $X$, $\mathcal{O}_1$, $\mathcal{O}_2$, $\mathcal{F}$ et $\mathcal{G}$
comme ci-dessus, il existe une bijection canonique
$$
\Hom_{\mathcal{O}_1}(\mathcal{G}, \mathcal{F}_{\mathcal{O}_1})
=
\Hom_{\mathcal{O}_2}(
\mathcal{O}_2 \otimes_{\mathcal{O}_1} \mathcal{G},
\mathcal{F}
)
$$
Autrement dit, les foncteurs de restriction et d'extension des scalaires sont
adjoints l'un de l'autre.
\end{lemma}

\begin{proof}
Cela résulte du Lemme \ref{lemma-adjointness-tensor-restrict-presheaves} et du
fait que
$\Hom_{\mathcal{O}_2}(
\mathcal{O}_2 \otimes_{\mathcal{O}_1} \mathcal{G},
\mathcal{F}
)
=
\Hom_{\mathcal{O}_2}(
\mathcal{O}_2 \otimes_{p, \mathcal{O}_1} \mathcal{G},
\mathcal{F}
)$
puisque $\mathcal{F}$ est un faisceau.
\end{proof}

\begin{lemma}
\label{lemma-epimorphism-modules}
Soit $\mathcal{C}$ un site. Soit $\mathcal{O} \to \mathcal{O}'$ un
épimorphisme de faisceaux d'anneaux. Soient
$\mathcal{G}_1, \mathcal{G}_2$ des $\mathcal{O}'$-modules. Alors
$$
\Hom_{\mathcal{O}'}(\mathcal{G}_1, \mathcal{G}_2) =
\Hom_\mathcal{O}(\mathcal{G}_1, \mathcal{G}_2).
$$
Autrement dit, le foncteur de restriction
$\textit{Mod}(\mathcal{O}') \to \textit{Mod}(\mathcal{O})$ est pleinement
fidèle.
\end{lemma}

\begin{proof}
C'est la version faisceautique d'Algèbre, Lemme
\ref{algebra-lemma-epimorphism-modules}, et elle se démontre exactement de la
même manière.
\end{proof}




\section{Morphismes de topos et faisceaux de modules}
\label{section-sheaves-modules-functorial}

\noindent
Tout ce qui suit est immédiat. Nous formulons les résultats pour les morphismes
de topos, mais ils valent bien entendu aussi pour les morphismes de sites.

\begin{lemma}
\label{lemma-pushforward-module}
Soient $\mathcal{C}$ et $\mathcal{D}$ des sites. Soit
$f : \Sh(\mathcal{C}) \to \Sh(\mathcal{D})$ un morphisme de topos. Soit
$\mathcal{O}$ un faisceau d'anneaux sur $\mathcal{C}$ et soit $\mathcal{F}$ un
faisceau de $\mathcal{O}$-modules. Il existe une application naturelle de
faisceaux d'ensembles
$$
f_*\mathcal{O} \times f_*\mathcal{F}
\longrightarrow
f_*\mathcal{F}
$$
qui munit $f_*\mathcal{F}$ d'une structure de faisceau de
$f_*\mathcal{O}$-modules. Cette construction est fonctorielle en $\mathcal{F}$.
\end{lemma}

\begin{proof}
Notons $\mu : \mathcal{O} \times \mathcal{F} \to \mathcal{F}$ l'application
de multiplication. Rappelons que $f_*$ (sur les faisceaux d'ensembles) est
exact à gauche et commute donc aux produits. Ainsi $f_*\mu$ est une application
du type indiqué, ce qui démontre le lemme.
\end{proof}

\begin{lemma}
\label{lemma-pullback-module}
Soient $\mathcal{C}$ et $\mathcal{D}$ des sites. Soit
$f : \Sh(\mathcal{C}) \to \Sh(\mathcal{D})$ un morphisme de topos. Soit
$\mathcal{O}$ un faisceau d'anneaux sur $\mathcal{D}$ et soit $\mathcal{G}$ un
faisceau de $\mathcal{O}$-modules. Il existe une application naturelle de
faisceaux d'ensembles
$$
f^{-1}\mathcal{O} \times f^{-1}\mathcal{G}
\longrightarrow
f^{-1}\mathcal{G}
$$
qui munit $f^{-1}\mathcal{G}$ d'une structure de faisceau de
$f^{-1}\mathcal{O}$-modules. Cette construction est fonctorielle en
$\mathcal{G}$.
\end{lemma}

\begin{proof}
Notons $\mu : \mathcal{O} \times \mathcal{G} \to \mathcal{G}$ l'application
de multiplication. Rappelons que $f^{-1}$ (sur les faisceaux d'ensembles) est
exact et commute donc aux produits. Ainsi $f^{-1}\mu$ est une application du
type indiqué, ce qui démontre le lemme.
\end{proof}

\begin{lemma}
\label{lemma-adjoint-push-pull-modules}
Soient $\mathcal{C}$ et $\mathcal{D}$ des sites. Soit
$f : \Sh(\mathcal{C}) \to \Sh(\mathcal{D})$ un morphisme de topos. Soit
$\mathcal{O}$ un faisceau d'anneaux sur $\mathcal{D}$. Soit $\mathcal{G}$ un
faisceau de $\mathcal{O}$-modules et soit $\mathcal{F}$ un faisceau de
$f^{-1}\mathcal{O}$-modules. Alors
$$
\Mor_{\textit{Mod}(f^{-1}\mathcal{O})}(f^{-1}\mathcal{G}, \mathcal{F})
=
\Mor_{\textit{Mod}(\mathcal{O})}(\mathcal{G}, f_*\mathcal{F}).
$$
Nous utilisons ici les Lemmes \ref{lemma-pullback-module} et
\ref{lemma-pushforward-module}, et nous considérons $f_*\mathcal{F}$ comme un
$\mathcal{O}$-module par restriction des scalaires suivant
$\mathcal{O} \to f_*f^{-1}\mathcal{O}$.
\end{lemma}

\begin{proof}
Remarquons d'abord que
$$
\Mor_{\textit{Ab}(\mathcal{C})}(f^{-1}\mathcal{G}, \mathcal{F})
=
\Mor_{\textit{Ab}(\mathcal{D})}(\mathcal{G}, f_*\mathcal{F}).
$$
d'après Sites, Proposition
\ref{sites-proposition-functoriality-algebraic-structures-topoi}. Supposons que
$\alpha : f^{-1}\mathcal{G} \to \mathcal{F}$ et
$\beta : \mathcal{G} \to f_*\mathcal{F}$ soient des morphismes de faisceaux
abéliens qui se correspondent par la formule ci-dessus. Il faut montrer que
$\alpha$ est $f^{-1}\mathcal{O}$-linéaire si et seulement si $\beta$ est
$\mathcal{O}$-linéaire. Supposons par exemple que $\alpha$ soit
$f^{-1}\mathcal{O}$-linéaire. Alors $f_*\alpha$ est manifestement
$f_*f^{-1}\mathcal{O}$-linéaire, donc $\mathcal{O}$-linéaire puisque la
restriction des scalaires est fonctorielle. Il suffit donc de montrer que
l'application d'adjonction $\mathcal{G} \to f_*f^{-1}\mathcal{G}$ est
$\mathcal{O}$-linéaire. Comme $f_*$ et $f^{-1}$ commutent tous deux aux
produits (sur les faisceaux d'ensembles), cela revient à montrer que
$$
\xymatrix{
\mathcal{O} \times \mathcal{G} \ar[r] \ar[d] &
f_*f^{-1}(\mathcal{O} \times \mathcal{G}) \ar[d] \\
\mathcal{G} \ar[r] & f_*f^{-1}\mathcal{G}
}
$$
est commutatif. Cela résulte de ce que l'application d'adjonction
$\text{id}_{\Sh(\mathcal{D})} \to f_*f^{-1}$ est une transformation de
foncteurs. Nous omettons la démonstration de l'implication « $\beta$ linéaire
$\Rightarrow$ $\alpha$ linéaire ».
\end{proof}

\begin{lemma}
\label{lemma-adjoint-pull-push-modules}
Soient $\mathcal{C}$ et $\mathcal{D}$ des sites. Soit
$f : \Sh(\mathcal{C}) \to \Sh(\mathcal{D})$ un morphisme de topos. Soit
$\mathcal{O}$ un faisceau d'anneaux sur $\mathcal{C}$. Soit $\mathcal{F}$ un
faisceau de $\mathcal{O}$-modules et soit $\mathcal{G}$ un faisceau de
$f_*\mathcal{O}$-modules. Alors
$$
\Mor_{\textit{Mod}(\mathcal{O})}(
\mathcal{O} \otimes_{f^{-1}f_*\mathcal{O}} f^{-1}\mathcal{G}, \mathcal{F})
=
\Mor_{\textit{Mod}(f_*\mathcal{O})}(\mathcal{G}, f_*\mathcal{F}).
$$
Nous utilisons ici les Lemmes \ref{lemma-pullback-module} et
\ref{lemma-pushforward-module}, ainsi que l'application canonique
$f^{-1}f_*\mathcal{O} \to \mathcal{O}$ dans la définition du produit tensoriel.
\end{lemma}

\begin{proof}
Remarquons que
$$
\Mor_{\textit{Mod}(\mathcal{O})}(
\mathcal{O} \otimes_{f^{-1}f_*\mathcal{O}} f^{-1}\mathcal{G}, \mathcal{F})
=
\Mor_{\textit{Mod}(f^{-1}f_*\mathcal{O})}(
f^{-1}\mathcal{G}, \mathcal{F}_{f^{-1}f_*\mathcal{O}})
$$
d'après le Lemme \ref{lemma-adjointness-tensor-restrict}. Le résultat découle
donc du Lemme \ref{lemma-adjoint-push-pull-modules}.
\end{proof}






\section{Morphismes de topos annelés et modules}
\label{section-functoriality-modules}

\noindent
Nous avons maintenant introduit assez de notations pour définir l'image inverse
et l'image directe des modules le long d'un morphisme de topos annelés.

\begin{definition}
\label{definition-pushforward}
Soit
$(f, f^\sharp) :
(\Sh(\mathcal{C}), \mathcal{O}_\mathcal{C})
\to
(\Sh(\mathcal{D}), \mathcal{O}_\mathcal{D})$
un morphisme de topos annelés ou de sites annelés.
\begin{enumerate}
\item Soit $\mathcal{F}$ un faisceau de $\mathcal{O}_\mathcal{C}$-modules. Nous
définissons l'{\it image directe} de $\mathcal{F}$ comme le faisceau de
$\mathcal{O}_\mathcal{D}$-modules dont le faisceau de groupes abéliens
sous-jacent est $f_*\mathcal{F}$ et dont la structure de module est obtenue par
restriction des scalaires suivant
$f^\sharp : \mathcal{O}_\mathcal{D} \to f_*\mathcal{O}_\mathcal{C}$ à partir
de la structure de module
$$
f_*\mathcal{O}_\mathcal{C} \times f_*\mathcal{F}
\longrightarrow
f_*\mathcal{F}
$$
du Lemme \ref{lemma-pushforward-module}.
\item Soit $\mathcal{G}$ un faisceau de $\mathcal{O}_\mathcal{D}$-modules. Nous
définissons l'{\it image inverse} $f^*\mathcal{G}$ comme le faisceau de
$\mathcal{O}_\mathcal{C}$-modules donné par la formule
$$
f^*\mathcal{G}
=
\mathcal{O}_\mathcal{C} \otimes_{f^{-1}\mathcal{O}_\mathcal{D}}
f^{-1}\mathcal{G}
$$
où l'homomorphisme d'anneaux
$f^{-1}\mathcal{O}_\mathcal{D} \to \mathcal{O}_\mathcal{C}$
est $f^\sharp$, et où la structure de module est celle du Lemme
\ref{lemma-pullback-module}.
\end{enumerate}
\end{definition}

\noindent
Nous avons ainsi défini des foncteurs
\begin{eqnarray*}
f_* : \textit{Mod}(\mathcal{O}_\mathcal{C})
& \longrightarrow &
\textit{Mod}(\mathcal{O}_\mathcal{D}) \\
f^* : \textit{Mod}(\mathcal{O}_\mathcal{D})
& \longrightarrow &
\textit{Mod}(\mathcal{O}_\mathcal{C})
\end{eqnarray*}
Le dernier résultat sur ces foncteurs affirme qu'ils sont bien adjoints, comme
prévu.

\begin{lemma}
\label{lemma-adjoint-pullback-pushforward-modules}
Soit
$(f, f^\sharp) :
(\Sh(\mathcal{C}), \mathcal{O}_\mathcal{C})
\to
(\Sh(\mathcal{D}), \mathcal{O}_\mathcal{D})$
un morphisme de topos annelés ou de sites annelés. Soit $\mathcal{F}$ un
faisceau de $\mathcal{O}_\mathcal{C}$-modules et soit $\mathcal{G}$ un faisceau
de $\mathcal{O}_\mathcal{D}$-modules. Il existe une bijection canonique
$$
\Hom_{\mathcal{O}_\mathcal{C}}(f^*\mathcal{G}, \mathcal{F})
=
\Hom_{\mathcal{O}_\mathcal{D}}(\mathcal{G}, f_*\mathcal{F}).
$$
Autrement dit, le foncteur $f^*$ est adjoint à gauche de $f_*$.
\end{lemma}

\begin{proof}
Cela résulte des calculs précédents :
\begin{eqnarray*}
\Hom_{\mathcal{O}_\mathcal{C}}(f^*\mathcal{G}, \mathcal{F})
& = &
\Mor_{\textit{Mod}(\mathcal{O}_\mathcal{C})}(
\mathcal{O}_\mathcal{C}
\otimes_{f^{-1}\mathcal{O}_\mathcal{D}} f^{-1}\mathcal{G},
\mathcal{F}) \\
& = &
\Mor_{\textit{Mod}(f^{-1}\mathcal{O}_\mathcal{D})}(
f^{-1}\mathcal{G}, \mathcal{F}_{f^{-1}\mathcal{O}_\mathcal{D}}) \\
& = &
\Hom_{\mathcal{O}_\mathcal{D}}(\mathcal{G}, f_*\mathcal{F}).
\end{eqnarray*}
Nous utilisons ici les Lemmes \ref{lemma-adjointness-tensor-restrict} et
\ref{lemma-adjoint-push-pull-modules}.
\end{proof}

\begin{lemma}
\label{lemma-push-pull-composition-modules}
$(f, f^\sharp) :
(\Sh(\mathcal{C}_1), \mathcal{O}_1)
\to (\Sh(\mathcal{C}_2), \mathcal{O}_2)$ and
$(g, g^\sharp) :
(\Sh(\mathcal{C}_2), \mathcal{O}_2) \to
(\Sh(\mathcal{C}_3), \mathcal{O}_3)$
soient des morphismes de topos annelés. Il existe des isomorphismes canoniques
de foncteurs
$(g \circ f)_* \cong g_* \circ f_*$ and
$(g \circ f)^* \cong f^* \circ g^*$.
\end{lemma}

\begin{proof}
Cela résulte immédiatement des définitions.
\end{proof}





\section{La catégorie abélienne des faisceaux de modules}
\label{section-kernels}

\noindent
Soit $(\Sh(\mathcal{C}), \mathcal{O})$ un topos annelé. Soient
$\mathcal{F}$ et $\mathcal{G}$ des faisceaux de $\mathcal{O}$-modules, voir
Faisceaux, Définition \ref{sheaves-definition-sheaf-modules}. Soient
$\varphi, \psi : \mathcal{F} \to \mathcal{G}$ des morphismes de faisceaux de
$\mathcal{O}$-modules. Nous définissons
$\varphi + \psi : \mathcal{F} \to \mathcal{G}$ comme la somme de $\varphi$ et
$\psi$ en tant que morphismes de faisceaux abéliens. C'est manifestement encore
une application de $\mathcal{O}$-modules. La composition des applications de
$\mathcal{O}$-modules est en outre bilinéaire pour cette addition. Ainsi,
$\textit{Mod}(\mathcal{O})$ est une catégorie préadditive, voir Homologie,
Définition \ref{homology-definition-preadditive}.

\medskip\noindent
Nous noterons $0$ le faisceau de $\mathcal{O}$-modules qui prend la valeur
constante $\{0\}$ sur tout objet $U$ de $\mathcal{C}$. C'est manifestement à
la fois un objet final et un objet initial de $\textit{Mod}(\mathcal{O})$.
Étant donné un morphisme de $\mathcal{O}$-modules
$\varphi : \mathcal{F} \to \mathcal{G}$, les assertions suivantes sont
équivalentes : (a) $\varphi$ est nul, (b) $\varphi$ se factorise par $0$, et
(c) $\varphi$ est nul sur les sections au-dessus de tout objet $U$.

\medskip\noindent
En outre, étant donnés deux faisceaux $\mathcal{F}$ et $\mathcal{G}$ de
$\mathcal{O}$-modules, nous pouvons définir leur somme directe par
$$
\mathcal{F} \oplus \mathcal{G} = \mathcal{F} \times \mathcal{G}
$$
munie des applications évidentes $(i, j, p, q)$ comme dans Homologie,
Définition \ref{homology-definition-direct-sum}. Ainsi
$\textit{Mod}(\mathcal{O})$ est une catégorie additive, voir Homologie,
Définition \ref{homology-definition-additive-category}.

\medskip\noindent
Soit $\varphi : \mathcal{F} \to \mathcal{G}$ un morphisme de
$\mathcal{O}$-modules. Nous pouvons définir $\Ker(\varphi)$ comme le noyau de
$\varphi$ en tant qu'application de faisceaux abéliens. D'après la Section
\ref{section-abelian-sheaves}, c'est le sous-faisceau de $\mathcal{F}$ dont les
sections sont
$$
\Ker(\varphi)(U) =
\{ s \in \mathcal{F}(U) \mid \varphi(s) = 0 \text{ dans } \mathcal{G}(U)\}
$$
pour tout objet $U$ de $\mathcal{C}$. On vérifie aisément qu'il s'agit bien
d'un noyau dans la catégorie des $\mathcal{O}$-modules. Autrement dit, un
morphisme $\alpha : \mathcal{H} \to \mathcal{F}$ se factorise par
$\Ker(\varphi)$ si et seulement si $\varphi \circ \alpha = 0$.

\medskip\noindent
De même, nous définissons $\Coker(\varphi)$ comme le conoyau de $\varphi$ en
tant qu'application de faisceaux abéliens. Il existe une unique application de
multiplication
$$
\mathcal{O} \times \Coker(\varphi) \longrightarrow \Coker(\varphi)
$$
qui fait de $\mathcal{G} \to \Coker(\varphi)$ un morphisme de
$\mathcal{O}$-modules (vérification omise). L'application
$\mathcal{G} \to \Coker(\varphi)$ est surjective (comme application de
faisceaux d'ensembles, voir Section \ref{section-abelian-sheaves}). Pour montrer
que $\Coker(\varphi)$ est un conoyau dans $\textit{Mod}(\mathcal{O})$,
remarquons que si $\beta : \mathcal{G} \to \mathcal{H}$ est un morphisme de
$\mathcal{O}$-modules tel que $\beta \circ \varphi$ soit nul, alors, pour tout
objet $U$ de $\mathcal{C}$, $\beta$ induit une application de
$\mathcal{G}(U)/\varphi(\mathcal{F}(U))$ vers $\mathcal{H}(U)$. Par la
propriété universelle de la faisceautisation (voir Faisceaux, Lemme
\ref{sheaves-lemma-sheafification-presheaf-modules}), nous obtenons une
application canonique $\Coker(\varphi) \to \mathcal{H}$ telle que $\beta$ soit
la composée $\mathcal{G} \to \Coker(\varphi) \to \mathcal{H}$. Le morphisme
$\Coker(\varphi) \to \mathcal{H}$ est unique en raison de la surjectivité
mentionnée ci-dessus.

\begin{lemma}
\label{lemma-abelian}
Soit $(\Sh(\mathcal{C}), \mathcal{O})$ un topos annelé. La catégorie
$\textit{Mod}(\mathcal{O})$ est abélienne. Le foncteur d'oubli
$\textit{Mod}(\mathcal{O}) \to \textit{Ab}(\mathcal{C})$
est exact ; les noyaux, les conoyaux et l'exactitude des
$\mathcal{O}$-modules correspondent donc aux notions analogues pour les
faisceaux abéliens.
\end{lemma}

\begin{proof}
Nous avons vu ci-dessus que $\textit{Mod}(\mathcal{O})$ est une catégorie
additive, qu'elle possède noyaux et conoyaux, et que
$\textit{Mod}(\mathcal{O}) \to \textit{Ab}(\mathcal{C})$ les préserve. D'après
Homologie, Définition \ref{homology-definition-abelian-category}, il reste à
montrer que l'image et la coimage coïncident. Cela est clair puisque c'est vrai
dans $\textit{Ab}(\mathcal{C})$. Le lemme en résulte.
\end{proof}

\begin{lemma}
\label{lemma-limits-colimits}
Soit $(\Sh(\mathcal{C}), \mathcal{O})$ un topos annelé. Toutes les limites et
toutes les colimites existent dans $\textit{Mod}(\mathcal{O})$, et le foncteur
d'oubli
$\textit{Mod}(\mathcal{O}) \to \textit{Ab}(\mathcal{C})$
commute avec elles. De plus, les colimites filtrantes sont exactes.
\end{lemma}

\begin{proof}
La dernière assertion résulte de la première, car les colimites filtrantes sont
exactes dans $\textit{Ab}(\mathcal{C})$ d'après le Lemme
\ref{lemma-limits-colimits-abelian-sheaves}. Soit
$\mathcal{I} \to \textit{Mod}(\mathcal{C})$, $i \mapsto \mathcal{F}_i$, un
diagramme. Soit $\lim_i \mathcal{F}_i$ sa limite dans
$\textit{Ab}(\mathcal{C})$. La description de cette limite au Lemme
\ref{lemma-limits-colimits-abelian-sheaves} montre immédiatement qu'il existe
une multiplication
$$
\mathcal{O} \times \lim_i \mathcal{F}_i
\longrightarrow
\lim_i \mathcal{F}_i
$$
qui munit $\lim_i \mathcal{F}_i$ d'une structure de faisceau de
$\mathcal{O}$-modules. On voit aisément qu'il s'agit de la limite du diagramme
dans $\textit{Mod}(\mathcal{C})$. Soit $\colim_i \mathcal{F}_i$ la colimite du
diagramme dans $\textit{PAb}(\mathcal{C})$. La description de cette colimite
dans la démonstration du Lemme
\ref{lemma-limits-colimits-abelian-presheaves} montre immédiatement qu'il
existe une multiplication
$$
\mathcal{O} \times \colim_i \mathcal{F}_i
\longrightarrow
\colim_i \mathcal{F}_i
$$
qui munit $\colim_i \mathcal{F}_i$ d'une structure de préfaisceau de
$\mathcal{O}$-modules. En faisceautisant, nous obtenons le faisceau de
$\mathcal{O}$-modules $(\colim_i \mathcal{F}_i)^\#$, voir Lemme
\ref{lemma-sheafification-presheaf-modules}. On voit aisément que
$(\colim_i \mathcal{F}_i)^\#$ est la colimite du diagramme dans
$\textit{Mod}(\mathcal{O})$ et que, d'après le Lemme
\ref{lemma-limits-colimits-abelian-sheaves}, l'oubli de sa structure de
$\mathcal{O}$-module donne la colimite dans $\textit{Ab}(\mathcal{C})$.
\end{proof}

\noindent
L'existence des limites et des colimites permet d'exprimer par celles-ci les
propriétés d'exactitude des foncteurs définis sur la catégorie des
$\mathcal{O}$-modules, comme dans Catégories, Section
\ref{categories-section-exact-functor}. Voir Homologie, Lemme
\ref{homology-lemma-exact-functor}, pour une description de ces propriétés en
termes de suites exactes courtes.

\begin{lemma}
\label{lemma-exactness-pushforward-pullback}
Soit $f : (\Sh(\mathcal{C}), \mathcal{O}_\mathcal{C})
\to (\Sh(\mathcal{D}), \mathcal{O}_\mathcal{D})$
un morphisme de topos annelés.
\begin{enumerate}
\item Le foncteur $f_*$ est exact à gauche. En fait, il commute à toutes les
limites.
\item Le foncteur $f^*$ est exact à droite. En fait, il commute à toutes les
colimites.
\end{enumerate}
\end{lemma}

\begin{proof}
Cela est vrai parce que $(f^*, f_*)$ est un couple de foncteurs adjoints, voir
Lemme \ref{lemma-adjoint-pullback-pushforward-modules}. Voir Catégories,
Section \ref{categories-section-adjoint}.
\end{proof}

\begin{lemma}
\label{lemma-check-exactness-stalks}
Soit $\mathcal{C}$ un site. Si $\{p_i\}_{i \in I}$ est une famille conservative
de points, l'exactitude d'une suite de faisceaux abéliens peut être vérifiée sur
les germes aux points $p_i$, $i \in I$. Si $\mathcal{C}$ possède suffisamment
de points, l'exactitude d'une suite de faisceaux abéliens peut être vérifiée sur
les germes.
\end{lemma}

\begin{proof}
Cela résulte immédiatement de Sites, Lemme
\ref{sites-lemma-exactness-stalks}.
\end{proof}





\section{Exactitude de l'image directe}
\label{section-pushforward}

\noindent
Quelques lemmes techniques sur les propriétés d'exactitude de l'image directe.

\begin{lemma}
\label{lemma-reflect-surjections}
Soit $f : \Sh(\mathcal{C}) \to \Sh(\mathcal{D})$ un morphisme de topos. Les
assertions suivantes sont équivalentes :
\begin{enumerate}
\item $f^{-1}f_*\mathcal{F} \to \mathcal{F}$ est surjectif pour tout
$\mathcal{F}$ dans $\textit{Ab}(\mathcal{C})$, et
\item $f_* : \textit{Ab}(\mathcal{C}) \to \textit{Ab}(\mathcal{D})$
reflète les surjections.
\end{enumerate}
Dans ce cas, le foncteur
$f_* : \textit{Ab}(\mathcal{C}) \to \textit{Ab}(\mathcal{D})$
est fidèle.
\end{lemma}

\begin{proof}
Supposons (1). Soit $a : \mathcal{F} \to \mathcal{F}'$ une application de
faisceaux abéliens sur $\mathcal{C}$ telle que $f_*a$ soit surjective. Puisque
$f^{-1}$ est exact, il s'ensuit que
$f^{-1}f_*a : f^{-1}f_*\mathcal{F} \to f^{-1}f_*\mathcal{F}'$
est surjectif. Avec (1), cela implique que $a$ est surjectif. Ainsi (2) est
satisfaite.

\medskip\noindent
Supposons (2). Soit $\mathcal{F}$ un faisceau abélien sur $\mathcal{C}$. Nous
devons montrer que $f^{-1}f_*\mathcal{F} \to \mathcal{F}$ est surjectif.
D'après (2), il suffit de montrer que
$f_*f^{-1}f_*\mathcal{F} \to f_*\mathcal{F}$ est surjectif. Or cela est vrai,
car il existe une application canonique
$f_*\mathcal{F} \to f_*f^{-1}f_*\mathcal{F}$ qui en est un inverse d'un côté.

\medskip\noindent
Nous omettons la démonstration de la dernière assertion.
\end{proof}

\begin{lemma}
\label{lemma-exactness}
Soit $f : \Sh(\mathcal{C}) \to \Sh(\mathcal{D})$ un morphisme de topos.
Supposons qu'au moins l'une des propriétés suivantes soit satisfaite :
\begin{enumerate}
\item $f_*$ transforme les surjections de faisceaux d'ensembles en
surjections,
\item $f_*$ transforme les surjections de faisceaux abéliens en surjections,
\item $f_*$ commute aux coégalisateurs de faisceaux d'ensembles,
\item $f_*$ commute aux sommes amalgamées de faisceaux d'ensembles,
\end{enumerate}
Alors $f_* : \textit{Ab}(\mathcal{C}) \to \textit{Ab}(\mathcal{D})$ est
exact.
\end{lemma}

\begin{proof}
Puisque $f_* : \textit{Ab}(\mathcal{C}) \to \textit{Ab}(\mathcal{D})$ est un
adjoint à droite, nous savons déjà qu'il transforme une suite exacte courte
$0 \to \mathcal{F}_1 \to \mathcal{F}_2 \to \mathcal{F}_3 \to 0$
de faisceaux abéliens sur $\mathcal{C}$ en une suite exacte
$$
0 \to f_*\mathcal{F}_1 \to f_*\mathcal{F}_2 \to f_*\mathcal{F}_3
$$
voir Catégories, Sections \ref{categories-section-exact-functor} et
\ref{categories-section-adjoint}, ainsi qu'Homologie, Section
\ref{homology-section-functors}. Il suffit donc de démontrer que l'application
$f_*\mathcal{F}_2 \to f_*\mathcal{F}_3$ est surjective. Si (1) ou (2) est
satisfaite, cela découle des définitions. D'après Sites, Lemme
\ref{sites-lemma-exactness-properties}, chacune des propriétés (3) et (4)
implique formellement (1), ce qui conclut également dans ces cas.
\end{proof}

\begin{lemma}
\label{lemma-morphism-ringed-sites-almost-cocontinuous}
Soit $f : \mathcal{D} \to \mathcal{C}$ un morphisme de sites associé au
foncteur continu $u : \mathcal{C} \to \mathcal{D}$. Supposons que $u$ soit
presque cocontinu. Alors
\begin{enumerate}
\item $f_* : \textit{Ab}(\mathcal{D}) \to \textit{Ab}(\mathcal{C})$ est
exact.
\item si l'on se donne
$f^\sharp : f^{-1}\mathcal{O}_\mathcal{C} \to \mathcal{O}_\mathcal{D}$ de
sorte que $f$ devienne un morphisme de sites annelés, alors
$f_* : \textit{Mod}(\mathcal{O}_\mathcal{D}) \to
\textit{Mod}(\mathcal{O}_\mathcal{C})$ est exact.
\end{enumerate}
\end{lemma}

\begin{proof}
La partie (2) résulte de la partie (1) par le Lemme
\ref{lemma-limits-colimits}. La partie (1) résulte de Sites, Lemmes
\ref{sites-lemma-morphism-of-sites-almost-cocontinuous} et
\ref{sites-lemma-exactness-properties}.
\end{proof}





\section{Exactitude de l'adjoint à gauche}
\label{section-exactness-lower-shriek}

\noindent
Soit $u : \mathcal{C} \to \mathcal{D}$ un foncteur entre sites. Supposons que
\begin{enumerate}
\item[(a)] $u$ soit cocontinu, et
\item[(b)] $u$ soit continu.
\end{enumerate}
Soit $g : \Sh(\mathcal{C}) \to \Sh(\mathcal{D})$ le morphisme de topos associé
à $u$, voir Sites, Lemme \ref{sites-lemma-cocontinuous-morphism-topoi}.
Rappelons que $g^{-1} = u^p$, c'est-à-dire que $g^{-1}$ est donné par la
formule simple
$(g^{-1}\mathcal{G})(U) = \mathcal{G}(u(U))$, voir Sites, Lemme
\ref{sites-lemma-when-shriek}.
Nous voulons montrer que
$g^{-1} : \textit{Ab}(\mathcal{D}) \to \textit{Ab}(\mathcal{C})$
admet un adjoint à gauche $g_!$. D'après Sites, Lemme
\ref{sites-lemma-when-shriek}, le foncteur
$g^{Sh}_! = (u_p\ )^\#$ est un adjoint à gauche sur les faisceaux d'ensembles.
De plus, nous savons que $g^{Sh}_!\mathcal{F}$ est le faisceau associé au
préfaisceau
$$
V \longmapsto \colim_{V \to u(U)} \mathcal{F}(U)
$$
où la colimite est indexée par $(\mathcal{I}_V^u)^{opp}$ et prise dans la
catégorie des ensembles. La définition suivante est donc naturelle.

\begin{definition}
\label{definition-g-shriek}
Soit $u : \mathcal{C} \to \mathcal{D}$ satisfaisant (a) et (b) ci-dessus. Pour
$\mathcal{F} \in \textit{PAb}(\mathcal{C})$, nous définissons
{\it $g_{p!}\mathcal{F}$} comme le préfaisceau
$$
V \longmapsto \colim_{V \to u(U)} \mathcal{F}(U)
$$
où les colimites sur $(\mathcal{I}_V^u)^{opp}$ sont prises dans
$\textit{Ab}$. Pour $\mathcal{F} \in \textit{PAb}(\mathcal{C})$, nous posons
{\it $g_!\mathcal{F} = (g_{p!}\mathcal{F})^\#$}.
\end{definition}

\noindent
Nous sommes si explicites parce que les foncteurs $g^{Sh}_!$ et $g_!$ sont
différents. Lorsque nous les utilisons tous deux, il faut soigneusement les
distinguer.

\begin{lemma}
\label{lemma-g-shriek-adjoint}
Le foncteur $g_{p!}$ est adjoint à gauche du foncteur $u^p$. Le foncteur $g_!$
est adjoint à gauche du foncteur $g^{-1}$. Autrement dit, les formules
\begin{align*}
\Mor_{\textit{PAb}(\mathcal{C})}(\mathcal{F}, u^p\mathcal{G})
& =
\Mor_{\textit{PAb}(\mathcal{D})}(g_{p!}\mathcal{F}, \mathcal{G}), \\
\Mor_{\textit{Ab}(\mathcal{C})}(\mathcal{F}, g^{-1}\mathcal{G})
& =
\Mor_{\textit{Ab}(\mathcal{D})}(g_!\mathcal{F}, \mathcal{G})
\end{align*}
sont bifonctorielles en $\mathcal{F}$ et $\mathcal{G}$.
\end{lemma}

\begin{proof}
La seconde formule résulte formellement de la première, car si $\mathcal{F}$ et
$\mathcal{G}$ sont des faisceaux abéliens, alors
\begin{align*}
\Mor_{\textit{Ab}(\mathcal{C})}(\mathcal{F}, g^{-1}\mathcal{G})
& =
\Mor_{\textit{PAb}(\mathcal{D})}(g_{p!}\mathcal{F}, \mathcal{G}) \\
& =
\Mor_{\textit{Ab}(\mathcal{D})}(g_!\mathcal{F}, \mathcal{G})
\end{align*}
par la propriété universelle de la faisceautisation.

\medskip\noindent
Pour démontrer la première formule, soient $\mathcal{F}$ et $\mathcal{G}$ des
préfaisceaux abéliens. Nous construirons des applications du groupe de gauche
vers celui de droite et omettrons de vérifier qu'elles sont inverses l'une de
l'autre.

\medskip\noindent
Il existe une application canonique de préfaisceaux abéliens
$\mathcal{F} \to u^pg_{p!}\mathcal{F}$ qui, sur les sections au-dessus de $U$,
est l'application naturelle
$\mathcal{F}(U) \to \colim_{u(U) \to u(U')} \mathcal{F}(U')$, voir Sites,
Lemme \ref{sites-lemma-recover}.
Étant donnée une application $\alpha : g_{p!}\mathcal{F} \to \mathcal{G}$,
nous obtenons $u^p\alpha : u^pg_{p!}\mathcal{F} \to u^p\mathcal{G}$, que nous
pouvons précomposer par l'application
$\mathcal{F} \to u^pg_{p!}\mathcal{F}$.

\medskip\noindent
Il existe une application canonique de préfaisceaux abéliens
$g_{p!}u^p\mathcal{G} \to \mathcal{G}$ qui, sur les sections au-dessus de $V$,
est l'application naturelle
$\colim_{V \to u(U)} \mathcal{G}(u(U)) \to \mathcal{G}(V)$.
Elle envoie une section $s \in \mathcal{G}(u(U))$ dans le facteur correspondant
à $t : V \to u(U)$ sur $t^*s \in \mathcal{G}(V)$. Ainsi, étant donnée une
application $\beta : \mathcal{F} \to u^p\mathcal{G}$, nous obtenons une
application $g_{p!}\beta : g_{p!}\mathcal{F} \to g_{p!}u^p\mathcal{G}$, que
nous pouvons postcomposer par l'application
$g_{p!}u^p\mathcal{G} \to \mathcal{G}$ ci-dessus.
\end{proof}

\begin{lemma}
\label{lemma-exactness-lower-shriek}
Soient $\mathcal{C}$ et $\mathcal{D}$ des sites. Soit
$u : \mathcal{C} \to \mathcal{D}$ un foncteur. Supposons que
\begin{enumerate}
\item[(a)] $u$ soit cocontinu,
\item[(b)] $u$ soit continu, et
\item[(c)] les produits fibrés et les égalisateurs existent dans
$\mathcal{C}$ et $u$ commute à ces constructions.
\end{enumerate}
Dans ce cas, le foncteur
$g_! : \textit{Ab}(\mathcal{C}) \to \textit{Ab}(\mathcal{D})$
est exact.
\end{lemma}

\begin{proof}
Comparer avec Sites, Lemme \ref{sites-lemma-preserve-equalizers}. Supposons
(a), (b) et (c). Nous savons déjà que $g_!$ est exact à droite puisqu'il est un
adjoint à gauche, voir Catégories, Lemme
\ref{categories-lemma-exact-adjoint}, et Homologie, Section
\ref{homology-section-functors}. Nous avons $g_! = (g_{p!}\ )^\#$. Il faut
montrer que $g_!$ transforme les applications injectives de faisceaux abéliens
en applications injectives de préfaisceaux abéliens. Rappelons que la
faisceautisation des préfaisceaux abéliens est exacte, voir Lemme
\ref{lemma-limits-colimits-abelian-sheaves}. Il suffit donc de montrer que
$g_{p!}$ transforme les applications injectives de préfaisceaux abéliens en
applications injectives de préfaisceaux abéliens. Pour cela, il suffit que les
colimites sur les catégories
$(\mathcal{I}_V^u)^{opp}$ de Sites, Section
\ref{sites-section-functoriality-PSh}
transforment les applications injectives entre diagrammes en injections. Cela
résulte de Sites, Lemme \ref{sites-lemma-almost-directed}, et d'Algèbre, Lemme
\ref{algebra-lemma-almost-directed-colimit-exact}.
\end{proof}

\begin{lemma}
\label{lemma-back-and-forth}
Soient $\mathcal{C}$ et $\mathcal{D}$ des sites. Soit
$u : \mathcal{C} \to \mathcal{D}$ un foncteur. Supposons que
\begin{enumerate}
\item[(a)] $u$ soit cocontinu,
\item[(b)] $u$ soit continu, et
\item[(c)] $u$ soit pleinement fidèle.
\end{enumerate}
Pour $g_!, g^{-1}, g_*$ comme ci-dessus, les applications canoniques
$\mathcal{F} \to g^{-1}g_!\mathcal{F}$ and
$g^{-1}g_*\mathcal{F} \to \mathcal{F}$
sont des isomorphismes pour tout faisceau abélien $\mathcal{F}$ sur
$\mathcal{C}$.
\end{lemma}

\begin{proof}
L'application $g^{-1}g_*\mathcal{F} \to \mathcal{F}$ est un isomorphisme
d'après Sites, Lemme \ref{sites-lemma-back-and-forth}, et puisque les images
inverse et directe des faisceaux abéliens coïncident avec celles de leurs
faisceaux d'ensembles sous-jacents.

\medskip\noindent
Choisissons $U \in \Ob(\mathcal{C})$. Nous allons montrer que
$g^{-1}g_!\mathcal{F}(U) = \mathcal{F}(U)$. Remarquons d'abord que
$g^{-1}g_!\mathcal{F}(U) = g_!\mathcal{F}(u(U))$. Il suffit donc de montrer que
$g_!\mathcal{F}(u(U)) = \mathcal{F}(U)$. Nous savons que $g_!\mathcal{F}$ est
le faisceau abélien associé au préfaisceau $g_{p!}\mathcal{F}$ défini par la
règle
$$
V \longmapsto \colim_{V \to u(U')} \mathcal{F}(U')
$$
où la colimite est prise dans $\textit{Ab}$. Si $V = u(U)$, alors, puisque $u$
est pleinement fidèle, cette colimite est indexée par les $U \to U'$. Nous en
concluons que $g_{p!}\mathcal{F}(u(U)) = \mathcal{F}(U)$. Comme $u$ est
cocontinu et continu, tout recouvrement de $u(U)$ dans $\mathcal{D}$ peut être
raffiné par un recouvrement (!) $\{u(U_i) \to u(U)\}$ de $\mathcal{D}$, où
$\{U_i \to U\}$ est un recouvrement dans $\mathcal{C}$. Il s'ensuit également
que $(g_{p!}\mathcal{F})^+(u(U)) = \mathcal{F}(U)$, car, dans la colimite qui
définit la valeur de $(g_{p!}\mathcal{F})^+$ en $u(U)$, nous pouvons nous
restreindre au système cofinal des recouvrements
$\{u(U_i) \to u(U)\}$ ci-dessus. Ainsi
$(g_{p!}\mathcal{F})^+(u(U)) = \mathcal{F}(U)$ pour tout objet $U$ de
$\mathcal{C}$. En répétant encore une fois cet argument, nous obtenons
$(g_{p!}\mathcal{F})^\#(u(U)) = \mathcal{F}(U)$ pour tout objet $U$ de
$\mathcal{C}$, d'où l'égalité cherchée $g^{-1}g_!\mathcal{F} = \mathcal{F}$.
\end{proof}

\begin{remark}
\label{remark-no-extension}
En général, le foncteur $g_!$ ne s'étend pas aux catégories de modules lorsque
$g$ est (une partie d')un morphisme de topos annelés. En effet, pour tout
homomorphisme d'anneaux $A \to B$, le foncteur
$M \mapsto B \otimes_A M$ admet un adjoint à droite (la restriction des
scalaires), mais n'admet pas en général d'adjoint à gauche, car l'existence de
celui-ci impliquerait que $A \to B$ est plat. Nous verrons à la Section
\ref{section-localize} qu'il est possible de définir $j_!$ sur les faisceaux de
modules dans le cas d'une localisation de sites. Nous en discuterons avec une
plus grande généralité à la Section \ref{section-lower-shriek-modules}.
\end{remark}

\begin{lemma}
\label{lemma-have-left-adjoint}
Soient $\mathcal{C}$ et $\mathcal{D}$ des sites. Soit
$g : \Sh(\mathcal{C}) \to \Sh(\mathcal{D})$ le morphisme de topos associé à un
foncteur continu et cocontinu $u : \mathcal{C} \to \mathcal{D}$.
\begin{enumerate}
\item Si $u$ admet un adjoint à gauche $w$, alors $g_!$ et $g_!^{\Sh}$
coïncident sur les faisceaux d'ensembles sous-jacents et $g_!$ est exact.
\item Si, de plus, $w$ est cocontinu, alors $g_! = h^{-1}$ et
$g^{-1} = h_*$, où $h : \Sh(\mathcal{D}) \to \Sh(\mathcal{C})$ est le
morphisme de topos associé à $w$.
\end{enumerate}
\end{lemma}

\begin{proof}
Ce lemme est l'analogue de Sites, Lemme
\ref{sites-lemma-have-left-adjoint}. D'après Sites, Lemme
\ref{sites-lemma-adjoint-functors}, les catégories $\mathcal{I}_V^u$ possèdent
un objet initial. L'ensemble sous-jacent à la colimite d'un système de groupes
abéliens sur $(\mathcal{I}_V^u)^{opp}$ est donc la colimite des ensembles
sous-jacents. Il s'ensuit que $g_!^{\Sh}$ et $g_!$ coïncident, compte tenu de
notre construction de $g_!$ dans la Définition \ref{definition-g-shriek}.
L'exactitude et (2) résultent immédiatement des assertions correspondantes de
Sites, Lemme \ref{sites-lemma-have-left-adjoint}.
\end{proof}





\section{Types globaux de modules}
\label{section-global}

\begin{definition}
\label{definition-global}
Soit $(\Sh(\mathcal{C}), \mathcal{O})$ un topos annelé. Soit $\mathcal{F}$ un
faisceau de $\mathcal{O}$-modules.
\begin{enumerate}
\item Nous disons que $\mathcal{F}$ est un {\it $\mathcal{O}$-module libre}
si $\mathcal{F}$ est isomorphe, comme $\mathcal{O}$-module, à un faisceau de la
forme
$\bigoplus_{i \in I} \mathcal{O}$.
\item Nous disons que $\mathcal{F}$ est {\it libre de rang fini} si
$\mathcal{F}$ est isomorphe, comme $\mathcal{O}$-module, à un faisceau de la forme
$\bigoplus_{i \in I} \mathcal{O}$ avec $I$ fini.
\item Nous disons que $\mathcal{F}$ est {\it engendré par ses sections
globales} s'il existe une surjection
$$
\bigoplus\nolimits_{i \in I} \mathcal{O} \longrightarrow \mathcal{F}
$$
depuis un $\mathcal{O}$-module libre vers $\mathcal{F}$.
\item Pour $r \geq 0$, nous disons que $\mathcal{F}$ est {\it engendré par
$r$ sections globales} s'il existe une surjection
$\mathcal{O}^{\oplus r} \to \mathcal{F}$.
\item Nous disons que $\mathcal{F}$ est {\it engendré par un nombre fini de
sections globales} s'il est engendré par $r$ sections globales pour un certain
$r \geq 0$.
\item Nous disons que $\mathcal{F}$ admet une {\it présentation globale} s'il
existe une suite exacte
$$
\bigoplus\nolimits_{j \in J} \mathcal{O} \longrightarrow
\bigoplus\nolimits_{i \in I} \mathcal{O} \longrightarrow
\mathcal{F} \longrightarrow 0
$$
de $\mathcal{O}$-modules.
\item Nous disons que $\mathcal{F}$ admet une {\it présentation globale finie}
s'il existe une suite exacte
$$
\bigoplus\nolimits_{j \in J} \mathcal{O} \longrightarrow
\bigoplus\nolimits_{i \in I} \mathcal{O} \longrightarrow
\mathcal{F} \longrightarrow 0
$$
de $\mathcal{O}$-modules avec $I$ et $J$ finis.
\end{enumerate}
\end{definition}

\noindent
Remarquons que, pour tout ensemble $I$, la somme directe
$\bigoplus_{i \in I} \mathcal{O}$ existe (Lemme
\ref{lemma-limits-colimits}) et est le faisceau associé au préfaisceau
$U \mapsto \bigoplus_{i \in I} \mathcal{O}(U)$. Ce module est appelé le
{\it $\mathcal{O}$-module libre sur l'ensemble $I$}.

\begin{lemma}
\label{lemma-global-pullback}
Soit
$(f, f^\sharp) :
(\Sh(\mathcal{C}), \mathcal{O}_\mathcal{C})
\to
(\Sh(\mathcal{D}), \mathcal{O}_\mathcal{D})$
un morphisme de topos annelés. Soit $\mathcal{F}$ un
$\mathcal{O}_\mathcal{D}$-module.
\begin{enumerate}
\item Si $\mathcal{F}$ est libre, alors $f^*\mathcal{F}$ est libre.
\item Si $\mathcal{F}$ est libre de rang fini, alors $f^*\mathcal{F}$ est
libre de rang fini.
\item Si $\mathcal{F}$ est engendré par ses sections globales, alors
$f^*\mathcal{F}$ l'est aussi.
\item Si $r \geq 0$ et si $\mathcal{F}$ est engendré par $r$ sections
globales, alors $f^*\mathcal{F}$ est engendré par $r$ sections globales.
\item Si $\mathcal{F}$ est engendré par un nombre fini de sections globales,
alors $f^*\mathcal{F}$ l'est aussi.
\item Si $\mathcal{F}$ admet une présentation globale, alors
$f^*\mathcal{F}$ en admet une.
\item Si $\mathcal{F}$ admet une présentation globale finie, alors
$f^*\mathcal{F}$ en admet une.
\end{enumerate}
\end{lemma}

\begin{proof}
Cela est vrai parce que $f^*$ commute aux colimites quelconques (Lemme
\ref{lemma-exactness-pushforward-pullback}) et que
$f^*\mathcal{O}_\mathcal{D} = \mathcal{O}_\mathcal{C}$.
\end{proof}






\section{Propriétés intrinsèques des modules}
\label{section-intrinsic}

\noindent
Soit $\mathcal{P}$ une propriété des faisceaux de modules sur les topos
annelés. Nous disons que $\mathcal{P}$ est une {\it propriété intrinsèque} si
$\mathcal{P}(\mathcal{F}) \Leftrightarrow \mathcal{P}(f^*\mathcal{F})$
chaque fois que $(f, f^\sharp) :
(\Sh(\mathcal{C}'), \mathcal{O}')
\to
(\Sh(\mathcal{C}), \mathcal{O})$
est une équivalence de topos annelés. Par exemple, la propriété d'être libre est
intrinsèque. En effet, le $\mathcal{O}$-module libre sur l'ensemble $I$ est
caractérisé par la propriété
$$
\Mor_{\textit{Mod}(\mathcal{O})}(
\bigoplus\nolimits_{i \in I} \mathcal{O},
\mathcal{F})
=
\prod\nolimits_{i \in I} \Mor_{\Sh(\mathcal{C})}(\{*\},
\mathcal{F})
$$
pour $\mathcal{F}$ variable dans $\textit{Mod}(\mathcal{O})$. On peut aussi
utiliser le Lemme \ref{lemma-global-pullback} pour voir que la liberté est
intrinsèque. En fait, chacune des propriétés définies dans la Définition
\ref{definition-global} est intrinsèque pour la même raison. Comment
définirons-nous d'autres propriétés intrinsèques des $\mathcal{O}$-modules ?

\medskip\noindent
La conséquence du Lemme
\ref{lemma-morphism-ringed-topoi-comes-from-morphism-ringed-sites} est la
suivante. Supposons que l'on veuille définir une propriété intrinsèque
$\mathcal{P}$ d'un $\mathcal{O}$-module sur un topos. On peut procéder comme
suit :
\begin{enumerate}
\item Pour tout site $\mathcal{C}$, tout faisceau d'anneaux $\mathcal{O}$ sur
$\mathcal{C}$ et tout $\mathcal{O}$-module $\mathcal{F}$, définir la propriété
correspondante $\mathcal{P}(\mathcal{C}, \mathcal{O}, \mathcal{F})$.
\item Pour toute paire de sites $\mathcal{C}$, $\mathcal{C}'$, tout foncteur
cocontinu spécial $u : \mathcal{C} \to \mathcal{C}'$, tout faisceau d'anneaux
$\mathcal{O}$ sur $\mathcal{C}$ et tout $\mathcal{O}$-module $\mathcal{F}$,
montrer que
$$
\mathcal{P}(\mathcal{C}, \mathcal{O}, \mathcal{F})
\Leftrightarrow
\mathcal{P}(\mathcal{C}', g_*\mathcal{O}, g_*\mathcal{F})
$$
où $g : \Sh(\mathcal{C}) \to \Sh(\mathcal{C}')$ est l'équivalence de topos
associée à $u$.
\end{enumerate}
Dans ce cas, étant donnés un topos annelé
$(\Sh(\mathcal{C}), \mathcal{O})$ et un faisceau de $\mathcal{O}$-modules
$\mathcal{F}$, nous disons simplement que $\mathcal{F}$ possède la
propriété $\mathcal{P}$ si
$\mathcal{P}(\mathcal{C}, \mathcal{O}, \mathcal{F})$ est vraie. Le Lemme
\ref{lemma-morphism-ringed-topoi-comes-from-morphism-ringed-sites}, combiné à
(2), garantit que cette définition est indépendante des choix.

\medskip\noindent
De plus, le même Lemme
\ref{lemma-morphism-ringed-topoi-comes-from-morphism-ringed-sites} garantit que
si, en outre,
\begin{enumerate}
\item[(3)] pour tout morphisme de sites annelés
$(f, f^\sharp) :
(\mathcal{C}, \mathcal{O}_\mathcal{C})
\to
(\mathcal{D}, \mathcal{O}_\mathcal{D})$
tel que $f$ soit donné par un foncteur
$u : \mathcal{D} \to \mathcal{C}$ satisfaisant les hypothèses de Sites,
Proposition \ref{sites-proposition-get-morphism}, et pour tout
$\mathcal{O}_\mathcal{D}$-module $\mathcal{G}$, on ait
$$
\mathcal{P}(\mathcal{D}, \mathcal{O}_\mathcal{D}, \mathcal{G})
\Rightarrow
\mathcal{P}(\mathcal{C}, \mathcal{O}_\mathcal{C}, f^*\mathcal{G})
$$
\end{enumerate}
alors $\mathcal{P}$ est préservée par image inverse des modules suivant les
morphismes arbitraires de topos annelés.

\medskip\noindent
Dans les sections suivantes, nous utiliserons cette méthode pour montrer que les
propriétés suivantes sont intrinsèques : être localement libre, être localement
engendré par des sections, être localement engendré par $r$ sections, être de
type fini, être de présentation finie, être quasi-cohérent, et être cohérent.

\medskip\noindent
Une méthode peut-être plus satisfaisante consisterait à donner une définition
intrinsèque de ces notions, plutôt qu'à suivre la procédure laborieuse esquissée
ici. Cependant, dans de nombreuses situations géométriques où nous voulons
appliquer ces définitions, un site annelé précis et un faisceau de modules précis
nous sont donnés ; il est alors utile de disposer d'une définition déjà adaptée
à ce langage.
\section{Localisation de sites annelés}
\label{section-localize}

\noindent
Soit $(\mathcal{C}, \mathcal{O})$ un site annelé.
Soit $U \in \Ob(\mathcal{C})$.
Nous exposons dans ce cadre les analogues des résultats de
Sites, Section \ref{sites-section-localize}.

\medskip\noindent
Notons
$\mathcal{O}_U = j_U^{-1}\mathcal{O}$ la restriction de $\mathcal{O}$
au site $\mathcal{C}/U$. Elle est décrite par la règle simple
$\mathcal{O}_U(V/U) = \mathcal{O}(V)$. Avec cette notation,
le morphisme de localisation $j_U$ devient un morphisme de topos annelés
$$
(j_U, j_U^\sharp) :
(\Sh(\mathcal{C}/U), \mathcal{O}_U)
\longrightarrow
(\Sh(\mathcal{C}), \mathcal{O})
$$
à savoir que l'on prend pour
$j_U^\sharp : j_U^{-1}\mathcal{O} \to \mathcal{O}_U$
l'application identique.
On obtient en outre les descriptions suivantes de l'image directe
et de l'image inverse des modules.

\begin{definition}
\label{definition-localize-ringed-site}
Soit $(\mathcal{C}, \mathcal{O})$ un site annelé.
Soit $U \in \Ob(\mathcal{C})$.
\begin{enumerate}
\item Le site annelé $(\mathcal{C}/U, \mathcal{O}_U)$ est appelé la
{\it localisation du site annelé $(\mathcal{C}, \mathcal{O})$
en l'objet $U$}.
\item Le morphisme de topos annelés
$(j_U, j_U^\sharp) :
(\Sh(\mathcal{C}/U), \mathcal{O}_U)
\to
(\Sh(\mathcal{C}), \mathcal{O})$
est appelé le {\it morphisme de localisation}.
\item Le foncteur
$j_{U*} : \textit{Mod}(\mathcal{O}_U) \to \textit{Mod}(\mathcal{O})$
est appelé le {\it foncteur image directe}.
\item Pour un faisceau de $\mathcal{O}$-modules $\mathcal{F}$ sur $\mathcal{C}$,
le faisceau $j_U^*\mathcal{F}$ est appelé la
{\it restriction de $\mathcal{F}$ à $\mathcal{C}/U$}.
Nous le noterons parfois
$\mathcal{F}|_{\mathcal{C}/U}$, voire $\mathcal{F}|_U$.
Il est décrit par la règle simple
$j_U^*(\mathcal{F})(X/U) = \mathcal{F}(X)$.
\item L'adjoint à gauche
$j_{U!} : \textit{Mod}(\mathcal{O}_U) \to \textit{Mod}(\mathcal{O})$
de la restriction est appelé {\it prolongement par zéro}. Il existe et est
exact d'après les
Lemmes \ref{lemma-extension-by-zero} et
\ref{lemma-extension-by-zero-exact}.
\end{enumerate}
\end{definition}

\noindent
Comme dans le cas topologique, voir
Sheaves, Section \ref{sheaves-section-open-immersions},
le foncteur de prolongement par zéro $j_{U!}$ diffère du
prolongement par l'ensemble vide $j_{U!}$ défini sur les faisceaux
d'ensembles. Voici le lemme qui définit le prolongement par zéro.

\begin{lemma}
\label{lemma-extension-by-zero}
Soit $(\mathcal{C}, \mathcal{O})$ un site annelé.
Soit $U \in \Ob(\mathcal{C})$.
Le foncteur de restriction
$j_U^* : \textit{Mod}(\mathcal{O}) \to \textit{Mod}(\mathcal{O}_U)$
a un adjoint à gauche
$j_{U!} : \textit{Mod}(\mathcal{O}_U) \to \textit{Mod}(\mathcal{O})$.
Ainsi,
$$
\Mor_{\textit{Mod}(\mathcal{O}_U)}(\mathcal{G}, j_U^*\mathcal{F})
=
\Mor_{\textit{Mod}(\mathcal{O})}(j_{U!}\mathcal{G}, \mathcal{F})
$$
pour $\mathcal{F} \in \Ob(\textit{Mod}(\mathcal{O}))$
et $\mathcal{G} \in \Ob(\textit{Mod}(\mathcal{O}_U))$.
De plus, le prolongement par zéro $j_{U!}\mathcal{G}$ de $\mathcal{G}$
est le faisceau associé au préfaisceau
$$
V
\longmapsto
\bigoplus\nolimits_{\varphi \in \Mor_\mathcal{C}(V, U)}
\mathcal{G}(V \xrightarrow{\varphi} U)
$$
muni des applications de restriction évidentes et de la structure
évidente de $\mathcal{O}$-module.
\end{lemma}

\begin{proof}
La structure de $\mathcal{O}$-module sur le préfaisceau est définie
comme suit. Si $f \in \mathcal{O}(V)$ et
$s \in \mathcal{G}(V \xrightarrow{\varphi} U)$, on pose
$f \cdot s = fs$, où
$f \in \mathcal{O}_U(\varphi : V \to U) = \mathcal{O}(V)$
(car $\mathcal{O}_U$ est la restriction de $\mathcal{O}$ à
$\mathcal{C}/U$).

\medskip\noindent
De même, soit $\alpha : \mathcal{G} \to \mathcal{F}|_U$ un
morphisme de $\mathcal{O}_U$-modules. On peut alors définir une
application du préfaisceau du lemme dans $\mathcal{F}$ en envoyant
$$
\bigoplus\nolimits_{\varphi \in \Mor_\mathcal{C}(V, U)}
\mathcal{G}(V \xrightarrow{\varphi} U)
\longrightarrow
\mathcal{F}(V)
$$
selon la règle qui à $s \in \mathcal{G}(V \xrightarrow{\varphi} U)$
associe $\alpha(s) \in \mathcal{F}(V)$. Cette application est clairement
$\mathcal{O}$-linéaire et induit donc un morphisme de
$\mathcal{O}$-modules $\alpha' : j_{U!}\mathcal{G} \to \mathcal{F}$
par les propriétés de la faisceautisation des modules
(Lemme \ref{lemma-sheafification-presheaf-modules}).

\medskip\noindent
Réciproquement, soit $\beta : j_{U!}\mathcal{G} \to \mathcal{F}$
un morphisme de $\mathcal{O}$-modules.
Rappelons d'après Sites, Section \ref{sites-section-localize},
qu'il existe un prolongement par l'ensemble vide
$j^{Sh}_{U!} : \Sh(\mathcal{C}/U) \to \Sh(\mathcal{C})$
sur les faisceaux d'ensembles, adjoint à gauche de $j_U^{-1}$.
De plus, $j^{Sh}_{U!}\mathcal{G}$ est le faisceau associé au préfaisceau
$$
V
\longmapsto
\coprod\nolimits_{\varphi \in \Mor_\mathcal{C}(V, U)}
\mathcal{G}(V \xrightarrow{\varphi} U)
$$
Il existe donc une application naturelle
$j^{Sh}_{U!}\mathcal{G} \to j_{U!}\mathcal{G}$ de faisceaux d'ensembles.
En précomposant $\beta$ avec cette application, on obtient une
application de faisceaux d'ensembles
$j^{Sh}_{U!}\mathcal{G} \to \mathcal{F}$ qui correspond par adjonction
à une application de faisceaux d'ensembles
$\beta' : \mathcal{G} \to \mathcal{F}|_U$.
Nous affirmons que $\beta'$ est $\mathcal{O}_U$-linéaire. Supposons en effet
que $\varphi : V \to U$ soit un objet de $\mathcal{C}/U$ et que
$s, s' \in \mathcal{G}(\varphi : V \to U)$, and
$f \in \mathcal{O}(V) = \mathcal{O}_U(\varphi : V \to U)$.
La discussion précédente montre alors que
$\beta'(s + s')$, resp.\ $\beta'(fs)$ dans
$\mathcal{F}|_U(\varphi : V \to U)$, correspond à
$\beta(s + s')$, resp.\ $\beta(fs)$ dans
$\mathcal{F}(V)$. On conclut puisque $\beta$ est un homomorphisme.

\medskip\noindent
Pour achever la preuve du lemme, il reste à montrer que les constructions
$\alpha \mapsto \alpha'$ et $\beta \mapsto \beta'$ sont inverses l'une de
l'autre. Nous omettons les vérifications.
\end{proof}

\noindent
Notons que, dans la situation de la
Définition \ref{definition-localize-ringed-site}, on a
\begin{equation}
\label{equation-map-lower-shriek-OU-into-module}
\Hom_\mathcal{O}(j_{U!}\mathcal{O}_U, \mathcal{F}) =
\Hom_{\mathcal{O}_U}(\mathcal{O}_U, j_U^*\mathcal{F}) =
\mathcal{F}(U)
\end{equation}
pour tout $\mathcal{O}$-module $\mathcal{F}$. En effet, la première égalité
résulte de l'adjonction entre $j_{U!}$ et $j_U^*$, et la seconde du fait que
$\Hom_{\mathcal{O}_U}(\mathcal{O}_U, j_U^*\mathcal{F}) =
j_U^*\mathcal{F}(U/U) = \mathcal{F}|_U(U/U) = \mathcal{F}(U)$.

\begin{lemma}
\label{lemma-extension-by-zero-exact}
Soit $(\mathcal{C}, \mathcal{O})$ un site annelé.
Soit $U \in \Ob(\mathcal{C})$.
Le foncteur
$j_{U!} : \textit{Mod}(\mathcal{O}_U) \to \textit{Mod}(\mathcal{O})$
est exact.
\end{lemma}

\begin{proof}
Puisque $j_{U!}$ est adjoint à gauche de $j_U^*$, il est exact à droite
(voir
Categories, Lemme \ref{categories-lemma-exact-adjoint}
et
Homology, Section \ref{homology-section-functors}).
Il suffit donc de montrer que, si $\mathcal{G}_1 \to \mathcal{G}_2$
est un morphisme injectif de $\mathcal{O}_U$-modules, alors
$j_{U!}\mathcal{G}_1 \to j_{U!}\mathcal{G}_2$ est injectif.
Le morphisme sur les sections des préfaisceaux au-dessus d'un objet $V$
(comme dans le Lemme \ref{lemma-extension-by-zero}) est
$$
\bigoplus\nolimits_{\varphi \in \Mor_\mathcal{C}(V, U)}
\mathcal{G}_1(V \xrightarrow{\varphi} U)
\longrightarrow
\bigoplus\nolimits_{\varphi \in \Mor_\mathcal{C}(V, U)}
\mathcal{G}_2(V \xrightarrow{\varphi} U)
$$
qui est injectif par hypothèse. La faisceautisation étant exacte d'après le
Lemme \ref{lemma-sheafification-exact},
on en conclut que $j_{U!}\mathcal{G}_1 \to j_{U!}\mathcal{G}_2$ est injectif.
\end{proof}

\begin{lemma}
\label{lemma-j-shriek-reflects-exactness}
Soit $(\mathcal{C}, \mathcal{O})$ un site annelé.
Soit $U \in \Ob(\mathcal{C})$. Un complexe de $\mathcal{O}_U$-modules
$\mathcal{G}_1 \to \mathcal{G}_2 \to \mathcal{G}_3$ est exact
si et seulement si
$j_{U!}\mathcal{G}_1 \to j_{U!}\mathcal{G}_2 \to j_{U!}\mathcal{G}_3$
est exact comme suite de $\mathcal{O}$-modules.
\end{lemma}

\begin{proof}
Nous savons déjà que $j_{U!}$ est exact, voir le
Lemme \ref{lemma-extension-by-zero-exact}.
Il suffit donc de montrer que
$j_{U!} :  \textit{Mod}(\mathcal{O}_U) \to \textit{Mod}(\mathcal{O})$
reflète les injections et les surjections.

\medskip\noindent
Pour tout $\mathcal{G}$ dans $\textit{Mod}(\mathcal{O}_U)$,
on dispose de l'unité $\mathcal{G} \to j_U^*j_{U!}\mathcal{G}$
de l'adjonction. Nous affirmons que cette application est une injection
de faisceaux. En examinant la construction du
Lemme \ref{lemma-extension-by-zero}, on voit en effet qu'elle est la
faisceautisation de la règle qui associe à l'objet $V/U$ de
$\mathcal{C}/U$ l'application injective
$$
\mathcal{G}(V/U) \longrightarrow
\bigoplus\nolimits_{\varphi \in \Mor_\mathcal{C}(V, U)}
\mathcal{G}(V \xrightarrow{\varphi} U)
$$
donnée par l'inclusion du facteur correspondant au morphisme structural
$V \to U$. L'affirmation résulte alors de l'exactitude de la faisceautisation.
Nous omettons quelques détails.

\medskip\noindent
Si $\mathcal{G} \to \mathcal{G}'$ est un morphisme de
$\mathcal{O}_U$-modules tel que
$j_{U!}\mathcal{G} \to j_{U!}\mathcal{G}'$ soit injectif, alors
$j_U^*j_{U!}\mathcal{G} \to j_U^*j_{U!}\mathcal{G}'$ est injectif
(la restriction est exacte), donc
$\mathcal{G} \to j_U^*j_{U!}\mathcal{G}'$ est injectif, et par suite
$\mathcal{G} \to \mathcal{G}'$ est injectif.
Ainsi $j_{U!}$ reflète les injections.

\medskip\noindent
Soit $a : \mathcal{G} \to \mathcal{G}'$ un morphisme de
$\mathcal{O}_U$-modules tel que
$j_{U!}\mathcal{G} \to j_{U!}\mathcal{G}'$ soit surjectif.
Soit $\mathcal{H}$ le conoyau de $a$.
Alors $j_{U!}\mathcal{H} = 0$, puisque $j_{U!}$ est exact.
D'après ce qui précède, l'application
$\mathcal{H} \to j^*_U j_{U!}\mathcal{H}$ est injective.
Ainsi $\mathcal{H} = 0$, comme voulu.
\end{proof}

\begin{lemma}
\label{lemma-relocalize}
Soit $(\mathcal{C}, \mathcal{O})$ un site annelé.
Soit $f : V \to U$ un morphisme de $\mathcal{C}$.
Il existe alors un diagramme commutatif
$$
\xymatrix{
(\Sh(\mathcal{C}/V), \mathcal{O}_V)
\ar[rd]_{(j_V, j_V^\sharp)} \ar[rr]_{(j, j^\sharp)} & &
(\Sh(\mathcal{C}/U), \mathcal{O}_U)
\ar[ld]^{(j_U, j_U^\sharp)} \\
& (\Sh(\mathcal{C}), \mathcal{O}) &
}
$$
dans la catégorie des topos annelés. Ici $(j, j^\sharp)$ est le morphisme
de localisation associé à l'objet $V/U$ du site annelé
$(\mathcal{C}/V, \mathcal{O}_V)$.
\end{lemma}

\begin{proof}
Le seul point à vérifier est que
$j_V^\sharp = j^\sharp \circ j^{-1}(j_U^\sharp)$,
car tout le reste résulte directement de
Sites, Lemme \ref{sites-lemma-relocalize}, et
Sites, Équation (\ref{sites-equation-relocalize}).
Nous omettons la vérification de cette égalité.
\end{proof}

\begin{remark}
\label{remark-localize-shriek-equal}
Dans la situation du Lemme \ref{lemma-extension-by-zero},
le diagramme
$$
\xymatrix{
\textit{Mod}(\mathcal{O}_U) \ar[r]_{j_{U!}} \ar[d]_{forget} &
\textit{Mod}(\mathcal{O}_\mathcal{C}) \ar[d]^{forget} \\
\textit{Ab}(\mathcal{C}/U) \ar[r]^{j^{Ab}_{U!}} &
\textit{Ab}(\mathcal{C})
}
$$
est commutatif. Cela résulte clairement de la description explicite du
foncteur $j_{U!}$ donnée dans le lemme.
\end{remark}

\begin{remark}
\label{remark-localize-presheaves}
Localisation et préfaisceaux de modules ; voir
Sites, Remarque \ref{sites-remark-localize-presheaves}.
Soit $\mathcal{C}$ une catégorie.
Soit $\mathcal{O}$ un préfaisceau d'anneaux.
Soit $U$ un objet de $\mathcal{C}$.
À proprement parler, les foncteurs $j_U^*$, $j_{U*}$ et $j_{U!}$
n'ont pas été définis pour les préfaisceaux de $\mathcal{O}$-modules.
On peut cependant considérer un préfaisceau comme un faisceau pour la
topologie chaotique sur $\mathcal{C}$ (voir
Sites, Examples \ref{sites-example-indiscrete}).
On obtient donc aussi un foncteur
$$
j_U^* :
\textit{PMod}(\mathcal{O})
\longrightarrow
\textit{PMod}(\mathcal{O}_U)
$$
et des foncteurs
$$
j_{U*}, j_{U!} :
\textit{PMod}(\mathcal{O}_U)
\longrightarrow
\textit{PMod}(\mathcal{O})
$$
qui sont respectivement adjoint à droite et adjoint à gauche de $j_U^*$.
En examinant la preuve du Lemme \ref{lemma-extension-by-zero}, on voit que
$j_{U!}\mathcal{G}$ est le préfaisceau
$$
V \longmapsto
\bigoplus\nolimits_{\varphi \in \Mor_\mathcal{C}(V, U)}
\mathcal{G}(V \xrightarrow{\varphi} U)
$$
En outre, le foncteur $j_{U!}$ est exact (d'après le
Lemme \ref{lemma-extension-by-zero-exact}, appliqué aux topologies
discrètes). Si de plus $\mathcal{C}$ est un site et $\mathcal{O}$ un faisceau
d'anneaux, alors le diagramme
$$
\xymatrix{
\textit{Mod}(\mathcal{O}_U) \ar[r]_{j_{U!}} \ar[d]_{forget} &
\textit{Mod}(\mathcal{O}) \\
\textit{PMod}(\mathcal{O}_U) \ar[r]^{j_{U!}} &
\textit{PMod}(\mathcal{O}) \ar[u]_{(\ )^\#}
}
$$
est commutatif.
\end{remark}

\begin{lemma}
\label{lemma-restrict-back}
Soit $\mathcal{C}$ un site et $U \in \Ob(\mathcal{C})$.
Supposons que tout $X$ de $\mathcal{C}$ admette au plus un morphisme vers $U$.
Soit $\mathcal{F}$ un faisceau abélien sur $\mathcal{C}/U$.
Les applications canoniques $\mathcal{F} \to j_U^{-1}j_{U!}\mathcal{F}$
et $j_U^{-1}j_{U*}\mathcal{F} \to \mathcal{F}$ sont des isomorphismes.
\end{lemma}

\begin{proof}
C'est un cas particulier du Lemme \ref{lemma-back-and-forth}, car l'hypothèse
sur $U$ équivaut à la pleine fidélité du foncteur de localisation
$\mathcal{C}/U \to \mathcal{C}$.
\end{proof}









\section{Localisation des morphismes de sites annelés}
\label{section-localize-morphisms}

\noindent
Cette section est l'analogue de
Sites, Section \ref{sites-section-localize-morphisms}.

\begin{lemma}
\label{lemma-localize-morphism-ringed-sites}
Soit
$(f, f^\sharp) :
(\mathcal{C}, \mathcal{O})
\longrightarrow
(\mathcal{D}, \mathcal{O}')$
un morphisme de sites annelés, où $f$ est donné par le foncteur continu
$u : \mathcal{D} \to \mathcal{C}$.
Soit $V$ un objet de $\mathcal{D}$ et posons $U = u(V)$.
Il existe alors une application canonique de faisceaux d'anneaux
$(f')^\sharp$ telle que le diagramme du
Lemme \ref{sites-lemma-localize-morphism} de Sites
devienne un diagramme commutatif de topos annelés
$$
\xymatrix{
(\Sh(\mathcal{C}/U), \mathcal{O}_U)
\ar[rr]_{(j_U, j_U^\sharp)} \ar[d]_{(f', (f')^\sharp)} & &
(\Sh(\mathcal{C}), \mathcal{O})
\ar[d]^{(f, f^\sharp)} \\
(\Sh(\mathcal{D}/V), \mathcal{O}'_V)
\ar[rr]^{(j_V, j_V^\sharp)} & &
(\Sh(\mathcal{D}), \mathcal{O}').
}
$$
De plus, dans cette situation, on a $f'_*j_U^{-1} = j_V^{-1}f_*$
et $f'_*j_U^* = j_V^*f_*$.
\end{lemma}

\begin{proof}
Il suffit de prendre pour $(f')^\sharp$ l'application
$$
(f')^{-1}\mathcal{O}'_V =
(f')^{-1}j_V^{-1}\mathcal{O}' =
j_U^{-1}f^{-1}\mathcal{O}' \xrightarrow{j_U^{-1}f^\sharp}
j_U^{-1}\mathcal{O} = \mathcal{O}_U
$$
et tout le reste résulte de
Sites, Lemme \ref{sites-lemma-localize-morphism}.
(Notons que $j^{-1} = j^*$ sur les faisceaux de modules lorsque $j$ est un
morphisme de localisation ; la première égalité de foncteurs implique donc
la seconde.)
\end{proof}

\begin{lemma}
\label{lemma-relocalize-morphism-ringed-sites}
Soit
$(f, f^\sharp) :
(\mathcal{C}, \mathcal{O})
\longrightarrow
(\mathcal{D}, \mathcal{O}')$
un morphisme de sites annelés, où $f$ est donné par le foncteur continu
$u : \mathcal{D} \to \mathcal{C}$.
Soient $V \in \Ob(\mathcal{D})$, $U \in \Ob(\mathcal{C})$ et
$c : U \to u(V)$ un morphisme de $\mathcal{C}$.
Il existe un diagramme commutatif de topos annelés
$$
\xymatrix{
(\Sh(\mathcal{C}/U), \mathcal{O}_U)
\ar[rr]_{(j_U, j_U^\sharp)} \ar[d]_{(f_c, f_c^\sharp)} & &
(\Sh(\mathcal{C}), \mathcal{O}) \ar[d]^{(f, f^\sharp)} \\
(\Sh(\mathcal{D}/V), \mathcal{O}'_V)
\ar[rr]^{(j_V, j_V^\sharp)} & &
(\Sh(\mathcal{D}), \mathcal{O}').
}
$$
Le morphisme $(f_c, f_c^\sharp)$ est égal au composé du morphisme
$$
(f', (f')^\sharp) :
(\Sh(\mathcal{C}/u(V)), \mathcal{O}_{u(V)})
\longrightarrow
(\Sh(\mathcal{D}/V), \mathcal{O}'_V)
$$
du Lemme \ref{lemma-localize-morphism-ringed-sites} et du morphisme
$$
(j, j^\sharp) :
(\Sh(\mathcal{C}/U), \mathcal{O}_U)
\to
(\Sh(\mathcal{C}/u(V)), \mathcal{O}_{u(V)})
$$
du Lemme \ref{lemma-relocalize}.
Étant donnés des morphismes $b : V' \to V$, $a : U' \to U$ et
$c' : U' \to u(V')$ tels que
$$
\xymatrix{
U' \ar[r]_-{c'} \ar[d]_a & u(V') \ar[d]^{u(b)} \\
U \ar[r]^-c & u(V)
}
$$
soit commutatif, le diagramme suivant de topos annelés
$$
\xymatrix{
(\Sh(\mathcal{C}/U'), \mathcal{O}_{U'})
\ar[rr]_{(j_{U'/U}, j_{U'/U}^\sharp)} \ar[d]_{(f_{c'}, f_{c'}^\sharp)} & &
(\Sh(\mathcal{C}/U), \mathcal{O}_U)
\ar[d]^{(f_c, f_c^\sharp)} \\
(\Sh(\mathcal{D}/V'), \mathcal{O}'_{V'})
\ar[rr]^{(j_{V'/V}, j_{V'/V}^\sharp)} & &
(\Sh(\mathcal{D}/V), \mathcal{O}'_{V'})
}
$$
est commutatif.
\end{lemma}

\begin{proof}
Au niveau des morphismes de topos, c'est le
Lemme \ref{sites-lemma-relocalize-morphism} de Sites.
Pour vérifier que les diagrammes commutent comme diagrammes de morphismes de
topos annelés, on utilise les Lemmes \ref{lemma-relocalize} et
\ref{lemma-localize-morphism-ringed-sites}, exactement comme dans la preuve du
Lemme \ref{sites-lemma-relocalize-morphism} de Sites.
\end{proof}
















\section{Localisation de topos annelés}
\label{section-localize-ringed-topoi}

\noindent
Cette section est l'analogue de
Sites, Section \ref{sites-section-localize-topoi}
dans le cadre des topos annelés.

\begin{lemma}
\label{lemma-localize-ringed-topos}
Soit $(\Sh(\mathcal{C}), \mathcal{O})$ un topos annelé.
Soit $\mathcal{F} \in \Sh(\mathcal{C})$ un faisceau.
Pour tout faisceau $\mathcal{H}$ sur $\mathcal{C}$, notons
$\mathcal{H}_\mathcal{F}$ le faisceau $\mathcal{H} \times \mathcal{F}$,
considéré comme objet de la catégorie
$\Sh(\mathcal{C})/\mathcal{F}$. Le couple
$(\Sh(\mathcal{C})/\mathcal{F}, \mathcal{O}_\mathcal{F})$
est un topos annelé et il existe un morphisme canonique de topos annelés
$$
(j_\mathcal{F}, j_\mathcal{F}^\sharp) :
(\Sh(\mathcal{C})/\mathcal{F}, \mathcal{O}_\mathcal{F})
\longrightarrow
(\Sh(\mathcal{C}), \mathcal{O})
$$
qui est une localisation au sens de la
Section \ref{section-localize} et tel que
\begin{enumerate}
\item le foncteur $j_\mathcal{F}^{-1}$ est le foncteur
$\mathcal{H} \mapsto \mathcal{H}_\mathcal{F}$,
\item le foncteur $j_\mathcal{F}^*$ est le foncteur
$\mathcal{H} \mapsto \mathcal{H}_\mathcal{F}$,
\item le foncteur $j_{\mathcal{F}!}$ sur les faisceaux d'ensembles est le
foncteur d'oubli $\mathcal{G}/\mathcal{F} \mapsto \mathcal{G}$,
\item le foncteur $j_{\mathcal{F}!}$ sur les faisceaux de modules associe au
$\mathcal{O}_\mathcal{F}$-module
$\varphi : \mathcal{G} \to \mathcal{F}$ le $\mathcal{O}$-module qui est la
faisceautisation du préfaisceau
$$
V \longmapsto
\bigoplus\nolimits_{s \in \mathcal{F}(V)}
\{\sigma \in \mathcal{G}(V) \mid \varphi(\sigma) = s \}
$$
pour $V \in \Ob(\mathcal{C})$.
\end{enumerate}
\end{lemma}

\begin{proof}
D'après le Lemme \ref{sites-lemma-localize-topos} de Sites,
$\Sh(\mathcal{C})/\mathcal{F}$ est un topos et (1) et (3) sont vrais.
En particulier, cela montre que
$j_\mathcal{F}^{-1}\mathcal{O} = \mathcal{O}_\mathcal{F}$ et que
$\mathcal{O}_\mathcal{F}$ est un faisceau d'anneaux.
On peut donc choisir pour $j_\mathcal{F}^\sharp$ l'application identique ;
en particulier, (2) est vrai. De plus, la preuve du
Lemme \ref{sites-lemma-localize-topos} de Sites montre que l'on peut supposer
que $\mathcal{C}$ est un site admettant toutes les limites finies et muni
d'une topologie sous-canonique, et que $\mathcal{F} = h_U$ pour un objet
$U$ de $\mathcal{C}$. Alors (4) résulte de la description de
$j_{\mathcal{F}!}$ au Lemme \ref{lemma-extension-by-zero}.
On pourrait aussi montrer directement que le foncteur décrit en (4) est
adjoint à gauche de $j_\mathcal{F}^*$.
\end{proof}

\begin{definition}
\label{definition-localize-ringed-topos}
Soit $(\Sh(\mathcal{C}), \mathcal{O})$ un topos annelé.
Soit $\mathcal{F} \in \Sh(\mathcal{C})$.
\begin{enumerate}
\item Le topos annelé
$(\Sh(\mathcal{C})/\mathcal{F}, \mathcal{O}_\mathcal{F})$
est appelé la
{\it localisation du topos annelé
$(\Sh(\mathcal{C}), \mathcal{O})$ en $\mathcal{F}$}.
\item Le morphisme de topos annelés
$(j_\mathcal{F}, j_\mathcal{F}^\sharp) :
(\Sh(\mathcal{C})/\mathcal{F}, \mathcal{O}_\mathcal{F})
\to
(\Sh(\mathcal{C}), \mathcal{O})$ of
Lemme \ref{lemma-localize-ringed-topos}
est appelé le {\it morphisme de localisation}.
\end{enumerate}
\end{definition}

\noindent
Conformément à l'usage établi dans le chapitre sur les sites, nous
vérifions, chaque fois que cela a un sens, que les constructions de
localisation des topos sont compatibles avec celles des sites.

\begin{lemma}
\label{lemma-localize-compare}
Soient
$(\Sh(\mathcal{C}), \mathcal{O})$ and
$\mathcal{F} \in \Sh(\mathcal{C})$ comme dans le
Lemme \ref{lemma-localize-ringed-topos}.
Si $\mathcal{F} = h_U^\#$ pour un objet $U$ de $\mathcal{C}$, alors, via
l'identification
$\Sh(\mathcal{C}/U) = \Sh(\mathcal{C})/h_U^\#$ of
Sites, Lemme \ref{sites-lemma-essential-image-j-shriek}
on a
\begin{enumerate}
\item canoniquement $\mathcal{O}_U = \mathcal{O}_\mathcal{F}$, et
\item sous ces identifications,
$(j_\mathcal{F}, j_\mathcal{F}^\sharp) = (j_U, j_U^\sharp)$.
\end{enumerate}
\end{lemma}

\begin{proof}
L'assertion concernant les topos sous-jacents est le
Lemme \ref{sites-lemma-localize-compare} de Sites.
Notons que $\mathcal{O}_U$ est la restriction de $\mathcal{O}$ qui, d'après
Sites, Lemme \ref{sites-lemma-compute-j-shriek-restrict}
correspond à $\mathcal{O} \times h_U^\#$ sous l'équivalence fournie par
Sites, Lemme \ref{sites-lemma-essential-image-j-shriek}.
La définition de $\mathcal{O}_\mathcal{F}$ donne (1).
Il reste à prouver que $j_\mathcal{F}^\sharp = j_U^\sharp$ sous cette
identification. Nous omettons la vérification.
\end{proof}

\noindent
La localisation est fonctorielle de deux manières : on peut
« relocaliser » une localisation (voir le
Lemme \ref{lemma-relocalize-ringed-topos}) ou, étant donné un morphisme de
topos annelés, localiser en haut en l'image inverse d'un faisceau d'en bas et
obtenir un diagramme commutatif de topos annelés (voir le
Lemme \ref{lemma-localize-morphism-ringed-topoi}).

\begin{lemma}
\label{lemma-relocalize-ringed-topos}
Soit $(\Sh(\mathcal{C}), \mathcal{O})$ un topos annelé.
Si $s : \mathcal{G} \to \mathcal{F}$ est un morphisme de faisceaux sur
$\mathcal{C}$, il existe un diagramme commutatif naturel de morphismes de
topos annelés
$$
\xymatrix{
(\Sh(\mathcal{C})/\mathcal{G}, \mathcal{O}_\mathcal{G})
\ar[rd]_{(j_\mathcal{G}, j_\mathcal{G}^\sharp)} \ar[rr]_{(j, j^\sharp)} & &
(\Sh(\mathcal{C})/\mathcal{F}, \mathcal{O}_\mathcal{F})
\ar[ld]^{(j_\mathcal{F}, j_\mathcal{F}^\sharp)} \\
& (\Sh(\mathcal{C}), \mathcal{O}) &
}
$$
où $(j, j^\sharp)$ est le morphisme de localisation du topos annelé
$(\Sh(\mathcal{C})/\mathcal{F}, \mathcal{O}_\mathcal{F})$
en l'objet $\mathcal{G}/\mathcal{F}$.
\end{lemma}

\begin{proof}
Toutes les assertions résultent de
Sites, Lemme \ref{sites-lemma-relocalize-topos}
excepté l'assertion
$j_\mathcal{G}^\sharp = j^\sharp \circ j^{-1}(j_\mathcal{F}^\sharp)$.
Nous omettons la vérification.
\end{proof}

\begin{lemma}
\label{lemma-relocalize-compare}
Soient $(\Sh(\mathcal{C}), \mathcal{O})$ et
$s : \mathcal{G} \to \mathcal{F}$ comme dans le
Lemme \ref{lemma-relocalize-ringed-topos}.
S'il existe un morphisme $f : V \to U$ de $\mathcal{C}$ tel que
$\mathcal{G} = h_V^\#$, $\mathcal{F} = h_U^\#$ et que $s$ soit induit par
$f$, alors les diagrammes du
Lemme \ref{lemma-relocalize}
et du
Lemme \ref{lemma-relocalize-ringed-topos}
coïncident via les identifications
$(j_\mathcal{F}, j_\mathcal{F}^\sharp) = (j_U, j_U^\sharp)$
and
$(j_\mathcal{G}, j_\mathcal{G}^\sharp) = (j_V, j_V^\sharp)$
du Lemme \ref{lemma-localize-compare}.
\end{lemma}

\begin{proof}
Toutes les assertions résultent de
Sites, Lemme \ref{sites-lemma-relocalize-compare}
excepté l'assertion selon laquelle les deux applications $j^\sharp$
coïncident. Elle est vraie puisque, dans les deux cas, $j^\sharp$ est
simplement l'identité. Nous omettons quelques détails.
\end{proof}










\section{Localisation des morphismes de topos annelés}
\label{section-localize-morphisms-ringed-topoi}

\noindent
Cette section est l'analogue de
Sites, Section \ref{sites-section-localize-morphisms-topoi}.

\begin{lemma}
\label{lemma-localize-morphism-ringed-topoi}
Soit
$$
f :
(\Sh(\mathcal{C}), \mathcal{O})
\longrightarrow
(\Sh(\mathcal{D}), \mathcal{O}')
$$
soit un morphisme de topos annelés. Soit $\mathcal{G}$ un faisceau sur
$\mathcal{D}$. Posons $\mathcal{F} = f^{-1}\mathcal{G}$.
Il existe alors un diagramme commutatif de topos annelés
$$
\xymatrix{
(\Sh(\mathcal{C})/\mathcal{F}, \mathcal{O}_\mathcal{F})
\ar[rr]_{(j_\mathcal{F}, j_\mathcal{F}^\sharp)}
\ar[d]_{(f', (f')^\sharp)} & &
(\Sh(\mathcal{C}), \mathcal{O}) \ar[d]^{(f, f^\sharp)} \\
(\Sh(\mathcal{D})/\mathcal{G}, \mathcal{O}'_\mathcal{G})
\ar[rr]^{(j_\mathcal{G}, j_\mathcal{G}^\sharp)} & &
(\Sh(\mathcal{D}), \mathcal{O}')
}
$$
On a $f'_*j_\mathcal{F}^{-1} = j_\mathcal{G}^{-1}f_*$ et
$f'_*j_\mathcal{F}^* = j_\mathcal{G}^*f_*$. De plus, le morphisme $f'$ est
caractérisé par la règle
$$
(f')^{-1}(\mathcal{H} \xrightarrow{\varphi} \mathcal{G})
=
(f^{-1}\mathcal{H} \xrightarrow{f^{-1}\varphi} \mathcal{F}).
$$
\end{lemma}

\begin{proof}
Le Lemme \ref{sites-lemma-localize-morphism-topoi} de Sites donne le
diagramme des topos sous-jacents, l'égalité
$f'_*j_\mathcal{F}^{-1} = j_\mathcal{G}^{-1}f_*$ et la description de
$(f')^{-1}$. Pour définir $(f')^\sharp$, on utilise l'application
$$
(f')^\sharp :
\mathcal{O}'_\mathcal{G} =
j_\mathcal{G}^{-1} \mathcal{O}'
\xrightarrow{j_\mathcal{G}^{-1}f^\sharp}
j_\mathcal{G}^{-1} f_*\mathcal{O} =
f'_* j_\mathcal{F}^{-1}\mathcal{O} =
f'_* \mathcal{O}_\mathcal{F}
$$
ou, de manière équivalente, l'application
$$
(f')^\sharp :
(f')^{-1}\mathcal{O}'_\mathcal{G} =
(f')^{-1}j_\mathcal{G}^{-1} \mathcal{O}' =
j_\mathcal{F}^{-1}f^{-1}\mathcal{O}'
\xrightarrow{j_\mathcal{F}^{-1}f^\sharp}
j_\mathcal{F}^{-1} \mathcal{O} =
\mathcal{O}_\mathcal{F}.
$$
Nous omettons la vérification que ces deux applications sont bien adjointes.
La seconde construction de $(f')^\sharp$ montre que le diagramme commute
dans la $2$-catégorie des topos annelés (car les applications
$j_\mathcal{F}^\sharp$ et $j_\mathcal{G}^\sharp$ sont des identités).
Enfin, l'égalité $f'_*j_\mathcal{F}^* = j_\mathcal{G}^*f_*$ résulte de
$f'_*j_\mathcal{F}^{-1} = j_\mathcal{G}^{-1}f_*$
et du fait que les images inverses de faisceaux de modules et de faisceaux
d'ensembles coïncident, voir le Lemme \ref{lemma-localize-ringed-topos}.
\end{proof}

\begin{lemma}
\label{lemma-localize-morphism-compare}
Soit
$$
f :
(\Sh(\mathcal{C}), \mathcal{O})
\longrightarrow
(\Sh(\mathcal{D}), \mathcal{O}')
$$
soit un morphisme de topos annelés.
Soit $\mathcal{G}$ un faisceau sur $\mathcal{D}$.
Posons $\mathcal{F} = f^{-1}\mathcal{G}$.
Si $f$ est donné par un foncteur continu $u : \mathcal{D} \to \mathcal{C}$
et $\mathcal{G} = h_V^\#$, alors les diagrammes commutatifs du
Lemme \ref{lemma-localize-morphism-ringed-sites}
et du
Lemme \ref{lemma-localize-morphism-ringed-topoi}
coïncident via les identifications du
Lemme \ref{lemma-localize-compare}.
\end{lemma}

\begin{proof}
Au niveau des morphismes de topos, c'est le
Lemme \ref{sites-lemma-localize-morphism-compare} de Sites.
Cela vaut aussi au niveau des morphismes de topos annelés puisque les
formules qui définissent $(f')^\sharp$ dans les preuves du
Lemme \ref{lemma-localize-morphism-ringed-sites}
et du
Lemme \ref{lemma-localize-morphism-ringed-topoi}
coïncident.
\end{proof}

\begin{lemma}
\label{lemma-relocalize-morphism-ringed-topoi}
Soit
$(f, f^\sharp) :
(\Sh(\mathcal{C}), \mathcal{O})
\to
(\Sh(\mathcal{D}), \mathcal{O}')$
soit un morphisme de topos annelés.
Soient $\mathcal{G}$ un faisceau sur $\mathcal{D}$,
$\mathcal{F}$ un faisceau sur $\mathcal{C}$ et
$s : \mathcal{F} \to f^{-1}\mathcal{G}$ un morphisme de faisceaux.
Il existe un diagramme commutatif de topos annelés
$$
\xymatrix{
(\Sh(\mathcal{C})/\mathcal{F}, \mathcal{O}_\mathcal{F})
\ar[rr]_{(j_\mathcal{F}, j_\mathcal{F}^\sharp)}
\ar[d]_{(f_c, f_c^\sharp)} & &
(\Sh(\mathcal{C}), \mathcal{O})
\ar[d]^{(f, f^\sharp)} \\
(\Sh(\mathcal{D})/\mathcal{G}, \mathcal{O}'_\mathcal{G})
\ar[rr]^{(j_\mathcal{G}, j_\mathcal{G}^\sharp)} & &
(\Sh(\mathcal{D}), \mathcal{O}').
}
$$
Le morphisme $(f_s, f_s^\sharp)$ est égal au composé du morphisme
$$
(f', (f')^\sharp) :
(\Sh(\mathcal{C})/f^{-1}\mathcal{G}, \mathcal{O}_{f^{-1}\mathcal{G}})
\longrightarrow
(\Sh(\mathcal{D})/{\mathcal{G}}, \mathcal{O}'_\mathcal{G})
$$
du Lemme \ref{lemma-localize-morphism-ringed-topoi} et du morphisme
$$
(j, j^\sharp) :
(\Sh(\mathcal{C})/\mathcal{F}, \mathcal{O}_\mathcal{F})
\to
(\Sh(\mathcal{C})/f^{-1}\mathcal{G}, \mathcal{O}_{f^{-1}\mathcal{G}})
$$
du Lemme \ref{lemma-relocalize-ringed-topos}.
Étant donnés des morphismes $b : \mathcal{G}' \to \mathcal{G}$,
$a : \mathcal{F}' \to \mathcal{F}$ et
$s' : \mathcal{F}' \to f^{-1}\mathcal{G}'$ tels que
$$
\xymatrix{
\mathcal{F}' \ar[r]_-{s'} \ar[d]_a &
f^{-1}\mathcal{G}' \ar[d]^{f^{-1}b} \\
\mathcal{F} \ar[r]^-s &
f^{-1}\mathcal{G}
}
$$
soit commutatif, le diagramme suivant de topos annelés
$$
\xymatrix{
(\Sh(\mathcal{C})/\mathcal{F}', \mathcal{O}_{\mathcal{F}'})
\ar[rr]_{(j_{\mathcal{F}'/\mathcal{F}}, j_{\mathcal{F}'/\mathcal{F}}^\sharp)}
\ar[d]_{(f_{s'}, f_{s'}^\sharp)} & &
(\Sh(\mathcal{C})/\mathcal{F}, \mathcal{O}_\mathcal{F})
\ar[d]^{(f_s, f_s^\sharp)} \\
(\Sh(\mathcal{D})/\mathcal{G}', \mathcal{O}'_{\mathcal{G}'})
\ar[rr]^{(j_{\mathcal{G}'/\mathcal{G}}, j_{\mathcal{G}'/\mathcal{G}}^\sharp)}
& &
(\Sh(\mathcal{D})/\mathcal{G}, \mathcal{O}'_{\mathcal{G}'})
}
$$
est commutatif.
\end{lemma}

\begin{proof}
Au niveau des morphismes de topos, c'est le
Lemme \ref{sites-lemma-relocalize-morphism-topoi} de Sites.
Pour vérifier que les diagrammes commutent comme diagrammes de morphismes de
topos annelés, on utilise les diagrammes commutatifs des
Lemmes \ref{lemma-relocalize-ringed-topos} et
\ref{lemma-localize-morphism-ringed-topoi}.
\end{proof}

\begin{lemma}
\label{lemma-relocalize-morphism-compare}
Soient
$(f, f^\sharp) :
(\Sh(\mathcal{C}), \mathcal{O})
\to
(\Sh(\mathcal{D}), \mathcal{O}')$,
$s : \mathcal{F} \to f^{-1}\mathcal{G}$ étant comme dans le
Lemme \ref{lemma-relocalize-morphism-ringed-topoi}.
Si $f$ est donné par un foncteur continu
$u : \mathcal{D} \to \mathcal{C}$
et si $\mathcal{G} = h_V^\#$, $\mathcal{F} = h_U^\#$ et $s$ provient d'un
morphisme $c : U \to u(V)$, alors les diagrammes commutatifs du
Lemme \ref{lemma-relocalize-morphism-ringed-sites}
et du
Lemme \ref{lemma-relocalize-morphism-ringed-topoi}
coïncident via les identifications du
Lemme \ref{lemma-localize-compare}.
\end{lemma}

\begin{proof}
Cela résulte formellement des
Lemmes \ref{lemma-relocalize-compare} et
\ref{lemma-localize-morphism-compare}.
\end{proof}

















\section{Propriétés locales des modules}
\label{section-local}

\noindent
Conformément à la stratégie générale expliquée à la
Section \ref{section-intrinsic}, nous définissons d'abord les propriétés
locales des faisceaux de modules sur un site annelé, puis montrons aussitôt
qu'elles sont intrinsèques et ont donc un sens pour les faisceaux de modules
sur les topos annelés.

\begin{definition}
\label{definition-site-local}
Soit $(\mathcal{C}, \mathcal{O})$ un site annelé.
Soit $\mathcal{F}$ un faisceau de $\mathcal{O}$-modules.
Nous utiliserons librement les notions définies dans la
Définition \ref{definition-global}.
\begin{enumerate}
\item On dit que $\mathcal{F}$ est {\it localement libre} si, pour tout objet
$U$ de $\mathcal{C}$, il existe un recouvrement
$\{U_i \to U\}_{i \in I}$ de $\mathcal{C}$ tel que chaque restriction
$\mathcal{F}|_{\mathcal{C}/U_i}$ soit un
$\mathcal{O}_{U_i}$-module libre.
\item On dit que $\mathcal{F}$ est {\it localement libre de rang fini} si,
pour tout objet $U$ de $\mathcal{C}$, il existe un recouvrement
$\{U_i \to U\}_{i \in I}$ de $\mathcal{C}$ tel que chaque restriction
$\mathcal{F}|_{\mathcal{C}/U_i}$ soit un
$\mathcal{O}_{U_i}$-module libre de rang fini.
\item On dit que $\mathcal{F}$ est {\it localement engendré par des sections}
si, pour tout objet $U$ de $\mathcal{C}$, il existe un recouvrement
$\{U_i \to U\}_{i \in I}$ de $\mathcal{C}$ tel que chaque restriction
$\mathcal{F}|_{\mathcal{C}/U_i}$ soit un $\mathcal{O}_{U_i}$-module engendré
par ses sections globales.
\item Étant donné $r \geq 0$, on dit que $\mathcal{F}$ est
{\it localement engendré par $r$ sections} si, pour tout objet $U$ de
$\mathcal{C}$, il existe un recouvrement
$\{U_i \to U\}_{i \in I}$ de $\mathcal{C}$ tel que chaque restriction
$\mathcal{F}|_{\mathcal{C}/U_i}$ soit un $\mathcal{O}_{U_i}$-module engendré
par $r$ sections globales.
\item On dit que $\mathcal{F}$ est {\it de type fini} si, pour tout objet $U$
de $\mathcal{C}$, il existe un recouvrement
$\{U_i \to U\}_{i \in I}$ de $\mathcal{C}$ tel que chaque restriction
$\mathcal{F}|_{\mathcal{C}/U_i}$ soit un $\mathcal{O}_{U_i}$-module engendré
par un nombre fini de sections globales.
\item On dit que $\mathcal{F}$ est {\it quasi-cohérent} si, pour tout objet
$U$ de $\mathcal{C}$, il existe un recouvrement
$\{U_i \to U\}_{i \in I}$ de $\mathcal{C}$ tel que chaque restriction
$\mathcal{F}|_{\mathcal{C}/U_i}$ soit un $\mathcal{O}_{U_i}$-module admettant
une présentation globale.
\item On dit que $\mathcal{F}$ est {\it de présentation finie} si, pour tout
objet $U$ de $\mathcal{C}$, il existe un recouvrement
$\{U_i \to U\}_{i \in I}$ de $\mathcal{C}$ tel que chaque restriction
$\mathcal{F}|_{\mathcal{C}/U_i}$ soit un $\mathcal{O}_{U_i}$-module admettant
une présentation globale finie.
\item On dit que $\mathcal{F}$ est {\it cohérent} si et seulement si
$\mathcal{F}$ est de type fini et si, pour tout objet $U$ de $\mathcal{C}$
et tous $s_1, \ldots, s_n \in \mathcal{F}(U)$, le noyau de l'application
$\bigoplus_{i = 1, \ldots, n} \mathcal{O}_U \to \mathcal{F}|_U$
est de type fini sur $(\mathcal{C}/U, \mathcal{O}_U)$.
\end{enumerate}
\end{definition}

\begin{lemma}
\label{lemma-special-locally-free}
Chacune des propriétés (1) -- (8) de la
Définition \ref{definition-site-local} est intrinsèque (voir la discussion
de la Section \ref{section-intrinsic}).
\end{lemma}

\begin{proof}
Soient $\mathcal{C}$ et $\mathcal{D}$ des sites.
Soit $u : \mathcal{C} \to \mathcal{D}$ un foncteur cocontinu spécial.
Soit $\mathcal{O}$ un faisceau d'anneaux sur $\mathcal{C}$.
Soit $\mathcal{F}$ un faisceau de $\mathcal{O}$-modules sur $\mathcal{C}$.
Soit $g : \Sh(\mathcal{C}) \to \Sh(\mathcal{D})$ l'équivalence de topos
associée à $u$.
Posons $\mathcal{O}' = g_*\mathcal{O}$ et prenons pour
$g^\sharp : \mathcal{O}' \to g_*\mathcal{O}$ l'identité.
Enfin, posons $\mathcal{F}' = g_*\mathcal{F}$.
Soit $\mathcal{P}_l$ l'une des propriétés (1) -- (7) de la
Définition \ref{definition-site-local}.
(Nous traiterons le cas cohérent à la fin de la preuve.)
Notons $\mathcal{P}_g$ la propriété correspondante de la
Définition \ref{definition-global}. Nous avons déjà vu que
$\mathcal{P}_g$ est intrinsèque. Il faut montrer que
$\mathcal{P}_l(\mathcal{C}, \mathcal{O}, \mathcal{F})$
est vraie si et seulement si
$\mathcal{P}_l(\mathcal{D}, \mathcal{O}', \mathcal{F}')$
est vraie.

\medskip\noindent
Supposons que $\mathcal{F}$ possède la propriété $\mathcal{P}_l$.
Soit $V$ un objet de $\mathcal{D}$.
L'une des propriétés d'un foncteur cocontinu spécial assure l'existence
d'un recouvrement $\{u(U_i) \to V\}_{i \in I}$ dans le site $\mathcal{D}$.
Par hypothèse, pour chaque $i$, il existe un recouvrement
$\{U_{ij} \to U_i\}_{j \in J_i}$ dans $\mathcal{C}$ tel que chaque restriction
$\mathcal{F}|_{U_{ij}}$ possède la propriété $\mathcal{P}_g$. D'après
Sites, Lemme \ref{sites-lemma-localize-special-cocontinuous}
on a des diagrammes commutatifs de topos annelés
$$
\xymatrix{
(\Sh(\mathcal{C}/U_{ij}), \mathcal{O}_{U_{ij}}) \ar[r] \ar[d] &
(\Sh(\mathcal{C}), \mathcal{O}) \ar[d] \\
(\Sh(\mathcal{D}/u(U_{ij})), \mathcal{O}'_{u(U_{ij})}) \ar[r] &
(\Sh(\mathcal{D}), \mathcal{O}')
}
$$
où les flèches verticales sont des équivalences. Ainsi
$\mathcal{F}'|_{u(U_{ij})}$ possède aussi la propriété $\mathcal{P}_g$.
De plus, $\{u(U_{ij}) \to V\}_{i \in I, j \in J_i}$ est un recouvrement du
site $\mathcal{D}$. Par conséquent, $\mathcal{F}'$ possède la propriété
$\mathcal{P}_l$.

\medskip\noindent
Supposons que $\mathcal{F}'$ possède la propriété $\mathcal{P}_l$.
Soit $U$ un objet de $\mathcal{C}$.
Par hypothèse, il existe un recouvrement
$\{V_i \to u(U)\}_{i \in I}$ tel que $\mathcal{F}'|_{V_i}$ possède la
propriété $\mathcal{P}_g$. Puisque $u$ est cocontinu, on peut raffiner ce
recouvrement par une famille $\{u(U_j) \to u(U)\}_{j \in J}$, où
$\{U_j \to U\}_{j \in J}$ est un recouvrement dans $\mathcal{C}$.
Notons que le raffinement est donné par $\alpha : J \to I$ et par des
morphismes $u(U_j) \to V_{\alpha(j)}$.
La restriction est transitive, c'est-à-dire
$(\mathcal{F}'|_{V_{\alpha(j)}})|_{u(U_j)} = \mathcal{F}'|_{u(U_j)}$.
Le Lemme \ref{lemma-global-pullback} montre donc que
$\mathcal{F}'|_{u(U_j)}$ possède la propriété $\mathcal{P}_g$.
Ainsi, le diagramme
$$
\xymatrix{
(\Sh(\mathcal{C}/U_j), \mathcal{O}_{U_j}) \ar[r] \ar[d] &
(\Sh(\mathcal{C}), \mathcal{O}) \ar[d] \\
(\Sh(\mathcal{D}/u(U_j)), \mathcal{O}'_{u(U_j)})
\ar[r] &
(\Sh(\mathcal{D}), \mathcal{O}')
}
$$
où les flèches verticales sont des équivalences montre que
$\mathcal{F}|_{U_j}$ possède aussi la propriété $\mathcal{P}_g$.
Ainsi $\mathcal{F}$ possède la propriété $\mathcal{P}_l$, comme voulu.

\medskip\noindent
Enfin, prouvons le lemme dans le cas
$\mathcal{P}_l = coherent$\footnote{Les détails techniques sont quelque peu
malcommodes ; nous conseillons au lecteur de sauter cette partie de la preuve.}.
Supposons $\mathcal{F}$ cohérent. Alors $\mathcal{F}$ est de type fini, et la
première partie de la preuve montre que $\mathcal{F}'$ l'est aussi.
Soit $V$ un objet de $\mathcal{D}$ et soient
$s_1, \ldots, s_n \in \mathcal{F}'(V)$. Il faut montrer que le noyau
$\mathcal{K}'$ of
$\bigoplus_{j = 1, \ldots, n} \mathcal{O}_V \to \mathcal{F}'|_V$
soit de type fini sur $\mathcal{D}/V$. Il faut donc montrer que, pour tout
$V'/V$, il existe un recouvrement $\{V'_i \to V'\}$ tel que
$\mathcal{F}'|_{V'_i}$ soit engendré par un nombre fini de sections.
En remplaçant $V$ par $V'$ (et en restreignant les sections $s_j$ à $V'$),
on se ramène au cas $V' = V$. Puisque $u$ est un foncteur cocontinu spécial,
il existe un recouvrement $\{u(U_i) \to V\}_{i \in I}$ dans le site
$\mathcal{D}$. En utilisant l'isomorphisme de topos
$\Sh(\mathcal{C}/U_i) = \Sh(\mathcal{D}/u(U_i))$
on voit que $\mathcal{K}'|_{u(U_i)}$ correspond au noyau
$\mathcal{K}_i$ d'un morphisme
$\bigoplus_{j = 1, \ldots, n} \mathcal{O}_{U_i} \to \mathcal{F}|_{U_i}$.
La cohérence de $\mathcal{F}$ montre que $\mathcal{K}_i$ est de type fini.
On en conclut (en appliquant à nouveau la première partie de la preuve) que
$\mathcal{K}|_{u(U_i)}$ est de type fini. Il existe donc des recouvrements
$\{V_{il} \to u(U_i)\}$ tels que $\mathcal{K}|_{V_{il}}$ soit engendré par un
nombre fini de sections globales. Comme $\{V_{il} \to V\}$ est un
recouvrement de $\mathcal{D}$, on conclut que $\mathcal{K}$ est de type fini,
comme voulu.

\medskip\noindent
Supposons $\mathcal{F}'$ cohérent. Alors $\mathcal{F}'$ est de type fini et,
d'après la première partie de la preuve, $\mathcal{F}$ l'est aussi.
Soit $U$ un objet de $\mathcal{C}$ et soient
$s_1, \ldots, s_n \in \mathcal{F}(U)$. Il faut montrer que le noyau
$\mathcal{K}$ of
$\bigoplus_{j = 1, \ldots, n} \mathcal{O}_U \to \mathcal{F}|_U$
est de type fini sur $\mathcal{C}/U$. En utilisant l'isomorphisme de topos
$\Sh(\mathcal{C}/U) = \Sh(\mathcal{D}/u(U))$
on voit que $\mathcal{K}|_U$ correspond au noyau
$\mathcal{K}'$ d'un morphisme
$\bigoplus_{j = 1, \ldots, n} \mathcal{O}_{u(U)} \to \mathcal{F}'|_{u(U)}$.
Comme $\mathcal{F}'$ est cohérent, $\mathcal{K}'$ est de type fini.
La première partie de la preuve montre alors que $\mathcal{K}$ est de type
fini.
\end{proof}

\noindent
Désormais, nous pouvons donc parler des propriétés des
$\mathcal{O}$-modules définies dans la Définition \ref{definition-site-local}
sans préciser de site.

\begin{lemma}
\label{lemma-local-final-object}
Soit $(\Sh(\mathcal{C}), \mathcal{O})$ un topos annelé.
Soit $\mathcal{F}$ un $\mathcal{O}$-module.
Supposons que le site $\mathcal{C}$ admette un objet final $X$.
Alors
\begin{enumerate}
\item Les conditions suivantes sont équivalentes :
\begin{enumerate}
\item $\mathcal{F}$ est localement libre ;
\item il existe un recouvrement $\{X_i \to X\}$ dans $\mathcal{C}$ tel que
chaque restriction $\mathcal{F}|_{\mathcal{C}/X_i}$ soit un
$\mathcal{O}_{X_i}$-module localement libre ;
\item il existe un recouvrement $\{X_i \to X\}$ dans $\mathcal{C}$ tel que
chaque restriction $\mathcal{F}|_{\mathcal{C}/X_i}$ soit un
$\mathcal{O}_{X_i}$-module libre.
\end{enumerate}
\item Les conditions suivantes sont équivalentes :
\begin{enumerate}
\item $\mathcal{F}$ est localement libre de rang fini ;
\item il existe un recouvrement $\{X_i \to X\}$ dans $\mathcal{C}$ tel que
chaque restriction $\mathcal{F}|_{\mathcal{C}/X_i}$ soit un
$\mathcal{O}_{X_i}$-module localement libre de rang fini ;
\item il existe un recouvrement $\{X_i \to X\}$ dans $\mathcal{C}$ tel que
chaque restriction $\mathcal{F}|_{\mathcal{C}/X_i}$ soit un
$\mathcal{O}_{X_i}$-module libre de rang fini.
\end{enumerate}
\item Les conditions suivantes sont équivalentes :
\begin{enumerate}
\item $\mathcal{F}$ est localement engendré par des sections ;
\item il existe un recouvrement $\{X_i \to X\}$ dans $\mathcal{C}$ tel que
chaque restriction $\mathcal{F}|_{\mathcal{C}/X_i}$ soit un
$\mathcal{O}_{X_i}$-module localement engendré par des sections ;
\item il existe un recouvrement $\{X_i \to X\}$ dans $\mathcal{C}$ tel que
chaque restriction $\mathcal{F}|_{\mathcal{C}/X_i}$ soit un
$\mathcal{O}_{X_i}$-module engendré par ses sections globales.
\end{enumerate}
\item Étant donné $r \geq 0$, les conditions suivantes sont équivalentes :
\begin{enumerate}
\item $\mathcal{F}$ est localement engendré par $r$ sections ;
\item il existe un recouvrement $\{X_i \to X\}$ dans $\mathcal{C}$ tel que
chaque restriction $\mathcal{F}|_{\mathcal{C}/X_i}$ soit un
$\mathcal{O}_{X_i}$-module localement engendré par $r$ sections ;
\item il existe un recouvrement $\{X_i \to X\}$ dans $\mathcal{C}$ tel que
chaque restriction $\mathcal{F}|_{\mathcal{C}/X_i}$ soit un
$\mathcal{O}_{X_i}$-module engendré par $r$ sections globales.
\end{enumerate}
\item Les conditions suivantes sont équivalentes :
\begin{enumerate}
\item $\mathcal{F}$ est de type fini ;
\item il existe un recouvrement $\{X_i \to X\}$ dans $\mathcal{C}$ tel que
chaque restriction $\mathcal{F}|_{\mathcal{C}/X_i}$ soit un
$\mathcal{O}_{X_i}$-module de type fini ;
\item il existe un recouvrement $\{X_i \to X\}$ dans $\mathcal{C}$ tel que
chaque restriction $\mathcal{F}|_{\mathcal{C}/X_i}$ soit un
$\mathcal{O}_{X_i}$-module engendré par un nombre fini de sections globales.
\end{enumerate}
\item Les conditions suivantes sont équivalentes :
\begin{enumerate}
\item $\mathcal{F}$ est quasi-cohérent ;
\item il existe un recouvrement $\{X_i \to X\}$ dans $\mathcal{C}$ tel que
chaque restriction $\mathcal{F}|_{\mathcal{C}/X_i}$ soit un
$\mathcal{O}_{X_i}$-module quasi-cohérent ;
\item il existe un recouvrement $\{X_i \to X\}$ dans $\mathcal{C}$ tel que
chaque restriction $\mathcal{F}|_{\mathcal{C}/X_i}$ soit un
$\mathcal{O}_{X_i}$-module admettant une présentation globale.
\end{enumerate}
\item Les conditions suivantes sont équivalentes :
\begin{enumerate}
\item $\mathcal{F}$ est de présentation finie ;
\item il existe un recouvrement $\{X_i \to X\}$ dans $\mathcal{C}$ tel que
chaque restriction $\mathcal{F}|_{\mathcal{C}/X_i}$ soit un
$\mathcal{O}_{X_i}$-module de présentation finie ;
\item il existe un recouvrement $\{X_i \to X\}$ dans $\mathcal{C}$ tel que
chaque restriction $\mathcal{F}|_{\mathcal{C}/X_i}$ soit un
$\mathcal{O}_{X_i}$-module admettant une présentation globale finie.
\end{enumerate}
\item Les conditions suivantes sont équivalentes :
\begin{enumerate}
\item $\mathcal{F}$ est cohérent ;
\item il existe un recouvrement $\{X_i \to X\}$ dans $\mathcal{C}$ tel que
chaque restriction $\mathcal{F}|_{\mathcal{C}/X_i}$ soit un
$\mathcal{O}_{X_i}$-module cohérent.
\end{enumerate}
\end{enumerate}
\end{lemma}

\begin{proof}
Dans chaque cas, (a) $\Rightarrow (b)$. Dans chacun des cas (1) -- (6),
la condition (b) implique la condition (c) par l'axiome (2) d'un site
(voir Sites, Définition \ref{sites-definition-site}) et par la définition
des propriétés locales des modules.
Supposons que $\{X_i \to X\}$ soit un recouvrement.
Pour tout objet $U$ de $\mathcal{C}$, on obtient alors le recouvrement induit
$\{X_i \times_X U \to U\}$. De plus, la propriété globale de
$\mathcal{F}|_{\mathcal{C}/X_i}$ dans (c) implique, d'après le
Lemme \ref{lemma-global-pullback}, la propriété globale correspondante de
$\mathcal{F}|_{\mathcal{C}/X_i \times_X U}$ ; le faisceau possède donc par
définition la propriété (a). Nous omettons la preuve de
(b) $\Rightarrow$ (a) dans le cas (7).
\end{proof}

\begin{lemma}
\label{lemma-local-pullback}
Soit
$(f, f^\sharp) :
(\Sh(\mathcal{C}), \mathcal{O}_\mathcal{C})
\to
(\Sh(\mathcal{D}), \mathcal{O}_\mathcal{D})$
soit un morphisme de topos annelés.
Soit $\mathcal{F}$ un $\mathcal{O}_\mathcal{D}$-module.
\begin{enumerate}
\item Si $\mathcal{F}$ est localement libre, alors $f^*\mathcal{F}$ est
localement libre.
\item Si $\mathcal{F}$ est localement libre de rang fini, alors
$f^*\mathcal{F}$ est localement libre de rang fini.
\item Si $\mathcal{F}$ est localement engendré par des sections, alors
$f^*\mathcal{F}$ est localement engendré par des sections.
\item Si $\mathcal{F}$ est localement engendré par $r$ sections, alors
$f^*\mathcal{F}$ est localement engendré par $r$ sections.
\item Si $\mathcal{F}$ est de type fini, alors $f^*\mathcal{F}$ est de type
fini.
\item Si $\mathcal{F}$ est quasi-cohérent, alors $f^*\mathcal{F}$ est
quasi-cohérent.
\item Si $\mathcal{F}$ est de présentation finie, alors $f^*\mathcal{F}$ est
de présentation finie.
\end{enumerate}
\end{lemma}

\begin{proof}
Selon la discussion de la Section \ref{section-intrinsic}, il suffit de
vérifier la conservation par image inverse pour un morphisme de sites annelés
$(f, f^\sharp) :
(\mathcal{C}, \mathcal{O}_\mathcal{C})
\to
(\mathcal{D}, \mathcal{O}_\mathcal{D})$
tel que $f$ soit donné par un foncteur exact à gauche et continu
$u : \mathcal{D} \to \mathcal{C}$ entre des sites admettant toutes les
limites finies.
Soit $\mathcal{G}$ un faisceau de $\mathcal{O}_\mathcal{D}$-modules possédant
l'une des propriétés (1) -- (6) de la
Définition \ref{definition-site-local}.
Nous savons que $\mathcal{D}$ admet un objet final $Y$ et que $X = u(Y)$ est
un objet final de $\mathcal{C}$. Par hypothèse, il existe un recouvrement
$\{Y_i \to Y\}$ tel que $\mathcal{G}|_{\mathcal{D}/Y_i}$ possède la
propriété globale correspondante. Posons $X_i = u(Y_i)$ ; alors
$\{X_i \to X\}$ est un recouvrement dans $\mathcal{C}$.
On obtient un diagramme commutatif de morphismes de sites annelés
$$
\xymatrix{
(\mathcal{C}/X_i, \mathcal{O}_\mathcal{C}|_{X_i}) \ar[r] \ar[d] &
(\mathcal{C}, \mathcal{O}_\mathcal{C}) \ar[d] \\
(\mathcal{D}/Y_i, \mathcal{O}_\mathcal{D}|_{Y_i}) \ar[r] &
(\mathcal{D}, \mathcal{O}_\mathcal{D})
}
$$
d'après Sites, Lemme \ref{sites-lemma-localize-morphism-strong}.
Le Lemme \ref{lemma-global-pullback} montre donc que
$f^*\mathcal{G}|_{X_i}$ possède la propriété globale correspondante.
Le Lemme \ref{lemma-local-final-object} permet alors de conclure que
$\mathcal{G}$ possède la propriété locale considérée.
\end{proof}







\section{Résultats élémentaires sur les propriétés locales des modules}
\label{section-basics}

\noindent
Voici quelques lemmes élémentaires relatifs aux définitions précédentes.

\begin{lemma}
\label{lemma-cokernel-finite-finite-presentation}
Soit $(\mathcal{C},\mathcal{O})$ un site annelé. Soit $\mathcal{F}$ un
$\mathcal{O}$-module de présentation finie. Soit
$\varphi : \mathcal{G} \to \mathcal{F}$ un morphisme de
$\mathcal{O}$-modules. Si $\mathcal{G}$ est de type fini, alors
$\Coker(\varphi)$ est de présentation finie.
\end{lemma}

\begin{proof}
La preuve est omise. Indication : voir Modules, Lemme
\ref{modules-lemma-cokernel-finite-finite-presentation}.
\end{proof}

\begin{lemma}
\label{lemma-kernel-surjection-finite-onto-finite-presentation}
Soit $(\mathcal{C}, \mathcal{O})$ un site annelé.
Soit $\theta : \mathcal{G} \to \mathcal{F}$ un morphisme surjectif de
$\mathcal{O}$-modules, où $\mathcal{F}$ est de présentation finie et
$\mathcal{G}$ de type fini. Alors $\Ker(\theta)$ est de type fini.
\end{lemma}

\begin{proof}
La preuve est omise. Indication : voir Modules, Lemme
\ref{modules-lemma-kernel-surjection-finite-free-onto-finite-presentation}.
\end{proof}

\begin{lemma}
\label{lemma-chaotic-quasi-coherent}
Soit $\mathcal{C}$ une catégorie considérée comme un site muni de la
topologie chaotique, voir Sites, Exemple \ref{sites-example-indiscrete}.
Soient $\mathcal{O}$ un faisceau d'anneaux sur $\mathcal{C}$ et
$\mathcal{F}$ un faisceau de $\mathcal{O}$-modules. Alors $\mathcal{F}$ est
quasi-cohérent si et seulement si, pour tout $U \to V$ dans $\mathcal{C}$,
l'application canonique
$$
\mathcal{F}(V) \otimes_{\mathcal{O}(V)} \mathcal{O}(U)
\longrightarrow
\mathcal{F}(U)
$$
est un isomorphisme.
\end{lemma}

\begin{proof}
Supposons $\mathcal{F}$ quasi-cohérent et soit $U \to V$ un morphisme de
$\mathcal{C}$. Comme tout recouvrement de $V$ est donné par un isomorphisme,
la Définition \ref{definition-site-local} donne une présentation
$$
\bigoplus\nolimits_{j \in J} \mathcal{O}_V \longrightarrow
\bigoplus\nolimits_{i \in I} \mathcal{O}_V \longrightarrow
\mathcal{F}|_{\mathcal{C}/V} \longrightarrow 0
$$
La topologie sur $\mathcal{C}$ étant chaotique, la prise des sections sur
tout objet de $\mathcal{C}$ est exacte. On obtient donc une présentation
$$
\bigoplus\nolimits_{j \in J} \mathcal{O}(V) \longrightarrow
\bigoplus\nolimits_{i \in I} \mathcal{O}(V) \longrightarrow
\mathcal{F}(V) \longrightarrow 0
$$
de $\mathcal{F}(V)$ comme $\mathcal{O}(V)$-module, et de même pour
$\mathcal{F}(U)$. Cela montre aisément que l'application affichée dans
l'énoncé du lemme est un isomorphisme.

\medskip\noindent
Supposons que l'application affichée dans l'énoncé du lemme soit un
isomorphisme pour tout morphisme $U \to V$ de $\mathcal{C}$. Fixons $V$ et
choisissons une présentation
$$
\bigoplus\nolimits_{j \in J} \mathcal{O}(V) \longrightarrow
\bigoplus\nolimits_{i \in I} \mathcal{O}(V) \longrightarrow
\mathcal{F}(V) \longrightarrow 0
$$
de $\mathcal{F}(V)$ comme $\mathcal{O}(V)$-module. L'hypothèse sur
$\mathcal{F}$ signifie exactement que la suite correspondante
$$
\bigoplus\nolimits_{j \in J} \mathcal{O}_V \longrightarrow
\bigoplus\nolimits_{i \in I} \mathcal{O}_V \longrightarrow
\mathcal{F}|_{\mathcal{C}/V} \longrightarrow 0
$$
est exacte, d'où l'on conclut que $\mathcal{F}$ est quasi-cohérent.
\end{proof}

\begin{lemma}
\label{lemma-chaotic-flat-quasi-coherent}
Soit $\mathcal{C}$ une catégorie considérée comme un site muni de la
topologie chaotique, voir Sites, Exemple \ref{sites-example-indiscrete}.
Soit $\mathcal{O}$ un faisceau d'anneaux sur $\mathcal{C}$.
Supposons que, pour tout $U \to V$ dans $\mathcal{C}$, l'application de
restriction $\mathcal{O}(V) \to \mathcal{O}(U)$ soit un homomorphisme plat
d'anneaux. Alors la catégorie des $\mathcal{O}$-modules quasi-cohérents est
une sous-catégorie de Serre faible de $\textit{Mod}(\mathcal{O})$.
\end{lemma}

\begin{proof}
Nous allons vérifier la définition d'une sous-catégorie de Serre faible,
voir Homology, Définition \ref{homology-definition-serre-subcategory}.
Pour cela, nous utiliserons la caractérisation des modules quasi-cohérents
donnée au Lemme \ref{lemma-chaotic-quasi-coherent}.
Considérons une suite exacte
$$
\mathcal{F}_0 \to \mathcal{F}_1 \to
\mathcal{F}_2 \to \mathcal{F}_3 \to
\mathcal{F}_4
$$
dans $\textit{Mod}(\mathcal{O})$, où $\mathcal{F}_0$, $\mathcal{F}_1$,
$\mathcal{F}_3$ et $\mathcal{F}_4$ sont quasi-cohérents. Soit $U \to V$ un
morphisme de $\mathcal{C}$ et considérons le diagramme commutatif
$$
\xymatrix@C=1em{
{\scriptstyle \mathcal{F}_0(V) \otimes_{\mathcal{O}(V)} \mathcal{O}(U)} \ar[r] \ar[d] &
{\scriptstyle \mathcal{F}_1(V) \otimes_{\mathcal{O}(V)} \mathcal{O}(U)} \ar[r] \ar[d] &
{\scriptstyle \mathcal{F}_2(V) \otimes_{\mathcal{O}(V)} \mathcal{O}(U)} \ar[r] \ar[d] &
{\scriptstyle \mathcal{F}_3(V) \otimes_{\mathcal{O}(V)} \mathcal{O}(U)} \ar[r] \ar[d] &
{\scriptstyle \mathcal{F}_4(V) \otimes_{\mathcal{O}(V)} \mathcal{O}(U)} \ar[d] \\
{\scriptstyle \mathcal{F}_0(U)} \ar[r] &
{\scriptstyle \mathcal{F}_1(U)} \ar[r] &
{\scriptstyle \mathcal{F}_2(U)} \ar[r] &
{\scriptstyle \mathcal{F}_3(U)} \ar[r] &
{\scriptstyle \mathcal{F}_4(U)}
}
$$
Par hypothèse, les flèches verticales d'indices $0$, $1$, $3$, $4$ sont des
isomorphismes. Comme la topologie sur $\mathcal{C}$ est chaotique, la prise
des sections sur un objet de $\mathcal{C}$ est exacte ; la ligne inférieure
est donc exacte. Comme $\mathcal{O}(V) \to \mathcal{O}(U)$ est plat, la ligne
supérieure est elle aussi exacte. Le lemme des $5$ permet alors de conclure
que la flèche du milieu est un isomorphisme
(Homology, Lemme \ref{homology-lemma-five-lemma}).
\end{proof}






\section{Immersions fermées de topos annelés}
\label{section-closed-immersion}

\noindent
Quand doit-on déclarer qu'un morphisme de topos annelés
$i : (\Sh(\mathcal{C}), \mathcal{O}) \to (\Sh(\mathcal{D}), \mathcal{O}')$
est une immersion fermée ? Par analogie avec la discussion de
Modules, Section \ref{modules-section-closed-immersion}, il semble naturel de
supposer au moins que :
\begin{enumerate}
\item le foncteur $i$ est une immersion fermée de topos
(Sites, Définition \ref{sites-definition-immersion-topoi}).
\item l'application associée $\mathcal{O}' \to i_*\mathcal{O}$ est surjective.
\end{enumerate}
Ces conditions impliquent déjà plusieurs résultats satisfaisants, que nous
examinons dans cette section. Il paraît toutefois prudent de ne pas définir
formellement la notion d'immersion fermée de topos annelés, car plusieurs
définitions différentes seraient possibles.

\begin{lemma}
\label{lemma-i-star-equivalence}
Soit $i : (\Sh(\mathcal{C}), \mathcal{O}) \to (\Sh(\mathcal{D}), \mathcal{O}')$
un morphisme de topos annelés. Supposons que $i$ soit une immersion fermée
de topos et que $i^\sharp : \mathcal{O}' \to i_*\mathcal{O}$ soit surjectif.
Notons $\mathcal{I} \subset \mathcal{O}'$ le noyau de $i^\sharp$.
Le foncteur
$$
i_* :
\textit{Mod}(\mathcal{O})
\longrightarrow
\textit{Mod}(\mathcal{O}')
$$
est exact et pleinement fidèle ; son image essentielle est formée des
$\mathcal{O}'$-modules $\mathcal{G}$ tels que $\mathcal{I}\mathcal{G} = 0$.
\end{lemma}

\begin{proof}
Le Lemme \ref{lemma-exactness} et Sites, Lemme
\ref{sites-lemma-closed-immersion} montrent que $i_*$ est exact.
Puisque $i_*$ est pleinement fidèle sur les faisceaux d'ensembles et que
$i^\sharp$ est surjectif, $i_*$ est pleinement fidèle comme foncteur
$\textit{Mod}(\mathcal{O}) \to \textit{Mod}(\mathcal{O}')$.
En effet, supposons que
$\alpha : i_*\mathcal{F}_1 \to i_*\mathcal{F}_2$ soit un morphisme de
$\mathcal{O}'$-modules. La pleine fidélité de $i_*$ fournit une application
$\beta : \mathcal{F}_1 \to \mathcal{F}_2$ de faisceaux d'ensembles.
Pour montrer que $\beta$ est un morphisme de modules, il faut vérifier que
$$
\xymatrix{
\mathcal{O} \times \mathcal{F}_1 \ar[r] \ar[d] &
\mathcal{F}_1 \ar[d] \\
\mathcal{O} \times \mathcal{F}_2 \ar[r] &
\mathcal{F}_2
}
$$
est commutatif. Il suffit de le vérifier après application de $i_*$.
Considérons
$$
\xymatrix{
\mathcal{O}' \times i_*\mathcal{F}_1 \ar[r] \ar[d] &
i_*\mathcal{O} \times i_*\mathcal{F}_1 \ar[r] \ar[d] &
i_*\mathcal{F}_1 \ar[d] \\
\mathcal{O}' \times i_*\mathcal{F}_2 \ar[r] &
i_*\mathcal{O} \times i_*\mathcal{F}_2 \ar[r] &
i_*\mathcal{F}_2
}
$$
Le rectangle extérieur est commutatif. La conclusion résulte de la
surjectivité de $i^\sharp$.

\medskip\noindent
Pour achever la preuve, il reste à établir l'assertion sur l'image
essentielle de $i_*$. Il est clair que $i_*\mathcal{F}$ est annulé par
$\mathcal{I}$ pour tout $\mathcal{O}$-module $\mathcal{F}$. Réciproquement,
soit $\mathcal{G}$ un $\mathcal{O}'$-module tel que
$\mathcal{I}\mathcal{G} = 0$. Par définition d'un sous-topos fermé, il existe
un sous-faisceau $\mathcal{U}$ de l'objet final de $\mathcal{D}$ tel que
l'image essentielle de $i_*$ sur les faisceaux d'ensembles soit la classe des
faisceaux d'ensembles $\mathcal{H}$ pour lesquels
$\mathcal{H} \times \mathcal{U} \to \mathcal{U}$ est un isomorphisme.
En particulier, $i_*\mathcal{O} \times \mathcal{U} = \mathcal{U}$, ce qui
implique $\mathcal{I} \times \mathcal{U} = \mathcal{O} \times \mathcal{U}$.
Notre module $\mathcal{G}$ satisfait donc
$\mathcal{G} \times \mathcal{U} = \{0\} \times \mathcal{U} = \mathcal{U}$
(car le module nul est isomorphe à l'objet final de la catégorie des
faisceaux d'ensembles). Il existe donc un faisceau d'ensembles $\mathcal{F}$
sur $\mathcal{C}$ tel que $i_*\mathcal{F} = \mathcal{G}$. Comme $i_*$ est
pleinement fidèle sur les faisceaux d'ensembles, pour définir l'addition
$\mathcal{F} \times \mathcal{F} \to \mathcal{F}$ et la multiplication
$\mathcal{O} \times \mathcal{F} \to \mathcal{F}$, il suffit d'utiliser
l'addition
$$
\mathcal{G} \times \mathcal{G} \longrightarrow \mathcal{G}
$$
(dont on dispose puisque $\mathcal{G}$ est un $\mathcal{O}'$-module) et la
multiplication
$$
i_*\mathcal{O} \times \mathcal{G} \to \mathcal{G}
$$
dont on dispose car la multiplication par $\mathcal{O}'$ annule
$\mathcal{I}$ par hypothèse et
$i_*\mathcal{O} = \mathcal{O}'/\mathcal{I}$. Par construction,
$\mathcal{G}$ est isomorphe à l'image directe du $\mathcal{O}$-module
$\mathcal{F}$ ainsi construit.
\end{proof}












\section{Produit tensoriel}
\label{section-tensor-product}

\noindent
Dans les Sections \ref{section-presheaves-modules} et
\ref{section-sheafification-presheaves-modules}, nous avons défini le
foncteur de changement d'anneaux au moyen d'un produit tensoriel.
Cette construction permet aussi de définir le produit tensoriel de
préfaisceaux de $\mathcal{O}$-modules. Plus précisément, supposons que
$\mathcal{C}$ soit une catégorie, $\mathcal{O}$ un préfaisceau d'anneaux et
$\mathcal{F}$, $\mathcal{G}$ des préfaisceaux de $\mathcal{O}$-modules.
On définit alors $\mathcal{F} \otimes_{p, \mathcal{O}} \mathcal{G}$ comme le
préfaisceau
$$
U
\longmapsto
(\mathcal{F} \otimes_{p, \mathcal{O}} \mathcal{G})(U)
=
\mathcal{F}(U) \otimes_{\mathcal{O}(U)} \mathcal{G}(U)
$$
Si $\mathcal{C}$ est un site, $\mathcal{O}$ un faisceau d'anneaux et
$\mathcal{F}$, $\mathcal{G}$ des faisceaux de $\mathcal{O}$-modules, on pose
$$
\mathcal{F} \otimes_\mathcal{O} \mathcal{G}
=
(\mathcal{F} \otimes_{p, \mathcal{O}} \mathcal{G})^\#
$$
qui est le faisceau de $\mathcal{O}$-modules associé au préfaisceau
$\mathcal{F} \otimes_{p, \mathcal{O}} \mathcal{G}$.

\medskip\noindent
Voici quelques formules que nous utiliserons plus loin sans autre mention :
$$
(\mathcal{F}
\otimes_{p, \mathcal{O}} \mathcal{G})
\otimes_{p, \mathcal{O}} \mathcal{H}
=
\mathcal{F}
\otimes_{p, \mathcal{O}} (\mathcal{G}
\otimes_{p, \mathcal{O}} \mathcal{H}),
$$
et de même pour les faisceaux.
Si $\mathcal{O}_1 \to \mathcal{O}_2$ est un morphisme de préfaisceaux
d'anneaux, alors
$$
(\mathcal{F} \otimes_{p, \mathcal{O}_1} \mathcal{G})
\otimes_{p, \mathcal{O}_1} \mathcal{O}_2 =
(\mathcal{F} \otimes_{p, \mathcal{O}_1} \mathcal{O}_2)
\otimes_{p, \mathcal{O}_2}
(\mathcal{G} \otimes_{p, \mathcal{O}_1} \mathcal{O}_2),
$$
et de même pour les faisceaux.
Ces formules résultent de leurs analogues algébriques et de la
faisceautisation.

\begin{lemma}
\label{lemma-sheafification-tensor}
Soit $\mathcal{C}$ un site. Soit $\mathcal{O}$ un préfaisceau d'anneaux.
Soient $\mathcal{F}$ et $\mathcal{G}$ des préfaisceaux de
$\mathcal{O}$-modules. Alors
$\mathcal{F}^\# \otimes_{\mathcal{O}^\#} \mathcal{G}^\#$
est égal à
$(\mathcal{F} \otimes_{p, \mathcal{O}} \mathcal{G})^\#$.
\end{lemma}

\begin{proof}
La preuve est omise. Indication : utiliser la caractérisation du produit
tensoriel par les applications bilinéaires donnée ci-dessous, ainsi que la
propriété universelle de la faisceautisation.
\end{proof}

\noindent
Soit $\mathcal{C}$ un site, soit $\mathcal{O}$ un faisceau d'anneaux et soient
$\mathcal{F}$, $\mathcal{G}$, $\mathcal{H}$ des faisceaux de
$\mathcal{O}$-modules. On pose
$$
\text{Bilin}_\mathcal{O}(\mathcal{F} \times \mathcal{G}, \mathcal{H})
=
\{\varphi \in
\Mor_{\Sh(\mathcal{C})}(
\mathcal{F} \times \mathcal{G}, \mathcal{H}) \mid
\varphi \text{ est }\mathcal{O}\text{-bilinéaire}\}.
$$
Avec cette définition, on a
$$
\Hom_\mathcal{O}
(\mathcal{F} \otimes_\mathcal{O} \mathcal{G}, \mathcal{H})
=
\text{Bilin}_\mathcal{O}(\mathcal{F} \times \mathcal{G}, \mathcal{H}).
$$
Autrement dit, $\mathcal{F} \otimes_\mathcal{O} \mathcal{G}$ représente le
foncteur qui associe à $\mathcal{H}$ l'ensemble des applications bilinéaires
$\mathcal{F} \times \mathcal{G} \to \mathcal{H}$.
En particulier, puisque la notion d'application bilinéaire a un sens pour
une paire de modules sur un topos annelé, le produit tensoriel de faisceaux de
modules est intrinsèque au topos (comparer la discussion de la
Section \ref{section-intrinsic}). On a en fait le résultat suivant.

\begin{lemma}
\label{lemma-tensor-product-pullback}
Soit $f : (\Sh(\mathcal{C}), \mathcal{O}_\mathcal{C})
\to (\Sh(\mathcal{D}), \mathcal{O}_\mathcal{D})$ un morphisme de topos
annelés. Soient $\mathcal{F}$ et $\mathcal{G}$ des
$\mathcal{O}_\mathcal{D}$-modules. Alors
$f^*(\mathcal{F} \otimes_{\mathcal{O}_\mathcal{D}} \mathcal{G})
= f^*\mathcal{F} \otimes_{\mathcal{O}_\mathcal{C}} f^*\mathcal{G}$
fonctoriellement en $\mathcal{F}$ et $\mathcal{G}$.
\end{lemma}

\begin{proof}
Pour tout faisceau $\mathcal{H}$ de $\mathcal{O}_\mathcal{C}$-modules, on a
\begin{align*}
\Hom_{\mathcal{O}_\mathcal{C}}(
f^*(\mathcal{F} \otimes_\mathcal{O} \mathcal{G}), \mathcal{H})
& =
\Hom_{\mathcal{O}_\mathcal{D}}(
\mathcal{F} \otimes_\mathcal{O} \mathcal{G}, f_*\mathcal{H}) \\
& =
\text{Bilin}_{\mathcal{O}_\mathcal{D}}(
\mathcal{F} \times \mathcal{G}, f_*\mathcal{H}) \\
& =
\text{Bilin}_{f^{-1}\mathcal{O}_\mathcal{D}}(
f^{-1}\mathcal{F} \times f^{-1}\mathcal{G}, \mathcal{H}) \\
& =
\Hom_{f^{-1}\mathcal{O}_\mathcal{D}}(
f^{-1}\mathcal{F} \otimes_{f^{-1}\mathcal{O}_\mathcal{D}} f^{-1}\mathcal{G},
\mathcal{H}) \\
& =
\Hom_{\mathcal{O}_\mathcal{C}}(
f^*\mathcal{F} \otimes_{f^*\mathcal{O}_\mathcal{D}} f^*\mathcal{G},
\mathcal{H})
\end{align*}
Le symbole « $=$ » intéressant de cette chaîne correspond à la
troisième égalité. Elle résulte
directement de la définition et de l'adjonction entre $f_*$ et $f^{-1}$
(discutée dans les sections précédentes).
\end{proof}

\begin{lemma}
\label{lemma-tensor-product-permanence}
Soit $(\mathcal{C}, \mathcal{O})$ un site annelé.
Soient $\mathcal{F}$ et $\mathcal{G}$ des faisceaux de $\mathcal{O}$-modules.
\begin{enumerate}
\item Si $\mathcal{F}$ et $\mathcal{G}$ sont localement libres, alors
$\mathcal{F} \otimes_\mathcal{O} \mathcal{G}$ l'est aussi.
\item Si $\mathcal{F}$ et $\mathcal{G}$ sont localement libres de rang fini,
alors $\mathcal{F} \otimes_\mathcal{O} \mathcal{G}$ l'est aussi.
\item Si $\mathcal{F}$ et $\mathcal{G}$ sont localement engendrés par des
sections, alors $\mathcal{F} \otimes_\mathcal{O} \mathcal{G}$ l'est aussi.
\item Si $\mathcal{F}$ et $\mathcal{G}$ sont de type fini, alors
$\mathcal{F} \otimes_\mathcal{O} \mathcal{G}$ l'est aussi.
\item Si $\mathcal{F}$ et $\mathcal{G}$ sont quasi-cohérents, alors
$\mathcal{F} \otimes_\mathcal{O} \mathcal{G}$ l'est aussi.
\item Si $\mathcal{F}$ et $\mathcal{G}$ sont de présentation finie, alors
$\mathcal{F} \otimes_\mathcal{O} \mathcal{G}$ l'est aussi.
\item Si $\mathcal{F}$ est de présentation finie et $\mathcal{G}$ cohérent,
alors $\mathcal{F} \otimes_\mathcal{O} \mathcal{G}$ est cohérent.
\item Si $\mathcal{F}$ et $\mathcal{G}$ sont cohérents, alors
$\mathcal{F} \otimes_\mathcal{O} \mathcal{G}$ l'est aussi.
\end{enumerate}
\end{lemma}

\begin{proof}
La preuve est omise. Indication : comparer avec
Sheaves of Modules, Lemme \ref{modules-lemma-tensor-product-permanence}.
\end{proof}



\section{Hom interne}
\label{section-internal-hom}

\noindent
Soit $\mathcal{C}$ une catégorie et soit $\mathcal{O}$ un préfaisceau
d'anneaux. Soient $\mathcal{F}$ et $\mathcal{G}$ des préfaisceaux de
$\mathcal{O}$-modules. Considérons la règle
$$
U \longmapsto \Hom_{\mathcal{O}_U}(\mathcal{F}|_U, \mathcal{G}|_U).
$$
Pour $\varphi : V \to U$ dans $\mathcal{C}$, on définit une application de
restriction
$$
\Hom_{\mathcal{O}_U}(\mathcal{F}|_U, \mathcal{G}|_U)
\longrightarrow
\Hom_{\mathcal{O}_V}(\mathcal{F}|_V, \mathcal{G}|_V)
$$
en restreignant par le morphisme de relocalisation
$j : \mathcal{C}/V \to \mathcal{C}/U$, voir Sites, Lemme
\ref{sites-lemma-relocalize}. Cela définit donc un préfaisceau
$\SheafHom_\mathcal{O}(\mathcal{F}, \mathcal{G})$.
De plus, étant donnés
$\varphi \in \Hom_{\mathcal{O}|_U}(\mathcal{F}|_U, \mathcal{G}|_U)$ et une
section $f \in \mathcal{O}(U)$, on peut définir
$f\varphi \in \Hom_{\mathcal{O}|_U}(\mathcal{F}|_U, \mathcal{G}|_U)$ soit
en précomposant par la multiplication par $f$ sur $\mathcal{F}|_U$, soit en
postcomposant par la multiplication par $f$ sur $\mathcal{G}|_U$ (le résultat
est le même). On obtient donc en fait un préfaisceau de
$\mathcal{O}$-modules. Il existe un morphisme canonique d'évaluation
$$
\mathcal{F}
\otimes_{p, \mathcal{O}}
\SheafHom_\mathcal{O}(\mathcal{F}, \mathcal{G})
\longrightarrow
\mathcal{G}.
$$

\begin{lemma}
\label{lemma-internal-hom}
Si $\mathcal{C}$ est un site, $\mathcal{O}$ un faisceau d'anneaux,
$\mathcal{F}$ un préfaisceau de $\mathcal{O}$-modules et
$\mathcal{G}$ un faisceau de $\mathcal{O}$-modules, alors
$\SheafHom_\mathcal{O}(\mathcal{F}, \mathcal{G})$
est un faisceau de $\mathcal{O}$-modules.
\end{lemma}

\begin{proof}
La preuve est omise. Indications : noter d'abord que
$\SheafHom_\mathcal{O}(\mathcal{F}, \mathcal{G})
= \SheafHom_\mathcal{O}(\mathcal{F}^\#, \mathcal{G})$, ce qui ramène
au cas où $\mathcal{F}$ et $\mathcal{G}$ sont tous deux des
faisceaux. Le résultat pour les faisceaux d'ensembles est Sites, Lemme
\ref{sites-lemma-glue-maps}.
\end{proof}

\begin{lemma}
\label{lemma-internal-hom-restriction}
Soit $(\mathcal{C}, \mathcal{O})$ un site annelé.
Soient $\mathcal{F}, \mathcal{G}$ des faisceaux de $\mathcal{O}$-modules.
Alors la formation de $\SheafHom_\mathcal{O}(\mathcal{F}, \mathcal{G})$
commute à la restriction à $U$, pour $U \in \Ob(\mathcal{C})$.
\end{lemma}

\begin{proof}
Cela résulte immédiatement de la définition.
\end{proof}

\begin{remark}
\label{remark-pullback-internal-hom}
Soit $f : (\mathcal{C}, \mathcal{O}_\mathcal{C}) \to
(\mathcal{D}, \mathcal{O}_\mathcal{D})$ un morphisme de sites annelés.
Soient $\mathcal{F}, \mathcal{G}$ des faisceaux de
$\mathcal{O}_\mathcal{D}$-modules. Il existe une application canonique
$$
f^*\SheafHom_{\mathcal{O}_\mathcal{D}}(\mathcal{F}, \mathcal{G})
\longrightarrow
\SheafHom_{\mathcal{O}_\mathcal{C}}(f^*\mathcal{F}, f^*\mathcal{G})
$$
Cette application est en effet adjointe à l'application
$$
\SheafHom_{\mathcal{O}_\mathcal{D}}(\mathcal{F}, \mathcal{G})
\longrightarrow
f_*\SheafHom_{\mathcal{O}_\mathcal{C}}(f^*\mathcal{F}, f^*\mathcal{G})
$$
définie comme suit. Supposons que $f$ soit donné par le foncteur continu
$u : \mathcal{D} \to \mathcal{C}$. Sur les sections au-dessus de
$V \in \Ob(\mathcal{D})$, on utilise l'application
\begin{align*}
\Gamma(V, \SheafHom_{\mathcal{O}_\mathcal{D}}(\mathcal{F}, \mathcal{G}))
& =
\Hom_{\mathcal{O}_V}(\mathcal{F}|_V, \mathcal{G}|_V) \\
& \longrightarrow
\Hom_{\mathcal{O}_{u(V)}}(f^*\mathcal{F}|_{u(V)}, f^*\mathcal{G}|_{u(V)}) \\
& =
\Gamma(u(V), \SheafHom_{\mathcal{O}_\mathcal{C}}(f^*\mathcal{F},
f^*\mathcal{G})) \\
& =
\Gamma(V, f_*\SheafHom_{\mathcal{O}_\mathcal{C}}(f^*\mathcal{F},
f^*\mathcal{G}))
\end{align*}
où, pour la flèche, on utilise l'image inverse par le morphisme
$(\mathcal{C}/u(V), \mathcal{O}_{u(V)}) \to
(\mathcal{D}/V, \mathcal{O}_V)$ induit par $f$.
\end{remark}

\noindent
Dans la situation du Lemme \ref{lemma-internal-hom}, le morphisme
d'évaluation se factorise par le produit tensoriel de faisceaux de
modules
$$
\mathcal{F}
\otimes_\mathcal{O}
\SheafHom_\mathcal{O}(\mathcal{F}, \mathcal{G})
\longrightarrow
\mathcal{G}.
$$

\begin{lemma}
\label{lemma-internal-hom-commute-limits}
Hom interne et (co)limites.
Soit $\mathcal{C}$ une catégorie et soit $\mathcal{O}$ un préfaisceau
d'anneaux.
\begin{enumerate}
\item Pour tout préfaisceau de $\mathcal{O}$-modules $\mathcal{F}$, le foncteur
$$
\textit{PMod}(\mathcal{O}) \longrightarrow \textit{PMod}(\mathcal{O})
, \quad
\mathcal{G} \longmapsto \SheafHom_\mathcal{O}(\mathcal{F}, \mathcal{G})
$$
commute aux limites arbitraires.
\item Pour tout préfaisceau de $\mathcal{O}$-modules $\mathcal{G}$, le foncteur
$$
\textit{PMod}(\mathcal{O}) \longrightarrow \textit{PMod}(\mathcal{O})^{opp}
, \quad
\mathcal{F} \longmapsto \SheafHom_\mathcal{O}(\mathcal{F}, \mathcal{G})
$$
commute aux colimites arbitraires ; en formule,
$$
\SheafHom_\mathcal{O}(\colim_i \mathcal{F}_i, \mathcal{G})
=
\lim_i \SheafHom_\mathcal{O}(\mathcal{F}_i, \mathcal{G}).
$$
\end{enumerate}
Supposons que $\mathcal{C}$ soit un site et $\mathcal{O}$ un faisceau d'anneaux.
\begin{enumerate}
\item[(3)] Pour tout faisceau de $\mathcal{O}$-modules $\mathcal{F}$, le foncteur
$$
\textit{Mod}(\mathcal{O}) \longrightarrow \textit{Mod}(\mathcal{O})
, \quad
\mathcal{G} \longmapsto \SheafHom_\mathcal{O}(\mathcal{F}, \mathcal{G})
$$
commute aux limites arbitraires.
\item[(4)] Pour tout faisceau de $\mathcal{O}$-modules $\mathcal{G}$, le foncteur
$$
\textit{Mod}(\mathcal{O}) \longrightarrow \textit{Mod}(\mathcal{O})^{opp}
, \quad
\mathcal{F} \longmapsto \SheafHom_\mathcal{O}(\mathcal{F}, \mathcal{G})
$$
commute aux colimites arbitraires ; en formule,
$$
\SheafHom_\mathcal{O}(\colim_i \mathcal{F}_i, \mathcal{G})
=
\lim_i \SheafHom_\mathcal{O}(\mathcal{F}_i, \mathcal{G}).
$$
\end{enumerate}
\end{lemma}

\begin{proof}
Soit $\mathcal{I} \to \textit{PMod}(\mathcal{O})$,
$i \mapsto \mathcal{G}_i$, un diagramme. Soit $U$ un objet de la catégorie
$\mathcal{C}$. Comme $j_U^*$ est à la fois adjoint à gauche et à droite,
on a
$\lim_i j_U^*\mathcal{G}_i = j_U^* \lim_i \mathcal{G}_i$.
Par conséquent,
\begin{align*}
\SheafHom_\mathcal{O}(\mathcal{F}, \lim_i \mathcal{G}_i)(U)
& =
\Hom_{\mathcal{O}_U}(\mathcal{F}|_U, \lim_i \mathcal{G}_i|_U) \\
& =
\lim_i \Hom_{\mathcal{O}_U}(\mathcal{F}|_U, \mathcal{G}_i|_U) \\
& = \lim_i \SheafHom_\mathcal{O}(\mathcal{F}, \mathcal{G}_i)(U)
\end{align*}
par définition d'une limite. Cela prouve (1). La preuve de (2) est exactement
la même. L'assertion (3) résulte de (1), car la limite d'un diagramme de
faisceaux est la même que sa limite dans la catégorie des préfaisceaux.
Enfin, (4) résulte du fait que, dans la formule, on a
$$
\Mor_{\textit{Mod}(\mathcal{O})}(
\colim_i \mathcal{F}_i, \mathcal{G})
=
\Mor_{\textit{PMod}(\mathcal{O})}(
\colim^{PSh}_i \mathcal{F}_i, \mathcal{G})
$$
car la colimite $\colim_i \mathcal{F}_i$ est la faisceautisation de la
colimite $\colim^{PSh}_i \mathcal{F}_i$ dans
$\textit{PMod}(\mathcal{O})$. Ainsi (4) résulte de (2) (en utilisant à
nouveau la remarque précédente sur les limites).
\end{proof}

\begin{lemma}
\label{lemma-internal-hom-exact}
Soit $(\mathcal{C}, \mathcal{O})$ un site annelé.
Soient $\mathcal{F}$ et $\mathcal{G}$ des $\mathcal{O}$-modules.
\begin{enumerate}
\item Si $\mathcal{F}_2 \to \mathcal{F}_1 \to \mathcal{F} \to 0$ est une
suite exacte de $\mathcal{O}$-modules, alors
$$
0 \to
\SheafHom_\mathcal{O}(\mathcal{F}, \mathcal{G}) \to
\SheafHom_\mathcal{O}(\mathcal{F}_1, \mathcal{G}) \to
\SheafHom_\mathcal{O}(\mathcal{F}_2, \mathcal{G})
$$
est exacte.
\item Si $0 \to \mathcal{G} \to \mathcal{G}_1 \to \mathcal{G}_2$ est une
suite exacte de $\mathcal{O}$-modules, alors
$$
0 \to
\SheafHom_\mathcal{O}(\mathcal{F}, \mathcal{G}) \to
\SheafHom_\mathcal{O}(\mathcal{F}, \mathcal{G}_1) \to
\SheafHom_\mathcal{O}(\mathcal{F}, \mathcal{G}_2)
$$
est exacte.
\end{enumerate}
\end{lemma}

\begin{proof}
Cela résulte du Lemme \ref{lemma-internal-hom-commute-limits} et de
Homology, Lemme \ref{homology-lemma-exact-functor}.
\end{proof}

\begin{lemma}
\label{lemma-internal-hom-adjoint-tensor}
Soit $\mathcal{C}$ une catégorie. Soit $\mathcal{O}$ un préfaisceau
d'anneaux.
\begin{enumerate}
\item Soient $\mathcal{F}$, $\mathcal{G}$ et $\mathcal{H}$ des préfaisceaux
de $\mathcal{O}$-modules. Il existe un isomorphisme canonique
$$
\SheafHom_\mathcal{O}
(\mathcal{F} \otimes_{p, \mathcal{O}} \mathcal{G}, \mathcal{H})
\longrightarrow
\SheafHom_\mathcal{O}
(\mathcal{F}, \SheafHom_\mathcal{O}(\mathcal{G}, \mathcal{H}))
$$
fonctoriel en ses trois arguments (le Hom est celui des préfaisceaux dans les
trois termes). En particulier,
$$
\Mor_{\textit{PMod}(\mathcal{O})}(
\mathcal{F} \otimes_{p, \mathcal{O}} \mathcal{G}, \mathcal{H})
=
\Mor_{\textit{PMod}(\mathcal{O})}(
\mathcal{F}, \SheafHom_\mathcal{O}(\mathcal{G}, \mathcal{H}))
$$
\item
Supposons que $\mathcal{C}$ soit un site, que $\mathcal{O}$ soit un faisceau
d'anneaux et que $\mathcal{F}$, $\mathcal{G}$, $\mathcal{H}$ soient des
faisceaux de $\mathcal{O}$-modules. Il existe un isomorphisme canonique
$$
\SheafHom_\mathcal{O}
(\mathcal{F} \otimes_\mathcal{O} \mathcal{G}, \mathcal{H})
\longrightarrow
\SheafHom_\mathcal{O}
(\mathcal{F}, \SheafHom_\mathcal{O}(\mathcal{G}, \mathcal{H}))
$$
fonctoriel en ses trois arguments (le Hom est celui des faisceaux dans les
trois termes). En particulier,
$$
\Mor_{\textit{Mod}(\mathcal{O})}(
\mathcal{F} \otimes_\mathcal{O} \mathcal{G}, \mathcal{H})
=
\Mor_{\textit{Mod}(\mathcal{O})}(
\mathcal{F}, \SheafHom_\mathcal{O}(\mathcal{G}, \mathcal{H}))
$$
\end{enumerate}
\end{lemma}

\begin{proof}
Il s'agit de l'analogue de
Algebra, Lemme \ref{algebra-lemma-hom-from-tensor-product}.
La preuve est la même et est omise.
\end{proof}

\begin{lemma}
\label{lemma-tensor-commute-colimits}
Produit tensoriel et colimites.
Soit $\mathcal{C}$ une catégorie et soit $\mathcal{O}$ un préfaisceau
d'anneaux.
\begin{enumerate}
\item Pour tout préfaisceau de $\mathcal{O}$-modules $\mathcal{F}$, le foncteur
$$
\textit{PMod}(\mathcal{O}) \longrightarrow \textit{PMod}(\mathcal{O})
, \quad
\mathcal{G} \longmapsto \mathcal{F} \otimes_{p, \mathcal{O}} \mathcal{G}
$$
commute aux colimites arbitraires.
\item
Supposons que $\mathcal{C}$ soit un site et $\mathcal{O}$ un faisceau
d'anneaux. Pour tout faisceau de $\mathcal{O}$-modules $\mathcal{F}$, le
foncteur
$$
\textit{Mod}(\mathcal{O}) \longrightarrow \textit{Mod}(\mathcal{O})
, \quad
\mathcal{G} \longmapsto \mathcal{F} \otimes_\mathcal{O} \mathcal{G}
$$
commute aux colimites arbitraires.
\end{enumerate}
\end{lemma}

\begin{proof}
Cela tient à ce que le produit tensoriel est adjoint au Hom interne d'après
le Lemme \ref{lemma-internal-hom-adjoint-tensor}.
Voir Categories, Lemme \ref{categories-lemma-adjoint-exact}.
\end{proof}

\begin{lemma}
\label{lemma-adjoint-hom-restrict}
Soit $\mathcal{C}$ une catégorie, resp.\ un site.
Soit $\mathcal{O} \to \mathcal{O}'$ un morphisme de préfaisceaux, resp.\ de
faisceaux d'anneaux. Alors
$$
\Hom_\mathcal{O}(\mathcal{G}, \mathcal{F}) =
\Hom_{\mathcal{O}'}(\mathcal{G},
\SheafHom_\mathcal{O}(\mathcal{O}', \mathcal{F}))
$$
pour tout $\mathcal{O}'$-module $\mathcal{G}$ et tout $\mathcal{O}$-module
$\mathcal{F}$.
\end{lemma}

\begin{proof}
Il s'agit de l'analogue de
Algebra, Lemme \ref{algebra-lemma-adjoint-hom-restrict}.
La preuve est la même et est omise.
\end{proof}

\begin{lemma}
\label{lemma-j-shriek-and-tensor}
Soit $(\mathcal{C}, \mathcal{O})$ un site annelé.
Soit $U \in \Ob(\mathcal{C})$.
Pour $\mathcal{G}$ dans $\textit{Mod}(\mathcal{O}_U)$ et $\mathcal{F}$ dans
$\textit{Mod}(\mathcal{O})$, on a
$j_{U!}\mathcal{G} \otimes_\mathcal{O} \mathcal{F} =
j_{U!}(\mathcal{G} \otimes_{\mathcal{O}_U} \mathcal{F}|_U)$.
\end{lemma}

\begin{proof}
Soit $\mathcal{H}$ un objet de $\textit{Mod}(\mathcal{O})$. Alors
\begin{align*}
\Hom_\mathcal{O}(j_{U!}(\mathcal{G} \otimes_{\mathcal{O}_U} \mathcal{F}|_U),
\mathcal{H})
& =
\Hom_{\mathcal{O}_U}(\mathcal{G} \otimes_{\mathcal{O}_U} \mathcal{F}|_U,
\mathcal{H}|_U) \\
& =
\Hom_{\mathcal{O}_U}(\mathcal{G},
\SheafHom_{\mathcal{O}_U}(\mathcal{F}|_U, \mathcal{H}|_U)) \\
& =
\Hom_{\mathcal{O}_U}(\mathcal{G},
\SheafHom_\mathcal{O}(\mathcal{F}, \mathcal{H})|_U) \\
& =
\Hom_\mathcal{O}(j_{U!}\mathcal{G},
\SheafHom_\mathcal{O}(\mathcal{F}, \mathcal{H})) \\
& =
\Hom_\mathcal{O}(j_{U!}\mathcal{G} \otimes_\mathcal{O} \mathcal{F},
\mathcal{H})
\end{align*}
La première égalité vient de ce que $j_{U!}$ est adjoint à gauche de la
restriction des modules. La deuxième résulte du
Lemme \ref{lemma-internal-hom-adjoint-tensor}.
La troisième résulte du Lemme \ref{lemma-internal-hom-restriction}.
La quatrième vient à nouveau de ce que $j_{U!}$ est adjoint à gauche de la
restriction des modules. La cinquième résulte du
Lemme \ref{lemma-internal-hom-adjoint-tensor}.
Le lemme découle de ces égalités et du lemme de Yoneda.
\end{proof}

\begin{remark}
\label{remark-j-shriek-tensor}
Soit $\mathcal{C}$ un site. Soit $\mathcal{F}$ un faisceau d'ensembles sur
$\mathcal{C}$ et considérons le morphisme de localisation
$j : \Sh(\mathcal{C})/\mathcal{F} \to \Sh(\mathcal{C})$.
Voir Sites, Définition \ref{sites-definition-localize-topos}.
Nous affirmons que (a) $j_!\mathbf{Z} = \mathbf{Z}_\mathcal{F}^\#$ et que
(b) $j_!(j^{-1}\mathcal{H}) = j_!\mathbf{Z} \otimes_\mathbf{Z} \mathcal{H}$
pour tout faisceau abélien $\mathcal{H}$ sur $\mathcal{C}$.
Soit $\mathcal{G}$ un faisceau abélien sur $\mathcal{C}$.
L'assertion (a) résulte du lemme de Yoneda, car
$$
\Hom(j_!\mathbf{Z}, \mathcal{G}) =
\Hom(\mathbf{Z}, j^{-1}\mathcal{G}) =
\Hom(\mathbf{Z}_\mathcal{F}^\#, \mathcal{G})
$$
où la seconde égalité vaut parce que ses deux membres s'identifient à
l'ensemble des applications $\mathcal{F} \to \mathcal{G}$, considéré comme
groupe abélien. Pour (b), on utilise le lemme de Yoneda et
\begin{align*}
\Hom(j_!(j^{-1}\mathcal{H}), \mathcal{G})
& =
\Hom(j^{-1}\mathcal{H}, j^{-1}\mathcal{G}) \\
& =
\Hom(\mathbf{Z}, \SheafHom(j^{-1}\mathcal{H}, j^{-1}\mathcal{G})) \\
& =
\Hom(\mathbf{Z}, j^{-1}\SheafHom(\mathcal{H}, \mathcal{G})) \\
& =
\Hom(j_!\mathbf{Z}, \SheafHom(\mathcal{H}, \mathcal{G})) \\
& =
\Hom(j_!\mathbf{Z} \otimes_\mathbf{Z} \mathcal{H}, \mathcal{G})
\end{align*}
On utilise ici l'adjonction, le fait que la formation de $\SheafHom$ commute à
la localisation, et le Lemme \ref{lemma-internal-hom-adjoint-tensor}.
\end{remark}

\begin{lemma}
\label{lemma-hom-finite-presentation-colimit}
Soit $(\mathcal{C}, \mathcal{O})$ un site annelé.
Soit $\mathcal{F}$ un $\mathcal{O}$-module de présentation finie. Soit
$\mathcal{G} = \colim_{\lambda \in \Lambda} \mathcal{G}_\lambda$ une
colimite filtrante de $\mathcal{O}$-modules. Alors l'application canonique
$$
\colim_\lambda \SheafHom_\mathcal{O}(\mathcal{F}, \mathcal{G}_\lambda)
\longrightarrow
\SheafHom_\mathcal{O}(\mathcal{F}, \mathcal{G})
$$
est un isomorphisme.
\end{lemma}

\begin{proof}
Il suffit de montrer que la flèche devient un isomorphisme après restriction
à $U$, pour tout $U$ de $\mathcal{C}$. La formation des colimites de faisceaux de
modules et celle du Hom interne commutent toutes deux à la restriction à
$U$. Voir par exemple les Lemmes \ref{lemma-exactness-pushforward-pullback}
et \ref{lemma-internal-hom-restriction}. Fixons $U$. Étant donné un
recouvrement $\{U_i \to U\}_{i \in I}$, il suffit de prouver que la
restriction à chaque $U_i$ est un isomorphisme. On peut donc supposer que
$\mathcal{F}$ admet une présentation globale
$$
\bigoplus\nolimits_{j = 1, \ldots, m}
\mathcal{O}
\longrightarrow
\bigoplus\nolimits_{i = 1, \ldots, n}
\mathcal{O}
\to
\mathcal{F}
\to
0
$$
Le foncteur $\SheafHom_\mathcal{O}(-, -)$ commute aux sommes directes finies
en chacune des variables, et $\SheafHom_\mathcal{O}(\mathcal{O}, -)$ est le
foncteur identité. Avec le Lemme \ref{lemma-internal-hom-exact}, cela donne
une suite exacte
$$
0 \to
\SheafHom_\mathcal{O}(\mathcal{F}, \mathcal{G}) \to
\bigoplus\nolimits_{i = 1, \ldots, n} \mathcal{G} \to
\bigoplus\nolimits_{j = 1, \ldots, m} \mathcal{G}
$$
Comme les colimites filtrantes sont exactes dans
$\textit{Mod}(\mathcal{O})$ d'après le Lemme \ref{lemma-limits-colimits},
la ligne supérieure du diagramme commutatif suivant est aussi exacte
$$
\xymatrix{
0 \ar[r] &
\colim_\lambda \SheafHom_\mathcal{O}(\mathcal{F}, \mathcal{G}_\lambda)
\ar[r] \ar[d] &
\colim_\lambda \bigoplus\nolimits_{i = 1, \ldots, n} \mathcal{G}_\lambda
\ar[r] \ar[d] &
\colim_\lambda \bigoplus\nolimits_{j = 1, \ldots, m} \mathcal{G}_\lambda
\ar[d] \\
0 \ar[r] &
\SheafHom_\mathcal{O}(\mathcal{F}, \mathcal{G}) \ar[r] &
\bigoplus\nolimits_{i = 1, \ldots, n} \mathcal{G} \ar[r] &
\bigoplus\nolimits_{j = 1, \ldots, m} \mathcal{G}
}
$$
On conclut puisque les deux flèches verticales de droite sont des
isomorphismes.
\end{proof}

\begin{lemma}
\label{lemma-finite-presentation-quasi-compact-colimit}
Soit $(\mathcal{C}, \mathcal{O})$ un site annelé.
Soit $\mathcal{G} = \colim_{\lambda \in \Lambda} \mathcal{G}_\lambda$ une
colimite filtrante de $\mathcal{O}$-modules.
Soit $\mathcal{F}$ un $\mathcal{O}$-module de présentation finie.
Alors on a
$$
\colim_\lambda \Hom_\mathcal{O}(\mathcal{F}, \mathcal{G}_\lambda)
=
\Hom_\mathcal{O}(\mathcal{F}, \mathcal{G}).
$$
si les hypothèses de Sites, Lemme
\ref{sites-lemma-directed-colimits-global-sections}, partie (4), sont
satisfaites pour le site $\mathcal{C}$ ; voir Sites, Remarque
\ref{sites-remark-stronger-conditions}.
\end{lemma}

\begin{proof}
Posons $\mathcal{H} =
\SheafHom_\mathcal{O}(\mathcal{F}, \colim \mathcal{G}_\lambda)$
et $\mathcal{H}_\lambda =
\SheafHom_\mathcal{O}(\mathcal{F}, \mathcal{G}_\lambda)$.
Rappelons que
$$
\Hom_\mathcal{O}(\mathcal{F}, \mathcal{G}) = \Gamma(\mathcal{C}, \mathcal{H})
\quad\text{and}\quad
\Hom_\mathcal{O}(\mathcal{F}, \mathcal{G}_\lambda) =
\Gamma(\mathcal{C}, \mathcal{H}_\lambda)
$$
par construction. Le Lemme \ref{lemma-hom-finite-presentation-colimit} donne
$\mathcal{H} = \colim \mathcal{H}_\lambda$. Le lemme résulte donc de Sites,
Lemme \ref{sites-lemma-directed-colimits-global-sections}.
\end{proof}
\section{Modules plats}

\label{section-flat}

\noindent
On peut définir les modules plats exactement comme dans le cas des modules sur un anneau.

\begin{definition}

\label{definition-flat}

Soit $\mathcal{C}$ une catégorie et soit $\mathcal{O}$ un préfaisceau d'anneaux.

\begin{enumerate}

\item Un préfaisceau $\mathcal{F}$ de $\mathcal{O}$-modules est dit
{\it plat} si le foncteur
$$
\textit{PMod}(\mathcal{O})
\longrightarrow
\textit{PMod}(\mathcal{O}), \quad
\mathcal{G} \mapsto \mathcal{G} \otimes_{p, \mathcal{O}} \mathcal{F}
$$
est exact.
\item Un morphisme $\mathcal{O} \to \mathcal{O}'$ de préfaisceaux d'anneaux
est dit {\it plat} si $\mathcal{O}'$ est plat comme préfaisceau de
$\mathcal{O}$-modules.
\item Si $\mathcal{C}$ est un site, si $\mathcal{O}$ est un faisceau
d'anneaux et si $\mathcal{F}$ est un faisceau de $\mathcal{O}$-modules,
on dit que $\mathcal{F}$ est {\it plat} si le foncteur
$$
\textit{Mod}(\mathcal{O})
\longrightarrow
\textit{Mod}(\mathcal{O}), \quad
\mathcal{G} \mapsto \mathcal{G} \otimes_\mathcal{O} \mathcal{F}
$$
est exact.
\item Un morphisme $\mathcal{O} \to \mathcal{O}'$ de faisceaux d'anneaux
sur un site est dit {\it plat} si $\mathcal{O}'$ est plat comme faisceau de
$\mathcal{O}$-modules.

\end{enumerate}

\end{definition}

\noindent
La notion de module plat ou de morphisme plat d'anneaux est intrinsèque
(section \ref{section-intrinsic}).

\begin{lemma}

\label{lemma-flatness-presheaves}
Soit $\mathcal{C}$ une catégorie, soit $\mathcal{O}$ un préfaisceau
d'anneaux et soit $\mathcal{F}$ un préfaisceau de $\mathcal{O}$-modules.
Si chaque $\mathcal{F}(U)$ est un $\mathcal{O}(U)$-module plat,
alors $\mathcal{F}$ est plat.
\end{lemma}

\begin{proof}
Cela résulte immédiatement des définitions.
\end{proof}

\begin{lemma}

\label{lemma-flatness-sheafification}
Soit $\mathcal{C}$ un site, soit $\mathcal{O}$ un préfaisceau d'anneaux
et soit $\mathcal{F}$ un préfaisceau de $\mathcal{O}$-modules.
Si $\mathcal{F}$ est plat sur $\mathcal{O}$, alors $\mathcal{F}^\#$
est plat sur $\mathcal{O}^\#$.
\end{lemma}

\begin{proof}

Démonstration omise. (Indication : la faisceautisation est exacte.)

\end{proof}

\begin{lemma}

\label{lemma-flatness-sheafification-refined}
Soit $\mathcal{C}$ un site, soit $\mathcal{O}$ un préfaisceau d'anneaux
et soit $\mathcal{F}$ un préfaisceau de $\mathcal{O}$-modules. Supposons
que tout objet $U$ de $\mathcal{C}$ admette un recouvrement
$\{U_i \to U\}_{i \in I}$ tel que $\mathcal{F}(U_i)$ soit un
$\mathcal{O}(U_i)$-module plat. Alors $\mathcal{F}^\#$ est un
$\mathcal{O}^\#$-module plat.
\end{lemma}

\begin{proof}

Soit $\mathcal{G} \subset \mathcal{G}'$ une inclusion de
$\mathcal{O}^\#$-modules. Il faut montrer que
$$
\mathcal{G} \otimes_{\mathcal{O}^\#} \mathcal{F}^\#
\longrightarrow
\mathcal{G}' \otimes_{\mathcal{O}^\#} \mathcal{F}^\#
$$
est injective. D'après le lemme \ref{lemma-sheafification-tensor},
la source de cette flèche est le faisceau associé au préfaisceau
$\mathcal{G} \otimes_{p, \mathcal{O}} \mathcal{F}$,
et il en va de même pour la cible. Si $U$ est un objet de $\mathcal{C}$ tel que
$\mathcal{F}(U)$ soit un $\mathcal{O}(U)$-module plat, alors
$$
(\mathcal{G} \otimes_{p, \mathcal{O}} \mathcal{F})(U) =
\mathcal{G}(U) \otimes_{\mathcal{O}(U)} \mathcal{F}(U)
\longrightarrow
\mathcal{G}'(U) \otimes_{\mathcal{O}(U)} \mathcal{F}(U) =
(\mathcal{G}' \otimes_{p, \mathcal{O}} \mathcal{F})(U)
$$
est injective. Il suffit donc d'établir l'assertion suivante : si un morphisme
de préfaisceaux $f : \mathcal{H} \to \mathcal{H}'$ sur $\mathcal{C}$ est tel
que tout objet $U$ de $\mathcal{C}$ admette un recouvrement
$\{U_i \to U\}_{i \in I}$ pour lequel
$\mathcal{H}(U_i) \to \mathcal{H}'(U_i)$ est injectif, alors $f^\#$ est
injectif. Nous laissons cette vérification au lecteur.

\end{proof}

\begin{lemma}

\label{lemma-colimits-flat}
Colimites et produit tensoriel.
\begin{enumerate}
\item Une colimite filtrante de préfaisceaux de modules plats est plate.
Une somme directe de préfaisceaux de modules plats est plate.
\item Une colimite filtrante de faisceaux de modules plats est plate.
Une somme directe de faisceaux de modules plats est plate.
\end{enumerate}

\end{lemma}

\begin{proof}

La partie (1) découle du lemme \ref{lemma-tensor-commute-colimits} et
d'Algèbre, lemme \ref{algebra-lemma-directed-colimit-exact}, en évaluant
sur les objets. Pour la partie (2), on utilise le lemme
\ref{lemma-tensor-commute-colimits} et le fait qu'une colimite filtrante de
complexes exacts est exacte ; cela repose sur l'exactitude de la
faisceautisation et sur sa commutation aux colimites. Certains détails sont
omis.

\end{proof}

\begin{lemma}

\label{lemma-restriction-flat}
Soit $(\mathcal{C}, \mathcal{O})$ un site annelé et soit $U$ un objet de
$\mathcal{C}$. Si $\mathcal{F}$ est un $\mathcal{O}$-module plat, alors
$\mathcal{F}|_U$ est un $\mathcal{O}_U$-module plat.
\end{lemma}

\begin{proof}
Soit $\mathcal{G}_1 \to \mathcal{G}_2 \to \mathcal{G}_3$ un complexe exact
de $\mathcal{O}_U$-modules. Comme $j_{U!}$ est exact (lemme
\ref{lemma-extension-by-zero-exact}) et que $\mathcal{F}$ est plat comme
$\mathcal{O}$-module, le complexe formé des modules
$$
j_{U!}(\mathcal{G}_i \otimes_{\mathcal{O}_U} \mathcal{F}|_U) =
j_{U!}\mathcal{G}_i \otimes_\mathcal{O} \mathcal{F}
$$
est exact d'après le lemme \ref{lemma-j-shriek-and-tensor}. On en déduit que
$\mathcal{G}_1 \otimes_{\mathcal{O}_U} \mathcal{F}|_U \to
\mathcal{G}_2 \otimes_{\mathcal{O}_U} \mathcal{F}|_U \to
\mathcal{G}_3 \otimes_{\mathcal{O}_U} \mathcal{F}|_U$ est exact par le lemme \ref{lemma-j-shriek-reflects-exactness}.
\end{proof}

\begin{lemma}

\label{lemma-j-shriek-flat}

Soient $\mathcal{C}$ une catégorie, $\mathcal{O}$ un préfaisceau d'anneaux
et $U$ un objet de $\mathcal{C}$. Considérons le foncteur
$j_U : \mathcal{C}/U \to \mathcal{C}$.

\begin{enumerate}

\item Le préfaisceau de $\mathcal{O}$-modules
$j_{U!}\mathcal{O}_U$ (voir la remarque
\ref{remark-localize-presheaves}) est plat.
\item Si $\mathcal{C}$ est un site, $\mathcal{O}$ est un faisceau d'anneaux,
$j_{U!}\mathcal{O}_U$ est un faisceau plat de $\mathcal{O}$-modules.

\end{enumerate}

\end{lemma}

\begin{proof}
Preuve de (1). La discussion de la remarque
\ref{remark-localize-presheaves} donne
$$
j_{U!}\mathcal{O}_U(V)
=
\bigoplus\nolimits_{\varphi \in \Mor_\mathcal{C}(V, U)}
\mathcal{O}(V)
$$
qui est un $\mathcal{O}(V)$-module plat. La partie (1) découle donc du lemme
\ref{lemma-flatness-presheaves}. La partie (2) résulte ensuite de l'égalité
$j_{U!}\mathcal{O}_U = (j_{U!}\mathcal{O}_U)^\#$ (le premier $j_{U!}$
étant pris sur les faisceaux et le second sur les préfaisceaux) et du lemme
\ref{lemma-flatness-sheafification}.
\end{proof}

\begin{lemma}

\label{lemma-module-quotient-flat}

Soient $\mathcal{C}$ une catégorie et $\mathcal{O}$ un préfaisceau d'anneaux.

\begin{enumerate}

\item Tout préfaisceau de $\mathcal{O}$-modules est quotient d'une somme
directe $\bigoplus j_{U_i!}\mathcal{O}_{U_i}$.
\item Tout préfaisceau de $\mathcal{O}$-modules est quotient d'un préfaisceau
plat de $\mathcal{O}$-modules.
\item Si $\mathcal{C}$ est un site et $\mathcal{O}$ un faisceau d'anneaux,
tout faisceau de $\mathcal{O}$-modules est quotient d'une somme directe
$\bigoplus j_{U_i!}\mathcal{O}_{U_i}$.
\item Si $\mathcal{C}$ est un site et $\mathcal{O}$ un faisceau d'anneaux,
tout faisceau de $\mathcal{O}$-modules est quotient d'un faisceau plat de
$\mathcal{O}$-modules.

\end{enumerate}

\end{lemma}

\begin{proof}
Preuve de (1). Pour tout objet $U$ de $\mathcal{C}$ et toute section
$s \in \mathcal{F}(U)$, on obtient un morphisme
$j_{U!}\mathcal{O}_U \to \mathcal{F}$, adjoint du morphisme
$\mathcal{O}_U \to \mathcal{F}|_U$, $1 \mapsto s$. L'application
$$
\bigoplus\nolimits_{(U, s)} j_{U!}\mathcal{O}_U
\longrightarrow
\mathcal{F}
$$
est surjective. Les lemmes \ref{lemma-colimits-flat} et
\ref{lemma-j-shriek-flat} montrent que sa source est plate, ce qui prouve
(2). Le cas des faisceaux s'en déduit soit par faisceautisation, soit en
répétant le même argument.
\end{proof}

\begin{lemma}

\label{lemma-flat-tor-zero}

Soient $\mathcal{C}$ une catégorie et $\mathcal{O}$ un préfaisceau d'anneaux.
$$
0 \to \mathcal{F}'' \to \mathcal{F}' \to \mathcal{F} \to 0
$$
une suite exacte courte de préfaisceaux de $\mathcal{O}$-modules, et soit
$\mathcal{G}$ un préfaisceau de $\mathcal{O}$-modules.

\begin{enumerate}

\item Si $\mathcal{F}$ est un préfaisceau plat de modules, alors la suite
$$
0 \to
\mathcal{F}'' \otimes_{p, \mathcal{O}} \mathcal{G} \to
\mathcal{F}' \otimes_{p, \mathcal{O}} \mathcal{G} \to
\mathcal{F} \otimes_{p, \mathcal{O}} \mathcal{G} \to 0
$$
est exacte.
\item Si $\mathcal{C}$ est un site, $\mathcal{O}$,
$\mathcal{F}$, $\mathcal{F}'$, $\mathcal{F}''$,
$\mathcal{G}$ sont des faisceaux et si $\mathcal{F}$ est plat comme faisceau
de modules, alors la suite
$$
0 \to
\mathcal{F}'' \otimes_\mathcal{O} \mathcal{G} \to
\mathcal{F}' \otimes_\mathcal{O} \mathcal{G} \to
\mathcal{F} \otimes_\mathcal{O} \mathcal{G} \to 0
$$
est exacte.

\end{enumerate}

\end{lemma}

\begin{proof}
Choisissons un préfaisceau plat de $\mathcal{O}$-modules $\mathcal{G}'$
surjectant sur $\mathcal{G}$ ; c'est possible d'après le lemme
\ref{lemma-module-quotient-flat}. Posons
$\mathcal{G}'' = \Ker(\mathcal{G}' \to \mathcal{G})$. Le résultat s'obtient
en appliquant le lemme du serpent au diagramme suivant :
$$
\begin{matrix}
 & & 0 & & 0 & & 0 & & \\
 & & \uparrow & & \uparrow & & \uparrow & & \\
 & & \mathcal{F}'' \otimes_{p, \mathcal{O}} \mathcal{G} & \to &
     \mathcal{F}' \otimes_{p, \mathcal{O}} \mathcal{G} & \to &
     \mathcal{F} \otimes_{p, \mathcal{O}} \mathcal{G} & \to & 0 \\
 & & \uparrow & & \uparrow & & \uparrow & & \\
0 & \to & \mathcal{F}'' \otimes_{p, \mathcal{O}} \mathcal{G}' & \to &
          \mathcal{F}' \otimes_{p, \mathcal{O}} \mathcal{G}' & \to &
          \mathcal{F} \otimes_{p, \mathcal{O}} \mathcal{G}' & \to & 0 \\
 & & \uparrow & & \uparrow & & \uparrow & & \\
 & & \mathcal{F}'' \otimes_{p, \mathcal{O}} \mathcal{G}'' & \to &
     \mathcal{F}' \otimes_{p, \mathcal{O}} \mathcal{G}'' & \to &
     \mathcal{F} \otimes_{p, \mathcal{O}} \mathcal{G}'' & \to & 0 \\
 & & & & & & \uparrow & & \\
 & & & & & & 0 & &
\end{matrix}
$$
dont les lignes et les colonnes sont exactes. La ligne médiane est exacte
parce que le produit tensoriel avec le module plat $\mathcal{G}'$ est exact.
La démonstration dans le cas des faisceaux est identique.
\end{proof}

\begin{lemma}

\label{lemma-flat-ses}

Soient $\mathcal{C}$ une catégorie et $\mathcal{O}$ un préfaisceau d'anneaux.
$$
0 \to
\mathcal{F}_2 \to
\mathcal{F}_1 \to
\mathcal{F}_0 \to 0
$$
une suite exacte courte de préfaisceaux de $\mathcal{O}$-modules.

\begin{enumerate}

\item Si $\mathcal{F}_2$ et $\mathcal{F}_0$ sont plats, alors
$\mathcal{F}_1$ l'est aussi.
\item Si $\mathcal{F}_1$ et $\mathcal{F}_0$ sont plats, alors
$\mathcal{F}_2$ l'est aussi.

\end{enumerate}

Si $\mathcal{C}$ est un site et $\mathcal{O}$ un faisceau d'anneaux, le même
résultat vaut dans $\textit{Mod}(\mathcal{O})$.

\end{lemma}

\begin{proof}
Soit $\mathcal{G}^\bullet$ un complexe exact quelconque de préfaisceaux de
$\mathcal{O}$-modules. Supposons $\mathcal{F}_0$ plat. Le lemme
\ref{lemma-flat-tor-zero} montre que
$$
0 \to
\mathcal{G}^\bullet \otimes_{p, \mathcal{O}} \mathcal{F}_2 \to
\mathcal{G}^\bullet \otimes_{p, \mathcal{O}} \mathcal{F}_1 \to
\mathcal{G}^\bullet \otimes_{p, \mathcal{O}} \mathcal{F}_0 \to 0
$$
est une suite exacte courte de complexes de préfaisceaux de
$\mathcal{O}$-modules. Les assertions (1) et (2) découlent donc du lemme du
serpent. Le cas des faisceaux de modules se démontre de la même manière.
\end{proof}

\begin{lemma}

\label{lemma-flat-resolution-of-flat}
Soient $\mathcal{C}$ une catégorie et $\mathcal{O}$ un préfaisceau d'anneaux.
Soit
$$
\ldots \to
\mathcal{F}_2 \to
\mathcal{F}_1 \to
\mathcal{F}_0 \to
\mathcal{Q} \to 0
$$
un complexe exact de préfaisceaux de $\mathcal{O}$-modules. Si
$\mathcal{Q}$ et tous les $\mathcal{F}_i$ sont des $\mathcal{O}$-modules
plats, alors, pour tout préfaisceau $\mathcal{G}$ de $\mathcal{O}$-modules,
le complexe
$$
\ldots \to
\mathcal{F}_2 \otimes_{p, \mathcal{O}} \mathcal{G} \to
\mathcal{F}_1 \otimes_{p, \mathcal{O}} \mathcal{G} \to
\mathcal{F}_0 \otimes_{p, \mathcal{O}} \mathcal{G} \to
\mathcal{Q} \otimes_{p, \mathcal{O}} \mathcal{G} \to 0
$$
est également exact. Si $\mathcal{C}$ est un site et $\mathcal{O}$ un
faisceau d'anneaux, le même résultat vaut dans
$\textit{Mod}(\mathcal{O})$.
\end{lemma}

\begin{proof}

Cela découle du lemme \ref{lemma-flat-tor-zero}, en décomposant le complexe
en suites exactes courtes et en appliquant le lemme \ref{lemma-flat-ses}
pour démontrer par récurrence que
$\Im(\mathcal{F}_{i + 1} \to \mathcal{F}_i)$ est plat.

\end{proof}

\begin{lemma}

\label{lemma-tensor-flats}
Soit $(\mathcal{C}, \mathcal{O})$ un site annelé. Si $\mathcal{G}$ et
$\mathcal{F}$ sont des $\mathcal{O}$-modules plats, alors
$\mathcal{G} \otimes_\mathcal{O} \mathcal{F}$ est un
$\mathcal{O}$-module plat.
\end{lemma}

\begin{proof}

C'est vrai parce que
$$
(\mathcal{G} \otimes_\mathcal{O} \mathcal{F}) \otimes_\mathcal{O} \mathcal{H}
=
\mathcal{G} \otimes_\mathcal{O} (\mathcal{F} \otimes_\mathcal{O} \mathcal{H})
$$
et une composition de foncteurs exacts est exacte.

\end{proof}

\begin{lemma}

\label{lemma-flat-change-of-rings}
Soit $\mathcal{O}_1 \to \mathcal{O}_2$ un morphisme de faisceaux d'anneaux
sur un site $\mathcal{C}$. Si $\mathcal{G}$ est un
$\mathcal{O}_1$-module plat, alors
$\mathcal{G} \otimes_{\mathcal{O}_1} \mathcal{O}_2$ est un
$\mathcal{O}_2$-module plat.
\end{lemma}

\begin{proof}

C'est vrai parce que
$$
(\mathcal{G} \otimes_{\mathcal{O}_1} \mathcal{O}_2)
\otimes_{\mathcal{O}_2} \mathcal{H}
=
\mathcal{G} \otimes_{\mathcal{O}_1} \mathcal{H}
$$
(par exemple comme faisceaux de groupes abéliens).

\end{proof}

\noindent
Le lemme suivant est l'analogue du critère d'équation de la planéité (Algèbre, lemme \ref{algebra-lemma-flat-eq}).

\begin{lemma}

\label{lemma-flat-eq}

Soit $(\mathcal{C}, \mathcal{O})$ un site annelé et soit $\mathcal{F}$ un
$\mathcal{O}$-module. Les conditions suivantes sont équivalentes :

\begin{enumerate}

\item $\mathcal{F}$ est un $\mathcal{O}$-module plat.
\item Soit $U$ un objet de $\mathcal{C}$ et soit
$$
\mathcal{O}_U \xrightarrow{(f_1, \ldots, f_n)}
\mathcal{O}_U^{\oplus n} \xrightarrow{(s_1, \ldots, s_n)}
\mathcal{F}|_U
$$
un complexe de $\mathcal{O}_U$-modules. Il existe alors un recouvrement
$\{U_i \to U\}$ et, pour chaque $i$, une factorisation
$$
\mathcal{O}_{U_i}^{\oplus n}
\xrightarrow{B_i}
\mathcal{O}_{U_i}^{\oplus l_i} \xrightarrow{(t_{i1}, \ldots, t_{il_i})}
\mathcal{F}|_{U_i}
$$
de $(s_1, \ldots, s_n)|_{U_i}$ telle que
$B_i \circ (f_1, \ldots, f_n)|_{U_i} = 0$.
\item Soit $U$ un objet de $\mathcal{C}$ et soit
$$
\mathcal{O}_U^{\oplus m} \xrightarrow{A}
\mathcal{O}_U^{\oplus n} \xrightarrow{(s_1, \ldots, s_n)}
\mathcal{F}|_U
$$
un complexe de $\mathcal{O}_U$-modules. Il existe alors un recouvrement
$\{U_i \to U\}$ et, pour chaque $i$, une factorisation
$$
\mathcal{O}_{U_i}^{\oplus n}
\xrightarrow{B_i}
\mathcal{O}_{U_i}^{\oplus l_i} \xrightarrow{(t_{i1}, \ldots, t_{il_i})}
\mathcal{F}|_{U_i}
$$
de $(s_1, \ldots, s_n)|_{U_i}$ telle que
$B_i \circ A|_{U_i} = 0$.

\end{enumerate}

\end{lemma}

\begin{proof}

Supposons (1). Soit $\mathcal{I} \subset \mathcal{O}_U$ le faisceau
d'idéaux engendré par $f_1, \ldots, f_n$. Alors
$\sum f_j \otimes s_j$ est une section de
$\mathcal{I} \otimes_{\mathcal{O}_U} \mathcal{F}|_U$
dont l'image dans $\mathcal{F}|_U$ est nulle. Comme $\mathcal{F}|_U$ est
plat (lemme \ref{lemma-restriction-flat}), l'application
$\mathcal{I} \otimes_{\mathcal{O}_U} \mathcal{F}|_U \to \mathcal{F}|_U$
est injective. $\mathcal{I} \otimes_{\mathcal{O}_U} \mathcal{F}|_U$
est le faisceau associé au préfaisceau produit tensoriel, il existe un
recouvrement $\{U_i \to U\}$ tel que
$\sum f_j|_{U_i} \otimes s_j|_{U_i}$ soit nul dans
$\mathcal{I}(U_i) \otimes_{\mathcal{O}(U_i)} \mathcal{F}(U_i)$. En
explicitant les définitions à l'aide d'Algèbre, lemme
\ref{algebra-lemma-relations}, on trouve
$t_{i1}, \ldots, t_{i l_i} \in \mathcal{F}(U_i)$ et
$a_{ijk} \in \mathcal{O}(U_i)$
de sorte que $\sum_j a_{ijk}f_j|_{U_i} = 0$ et
$s_j|_{U_i} = \sum_k a_{ijk}t_{ik}$. Ainsi (2) tient.

\medskip\noindent

Supposons (2), et donnons-nous $U$, $n$, $m$, $A$ et
$s_1, \ldots, s_n$ comme en (3). La matrice $A$ possède $m$ colonnes.
Nous allons démontrer l'assertion (3) par récurrence sur $m$. Pour
$m = 0$, on peut prendre le recouvrement $\{U \to U\}$ et la factorisation
$\mathcal{O}_U^n \to \mathcal{O}_U^n \to \mathcal{F}|_U$, la première
flèche étant l'identité et la seconde $(s_1, \ldots, s_n)$. Supposons
$m > 0$. Soit $(f_1, \ldots, f_n)$ la dernière colonne de $A$ et appliquons
(2). On obtient les nouveaux diagrammes
$$
\mathcal{O}_{U_i}^{\oplus m} \xrightarrow{B_i \circ A|_{U_i}}
\mathcal{O}_{U_i}^{\oplus l_i} \xrightarrow{(t_{i1}, \ldots, t_{il_i})}
\mathcal{F}|_{U_i}
$$
dont la dernière colonne de $A_i = B_i \circ A|_{U_i}$ est nulle.
On peut donc appliquer l'hypothèse de récurrence à
$U_i$, $l_i$, $m - 1$, la matrice composée de la première $m - 1$ colonnes $A_i$, $t_{i1}, \ldots, t_{il_i}$
pour obtenir des recouvrements $\{U_{ij} \to U_i\}$ et des factorisations
$$
\mathcal{O}_{U_{ij}}^{\oplus l_i}
\xrightarrow{C_{ij}}
\mathcal{O}_{U_{ij}}^{\oplus k_{ij}}
\xrightarrow{(v_{ij1}, \ldots, v_{ij k_{ij}})}
\mathcal{F}|_{U_{ij}}
$$
de $(t_{i1}, \ldots, t_{il_i})|_{U_{ij}}$ telles que
$C_{ij} \circ B_i|_{U_{ij}} \circ A|_{U_{ij}} = 0$. Alors
$\{U_{ij} \to U\}$ est un recouvrement, et les factorisations cherchées
s'obtiennent en prenant $B_{ij} = C_{ij} \circ B_i|_{U_{ij}}$ et
$v_{ija}$. Cela montre que (2) implique (3).

\medskip\noindent

Supposons (3). Soit $\mathcal{G} \to \mathcal{H}$ un morphisme injectif
de $\mathcal{O}$-modules. Il faut montrer que
$\mathcal{G} \otimes_\mathcal{O} \mathcal{F} \to
\mathcal{H} \otimes_\mathcal{O} \mathcal{F}$
est injective. Soit $U$ un objet de $\mathcal{C}$ et soit
$s \in (\mathcal{G} \otimes_\mathcal{O} \mathcal{F})(U)$ une section dont
l'image dans $\mathcal{H} \otimes_\mathcal{O} \mathcal{F}$ est nulle.
Il faut montrer que $s$ est nulle. Comme
$\mathcal{G} \otimes_\mathcal{O} \mathcal{F}$
est un faisceau, il suffit de trouver un recouvrement
$\{U_i \to U\}_{i \in I}$ de $\mathcal{C}$ tel que $s|_{U_i}$ soit nulle
pour tout $i \in I$. On peut donc remplacer $U$ par les membres d'un
recouvrement. En particulier, puisque
$\mathcal{G} \otimes_\mathcal{O} \mathcal{F}$
est la faisceautisation de $\mathcal{G} \otimes_{p, \mathcal{O}} \mathcal{F}$
on peut supposer que $s$ est l'image d'un élément $s' \in 
\mathcal{G}(U) \otimes_{\mathcal{O}(U)} \mathcal{F}(U)$. En raisonnant de
même pour $\mathcal{H} \otimes_\mathcal{O} \mathcal{F}$, on peut supposer
que l'image de $s'$ dans
$\mathcal{H}(U) \otimes_{\mathcal{O}(U)} \mathcal{F}(U)$ est nulle.
Écrivons $\mathcal{F}(U) = \colim M_\alpha$ comme colimite filtrante de
$\mathcal{O}(U)$-modules de présentation finie $M_\alpha$ (Algèbre, lemme
\ref{algebra-lemma-module-colimit-fp}). Puisque le produit tensoriel commute
aux colimites filtrantes (Algèbre, lemme
\ref{algebra-lemma-tensor-products-commute-with-limits}), on peut choisir un
$\alpha$ tel que $s'$ provienne d'un élément
$s'' \in \mathcal{G}(U) \otimes_{\mathcal{O}(U)} M_\alpha$ et que l'image
de $s''$ dans
$\mathcal{H}(U) \otimes_{\mathcal{O}(U)} M_\alpha$ soit nulle. Fixons
$\alpha$ et $s''$, puis choisissons une présentation
$$
\mathcal{O}(U)^{\oplus m} \xrightarrow{A} \mathcal{O}(U)^{\oplus n}
\to M_\alpha \to 0
$$
Appliquons (3) au complexe correspondant de $\mathcal{O}_U$-modules
$$
\mathcal{O}_U^{\oplus m} \xrightarrow{A}
\mathcal{O}_U^{\oplus n} \xrightarrow{(s_1, \ldots, s_n)}
\mathcal{F}|_U
$$
Après avoir remplacé $U$ par les membres du recouvrement $U_i$, on obtient
que l'application
$$
M_\alpha \to \mathcal{F}(U)
$$
se factorise par un module libre $\mathcal{O}(U)^{\oplus l}$ pour un certain
$l$. Comme $\mathcal{G}(U) \to \mathcal{H}(U)$ est injective, on en déduit que
$$
\mathcal{G}(U) \otimes_{\mathcal{O}(U)} \mathcal{O}(U)^{\oplus l}
\to
\mathcal{H}(U) \otimes_{\mathcal{O}(U)} \mathcal{O}(U)^{\oplus l}
$$
est également injective. Puisque l'image de $s''$ dans le module de droite
est nulle, son image dans celui de gauche l'est aussi ; autrement dit,
$s$ est nulle, comme voulu.

\end{proof}

\begin{lemma}

\label{lemma-flat-over-thickening}

Soit $\mathcal{C}$ un site. Soit $\mathcal{O}' \to \mathcal{O}$ une
surjection de faisceaux d'anneaux dont le noyau $\mathcal{I}$ est un idéal
de carré nul. Soit $\mathcal{F}'$ un $\mathcal{O}'$-module et posons
$\mathcal{F} = \mathcal{F}'/\mathcal{I}\mathcal{F}'$. Les conditions
suivantes sont équivalentes :

\begin{enumerate}

\item $\mathcal{F}'$ est un $\mathcal{O}'$-module plat ;
\item $\mathcal{F}$ est un $\mathcal{O}$-module plat et
$\mathcal{I} \otimes_\mathcal{O} \mathcal{F} \to \mathcal{F}'$
est injective.

\end{enumerate}

\end{lemma}

\begin{proof}

Si (1) est satisfaite, alors
$\mathcal{F} = \mathcal{F}' \otimes_{\mathcal{O}'} \mathcal{O}$ est plat
sur $\mathcal{O}$ d'après le lemme \ref{lemma-flat-change-of-rings}, et
l'application
$\mathcal{I} \otimes_\mathcal{O} \mathcal{F} \to \mathcal{F}'$
est injective lorsque l'on applique $- \otimes_{\mathcal{O}'} \mathcal{F}'$
à la suite exacte
$0 \to \mathcal{I} \to \mathcal{O}' \to \mathcal{O} \to 0$ ; voir le lemme
\ref{lemma-flat-tor-zero}. Supposons maintenant (2). Dans la suite, nous
utiliserons sans le rappeler l'égalité
$\mathcal{K} \otimes_{\mathcal{O}'} \mathcal{F}' =
\mathcal{K} \otimes_\mathcal{O} \mathcal{F}$
pour tout $\mathcal{O}'$-module $\mathcal{K}$ annulé par $\mathcal{I}$.
Soit $\alpha : \mathcal{G}' \to \mathcal{H}'$ un morphisme injectif de
$\mathcal{O}'$-modules. Soient $\mathcal{G} \subset \mathcal{G}'$,
respectivement $\mathcal{H} \subset \mathcal{H}'$, les sous-faisceaux des
sections annulées par $\mathcal{I}$. Considérons le diagramme
$$
\xymatrix{
\mathcal{G} \otimes_{\mathcal{O}'} \mathcal{F}' \ar[r] \ar[d] &
\mathcal{G}' \otimes_{\mathcal{O}'} \mathcal{F}' \ar[r] \ar[d] &
\mathcal{G}'/\mathcal{G} \otimes_{\mathcal{O}'} \mathcal{F}' \ar[r] \ar[d] &
0 \\
\mathcal{H} \otimes_{\mathcal{O}'} \mathcal{F}' \ar[r] &
\mathcal{H}' \otimes_{\mathcal{O}'} \mathcal{F}' \ar[r] &
\mathcal{H}'/\mathcal{H} \otimes_{\mathcal{O}'} \mathcal{F}' \ar[r] & 0
}
$$
Notez que $\mathcal{G}'/\mathcal{G}$ et $\mathcal{H}'/\mathcal{H}$
sont annulés par $\mathcal{I}$ et que
$\mathcal{G}'/\mathcal{G} \to \mathcal{H}'/\mathcal{H}$ est injectif.
La flèche verticale de droite est donc injective puisque $\mathcal{F}$ est
plat sur $\mathcal{O}$. Il en va de même de celle de gauche. La flèche
verticale médiane est alors injective, et $\mathcal{F}'$ est plat.

\end{proof}

\begin{lemma}

\label{lemma-neighbourhood-isomorphism}
Soit $\mathcal{C}$ un site. Soit $\mathcal{O} \to \mathcal{O}'$ un
morphisme plat de faisceaux d'anneaux. Soit
$\mathcal{I} \subset \mathcal{O}$ un faisceau d'idéaux tel que le morphisme
induit $\mathcal{O}/\mathcal{I} \to
\mathcal{O}'/\mathcal{I}\mathcal{O}'$ soit un isomorphisme. Pour tout
$\mathcal{O}$-module $\mathcal{F}$ annulé par $\mathcal{I}^n$ pour un
$n \geq 0$, l'application $\text{id} \otimes 1 :
\mathcal{F} \to \mathcal{F} \otimes_\mathcal{O} \mathcal{O}'$ est un isomorphisme.
\end{lemma}

\begin{proof}
Démonstration omise. Indication : voir Compléments sur l'algèbre, lemme
\ref{more-algebra-lemma-neighbourhood-isomorphism}.
\end{proof}





\section{Duaux}

\label{section-duals}

\noindent

Soit $(\mathcal{C}, \mathcal{O})$ un site annelé. La catégorie des
$\mathcal{O}$-modules, munie du produit tensoriel construit à la section
\ref{section-tensor-product}, est une catégorie monoïdale symétrique. Pour
un $\mathcal{O}$-module $\mathcal{F}$, les conditions suivantes sont
équivalentes :

\begin{enumerate}

\item $\mathcal{F}$ admet un dual à gauche dans la catégorie monoïdale des
$\mathcal{O}$-modules ;
\item pour tout objet $U$ de $\mathcal{C}$, il existe un recouvrement
$\{U_i \to U\}$ tel que $\mathcal{F}|_{U_i}$ soit facteur direct d'un
$\mathcal{O}|_{U_i}$-module libre de rang fini ;
\item $\mathcal{F}$ est de présentation finie et plate comme
$\mathcal{O}$-module.

\end{enumerate}

Ces équivalences sont établies par l'exemple \ref{example-dual} et les
lemmes \ref{lemma-left-dual-module} et
\ref{lemma-flat-locally-finite-presentation} de la présente section.

\begin{example}

\label{example-dual}

Soit $(\mathcal{C}, \mathcal{O})$ un site annelé. Soit $\mathcal{F}$ un
$\mathcal{O}$-module tel que, pour tout objet $U$ de $\mathcal{C}$, il
existe un recouvrement $\{U_i \to U\}$ pour lequel
$\mathcal{F}|_{U_i}$ est facteur direct d'un
$\mathcal{O}|_{U_i}$-module libre de rang fini. Alors l'application
$$
\mathcal{F} \otimes_\mathcal{O}
\SheafHom_\mathcal{O}(\mathcal{F}, \mathcal{O})
\longrightarrow
\SheafHom_\mathcal{O}(\mathcal{F}, \mathcal{F})
$$
est un isomorphisme. En effet, cette assertion est locale ; elle est vraie
lorsque $\mathcal{F}$ est libre de rang fini et reste vraie pour tout
facteur direct d'un module pour lequel elle est satisfaite. Nous omettons
les détails. Notons
$$
\eta :
\mathcal{O}
\longrightarrow
\mathcal{F} \otimes_\mathcal{O}
\SheafHom_\mathcal{O}(\mathcal{F}, \mathcal{O})
$$
le morphisme qui envoie $1$ sur la section correspondant à
$\text{id}_\mathcal{F}$ sous l'isomorphisme ci-dessus.
Notons également
$$
\epsilon : 
\SheafHom_\mathcal{O}(\mathcal{F}, \mathcal{O})
\otimes_\mathcal{O} \mathcal{F}
\longrightarrow
\mathcal{O}
$$
le morphisme d'évaluation. On vérifie alors que
$\SheafHom_\mathcal{O}(\mathcal{F}, \mathcal{O}), \eta, \epsilon$
est un dual à gauche de $\mathcal{F}$ au sens de Catégories, définition
\ref{categories-definition-dual}. Nous omettons la vérification des
identités
$(1 \otimes \epsilon) \circ (\eta \otimes 1) = \text{id}_\mathcal{F}$
et
$(\epsilon \otimes 1) \circ (1 \otimes \eta) =
\text{id}_{\SheafHom_\mathcal{O}(\mathcal{F}, \mathcal{O})}$.

\end{example}

\begin{lemma}

\label{lemma-left-dual-module}

Soit $(\mathcal{C}, \mathcal{O})$ un site annelé et soit $\mathcal{F}$ un
$\mathcal{O}$-module. Supposons que $\mathcal{G}, \eta, \epsilon$
soit un dual à gauche de $\mathcal{F}$ dans la catégorie monoïdale des
$\mathcal{O}$-modules ; voir Catégories, définition
\ref{categories-definition-dual}. Alors :

\begin{enumerate}

\item pour tout objet $U$ de $\mathcal{C}$, il existe un recouvrement
$\{U_i \to U\}$ tel que $\mathcal{F}|_{U_i}$ soit facteur direct d'un
$\mathcal{O}|_{U_i}$-module libre de rang fini ;
\item l'application
$e : \SheafHom_\mathcal{O}(\mathcal{F}, \mathcal{O}) \to \mathcal{G}$
qui envoie une section locale $\lambda$ sur $(\lambda \otimes 1)(\eta)$
est un isomorphisme ;
\item on a $\epsilon(f, g) = e^{-1}(g)(f)$ pour les sections locales
$f$ et $g$ de $\mathcal{F}$ et $\mathcal{G}$.

\end{enumerate}

\end{lemma}

\begin{proof}

Les hypothèses signifient que
$$
\mathcal{F} \xrightarrow{\eta \otimes 1}
\mathcal{F} \otimes_\mathcal{O} \mathcal{G}
\otimes_\mathcal{O} \mathcal{F}
\xrightarrow{1 \otimes \epsilon} \mathcal{F}
\quad\text{et}\quad
\mathcal{G} \xrightarrow{1 \otimes \eta}
\mathcal{G} \otimes_\mathcal{O} \mathcal{F}
\otimes_\mathcal{O} \mathcal{G}
\xrightarrow{\epsilon \otimes 1} \mathcal{G}
$$
sont les applications identiques. Soit $U$ un objet de $\mathcal{C}$.
Après avoir remplacé $U$ par les membres d'un recouvrement de $U$, on peut trouver
un nombre fini de sections $f_1, \ldots, f_n$ et
$g_1, \ldots, g_n$ de $\mathcal{F}$ et $\mathcal{G}$ sur $U$ telles que
$\eta(1) = \sum f_i g_i$. Considérons
$$
\mathcal{O}_U^{\oplus n} \to \mathcal{F}|_U
$$
l'application qui envoie le $i$-ième vecteur de base sur $f_i$. On peut
alors factoriser $\eta|_U$ par une application
$\tilde \eta : \mathcal{O}_U \to
\mathcal{O}_U^{\oplus n} \otimes_{\mathcal{O}_U} \mathcal{G}|_U$. Nous obtenons un diagramme commutatif
$$
\xymatrix{
\mathcal{F}|_U
\ar[rr]_-{\eta \otimes 1} \ar[rrd]_-{\tilde \eta \otimes 1} & &
\mathcal{F}|_U \otimes \mathcal{G}|_U \otimes \mathcal{F}|_U
\ar[r]_-{1 \otimes \epsilon} &
\mathcal{F}|_U \\
& &
\mathcal{O}_U^{\oplus n} \otimes \mathcal{G}|_U \otimes \mathcal{F}|_U
\ar[u] \ar[r]^-{1 \otimes \epsilon} &
\mathcal{O}_U^{\oplus n} \ar[u]
}
$$
Ce diagramme montre que l'identité de $\mathcal{F}|_U$ se factorise par un
$\mathcal{O}_U$-module libre de rang fini, ce qui prouve (1). La partie (2)
découle de Catégories, lemme \ref{categories-lemma-left-dual}, et de sa
démonstration. La partie (3) découle de la première identité de la preuve.
On peut aussi déduire (2) et (3) de l'unicité des duaux à gauche
(Catégories, remarque \ref{categories-remark-left-dual-adjoint}) et de la
construction du dual à gauche dans l'exemple \ref{example-dual}.

\end{proof}

\begin{lemma}

\label{lemma-flat-locally-finite-presentation}
Soit $(\mathcal{C}, \mathcal{O})$ un site annelé et soit $\mathcal{F}$
localement de présentation finie et plat. Pour tout
objet $U$ de $\mathcal{C}$, il existe alors un recouvrement
$\{U_i \to U\}$ tel que $\mathcal{F}|_{U_i}$ soit facteur direct d'un
$\mathcal{O}_{U_i}$-module libre de rang fini.
\end{lemma}

\begin{proof}

Choisissons un objet $U$ de $\mathcal{C}$. Après avoir remplacé $U$ par les
membres d'un recouvrement, on peut supposer qu'il existe une présentation
$$
\mathcal{O}_U^{\oplus r} \to
\mathcal{O}_U^{\oplus n} \to \mathcal{F}|_U \to 0
$$
D'après le lemme \ref{lemma-flat-eq}, après avoir de nouveau remplacé $U$
par les membres d'un recouvrement, on peut supposer qu'il existe une
factorisation
$$
\mathcal{O}_U^{\oplus n} \to
\mathcal{O}_U^{\oplus n_1} \to \mathcal{F}|_U
$$
tel que le composé
$\mathcal{O}_U^{\oplus r} \to \mathcal{O}_U^{\oplus n}
\to \mathcal{O}_U^{\oplus n_r}$ est zéro. Cela signifie que la surjection $\mathcal{O}_U^{\oplus n_r} \to \mathcal{F}|_U$
admet une section, ce qui conclut la preuve.

\end{proof}











\section{Vers les modules constructibles}

\label{section-constructible}

\noindent
Rappelons qu'un objet quasi-compact d'un site est, grossièrement, un objet
dont tout recouvrement admet un raffinement fini ; la définition précise est
un peu plus élaborée, voir Sites, section
\ref{sites-section-quasi-compact}. Si tout objet d'un site admet un
recouvrement par des objets quasi-compacts, les modules
$j_!\mathcal{O}_U$, avec $U$ quasi-compact, forment un ensemble de
générateurs particulièrement commode pour la catégorie de tous les modules.

\begin{lemma}

\label{lemma-covering-gives-surjection}
Soit $(\mathcal{C}, \mathcal{O})$ un site annelé et soit
$\{U_i \to U\}$ un recouvrement de $\mathcal{C}$. Alors la suite
$$
\bigoplus j_{U_i \times_U U_j!}\mathcal{O}_{U_i \times_U U_j} \to
\bigoplus j_{U_i!}\mathcal{O}_{U_i} \to j_!\mathcal{O}_U \to 0
$$ est exacte.
\end{lemma}

\begin{proof}
Pour tout $\mathcal{O}$-module $\mathcal{F}$, le foncteur
$\Hom_\mathcal{O}(-, \mathcal{F})$ transforme la suite de l'énoncé en la
suite exacte $0 \to \mathcal{F}(U) \to \prod \mathcal{F}(U_i) \to
\prod \mathcal{F}(U_i \times_U U_j)$ ; voir
(\ref{equation-map-lower-shriek-OU-into-module}). Le résultat découle alors
d'Homologie, lemme \ref{homology-lemma-check-exactness}.
\end{proof}

\begin{lemma}

\label{lemma-silly-quasi-compact}
Soit $(\mathcal{C}, \mathcal{O})$ un site annelé. Soit
$\mathcal{U} = \{U_i \to U\}_{i \in I}$ un recouvrement dans
$\mathcal{C}$. Si $U$ est quasi-compact, il existe un sous-ensemble fini
$I' \subset I$ tel que la suite
$$
\bigoplus\nolimits_{i, i' \in I'}
j_{U_i \times_U U_{i'}!}\mathcal{O}_{U_i \times_U U_{i'}} \to
\bigoplus\nolimits_{i \in I'}
j_{U_i!}\mathcal{O}_{U_i} \to
j_!\mathcal{O}_U \to 0
$$ est exacte.
\end{lemma}

\begin{proof}

Le lemme découle immédiatement du lemme
\ref{lemma-covering-gives-surjection} lorsque $U$ satisfait la condition
(3) de Sites, lemme \ref{sites-lemma-conclude-quasi-compact}. Le lecteur
peut ignorer la démonstration du cas général. Par définition, il existe un
recouvrement
$\mathcal{V} = \{V_j \to U\}_{j \in J}$ et un morphisme
$\mathcal{V} \to \mathcal{U}$ de familles de morphismes de but fixé, donné
par $\text{id} : U \to U$, $\alpha : J \to I$ et
$f_j : V_j \to U_{\alpha(j)}$
(voir Sites, définition \ref{sites-definition-morphism-coverings}), de sorte
que l'image $I' \subset I$ de $\alpha$ soit finie. D'après Homologie, lemme
\ref{homology-lemma-check-exactness}, il suffit de montrer que, pour tout
faisceau de $\mathcal{O}$-modules $\mathcal{F}$, le foncteur
$\Hom_\mathcal{O}(-, \mathcal{F})$ transforme la suite de l'énoncé en une
suite exacte. Par (\ref{equation-map-lower-shriek-OU-into-module}), on
obtient la suite usuelle
$$
0 \to
\mathcal{F}(U) \to
\prod\nolimits_{i \in I'} \mathcal{F}(U_i) \to
\prod\nolimits_{i, i' \in I'} \mathcal{F}(U_i \times_U U_{i'})
$$
Cette suite est exacte d'après Sites, lemme
\ref{sites-lemma-compare-sheaf-condition},
appliqué à la famille de cartes $\{U_i \to U\}_{i \in I'}$
qui est raffinée par le recouvrement $\mathcal{V}$.

\end{proof}

\begin{lemma}

\label{lemma-sections-over-quasi-compact}
Soit $\mathcal{C}$ un site et soit $W$ un objet quasi-compact de
$\mathcal{C}$.
\begin{enumerate}
\item Le foncteur $\Sh(\mathcal{C}) \to \textit{Sets}$,
$\mathcal{F} \mapsto \mathcal{F}(W)$ commute aux coproduits.
\item Soit $\mathcal{O}$ un faisceau d'anneaux sur $\mathcal{C}$. Le
foncteur $\textit{Mod}(\mathcal{O}) \to \textit{Ab}$,
$\mathcal{F} \mapsto \mathcal{F}(W)$ commute aux sommes directes.
\end{enumerate}

\end{lemma}

\begin{proof}
Preuve de (1). D'après Sites, lemme
\ref{sites-lemma-directed-colimits-sections}, la prise des sections sur $W$
commute aux colimites filtrantes dont les morphismes de transition sont
injectifs. Soit $\mathcal{F}_i$ une famille de faisceaux d'ensembles
indexée par un ensemble $I$. Le coproduit $\coprod \mathcal{F}_i$ est la
colimite filtrante, sur l'ensemble ordonné des sous-ensembles finis
$E \subset I$, des coproduits
$\mathcal{F}_E = \coprod_{i \in E} \mathcal{F}_i$. Les morphismes de
transition étant injectifs, on conclut.

\medskip\noindent

Preuve de (2). Soit $\mathcal{F}_i$ une famille de faisceaux de
$\mathcal{O}$-modules indexée par un ensemble $I$. Alors
$\bigoplus \mathcal{F}_i$ est la colimite filtrante, sur l'ensemble ordonné
des sous-ensembles finis $E \subset I$, des sommes directes
$\mathcal{F}_E = \bigoplus_{i \in E} \mathcal{F}_i$. Une colimite filtrante
de faisceaux abéliens se calcule dans la catégorie des faisceaux d'ensembles.
Si $E \subset E'$, le morphisme de transition
$\mathcal{F}_E \to \mathcal{F}_{E'}$ est injectif, car la faisceautisation
est exacte et l'injectivité est claire sur les préfaisceaux sous-jacents.
Sites, lemme \ref{sites-lemma-directed-colimits-sections}, ramène donc au
cas d'un ensemble fini d'indices. Ce cas est traité par le lemme
\ref{lemma-limits-colimits-abelian-sheaves}, qui vaut sur tout objet de
$\mathcal{C}$.

\end{proof}

\begin{lemma}

\label{lemma-quasi-compact-hom-from}
Soit $(\mathcal{C}, \mathcal{O})$ un site annelé et soit $U$ un objet
quasi-compact de $\mathcal{C}$. Alors le foncteur
$\Hom_\mathcal{O}(j_!\mathcal{O}_U, -)$ commute aux sommes directes.
\end{lemma}

\begin{proof}
En effet, $\Hom_\mathcal{O}(j_!\mathcal{O}_U, \mathcal{F}) =
\mathcal{F}(U)$ d'après
(\ref{equation-map-lower-shriek-OU-into-module}), et le foncteur
$\mathcal{F} \mapsto \mathcal{F}(U)$ commute aux sommes directes par le
lemme \ref{lemma-sections-over-quasi-compact}.
\end{proof}

\noindent
Afin d'indiquer les résultats les plus nets possibles dans ce qui suit, nous introduisons une notation.

\begin{situation}

\label{situation-quasi-compact-objects}
Soit $\mathcal{C}$ un site et soit
$\mathcal{B} \subset \text{Ob}(\mathcal{C})$ un ensemble d'objets.
Considérons les conditions suivantes :
\begin{enumerate}
\item
\label{item-enough}
Tout objet de $\mathcal{C}$ admet un recouvrement par des éléments de
$\mathcal{B}$. \item
\label{item-enough-qc}
Tout $U \in \mathcal{B}$ est quasi-compact (Sites, section
\ref{sites-section-quasi-compact}). \item
\label{item-enough-qc-qs}
Pour tout recouvrement $\{U_i \to U\}$ avec
$U_i, U \in \mathcal{B}$, les produits fibrés
$U_i \times_U U_j$ sont quasi-compacts.
\end{enumerate}

\end{situation}

\begin{lemma}

\label{lemma-module-quotient-direct-sum}
Dans la situation \ref{situation-quasi-compact-objects}, supposons
(\ref{item-enough}) satisfaite.
\begin{enumerate}
\item Tout faisceau d'ensembles est le but d'un morphisme surjectif dont la
source est un coproduit $\coprod h_{U_i}^\#$, avec $U_i$ dans
$\mathcal{B}$.
\item Si $\mathcal{O}$ est un faisceau d'anneaux, tout
$\mathcal{O}$-module est quotient d'une somme directe
$\bigoplus\nolimits j_{U_i!}\mathcal{O}_{U_i}$, avec $U_i$ dans
$\mathcal{B}$.
\end{enumerate}

\end{lemma}

\begin{proof}

La partie (1) découle de Sites, lemmes
\ref{sites-lemma-sheaf-coequalizer-representable} et
\ref{sites-lemma-covering-surjective-after-sheafification}. La partie (2)
découle des lemmes \ref{lemma-module-quotient-flat} et
\ref{lemma-covering-gives-surjection}.

\end{proof}

\begin{lemma}

\label{lemma-module-filtered-colimit-constructibles}

Dans la situation \ref{situation-quasi-compact-objects}, supposons
(\ref{item-enough}) et (\ref{item-enough-qc}) satisfaites.

\begin{enumerate}

\item Tout faisceau d'ensembles est une colimite filtrante de faisceaux de
la forme

\begin{equation}
\label{equation-towards-constructible-sets}
\text{Coégalisateur}\left(
\xymatrix{
\coprod\nolimits_{j = 1, \ldots, m} h_{V_j}^\#
\ar@<1ex>[r] \ar@<-1ex>[r] &
\coprod\nolimits_{i = 1, \ldots, n} h_{U_i}^\#
}
\right)
\end{equation}

avec $U_i$ et $V_j$ dans $\mathcal{B}$.
\item Si $\mathcal{O}$ est un faisceau d'anneaux, tout
$\mathcal{O}$-module est une colimite filtrante de faisceaux de la forme

\begin{equation}
\label{equation-towards-constructible}
\Coker\left(
\bigoplus\nolimits_{j = 1, \ldots, m} j_{V_j!}\mathcal{O}_{V_j}
\longrightarrow
\bigoplus\nolimits_{i = 1, \ldots, n} j_{U_i!}\mathcal{O}_{U_i}
\right)
\end{equation}

avec $U_i$ et $V_j$ dans $\mathcal{B}$.

\end{enumerate}

\end{lemma}

\begin{proof}

Preuve de (1). D'après le lemme
\ref{lemma-module-quotient-direct-sum}, tout faisceau d'ensembles
$\mathcal{F}$ est le but d'une surjection dont la source est un coproduit
$\mathcal{F}_0$ de faisceaux de la forme $h_{U}^\#$, avec
$U \in \mathcal{B}$. En appliquant ce résultat à
$\mathcal{F}_0 \times_\mathcal{F} \mathcal{F}_0$, on obtient que
$\mathcal{F}$ est le coégalisateur d'une paire de morphismes
$$
\xymatrix{
\coprod\nolimits_{j \in J} h_{V_j}^\#
\ar@<1ex>[r] \ar@<-1ex>[r] &
\coprod\nolimits_{i \in I} h_{U_i}^\#
}
$$
pour certains ensembles d'indices $I$, $J$, et certains $V_j$, $U_i$ dans
$\mathcal{B}$. Pour tout sous-ensemble fini $J' \subset J$, il existe un
sous-ensemble fini $I' \subset I$ tel que les deux morphismes envoient le
coproduit indexé par $j \in J'$ dans celui indexé par $i \in I'$ ; voir
Sites, lemme \ref{sites-lemma-directed-colimits-sections}. Nous omettons les
détails ; on utilise qu'un coproduit infini est la colimite filtrante de ses
sous-coproduits finis. Notre faisceau est donc la colimite des conoyaux de
ces morphismes entre coproduits finis.

\medskip\noindent

Preuve de (2). Par le lemme \ref{lemma-module-quotient-direct-sum}, tout
module est quotient d'une somme directe de modules de la forme
$j_{U!}\mathcal{O}_U$, avec $U \in \mathcal{B}$. Tout module est donc un
conoyau
$$
\Coker\left(
\bigoplus\nolimits_{j \in J} j_{V_j!}\mathcal{O}_{V_j}
\longrightarrow
\bigoplus\nolimits_{i \in I} j_{U_i!}\mathcal{O}_{U_i}
\right)
$$
pour certains ensembles d'indices $I$, $J$, et certains $V_j$, $U_i$ dans
$\mathcal{B}$. Pour tout sous-ensemble fini $J' \subset J$, il existe un
sous-ensemble fini $I' \subset I$ tel que la somme directe indexée par
$j \in J'$ s'envoie dans celle indexée par $i \in I'$ ; voir le lemme
\ref{lemma-quasi-compact-hom-from}. Notre module est donc la colimite des
conoyaux de ces morphismes entre sommes directes finies.

\end{proof}

\begin{lemma}

\label{lemma-cokernel-map-towards-constructibles}

Dans la situation \ref{situation-quasi-compact-objects}, supposons
(\ref{item-enough}) et (\ref{item-enough-qc}) satisfaites. Soit
$\mathcal{O}$ un faisceau d'anneaux. Le conoyau d'un morphisme entre deux
modules de la forme (\ref{equation-towards-constructible}) est encore un
module de la forme (\ref{equation-towards-constructible}).

\end{lemma}

\begin{proof}

Soit $\mathcal{F} = \Coker(\bigoplus j_{V_j!}\mathcal{O}_{V_j} \to
\bigoplus j_{U_i!}\mathcal{O}_{U_i})$
comme dans (\ref{equation-towards-constructible}). Il suffit de montrer que
le conoyau d'un morphisme
$\varphi : j_{W!}\mathcal{O}_W \to \mathcal{F}$, avec
$W \in \mathcal{B}$, est encore un module du même type. Le morphisme
$\varphi$ correspond à un élément $s \in \mathcal{F}(W)$. Comme
$\bigoplus j_{U_i!}\mathcal{O}_{U_i} \to \mathcal{F}$ est surjectif, on
peut, par (\ref{item-enough}), choisir un recouvrement
$\{W_k \to W\}_{k \in K}$ avec $W_k \in \mathcal{B}$
tel que $s|_{W_k}$ soit l'image d'une section
$s_k$ de $\bigoplus j_{U_i!}\mathcal{O}_{U_i})$. Par
(\ref{item-enough-qc}), l'objet $W$ est quasi-compact. D'après le lemme
\ref{lemma-silly-quasi-compact}, il existe un sous-ensemble fini
$K' \subset K$ tel que
$\bigoplus_{k \in K'} j_{W_k!}\mathcal{O}_{W_k} \to j_{W!}\mathcal{O}_W$
soit surjectif. On en déduit que $\Coker(\varphi)$ est égal à
$$
\Coker\left(
\bigoplus\nolimits_{k \in K'} j_{W_k!}\mathcal{O}_{W_k} \oplus
\bigoplus j_{V_j!}\mathcal{O}_{V_j}
\longrightarrow
\bigoplus j_{U_i!}\mathcal{O}_{U_i}
\right)
$$
où le morphisme $\bigoplus_{k \in K'} j_{W_k!}\mathcal{O}_{W_k}
\to \bigoplus j_{U_i!}\mathcal{O}_{U_i}$ correspond à
$\sum_{k \in K'} s_k$. Cela termine la preuve.

\end{proof}

\begin{lemma}

\label{lemma-change-presentation-towards-constructibles}

Dans la situation \ref{situation-quasi-compact-objects}, supposons
(\ref{item-enough}), (\ref{item-enough-qc}) et
(\ref{item-enough-qc-qs}) satisfaites. Soit $\mathcal{O}$ un faisceau
d'anneaux. Donnons-nous un morphisme
$$
\bigoplus\nolimits_{j = 1, \ldots, m} j_{V_j!}\mathcal{O}_{V_j}
\longrightarrow
\bigoplus\nolimits_{i = 1, \ldots, n} j_{U_i!}\mathcal{O}_{U_i}
$$
avec $U_i$ et $V_j$ dans $\mathcal{B}$, ainsi que des recouvrements
$\{U_{ik} \to U_i\}_{k \in K_i}$ avec $U_{ik} \in \mathcal{B}$. Il existe
alors des sous-ensembles finis $K'_i \subset K_i$, un ensemble fini $L$
d'objets $W_l \in \mathcal{B}$ et un diagramme commutatif
$$
\xymatrix{
\bigoplus_{l \in L} j_{W_l!}\mathcal{O}_{W_l} \ar[d] \ar[r] &
\bigoplus_{i = 1, \ldots, n} \bigoplus_{k \in K'_i}
j_{U_{ik}!}\mathcal{O}_{U_{ik}} \ar[d] \\
\bigoplus_{j = 1, \ldots, m} j_{V_j!}\mathcal{O}_{V_j} \ar[r] &
\bigoplus_{i = 1, \ldots, n} j_{U_i!}\mathcal{O}_{U_i}
}
$$
induisant un isomorphisme sur les conoyaux des morphismes horizontaux.

\end{lemma}

\begin{proof}

Puisque $U_i$ est quasi-compact, on peut choisir des sous-ensembles finis
$K'_i \subset K_i$ comme dans le lemme
\ref{lemma-silly-quasi-compact}. Puisque
$\bigoplus_{i = 1, \ldots, n}
\bigoplus_{k \in K'_i} j_{U_{ik}!}\mathcal{O}_{U_{ik}} \to
\bigoplus_{i = 1, \ldots, n} j_{U_i!}\mathcal{O}_{U_i}$ est surjectif, on
peut trouver des recouvrements
$\{V_{jm} \to V_j\}_{m \in M_j}$ avec $V_{jm} \in \mathcal{B}$ pour
lesquels il existe un diagramme commutatif
$$
\xymatrix{
\bigoplus_{j = 1, \ldots, m} \bigoplus_{m \in M_j}
j_{V_{jm}!}\mathcal{O}_{V_{jm}} \ar[d] \ar[r] &
\bigoplus_{i = 1, \ldots n} \bigoplus_{k \in K'_i}
j_{U_{ik}!}\mathcal{O}_{U_{ik}} \ar[d] \\
\bigoplus_{j = 1, \ldots, m} j_{V_j!}\mathcal{O}_{V_j} \ar[r] &
\bigoplus_{i = 1, \ldots, n} j_{U_i!}\mathcal{O}_{U_i}
}
$$
Puisque $V_j$ est quasi-compact, on peut choisir des sous-ensembles finis
$M'_j \subset M_j$ comme dans le lemme
\ref{lemma-silly-quasi-compact}. Posons
$$
L = \left(\coprod\nolimits_{i = 1, \ldots, n} K'_i \times K'_i \right)
\coprod
\left(\coprod\nolimits_{j = 1, \ldots, m} M'_j\right)
$$
et, pour $l = (k, k') \in K'_i \times K'_i \subset L$, posons
$W_l = U_{ik} \times_{U_i} U_{ik'}$ ; pour
$l = m \in M'_j \subset L$, posons $W_l = V_{jm}$. Les suites exactes du
lemme \ref{lemma-silly-quasi-compact}
pour les familles $\{U_{ik} \to U_i\}_{k \in K'_i}$
montrent que l'on obtient un diagramme comme dans l'énoncé du lemme ; nous
omettons les détails. Il reste seulement à assurer que
$W_l \in \mathcal{B}$. Or $W_l$ est quasi-compact pour tout $l \in L$.
Une nouvelle application du lemme \ref{lemma-silly-quasi-compact} fournit
des familles finies de morphismes
$\{W_{lt} \to W_l\}_{t \in T_l}$ avec $W_{lt} \in \mathcal{B}$ telles que
$\bigoplus j_{W_{lt}!}\mathcal{O}_{W_{lt}} \to j_{W_l!}\mathcal{O}_{W_l}$
soit surjectif. On remplace alors $L$ par $\coprod_{l \in L} T_l$, ce qui
achève la démonstration.

\end{proof}

\begin{lemma}

\label{lemma-extension-towards-constructibles}
Dans la situation \ref{situation-quasi-compact-objects}, supposons
(\ref{item-enough}), (\ref{item-enough-qc}) et
(\ref{item-enough-qc-qs}) satisfaites. Soit $\mathcal{O}$ un faisceau
d'anneaux. Toute extension de modules de la forme
(\ref{equation-towards-constructible}) est encore un module de la forme
(\ref{equation-towards-constructible}).
\end{lemma}

\begin{proof}

Soit $0 \to \mathcal{F}_1 \to \mathcal{F}_2 \to \mathcal{F}_3 \to 0$
une suite exacte courte de $\mathcal{O}$-modules dans laquelle
$\mathcal{F}_1$ et $\mathcal{F}_3$ sont de la forme
(\ref{equation-towards-constructible}). Choisissons des présentations
$$
\bigoplus A_{V_j} \to \bigoplus A_{U_i} \to \mathcal{F}_1 \to 0
\quad\text{et}\quad
\bigoplus A_{T_j} \to \bigoplus A_{W_i} \to \mathcal{F}_3 \to 0
$$
Dans cette preuve, toutes les sommes directes sont finies, et nous écrivons
$A_U = j_{U!}\mathcal{O}_U$ pour $U \in \mathcal{B}$. Puisque
$\mathcal{F}_2 \to \mathcal{F}_3$ est surjectif, on peut choisir des
recouvrements $\{W_{ik} \to W_i\}$ avec $W_{ik} \in \mathcal{B}$ tels que
$A_{W_{ik}} \to \mathcal{F}_3$ se relève en un morphisme
$A_{W_{ik}} \to \mathcal{F}_2$. D'après le lemme
\ref{lemma-change-presentation-towards-constructibles}, on peut remplacer
la famille $\{W_i\}$ par une sous-famille finie de la famille
$\{W_{ik}\}$ et supposer que le morphisme
$\bigoplus A_{W_i} \to \mathcal{F}_3$ se relève dans $\mathcal{F}_2$.
Considérons le noyau
$$
\mathcal{K}_2 = \Ker(\bigoplus A_{U_i} \oplus \bigoplus A_{W_i}
\longrightarrow \mathcal{F}_2)
$$
Le lemme du serpent montre que ce noyau surjecte sur
$\mathcal{K}_3 = \Ker(\bigoplus A_{W_i} \to \mathcal{F}_3)$. En raisonnant
comme ci-dessus, et après avoir remplacé chaque $T_j$ par une famille finie
d'éléments de $\mathcal{B}$, ce que permet le lemme
\ref{lemma-silly-quasi-compact}, on peut supposer qu'il existe un morphisme
$\bigoplus A_{T_j} \to \mathcal{K}_2$ relevant le morphisme donné
$\bigoplus A_{T_j} \to \mathcal{K}_3$. Alors $\bigoplus A_{V_j} \oplus \bigoplus A_{T_j} \to \mathcal{K}_2$
est surjectif, ce qui achève la démonstration.

\end{proof}

\begin{lemma}

\label{lemma-towards-constructible-when-serre-subcategory}

Dans la situation \ref{situation-quasi-compact-objects}, supposons
(\ref{item-enough}), (\ref{item-enough-qc}) et
(\ref{item-enough-qc-qs}) satisfaites. Soit $\mathcal{O}$ un faisceau
d'anneaux. Soit $\mathcal{A} \subset \textit{Mod}(\mathcal{O})$ la
sous-catégorie pleine des modules isomorphes à un conoyau de la forme
(\ref{equation-towards-constructible}). Si le noyau de tout morphisme de
$\mathcal{O}$-modules de la forme
$$
\bigoplus\nolimits_{j = 1, \ldots, m} j_{V_j!}\mathcal{O}_{V_j}
\longrightarrow
\bigoplus\nolimits_{i = 1, \ldots, n} j_{U_i!}\mathcal{O}_{U_i}
$$
avec $U_i$ et $V_j$ dans $\mathcal{B}$ appartient à $\mathcal{A}$, alors
$\mathcal{A}$ est une sous-catégorie de Serre faible de
$\textit{Mod}(\mathcal{O})$.

\end{lemma}

\begin{proof}

Nous utiliserons le critère d'Homologie, lemme
\ref{homology-lemma-characterize-weak-serre-subcategory}. D'après les
lemmes \ref{lemma-cokernel-map-towards-constructibles} et
\ref{lemma-extension-towards-constructibles}
il suffit de montrer que le noyau d'un morphisme
$\mathcal{F} \to \mathcal{G}$ entre objets de $\mathcal{A}$ appartient à
$\mathcal{A}$. Pour cela, choisissons des présentations
$$
\bigoplus A_{V_j} \to \bigoplus A_{U_i} \to \mathcal{F} \to 0
\quad\text{et}\quad
\bigoplus A_{T_j} \to \bigoplus A_{W_i} \to \mathcal{G} \to 0
$$
Dans cette preuve, toutes les sommes directes sont finies, et nous écrivons
$A_U = j_{U!}\mathcal{O}_U$ pour $U \in \mathcal{B}$. En appliquant les
lemmes \ref{lemma-covering-gives-surjection} et
\ref{lemma-change-presentation-towards-constructibles}
et en raisonnant comme dans la preuve du lemme
\ref{lemma-extension-towards-constructibles}, on peut supposer que le
morphisme $\mathcal{F} \to \mathcal{G}$ se relève en un morphisme de
présentations
$$
\xymatrix{
\bigoplus A_{V_j} \ar[r] \ar[d] &
\bigoplus A_{U_i} \ar[r] \ar[d] &
\mathcal{F} \ar[r] \ar[d] & 0 \\
\bigoplus A_{T_j} \ar[r] &
\bigoplus A_{W_i} \ar[r] &
\mathcal{G} \ar[r] & 0
}
$$
On obtient alors
$$
\Ker(\mathcal{F} \to \mathcal{G}) =
\Coker\left(\bigoplus A_{V_j} \to
\Ker\left(
\bigoplus A_{T_j} \oplus \bigoplus A_{U_i} \to \bigoplus A_{W_i}\right)\right)
$$
et le résultat découle de l'hypothèse et du lemme
\ref{lemma-cokernel-map-towards-constructibles}.

\end{proof}







\section{Morphismes plats}

\label{section-flat-morphisms}

\begin{definition}

\label{definition-flat-morphism}

Soit
$(f, f^\sharp) :
(\Sh(\mathcal{C}), \mathcal{O})
\longrightarrow
(\Sh(\mathcal{C}'), \mathcal{O}')$
un morphisme de topos annelés. On dit que $(f, f^\sharp)$ est
{\it plat} si le morphisme d'anneaux $f^\sharp : f^{-1}\mathcal{O}' \to \mathcal{O}$
est plat. Nous disons qu'un morphisme de sites annelés est {\it plat}
si le morphisme de topos annelés associé est plat.

\end{definition}

\begin{lemma}

\label{lemma-flat-pullback-exact}
Soit $f : \Sh(\mathcal{C}) \to \Sh(\mathcal{C}')$ un morphisme de topos
annelés. Alors
$$
f^{-1} : \textit{Ab}(\mathcal{C}') \longrightarrow \textit{Ab}(\mathcal{C}),
\quad
\mathcal{F} \longmapsto f^{-1}\mathcal{F}
$$ est exact. Si $(f, f^\sharp) :
(\Sh(\mathcal{C}), \mathcal{O})
\to
(\Sh(\mathcal{C}'), \mathcal{O}')$ est un morphisme de topos annelés plat, alors $$
f^* : \textit{Mod}(\mathcal{O}') \longrightarrow \textit{Mod}(\mathcal{O}),
\quad
\mathcal{F} \longmapsto f^*\mathcal{F}
$$ est exact.
\end{lemma}

\begin{proof}
Pour un faisceau abélien $\mathcal{G}$ sur $\mathcal{C}'$, le faisceau
d'ensembles sous-jacent à $f^{-1}\mathcal{G}$ est l'image inverse par
$f^{-1}$ du faisceau d'ensembles sous-jacent à $\mathcal{G}$ ; voir Sites,
section \ref{sites-section-sheaves-algebraic-structures}. L'exactitude de
$f^{-1}$ sur les faisceaux d'ensembles, qui fait partie de la définition
d'un morphisme de topos (voir Sites, définition
\ref{sites-definition-topos}), implique donc l'exactitude de $f^{-1}$ sur les
faisceaux abéliens.

\medskip\noindent

Pour voir l'énoncé sur les modules, rappelons que $f^*\mathcal{F}$ est défini comme le produit tensoriel
$f^{-1}\mathcal{F} \otimes_{f^{-1}\mathcal{O}', f^\sharp} \mathcal{O}$.
Ainsi, $f^*$ est le composé de deux foncteurs exacts.

\end{proof}

\begin{definition}

\label{definition-flat-module}
Soit $f : (\Sh(\mathcal{C}), \mathcal{O}) \to
(\Sh(\mathcal{D}), \mathcal{O}')$ un morphisme de topos annelés et soit
$\mathcal{F}$ un faisceau de $\mathcal{O}$-modules. On dit que
$\mathcal{F}$ est {\it plat sur
$(\Sh(\mathcal{D}), \mathcal{O}')$} si $\mathcal{F}$ est plat comme
$f^{-1}\mathcal{O}'$-module.
\end{definition}

\noindent
Ceci est compatible avec la notion telle que définie pour les morphismes des espaces annelés, voir Modules, définition \ref{modules-definition-flat-module} et la discussion suivante.

\begin{lemma}

\label{lemma-pullback-internal-hom}
Soit $f : (\mathcal{C}, \mathcal{O}_\mathcal{C}) \to
(\mathcal{D}, \mathcal{O}_\mathcal{D})$ un morphisme de sites annelés.
Soient $\mathcal{F}$ et $\mathcal{G}$ des
$\mathcal{O}_\mathcal{D}$-modules. Si $\mathcal{F}$ est de présentation
finie et si $f$ est plat, alors le morphisme canonique
$$
f^*\SheafHom_{\mathcal{O}_\mathcal{D}}(\mathcal{F}, \mathcal{G})
\longrightarrow
\SheafHom_{\mathcal{O}_\mathcal{C}}(f^*\mathcal{F}, f^*\mathcal{G})
$$
de la remarque \ref{remark-pullback-internal-hom} est un isomorphisme.
\end{lemma}

\begin{proof}

Supposons que $f$ soit donné par le foncteur continu
$u : \mathcal{D} \to \mathcal{C}$. Il faut montrer que la restriction du
morphisme à $\mathcal{C}/U$ est un isomorphisme pour tout
$U \in \Ob(\mathcal{C})$. On peut remplacer $U$ par les membres
d'un recouvrement de $U$. D'après Sites, lemme
\ref{sites-lemma-morphism-of-sites-covering}, on peut donc supposer qu'il
existe un morphisme $U \to u(V)$ pour un certain
$V \in \Ob(\mathcal{D})$. On peut alors remplacer $U$ par $u(V)$.
Puisque $u$ est continu, on peut remplacer $V$ par les membres d'un
recouvrement et supposer qu'il existe une présentation
$\mathcal{O}_V^{\oplus m} \to \mathcal{O}_V^{\oplus n} \to
\mathcal{F}|_V \to 0$
sur $\mathcal{D}/V$. Comme la formation de $\SheafHom$ commute à la
localisation (lemme \ref{lemma-internal-hom-restriction}), on peut remplacer
$f$ par le morphisme
$(\mathcal{C}/u(V), \mathcal{O}_{u(V)}) \to
(\mathcal{D}/V, \mathcal{O}_V)$ induit par $f$. On se ramène ainsi au cas
où $\mathcal{F}$ admet une présentation globale
$\mathcal{O}_\mathcal{D}^{\oplus m} \to \mathcal{O}_\mathcal{D}^{\oplus n} \to
\mathcal{F} \to 0$. Puisque $f$ est plat et que
$f^*\mathcal{O}_\mathcal{D} = \mathcal{O}_\mathcal{C}$, on obtient la
présentation correspondante
$\mathcal{O}_\mathcal{C}^{\oplus m} \to \mathcal{O}_\mathcal{C}^{\oplus n} \to
f^*\mathcal{F} \to 0$ ; voir le lemme
\ref{lemma-flat-pullback-exact}. En utilisant la commutation de
$\SheafHom$ aux sommes directes finies dans la première variable, le fait
que les deux foncteurs
$\SheafHom_{\mathcal{O}_\mathcal{C}}(\mathcal{O}_\mathcal{C}, -)$ et
$\SheafHom_{\mathcal{O}_\mathcal{D}}(\mathcal{O}_\mathcal{D}, -)$
sont le foncteur identité, et la fonctorialité de la construction de la
remarque \ref{remark-pullback-internal-hom}, on obtient un diagramme
commutatif
$$
\xymatrix{
0 \ar[r] &
f^*\SheafHom_{\mathcal{O}_\mathcal{D}}(\mathcal{F}, \mathcal{G})
\ar[d] \ar[r] &
f^*\mathcal{G}^{\oplus n} \ar[d] \ar[r] &
f^*\mathcal{G}^{\oplus m} \ar[d] \\
0 \ar[r] &
\SheafHom_{\mathcal{O}_\mathcal{C}}(f^*\mathcal{F}, f^*\mathcal{G})
\ar[r] &
f^*\mathcal{G}^{\oplus n} \ar[r] &
f^*\mathcal{G}^{\oplus m}
}
$$
où les deux flèches verticales de droite sont des isomorphismes. Les lignes
sont exactes d'après le lemme \ref{lemma-internal-hom-exact}. Le lemme des
cinq permet de conclure.

\end{proof}







\section{Modules inversibles}

\label{section-invertible}

\noindent
Voici la définition.

\begin{definition}

\label{definition-invertible-sheaf}

Soit $(\mathcal{C}, \mathcal{O})$ un site annelé.

\begin{enumerate}

\item Un $\mathcal{O}$-module $\mathcal{F}$ localement libre de rang fini
est dit de {\it rang $r$} si, pour tout objet $U$ de $\mathcal{C}$, il
existe un recouvrement $\{U_i \to U\}$ de $U$ tel que
$\mathcal{F}|_{U_i}$ soit isomorphe à
$\mathcal{O}_{U_i}^{\oplus r}$ comme $\mathcal{O}_{U_i}$-module.
\item Un $\mathcal{O}$-module $\mathcal{L}$ est dit {\it inversible}
si le foncteur
$$
\textit{Mod}(\mathcal{O}) \longrightarrow \textit{Mod}(\mathcal{O}),\quad
\mathcal{F}  \longmapsto \mathcal{F} \otimes_\mathcal{O} \mathcal{L}
$$
est une équivalence.
\item Le faisceau {\it $\mathcal{O}^*$} est le sous-faisceau de
$\mathcal{O}$ défini par
$$
U \longmapsto \mathcal{O}^*(U) = \{f \in \mathcal{O}(U) \mid
\exists g \in \mathcal{O}(U)\text{ tel que }fg = 1\}
$$
C'est un faisceau de groupes abéliens pour la multiplication.

\end{enumerate}

\end{definition}

\noindent
Le lemme \ref{lemma-invertible-is-locally-free-rank-1} ci-dessous précise
le lien avec les modules localement libres de rang $1$.

\begin{lemma}

\label{lemma-invertible}

Soit $(\mathcal{C}, \mathcal{O})$ un site annelé et soit $\mathcal{L}$ un
$\mathcal{O}$-module. Les conditions suivantes sont équivalentes :

\begin{enumerate}

\item $\mathcal{L}$ est inversible ;
\item il existe un $\mathcal{O}$-module $\mathcal{N}$
tel que
$\mathcal{L} \otimes_\mathcal{O} \mathcal{N} \cong \mathcal{O}$.

\end{enumerate}

Dans ce cas, nous avons

\begin{enumerate}

\item[(a)] $\mathcal{L}$ est un $\mathcal{O}$-module plat de présentation
finie ;
\item[b)] pour tout objet $U$ de $\mathcal{C}$, il existe un recouvrement
$\{U_i \to U\}$ tel que $\mathcal{L}|_{U_i}$ soit facteur direct d'un
module libre de type fini ;
\item[c)] le module $\mathcal{N}$ en (2) est isomorphe à
$\SheafHom_\mathcal{O}(\mathcal{L}, \mathcal{O})$.

\end{enumerate}

\end{lemma}

\begin{proof}
Supposons (1). Le foncteur $- \otimes_\mathcal{O} \mathcal{L}$ est
essentiellement surjectif ; il existe donc un $\mathcal{O}$-module
$\mathcal{N}$ comme en (2). Réciproquement, si (2) est satisfaite, le
foncteur $- \otimes_\mathcal{O} \mathcal{N}$ est un quasi-inverse du
foncteur $- \otimes_\mathcal{O} \mathcal{L}$, ce qui donne (1).

\medskip\noindent

Supposons (1) et (2). Puisque $- \otimes_\mathcal{O} \mathcal{L}$ est une
équivalence, ce foncteur est exact et $\mathcal{L}$ est donc plat. Soit
$\psi : \mathcal{L} \otimes_\mathcal{O} \mathcal{N} \to \mathcal{O}$
l'isomorphisme donné, et soit $U$ un objet de $\mathcal{C}$. Nous allons
montrer que la restriction de $\mathcal{L}$ aux membres d'un recouvrement de
$U$ est facteur direct d'un module libre ; cela impliquera en
particulier que $\mathcal{L}$ est de présentation finie. Par construction
du produit $\otimes$, après avoir
remplacé $U$ par les membres d'un recouvrement, on peut supposer qu'il
existe un entier $n \geq 1$, des sections $x_i \in \mathcal{L}(U)$ et
$y_i \in \mathcal{N}(U)$ telles que
$\psi(\sum x_i \otimes y_i) = 1$. Considérons les isomorphismes
$$
\mathcal{L}|_U \to
\mathcal{L}|_U \otimes_{\mathcal{O}_U}
\mathcal{L}|_U \otimes_{\mathcal{O}_U} \mathcal{N}|_U \to \mathcal{L}|_U
$$
où la première flèche envoie $x$ sur
$\sum x_i \otimes x \otimes y_i$ et la seconde envoie
$x \otimes x' \otimes y$ sur $\psi(x' \otimes y)x$. On en déduit que
$x \mapsto \sum \psi(x \otimes y_i)x_i$ est un automorphisme de
$\mathcal{L}|_U$. Cet automorphisme se factorise comme
$$
\mathcal{L}|_U \to \mathcal{O}_U^{\oplus n} \to \mathcal{L}|_U
$$
où la première flèche est donnée
$x \mapsto (\psi(x \otimes y_1), \ldots, \psi(x \otimes y_n))$
et la seconde par $(a_1, \ldots, a_n) \mapsto \sum a_i x_i$. Ainsi,
$\mathcal{L}|_U$ est facteur direct d'un $\mathcal{O}_U$-module libre de
rang fini.

\medskip\noindent

Toujours sous les hypothèses (1) et (2), considérons le morphisme
d'évaluation
$$
\mathcal{L} \otimes_\mathcal{O}
\SheafHom_\mathcal{O}(\mathcal{L}, \mathcal{O}_X)
\longrightarrow \mathcal{O}_X
$$
Pour achever la démonstration, montrons qu'il s'agit d'un isomorphisme. Le
lemme \ref{lemma-internal-hom-adjoint-tensor} donne
$$
\Hom_\mathcal{O}(\mathcal{O}, \mathcal{O}) =
\Hom_\mathcal{O}
(\mathcal{N} \otimes_\mathcal{O} \mathcal{L}, \mathcal{O})
\longrightarrow
\Hom_\mathcal{O}
(\mathcal{N}, \SheafHom_\mathcal{O}(\mathcal{L}, \mathcal{O}))
$$
L'image de $1$ définit un morphisme
$\mathcal{N} \to \SheafHom_\mathcal{O}(\mathcal{L}, \mathcal{O})$. En
prenant le produit tensoriel avec $\mathcal{L}$, on obtient
$$
\mathcal{O} = \mathcal{L} \otimes_\mathcal{O} \mathcal{N}
\longrightarrow
\mathcal{L} \otimes_\mathcal{O} \SheafHom_\mathcal{O}(\mathcal{L}, \mathcal{O})
$$
Ce morphisme est l'inverse du morphisme d'évaluation ; nous omettons le
calcul.

\end{proof}

\begin{lemma}

\label{lemma-pullback-invertible}

Soit $f : (\Sh(\mathcal{C}), \mathcal{O}_\mathcal{C}) \to
(\Sh(\mathcal{D}), \mathcal{O}_\mathcal{D})$ un morphisme de topos
annelés. L'image inverse $f^*\mathcal{L}$ d'un
$\mathcal{O}_\mathcal{D}$-module inversible est inversible.

\end{lemma}

\begin{proof}
D'après le lemme \ref{lemma-invertible}, il existe un
$\mathcal{O}_\mathcal{D}$-module $\mathcal{N}$ tel que
$\mathcal{L} \otimes_{\mathcal{O}_\mathcal{D}} \mathcal{N} \cong
\mathcal{O}_\mathcal{D}$. Par image inverse, on obtient
$f^*\mathcal{L} \otimes_{\mathcal{O}_\mathcal{C}} f^*\mathcal{N} \cong
\mathcal{O}_\mathcal{C}$ d'après le lemme
\ref{lemma-tensor-product-pullback}. Le lemme \ref{lemma-invertible} montre
donc que $f^*\mathcal{L}$ est inversible.
\end{proof}

\begin{lemma}

\label{lemma-constructions-invertible}
Soit $(\mathcal{C}, \mathcal{O})$ un site annelé.
\begin{enumerate}
\item Si $\mathcal{L}$ et $\mathcal{N}$ sont des $\mathcal{O}$-modules
inversibles, alors $\mathcal{L} \otimes_\mathcal{O} \mathcal{N}$ est
inversible.
\item Si $\mathcal{L}$ est un $\mathcal{O}$-module inversible, alors
$\SheafHom_\mathcal{O}(\mathcal{L}, \mathcal{O})$ est inversible et le
morphisme d'évaluation $\mathcal{L} \otimes_\mathcal{O}
\SheafHom_\mathcal{O}(\mathcal{L}, \mathcal{O}) \to \mathcal{O}$ est un isomorphisme.
\end{enumerate}

\end{lemma}

\begin{proof}

La partie (1) est immédiate d'après la définition, et la partie (2) découle
du lemme \ref{lemma-invertible} et de sa démonstration.

\end{proof}

\begin{lemma}

\label{lemma-pic-set}

Soit $(\mathcal{C}, \mathcal{O})$ un site annelé. Il existe un ensemble de
modules inversibles $\{\mathcal{L}_i\}_{i \in I}$ tel que tout module
inversible sur $(\mathcal{C}, \mathcal{O})$
est isomorphe à exactement l'un des $\mathcal{L}_i$.

\end{lemma}

\begin{proof}
Démonstration omise ; voir Faisceaux de modules, lemme
\ref{modules-lemma-pic-set}.
\end{proof}

\noindent
Le lemme \ref{lemma-pic-set} affirme que les classes d'isomorphisme des
faisceaux inversibles forment un ensemble. Le lemme
\ref{lemma-constructions-invertible} montre que le produit tensoriel munit
cet ensemble d'une structure de groupe abélien ; l'inverse de
$\mathcal{L}$ y est représenté par
$\SheafHom_\mathcal{O}(\mathcal{L}, \mathcal{O})$.

\medskip\noindent
Plus précisément, pour un $\mathcal{O}$-module inversible $\mathcal{L}$ et
$n \in \mathbf{Z}$, on définit la {\it $n$-ième puissance tensorielle}
$\mathcal{L}^{\otimes n}$ de $\mathcal{L}$ comme l'image de $\mathcal{O}$
par l'équivalence $\mathcal{F} \mapsto
\mathcal{F} \otimes_\mathcal{O} \mathcal{L}$ itérée $n$ fois. Cette
définition a également un sens pour $n$ négatif, puisqu'un
$\mathcal{O}$-module est inversible précisément lorsque le produit tensoriel
définit une équivalence.
Plus explicitement :
$$
\mathcal{L}^{\otimes n} =
\left\{
\begin{matrix}
\mathcal{O} & \text{si} & n = 0 \\
\SheafHom_\mathcal{O}(\mathcal{L}, \mathcal{O}) & \text{si} & n = -1\\
\mathcal{L} \otimes_\mathcal{O} \ldots \otimes_\mathcal{O} \mathcal{L}
& \text{si} & n > 0 \\
\mathcal{L}^{\otimes -1} \otimes_\mathcal{O} \ldots
\otimes_\mathcal{O} \mathcal{L}^{\otimes -1}
& \text{si} & n < -1
\end{matrix}
\right.
$$
voir le lemme \ref{lemma-constructions-invertible}. Avec cette définition,
on dispose d'isomorphismes canoniques
$\mathcal{L}^{\otimes n} \otimes_\mathcal{O}
\mathcal{L}^{\otimes m} \to
\mathcal{L}^{\otimes n + m}$ satisfaisant aux contraintes de commutativité
et d'associativité ; nous en omettons la formulation.

\begin{definition}

\label{definition-pic}
Soit $(\mathcal{C}, \mathcal{O})$ un site annelé. Le {\it groupe de Picard}
$\Pic(\mathcal{O})$ de ce site annelé est le groupe abélien des classes
d'isomorphisme de $\mathcal{O}$-modules inversibles, dont l'addition
correspond au produit tensoriel.
\end{definition}









\section{Modules de différentielles}

\label{section-differentials}

\noindent
Dans cette section, nous expliquons brièvement comment définir le module des
différentielles relatives d'un morphisme de topos annelés. Nous invitons le
lecteur à consulter la section correspondante du chapitre d'algèbre
commutative (Algèbre, section \ref{algebra-section-differentials}).

\begin{definition}

\label{definition-derivation}

Soit $\mathcal{C}$ un site. Soit
$\varphi : \mathcal{O}_1 \to \mathcal{O}_2$ un morphisme de faisceaux
d'anneaux, et soit $\mathcal{F}$ un $\mathcal{O}_2$-module. Une
{\it $\mathcal{O}_1$-dérivation}, ou plus précisément une
{\it $\varphi$-dérivation} à valeurs dans $\mathcal{F}$, est une application
$D : \mathcal{O}_2 \to \mathcal{F}$ additive, qui annule l'image de
$\mathcal{O}_1 \to \mathcal{O}_2$ et satisfait la {\it règle de Leibniz}

$$
D(ab) = aD(b) + D(a)b
$$
pour toutes sections locales $a, b$ de $\mathcal{O}_2$, là où elles sont
toutes deux définies. On note
$\text{Der}_{\mathcal{O}_1}(\mathcal{O}_2, \mathcal{F})$
l'ensemble des $\varphi$-dérivations à valeurs dans $\mathcal{F}$.

\end{definition}

\noindent
C'est l'analogue, pour les faisceaux, d'Algèbre, définition
\ref{definition-derivation}. Pour une dérivation
$D : \mathcal{O}_2 \to \mathcal{F}$ comme dans la définition, l'application
sur les sections globales
$$
D : \Gamma(\mathcal{O}_2) \longrightarrow \Gamma(\mathcal{F})
$$
est clairement une $\Gamma(\mathcal{O}_1)$-dérivation au sens algébrique.
Si $\alpha : \mathcal{F} \to \mathcal{G}$ est un morphisme de
$\mathcal{O}_2$-modules, on obtient un morphisme induit
$$
\text{Der}_{\mathcal{O}_1}(\mathcal{O}_2, \mathcal{F})
\longrightarrow
\text{Der}_{\mathcal{O}_1}(\mathcal{O}_2, \mathcal{G})
$$
donné par $D \mapsto \alpha \circ D$. On obtient ainsi un foncteur.

\begin{lemma}

\label{lemma-universal-module}

Soit $\mathcal{C}$ un site et soit
$\varphi : \mathcal{O}_1 \to \mathcal{O}_2$ un morphisme de faisceaux
d'anneaux. Le foncteur
$$
\textit{Mod}(\mathcal{O}_2) \longrightarrow \textit{Ab}, \quad
\mathcal{F} \longmapsto \text{Der}_{\mathcal{O}_1}(\mathcal{O}_2, \mathcal{F})
$$
est représentable.

\end{lemma}

\begin{proof}

La démonstration est identique à celle de l'énoncé analogue en algèbre.
Au cours de celle-ci, pour tout faisceau d'ensembles $\mathcal{F}$ sur
$\mathcal{C}$, notons $\mathcal{O}_2[\mathcal{F}]$ la faisceautisation du
préfaisceau $U \mapsto \mathcal{O}_2(U)[\mathcal{F}(U)]$, où le membre de
droite désigne le $\mathcal{O}_2(U)$-module libre sur l'ensemble
$\mathcal{F}(U)$. Pour $s \in \mathcal{F}(U)$, notons $[s]$ la section
correspondante de $\mathcal{O}_2[\mathcal{F}]$ au-dessus de $U$. Si
$\mathcal{F}$ est un faisceau de $\mathcal{O}_2$-modules, il existe un
morphisme canonique
$$
c : \mathcal{O}_2[\mathcal{F}] \longrightarrow \mathcal{F}
$$
qui, au niveau des préfaisceaux, est donné par la règle
$\sum f_s[s] \mapsto \sum f_s s$. Nous emploierons l'abréviation
$[s] \mapsto s$ pour décrire ce morphisme, et procéderons de même pour les
autres morphismes ci-dessous. Considérons le morphisme de
$\mathcal{O}_2$-modules

\begin{equation}
\label{equation-define-module-differentials}
\begin{matrix}
\mathcal{O}_2[\mathcal{O}_2 \times  \mathcal{O}_2] \oplus
\mathcal{O}_2[\mathcal{O}_2 \times  \mathcal{O}_2] \oplus
\mathcal{O}_2[\mathcal{O}_1] &
\longrightarrow &
\mathcal{O}_2[\mathcal{O}_2] \\
[(a, b)] \oplus [(f, g)] \oplus [h] & \longmapsto & [a + b] - [a] - [b] + \\
& & [fg] - g[f] - f[g] + \\
& & [\varphi(h)]
\end{matrix}
\end{equation}

où nous utilisons la notation abrégée précédente. Posons
$\Omega_{\mathcal{O}_2/\mathcal{O}_1}$ égal au conoyau de ce morphisme. Il
existe alors manifestement un morphisme de faisceaux d'ensembles
$$
\text{d} : \mathcal{O}_2 \longrightarrow \Omega_{\mathcal{O}_2/\mathcal{O}_1}
$$
qui envoie une section locale $f$ sur l'image de $[f]$ dans
$\Omega_{\mathcal{O}_2/\mathcal{O}_1}$. Par construction, $\text{d}$ est une
$\mathcal{O}_1$-dérivation. Soit maintenant $\mathcal{F}$ un faisceau de
$\mathcal{O}_2$-modules et soit $D : \mathcal{O}_2 \to \mathcal{F}$ une
$\mathcal{O}_1$-dérivation. On peut considérer le morphisme
$\mathcal{O}_2$-linéaire $\mathcal{O}_2[\mathcal{O}_2] \to \mathcal{F}$ qui
envoie $[g]$ sur $D(g)$. Par définition d'une dérivation, ce morphisme annule
les sections appartenant à l'image du morphisme
(\ref{equation-define-module-differentials}) et induit donc un morphisme
$$
\alpha_D : \Omega_{\mathcal{O}_2/\mathcal{O}_1} \longrightarrow \mathcal{F}
$$
Comme $D = \alpha_D \circ \text{d}$, le lemme est démontré.

\end{proof}

\begin{definition}

\label{definition-module-differentials}
Soit $\mathcal{C}$ un site. Soit
$\varphi : \mathcal{O}_1 \to \mathcal{O}_2$ un morphisme de faisceaux
d'anneaux. Le {\it module des différentielles} du morphisme d'anneaux
$\varphi$ est l'objet représentant le foncteur
$\mathcal{F} \mapsto \text{Der}_{\mathcal{O}_1}(\mathcal{O}_2, \mathcal{F})$,
dont l'existence résulte du lemme \ref{lemma-universal-module}. On le note
$\Omega_{\mathcal{O}_2/\mathcal{O}_1}$, et la {\it $\varphi$-dérivation
universelle} est notée
$\text{d} : \mathcal{O}_2 \to \Omega_{\mathcal{O}_2/\mathcal{O}_1}$.
\end{definition}

\noindent
Comme ce module et cette dérivation forment l'objet universel représentant un
foncteur, la notion est intrinsèque : elle ne dépend pas du choix du site
sous-jacent au topos annelé ; voir la section \ref{section-intrinsic}. Le
module $\Omega_{\mathcal{O}_2/\mathcal{O}_1}$ est le conoyau du morphisme
(\ref{equation-define-module-differentials}) de $\mathcal{O}_2$-modules. De plus,
$\text{d}$ est caractérisée par le fait que $\text{d}f$ est l'image de la
section locale $[f]$.

\begin{lemma}

\label{lemma-differentials-sheafify}
Soit $\mathcal{C}$ un site. Soit
$\varphi : \mathcal{O}_1 \to \mathcal{O}_2$ un morphisme de préfaisceaux
d'anneaux. Alors $\Omega_{\mathcal{O}_2^\#/\mathcal{O}_1^\#}$ est le
faisceau associé au préfaisceau
$U \mapsto \Omega_{\mathcal{O}_2(U)/\mathcal{O}_1(U)}$.
\end{lemma}

\begin{proof}
Considérons le morphisme (\ref{equation-define-module-differentials}). Il
existe un morphisme analogue de préfaisceaux dont la valeur en
$U \in \Ob(\mathcal{C})$ est
$$
\mathcal{O}_2(U)[\mathcal{O}_2(U) \times \mathcal{O}_2(U)] \oplus
\mathcal{O}_2(U)[\mathcal{O}_2(U) \times \mathcal{O}_2(U)] \oplus
\mathcal{O}_2(U)[\mathcal{O}_1(U)]
\longrightarrow
\mathcal{O}_2(U)[\mathcal{O}_2(U)]
$$
Le conoyau de ce morphisme vaut
$\Omega_{\mathcal{O}_2(U)/\mathcal{O}_1(U)}$ en $U$, par la construction du
module des différentielles dans Algèbre, définition
\ref{algebra-definition-differentials}. D'autre part, les faisceaux de
(\ref{equation-define-module-differentials}) sont les faisceautisations des
préfaisceaux ci-dessus. Le résultat découle donc de l'exactitude de la
faisceautisation.
\end{proof}

\begin{lemma}

\label{lemma-pullback-differentials}

Soit $f : \Sh(\mathcal{D}) \to \Sh(\mathcal{C})$ un morphisme de topos. Soit
$\varphi : \mathcal{O}_1 \to \mathcal{O}_2$ un morphisme de faisceaux
d'anneaux sur $\mathcal{C}$. Il existe alors une identification canonique
$f^{-1}\Omega_{\mathcal{O}_2/\mathcal{O}_1} =
\Omega_{f^{-1}\mathcal{O}_2/f^{-1}\mathcal{O}_1}$
compatible avec les dérivations universelles.

\end{lemma}

\begin{proof}
Cela résulte du fait que le faisceau
$\Omega_{\mathcal{O}_2/\mathcal{O}_1}$ est le conoyau du morphisme
(\ref{equation-define-module-differentials}), qu'il en va de même pour
$\Omega_{f^{-1}\mathcal{O}_2/f^{-1}\mathcal{O}_1}$, que le foncteur
$f^{-1}$ est exact et que $f^{-1}(\mathcal{O}_2[\mathcal{O}_2]) =
f^{-1}\mathcal{O}_2[f^{-1}\mathcal{O}_2]$, $f^{-1}(\mathcal{O}_2[\mathcal{O}_2 \times \mathcal{O}_2]) =
f^{-1}\mathcal{O}_2[f^{-1}\mathcal{O}_2 \times f^{-1}\mathcal{O}_2]$, et $f^{-1}(\mathcal{O}_2[\mathcal{O}_1]) =
f^{-1}\mathcal{O}_2[f^{-1}\mathcal{O}_1]$.
\end{proof}

\begin{lemma}

\label{lemma-localize-differentials}

Soit $\mathcal{C}$ un site. Soit
$\varphi : \mathcal{O}_1 \to \mathcal{O}_2$ un morphisme de faisceaux
d'anneaux. Pour tout objet $U$ de $\mathcal{C}$, il existe un isomorphisme
canonique
$$
\Omega_{\mathcal{O}_2/\mathcal{O}_1}|_U =
\Omega_{(\mathcal{O}_2|_U)/(\mathcal{O}_1|_U)}
$$
compatible avec les dérivations universelles.

\end{lemma}

\begin{proof}
C'est un cas particulier du lemme \ref{lemma-pullback-differentials}.
\end{proof}

\begin{lemma}

\label{lemma-functoriality-differentials}

Soit $\mathcal{C}$ un site. Soit
$$
\xymatrix{
\mathcal{O}_2 \ar[r]_\varphi & \mathcal{O}_2' \\
\mathcal{O}_1 \ar[r] \ar[u] & \mathcal{O}'_1 \ar[u]
}
$$
un diagramme commutatif de faisceaux d'anneaux sur $\mathcal{C}$. Le morphisme
$\mathcal{O}_2 \to \mathcal{O}'_2$, composé avec le morphisme
$\text{d} : \mathcal{O}'_2 \to \Omega_{\mathcal{O}'_2/\mathcal{O}'_1}$
est une $\mathcal{O}_1$-dérivation. On obtient donc un morphisme canonique de
$\mathcal{O}_2$-modules
$\Omega_{\mathcal{O}_2/\mathcal{O}_1} \to
\Omega_{\mathcal{O}'_2/\mathcal{O}'_1}$. Il est uniquement caractérisé par
la propriété que $\text{d}(f)$ s'envoie sur $\text{d}(\varphi(f))$ pour toute
section locale $f$ de $\mathcal{O}_2$. De cette manière, $\Omega_{-/-}$
devient un foncteur sur la catégorie des flèches de faisceaux d'anneaux.

\end{lemma}

\begin{proof}
Le lemme découle directement des définitions.
\end{proof}

\begin{lemma}

\label{lemma-differential-seq}
Dans le lemme \ref{lemma-functoriality-differentials}, supposons que
$\mathcal{O}_2 \to \mathcal{O}'_2$ soit surjectif, de noyau
$\mathcal{I} \subset \mathcal{O}_2$, et que
$\mathcal{O}_1 = \mathcal{O}'_1$. Il existe alors une suite exacte canonique
de $\mathcal{O}'_2$-modules
$$
\mathcal{I}/\mathcal{I}^2
\longrightarrow
\Omega_{\mathcal{O}_2/\mathcal{O}_1} \otimes_{\mathcal{O}_2} \mathcal{O}'_2
\longrightarrow
\Omega_{\mathcal{O}'_2/\mathcal{O}_1}
\longrightarrow
0
$$
Le morphisme de gauche est caractérisé par le fait qu'une section locale $f$
de $\mathcal{I}$ s'envoie sur $\text{d}f \otimes 1$.
\end{lemma}

\begin{proof}

Pour une section locale $f$ de $\mathcal{I}$, notons $\overline{f}$ l'image de
$f$ dans $\mathcal{I}/\mathcal{I}^2$. Pour montrer que le morphisme
$\overline{f} \mapsto \text{d}f \otimes 1$ est bien défini, il suffit de
vérifier que $\text{d} f_1f_2 \otimes 1 = 0$ lorsque $f_1, f_2$ sont des
sections locales de $\mathcal{I}$. Cela résulte immédiatement de la règle de
Leibniz :
$\text{d} f_1f_2 \otimes 1 =
(f_1 \text{d}f_2 + f_2 \text{d} f_1 )\otimes 1 =
\text{d}f_2 \otimes f_1 + \text{d}f_1 \otimes f_2 = 0$. Un calcul analogue
montre que ce morphisme est $\mathcal{O}'_2 = \mathcal{O}_2/\mathcal{I}$-linéaire.
Le morphisme de droite est celui du lemme
\ref{lemma-functoriality-differentials}.

\medskip\noindent
Pour établir l'exactitude, raisonnons comme suit. Soit
$\mathcal{O}''_2 \subset \mathcal{O}'_2$ le préfaisceau de
$\mathcal{O}_1$-algèbres dont la valeur en $U$ est l'image de
$\mathcal{O}_2(U) \to \mathcal{O}'_2(U)$. D'après Algèbre, lemme
\ref{algebra-lemma-differential-seq}, les suites
$$
\mathcal{I}(U)/\mathcal{I}(U)^2
\longrightarrow
\Omega_{\mathcal{O}_2(U)/\mathcal{O}_1(U)}
\otimes_{\mathcal{O}_2(U)} \mathcal{O}''_2(U)
\longrightarrow
\Omega_{\mathcal{O}''_2(U)/\mathcal{O}_1(U)}
\longrightarrow
0
$$
sont exactes pour tout objet $U$ de $\mathcal{C}$. L'exactitude de la
faisceautisation donne donc une suite exacte de faisceaux de
$(\mathcal{O}'_2)^\#$-modules. On conclut par le lemme
\ref{lemma-differentials-sheafify} et l'égalité
$(\mathcal{O}''_2)^\# = \mathcal{O}'_2$.
\end{proof}

\noindent
Voici une situation particulière où les dérivations apparaissent naturellement.

\begin{lemma}

\label{lemma-double-structure-gives-derivation}

Soit $\mathcal{C}$ un site. Soit
$\varphi : \mathcal{O}_1 \to \mathcal{O}_2$ un morphisme de faisceaux
d'anneaux. Considérons une suite exacte courte
$$
0 \to \mathcal{F} \to \mathcal{A} \to \mathcal{O}_2 \to 0
$$
Ici, $\mathcal{A}$ est un faisceau de $\mathcal{O}_1$-algèbres,
$\pi : \mathcal{A} \to \mathcal{O}_2$ est une surjection de faisceaux de
$\mathcal{O}_1$-algèbres et $\mathcal{F} = \Ker(\pi)$ est son noyau.
Supposons que $\mathcal{F}$ soit un faisceau d'idéaux de carré nul dans
$\mathcal{A}$. Le faisceau $\mathcal{F}$ possède alors une structure naturelle de
$\mathcal{O}_2$-module. Une section $s : \mathcal{O}_2 \to \mathcal{A}$ de
$\pi$ est un morphisme de $\mathcal{O}_1$-algèbres tel que
$\pi \circ s = \text{id}$. Étant données une section quelconque
$s : \mathcal{O}_2 \to \mathcal{A}$ de $\pi$ et une $\varphi$-dérivation
$D : \mathcal{O}_2 \to \mathcal{F}$, le morphisme
$$
s + D : \mathcal{O}_2 \to \mathcal{A}
$$
est une section de $\pi$, et toute section $s'$ s'écrit de manière unique
$s + D$ pour une $\varphi$-dérivation $D$.

\end{lemma}

\begin{proof}

Rappelons que la structure de $\mathcal{O}_2$-module sur $\mathcal{F}$ est
donnée par $h \tau = \tilde h \tau$ (le produit étant pris dans
$\mathcal{A}$), où $h$ est une section locale de $\mathcal{O}_2$,
$\tilde h$ un relèvement local de $h$ en une section locale de $\mathcal{A}$,
et $\tau$ une section locale de $\mathcal{F}$. En présence de $s$, on peut
prendre $\tilde h = s(h)$. Pour vérifier que $s + D$ est un morphisme de
faisceaux d'anneaux, calculons

\begin{eqnarray*}
(s + D)(ab) & = & s(ab) + D(ab) \\
& = & s(a)s(b) + aD(b) + D(a)b \\
& = & s(a) s(b) + s(a)D(b) + D(a)s(b) \\
& = & (s(a) + D(a))(s(b) + D(b))
\end{eqnarray*}

par la règle de Leibniz. De même, $s + D$ est un morphisme de
$\mathcal{O}_1$-algèbres parce que $D$ est une $\mathcal{O}_1$-dérivation.
Réciproquement, pour une section $s'$, posons $D = s' - s$. Les détails sont
laissés au lecteur.

\end{proof}

\begin{definition}

\label{definition-sheaf-differentials}

Soient $X = (\Sh(\mathcal{C}), \mathcal{O})$ et
$Y = (\Sh(\mathcal{C}'), \mathcal{O}')$ des topos annelés. Soit
$(f, f^\sharp) : X \to Y$ un morphisme de topos annelés. Dans cette situation,

\begin{enumerate}

\item pour un faisceau $\mathcal{F}$ de $\mathcal{O}$-modules, une
{\it $Y$-dérivation} $D : \mathcal{O} \to \mathcal{F}$ est simplement une
$f^\sharp$-dérivation, et
\item le {\it faisceau des différentielles $\Omega_{X/Y}$ de $X$ sur $Y$}
est le module des différentielles du morphisme
$f^\sharp : f^{-1}\mathcal{O}' \to \mathcal{O}$, voir définition \ref{definition-module-differentials}.

\end{enumerate}

Ainsi, $\Omega_{X/Y}$ est muni d'une {\it dérivation $Y$-universelle}

$\text{d}_{X/Y} : \mathcal{O} \longrightarrow \Omega_{X/Y}$. Nous écrivons parfois $\Omega_{X/Y} = \Omega_f$.

\end{definition}

\noindent
Rappelons que $f^\sharp : f^{-1}\mathcal{O}' \to \mathcal{O}$, ce qui donne
un sens à cette définition.

\begin{lemma}

\label{lemma-functoriality-differentials-ringed-topoi}

Soient
$X = (\Sh(\mathcal{C}_X), \mathcal{O}_X)$,
$Y = (\Sh(\mathcal{C}_Y), \mathcal{O}_Y)$,
$X' = (\Sh(\mathcal{C}_{X'}), \mathcal{O}_{X'})$,
$Y' = (\Sh(\mathcal{C}_{Y'}), \mathcal{O}_{Y'})$ des topos annelés. Soit
$$
\xymatrix{
X' \ar[d] \ar[r]_f & X \ar[d] \\
Y' \ar[r] & Y
}
$$
un diagramme commutatif de morphismes de topos annelés. Le morphisme
$f^\sharp : \mathcal{O}_X \to f_*\mathcal{O}_{X'}$, composé avec le morphisme
$f_*\text{d}_{X'/Y'} : f_*\mathcal{O}_{X'} \to f_*\Omega_{X'/Y'}$ est un
$Y$-dérivation. On obtient donc un morphisme canonique de
$\mathcal{O}_X$-modules $\Omega_{X/Y} \to f_*\Omega_{X'/Y'}$ et, par
adjonction entre $f_*$ et $f^*$, un homomorphisme canonique de
$\mathcal{O}_{X'}$-modules
$$
c_f : f^*\Omega_{X/Y} \longrightarrow \Omega_{X'/Y'}.
$$
Il est uniquement caractérisé par la propriété que
$f^*\text{d}_{X/Y}(t)$ s'envoie sur $\text{d}_{X'/Y'}(f^* t)$ pour toute
section locale $t$ de $\mathcal{O}_X$.

\end{lemma}

\begin{proof}

Tout est clair, sauf la dernière assertion. Précisons son sens. Soient
$U \in \Ob(\mathcal{C}_X)$ et $t \in \mathcal{O}_X(U)$ ; c'est en ce sens que
$t$ est une section locale de $\mathcal{O}_X$. On peut voir $t$ comme un
morphisme de faisceaux d'ensembles $t : h_U^\# \to \mathcal{O}_X$. On a alors
$f^{-1}t : f^{-1}h_U^\# \to f^{-1}\mathcal{O}_X$. Par $f^*t$, on entend le
composé
$$
\xymatrix{
f^{-1}h_U^\# \ar[rr]^{f^{-1}t} \ar@/^4ex/[rrrr]^{f^*t} & &
f^{-1}\mathcal{O}_X \ar[rr]^{f^\sharp} & &
\mathcal{O}_{X'}
}
$$
Notons que $\text{d}_{X/Y}(t) \in \Omega_{X/Y}(U)$. On peut donc voir
$\text{d}_{X/Y}(t)$ comme un morphisme
$\text{d}_{X/Y}(t) : h_U^\# \to \Omega_{X/Y}$. On a alors
$f^{-1}\text{d}_{X/Y}(t) : f^{-1}h_U^\# \to f^{-1}\Omega_{X/Y}$. Par
$f^*\text{d}_{X/Y}(t)$, on entend le composé
$$
\xymatrix{
f^{-1}h_U^\#
\ar[rr]^{f^{-1}\text{d}_{X/Y}(t)}
\ar@/^4ex/[rrrr]^{f^*\text{d}_{X/Y}(t)} & &
f^{-1}\Omega_{X/Y} \ar[rr]^{1 \otimes \text{id}} & &
f^*\Omega_{X/Y}
}
$$
L'énoncé du lemme signifie alors que l'on a
$$
c_f \circ f^*\text{d}_{X/Y}(t) = \text{d}_{X'/Y'}(f^*t)
$$
comme morphismes de $f^{-1}h_U^\#$ vers $\Omega_{X'/Y'}$. Nous omettons la
vérification que le morphisme $c_f$ défini dans le lemme satisfait cette
propriété. Indication : le premier morphisme
$c'_f : \Omega_{X/Y} \to f_*\Omega_{X'/Y'}$ satisfait
$c'_f(\text{d}_{X/Y}(t)) = f_*\text{d}_{X'/Y'}(f^\sharp(t))$ en tant que
sections de $f_*\Omega_{X'/Y'}$ au-dessus de $U$ ; il suffit de traduire cette
égalité au moyen de l'adjonction. Cette propriété caractérise uniquement
$c_f$, car les images de $f^*\text{d}_{X/Y}(t)$ engendrent le
$\mathcal{O}_{X'}$-module $f^*\Omega_{X/Y}$, puisque les sections locales
$\text{d}_{X/Y}(t)$ engendrent le $\mathcal{O}_X$-module $\Omega_{X/Y}$.

\end{proof}
\section{Opérateurs différentiels d'ordre fini}
\label{section-differential-operators}

\noindent
Dans cette section, nous introduisons les opérateurs différentiels d'ordre fini.
Nous conseillons au lecteur de consulter la section correspondante
du chapitre consacré à l'algèbre commutative
(Algèbre, section \ref{algebra-section-differential-operators}).

\begin{definition}
\label{definition-differential-operators}
Soit $\mathcal{C}$ un site. Soit $\varphi : \mathcal{O}_1 \to \mathcal{O}_2$
un homomorphisme de faisceaux d'anneaux. Soit $k \geq 0$ un entier.
Soient $\mathcal{F}$ et $\mathcal{G}$ des faisceaux de $\mathcal{O}_2$-modules.
Un {\it opérateur différentiel $D : \mathcal{F} \to \mathcal{G}$ d'ordre $k$}
est une application $\mathcal{O}_1$-linéaire telle que, pour toute section locale
$g$ de $\mathcal{O}_2$, l'application $s \mapsto D(gs) - gD(s)$ soit un
opérateur différentiel d'ordre $k - 1$. Pour le cas initial $k = 0$,
nous définissons un opérateur différentiel d'ordre $0$ comme une application
$\mathcal{O}_2$-linéaire.
\end{definition}

\noindent
Si $D : \mathcal{F} \to \mathcal{G}$ est un opérateur différentiel d'ordre $k$,
alors, pour toute section locale $g$ de $\mathcal{O}_2$, l'application $gD$ est
un opérateur différentiel d'ordre $k$. La somme de deux opérateurs différentiels
d'ordre $k$ en est encore un. Par conséquent, l'ensemble de tous ces opérateurs
$$
\text{Diff}^k(\mathcal{F}, \mathcal{G}) =
\text{Diff}^k_{\mathcal{O}_2/\mathcal{O}_1}(\mathcal{F}, \mathcal{G})
$$
est un $\Gamma(\mathcal{C}, \mathcal{O}_2)$-module. On a
$$
\text{Diff}^0(\mathcal{F}, \mathcal{G}) \subset
\text{Diff}^1(\mathcal{F}, \mathcal{G}) \subset
\text{Diff}^2(\mathcal{F}, \mathcal{G}) \subset \ldots
$$
La règle qui associe à $U \in \Ob(\mathcal{C})$ le module des
opérateurs différentiels $D : \mathcal{F}|_U \to \mathcal{G}|_U$
d'ordre $k$ est un faisceau de $\mathcal{O}_2$-modules sur le site $\mathcal{C}$.
On obtient ainsi un faisceau d'opérateurs différentiels (si nous en avons un
jour besoin, nous ajouterons ici une définition).

\begin{lemma}
\label{lemma-composition-differential-operators}
Soit $\mathcal{C}$ un site.
Soit $\mathcal{O}_1 \to \mathcal{O}_2$ un morphisme de faisceaux d'anneaux.
Soient $\mathcal{E}, \mathcal{F}, \mathcal{G}$ des faisceaux de
$\mathcal{O}_2$-modules.
Si $D : \mathcal{E} \to \mathcal{F}$ et $D' : \mathcal{F} \to \mathcal{G}$
sont des opérateurs différentiels d'ordres $k$ et $k'$, alors $D' \circ D$ est
un opérateur différentiel d'ordre $k + k'$.
\end{lemma}

\begin{proof}
Soit $g$ une section locale de $\mathcal{O}_2$.
L'application qui envoie une section locale $x$ de $\mathcal{E}$ sur
$$
D'(D(gx)) - gD'(D(x)) = D'(D(gx)) - D'(gD(x)) + D'(gD(x)) - gD'(D(x))
$$
est une somme de deux composées d'opérateurs différentiels d'ordre inférieur.
Le lemme en résulte par récurrence sur $k + k'$.
\end{proof}

\begin{lemma}
\label{lemma-module-principal-parts}
Soit $\mathcal{C}$ un site.
Soit $\mathcal{O}_1 \to \mathcal{O}_2$ un morphisme de faisceaux d'anneaux.
Soit $\mathcal{F}$ un faisceau de $\mathcal{O}_2$-modules.
Soit $k \geq 0$. Il existe un faisceau de $\mathcal{O}_2$-modules
$\mathcal{P}^k_{\mathcal{O}_2/\mathcal{O}_1}(\mathcal{F})$
et un isomorphisme canonique
$$
\text{Diff}^k_{\mathcal{O}_2/\mathcal{O}_1}(\mathcal{F}, \mathcal{G}) =
\Hom_{\mathcal{O}_2}(
\mathcal{P}^k_{\mathcal{O}_2/\mathcal{O}_1}(\mathcal{F}), \mathcal{G})
$$
fonctoriel par rapport au $\mathcal{O}_2$-module $\mathcal{G}$.
\end{lemma}

\begin{proof}
L'existence résulte d'arguments généraux de théorie des catégories
(insérer ici une référence ultérieure), mais nous donnerons également une
construction directe, car elle sera utile dans les démonstrations à venir.
Nous emploierons librement les notations introduites dans la démonstration du
lemme \ref{lemma-universal-module}.
Étant donné un opérateur différentiel $D : \mathcal{F} \to \mathcal{G}$,
nous obtenons une application $\mathcal{O}_2$-linéaire
$L_D : \mathcal{O}_2[\mathcal{F}] \to \mathcal{G}$
qui envoie $[m]$ sur $D(m)$. Si $D$ est d'ordre $0$,
alors $L_D$ annule les sections locales
$$
[m + m'] - [m] - [m'],\quad
g_0[m] - [g_0m]
$$
où $g_0$ est une section locale de $\mathcal{O}_2$ et $m, m'$
sont des sections locales de $\mathcal{F}$. Si $D$ est d'ordre $1$, alors $L_D$
annule les sections locales
$$
[m + m'] - [m] - [m'],\quad
f[m] - [fm], \quad
g_0g_1[m] - g_0[g_1m] - g_1[g_0m] + [g_1g_0m]
$$
où $f$ est une section locale de $\mathcal{O}_1$,
$g_0, g_1$ sont des sections locales de $\mathcal{O}_2$, et
$m, m'$ sont des sections locales de $\mathcal{F}$.
Si $D$ est d'ordre $k$, alors $L_D$ annule les sections locales
$[m + m'] - [m] - [m']$, $f[m] - [fm]$, ainsi que les sections locales
$$
g_0g_1\ldots g_k[m] - \sum g_0 \ldots \hat g_i \ldots g_k[g_im] + \ldots
+(-1)^{k + 1}[g_0\ldots g_km]
$$
Réciproquement, si $L : \mathcal{O}_2[\mathcal{F}] \to \mathcal{G}$ est une
application $\mathcal{O}_2$-linéaire qui annule toutes les sections locales
énumérées dans la phrase précédente, alors $m \mapsto L([m])$ est un
opérateur différentiel d'ordre $k$. On voit donc que
$\mathcal{P}^k_{\mathcal{O}_2/\mathcal{O}_1}(\mathcal{F})$
est le quotient de $\mathcal{O}_2[\mathcal{F}]$ par le
sous-$\mathcal{O}_2$-module engendré par ces sections locales.
\end{proof}

\begin{definition}
\label{definition-module-principal-parts}
Soit $\mathcal{C}$ un site.
Soit $\mathcal{O}_1 \to \mathcal{O}_2$ un morphisme de faisceaux d'anneaux.
Soit $\mathcal{F}$ un faisceau de $\mathcal{O}_2$-modules.
Le module $\mathcal{P}^k_{\mathcal{O}_2/\mathcal{O}_1}(\mathcal{F})$
construit dans le lemme \ref{lemma-module-principal-parts}
est appelé le {\it module des parties principales d'ordre $k$} de $\mathcal{F}$.
\end{definition}

\noindent
Remarquons que les inclusions
$$
\text{Diff}^0(\mathcal{F}, \mathcal{G}) \subset
\text{Diff}^1(\mathcal{F}, \mathcal{G}) \subset
\text{Diff}^2(\mathcal{F}, \mathcal{G}) \subset \ldots
$$
correspondent, par le lemme de Yoneda (Catégories, lemme
\ref{categories-lemma-yoneda}), aux surjections
$$
\ldots \to \mathcal{P}^2_{\mathcal{O}_2/\mathcal{O}_1}(\mathcal{F})
\to \mathcal{P}^1_{\mathcal{O}_2/\mathcal{O}_1}(\mathcal{F})
\to \mathcal{P}^0_{\mathcal{O}_2/\mathcal{O}_1}(\mathcal{F}) = \mathcal{F}
$$

\begin{lemma}
\label{lemma-differential-operators-sheafify}
Soit $\mathcal{C}$ un site. Soit $\mathcal{O}_1 \to \mathcal{O}_2$
un homomorphisme de préfaisceaux d'anneaux. Soit $\mathcal{F}$ un préfaisceau
de $\mathcal{O}_2$-modules. Alors
$\mathcal{P}^k_{\mathcal{O}_2^\#/\mathcal{O}_1^\#}(\mathcal{F}^\#)$
est le faisceau associé au préfaisceau
$U \mapsto P^k_{\mathcal{O}_2(U)/\mathcal{O}_1(U)}(\mathcal{F}(U))$.
\end{lemma}

\begin{proof}
On peut le démontrer exactement comme pour le faisceau des différentielles
dans le lemme \ref{lemma-differentials-sheafify}.
Une approche peut-être plus élégante consiste à appliquer directement la
propriété universelle du lemme \ref{lemma-module-principal-parts} pour obtenir
l'égalité. Nous omettons les détails.
\end{proof}

\begin{lemma}
\label{lemma-sequence-of-principal-parts}
Soit $\mathcal{C}$ un site. Soit $\mathcal{O}_1 \to \mathcal{O}_2$
un homomorphisme de faisceaux d'anneaux. Soit $\mathcal{F}$ un faisceau
de $\mathcal{O}_2$-modules. Il existe une suite exacte courte canonique
$$
0 \to
\Omega_{\mathcal{O}_2/\mathcal{O}_1} \otimes_{\mathcal{O}_2} \mathcal{F} \to
\mathcal{P}^1_{\mathcal{O}_2/\mathcal{O}_1}(\mathcal{F}) \to
\mathcal{F} \to 0
$$
fonctorielle en $\mathcal{F}$, appelée la {\it suite des parties principales}.
\end{lemma}

\begin{proof}
Cela résulte de la version d'algèbre commutative
(Algèbre, lemme \ref{algebra-lemma-sequence-of-principal-parts})
et des lemmes \ref{lemma-differentials-sheafify} et
\ref{lemma-differential-operators-sheafify}.
\end{proof}

\begin{remark}
\label{remark-functoriality-principal-parts}
Soit $\mathcal{C}$ un site. Supposons donné un diagramme commutatif de
faisceaux d'anneaux
$$
\xymatrix{
\mathcal{B} \ar[r] & \mathcal{B}' \\
\mathcal{A} \ar[u] \ar[r] & \mathcal{A}' \ar[u]
}
$$
un $\mathcal{B}$-module $\mathcal{F}$, un $\mathcal{B}'$-module $\mathcal{F}'$,
et une application $\mathcal{B}$-linéaire $\mathcal{F} \to \mathcal{F}'$.
On obtient alors un système compatible d'applications de modules
$$
\xymatrix{
\ldots \ar[r] &
\mathcal{P}^2_{\mathcal{B}'/\mathcal{A}'}(\mathcal{F}') \ar[r] &
\mathcal{P}^1_{\mathcal{B}'/\mathcal{A}'}(\mathcal{F}') \ar[r] &
\mathcal{P}^0_{\mathcal{B}'/\mathcal{A}'}(\mathcal{F}') \\
\ldots \ar[r] &
\mathcal{P}^2_{\mathcal{B}/\mathcal{A}}(\mathcal{F}) \ar[r] \ar[u] &
\mathcal{P}^1_{\mathcal{B}/\mathcal{A}}(\mathcal{F}) \ar[r] \ar[u] &
\mathcal{P}^0_{\mathcal{B}/\mathcal{A}}(\mathcal{F}) \ar[u]
}
$$
Ces applications sont compatibles avec toute composition ultérieure
d'applications de ce type. Le moyen le plus simple de le voir consiste à
utiliser la description des modules
$\mathcal{P}^k_{\mathcal{B}/\mathcal{A}}(\mathcal{M})$ en termes de
générateurs et relations (locaux) donnée dans la démonstration du
lemme \ref{lemma-module-principal-parts}; cela se voit aussi directement à
partir de la propriété universelle de ces modules. De plus, ces applications
sont compatibles avec les suites exactes courtes du
lemme \ref{lemma-sequence-of-principal-parts}.
\end{remark}








\section{Le complexe cotangent naïf}
\label{section-netherlander}

\noindent
Cette section est l'analogue de
Algèbre, section \ref{algebra-section-netherlander},
et de
Modules, section \ref{modules-section-netherlander}.
Nous conseillons au lecteur de lire d'abord ces sections.

\medskip\noindent
Soit $\mathcal{C}$ un site. Soit $\mathcal{A} \to \mathcal{B}$ un
homomorphisme de faisceaux d'anneaux sur $\mathcal{C}$. Dans cette section,
pour tout faisceau d'ensembles $\mathcal{E}$ sur $\mathcal{C}$, nous notons
$\mathcal{A}[\mathcal{E}]$ le faisceautisé du préfaisceau
$U \mapsto \mathcal{A}(U)[\mathcal{E}(U)]$. Ici,
$\mathcal{A}(U)[\mathcal{E}(U)]$
désigne l'algèbre de polynômes sur $\mathcal{A}(U)$ dont les variables
correspondent aux éléments de $\mathcal{E}(U)$.
Nous notons $[e] \in \mathcal{A}(U)[\mathcal{E}(U)]$ la variable
correspondant à $e \in \mathcal{E}(U)$.
Il existe une surjection canonique de $\mathcal{A}$-algèbres
\begin{equation}
\label{equation-canonical-presentation}
\mathcal{A}[\mathcal{B}] \longrightarrow \mathcal{B},\quad [b] \longmapsto b
\end{equation}
dont nous notons le noyau $\mathcal{I} \subset \mathcal{A}[\mathcal{B}]$.
On observe immédiatement que $\mathcal{I}$ est engendré par les
sections locales $[b][b'] - [bb']$ et $[a] - a$. D'après le
lemme \ref{lemma-differential-seq}, il existe une application canonique
\begin{equation}
\label{equation-naive-cotangent-complex}
\mathcal{I}/\mathcal{I}^2
\longrightarrow
\Omega_{\mathcal{A}[\mathcal{B}]/\mathcal{A}}
\otimes_{\mathcal{A}[\mathcal{B}]} \mathcal{B}
\end{equation}
dont le conoyau est canoniquement isomorphe à $\Omega_{\mathcal{B}/\mathcal{A}}$.

\begin{definition}
\label{definition-naive-cotangent-complex}
Soit $\mathcal{C}$ un site. Soit $\mathcal{A} \to \mathcal{B}$ un
homomorphisme de faisceaux d'anneaux sur $\mathcal{C}$.
Le {\it complexe cotangent naïf} $\NL_{\mathcal{B}/\mathcal{A}}$
est le complexe de chaînes (\ref{equation-naive-cotangent-complex})
$$
\NL_{\mathcal{B}/\mathcal{A}} =
\left(\mathcal{I}/\mathcal{I}^2
\longrightarrow
\Omega_{\mathcal{A}[\mathcal{B}]/\mathcal{A}}
\otimes_{\mathcal{A}[\mathcal{B}]} \mathcal{B}\right)
$$
où $\mathcal{I}/\mathcal{I}^2$ est placé en degré $-1$ et
$\Omega_{\mathcal{A}[\mathcal{B}]/\mathcal{A}}
\otimes_{\mathcal{A}[\mathcal{B}]} \mathcal{B}$
est placé en degré $0$.
\end{definition}

\noindent
Cette construction possède une fonctorialité analogue à celle examinée dans le
lemme \ref{lemma-functoriality-differentials} pour les modules de
différentielles. Plus précisément, étant donné un diagramme commutatif
\begin{equation}
\label{equation-commutative-square-sheaves}
\vcenter{
\xymatrix{
\mathcal{B} \ar[r] & \mathcal{B}' \\
\mathcal{A} \ar[u] \ar[r] & \mathcal{A}' \ar[u]
}
}
\end{equation}
de faisceaux d'anneaux sur $\mathcal{C}$, il existe une application canonique
$\mathcal{B}$-linéaire de complexes
$$
\NL_{\mathcal{B}/\mathcal{A}} \longrightarrow \NL_{\mathcal{B}'/\mathcal{A}'}
$$
En effet, les applications du diagramme commutatif donnent une application canonique
$\mathcal{A}[\mathcal{B}] \to \mathcal{A}'[\mathcal{B}']$
qui envoie $\mathcal{I}$ dans
$\mathcal{I}' = \Ker(\mathcal{A}'[\mathcal{B}'] \to \mathcal{B}')$.
On obtient ainsi une application
$\mathcal{I}/\mathcal{I}^2 \to \mathcal{I}'/(\mathcal{I}')^2$ et une
application entre les modules de différentielles; ensemble, elles donnent
l'application cherchée entre les complexes cotangents naïfs.

\medskip\noindent
On peut choisir une autre présentation de $\mathcal{B}$ comme quotient d'une
algèbre de polynômes sur $\mathcal{A}$ et obtenir néanmoins le même objet de
$D(\mathcal{B})$. Pour l'expliquer, supposons que $\mathcal{E}$ soit un
faisceau d'ensembles sur $\mathcal{C}$ et que
$\alpha : \mathcal{E} \to \mathcal{B}$ soit une application de faisceaux
d'ensembles. On obtient alors un homomorphisme de $\mathcal{A}$-algèbres
$\mathcal{A}[\mathcal{E}] \to \mathcal{B}$. Supposons cette application
surjective et notons $\mathcal{J} \subset \mathcal{A}[\mathcal{E}]$ son noyau.
Posons
$$
\NL(\alpha) = \left(
\mathcal{J}/\mathcal{J}^2
\longrightarrow
\Omega_{\mathcal{A}[\mathcal{E}]/\mathcal{A}}
\otimes_{\mathcal{A}[\mathcal{E}]} \mathcal{B}\right)
$$
Voici le résultat.

\begin{lemma}
\label{lemma-NL-up-to-qis}
Dans la situation ci-dessus, il existe un isomorphisme canonique
$\NL(\alpha) = \NL_{\mathcal{B}/\mathcal{A}}$ dans $D(\mathcal{B})$.
\end{lemma}

\begin{proof}
Remarquons que
$\NL_{\mathcal{B}/\mathcal{A}} = \NL(\text{id}_\mathcal{B})$.
Il suffit donc de montrer que, pour deux applications
$\alpha_i : \mathcal{E}_i \to \mathcal{B}$ comme ci-dessus, il existe un
quasi-isomorphisme canonique $\NL(\alpha_1) = \NL(\alpha_2)$ dans
$D(\mathcal{B})$. Pour cela, posons
$\mathcal{E} = \mathcal{E}_1 \amalg \mathcal{E}_2$ et
$\alpha = \alpha_1 \amalg \alpha_2 : \mathcal{E} \to \mathcal{B}$.
Posons
$\mathcal{J}_i = \Ker(\mathcal{A}[\mathcal{E}_i] \to \mathcal{B})$
et
$\mathcal{J} = \Ker(\mathcal{A}[\mathcal{E}] \to \mathcal{B})$.
On obtient des applications
$\mathcal{A}[\mathcal{E}_i] \to \mathcal{A}[\mathcal{E}]$ qui envoient
$\mathcal{J}_i$ dans $\mathcal{J}$, donc des applications canoniques de complexes
$$
\NL(\alpha_i) \longrightarrow \NL(\alpha)
$$
et il suffit de montrer que ces applications sont des quasi-isomorphismes.
Pour cela, raisonnons comme suit. Tout d'abord, remarquons que
$H^0(\NL(\alpha_i)) = \Omega_{\mathcal{B}/\mathcal{A}}$ et
$H^0(\NL(\alpha)) = \Omega_{\mathcal{B}/\mathcal{A}}$ d'après
le lemme \ref{lemma-differential-seq}; l'application est donc un isomorphisme
sur les faisceaux de cohomologie en degré $0$.
De même, nous affirmons que $H^{-1}(\NL(\alpha_i))$ et
$H^{-1}(\NL(\alpha))$ sont les faisceaux associés au préfaisceau
$U \mapsto H_1(L_{\mathcal{B}(U)/\mathcal{A}(U)})$, où
où $H_1(L_{-/-})$ est défini comme dans
Algèbre, définition \ref{algebra-definition-naive-cotangent-complex}.
Cette assertion achève la démonstration.

\medskip\noindent
Démonstration de l'assertion. Soit
$\alpha : \mathcal{E} \to \mathcal{B}$ comme ci-dessus. Soit
$\mathcal{B}' \subset \mathcal{B}$ le sous-préfaisceau de
$\mathcal{A}$-algèbres dont la valeur sur $U$ est l'image de
$\mathcal{A}(U)[\mathcal{E}(U)] \to \mathcal{B}(U)$. Soit $\mathcal{I}'$
le préfaisceau dont la valeur sur $U$ est le noyau de
$\mathcal{A}(U)[\mathcal{E}(U)] \to \mathcal{B}(U)$. Alors $\mathcal{I}$
est le faisceautisé de $\mathcal{I}'$, et $\mathcal{B}$ celui de
$\mathcal{B}'$. De même, $H^{-1}(\NL(\alpha))$ est le faisceautisé du préfaisceau
$$
U \longmapsto
\Ker(\mathcal{I}'(U)/\mathcal{I}'(U)^2 \to
\Omega_{\mathcal{A}(U)[\mathcal{E}(U)]/\mathcal{A}(U)}
\otimes_{\mathcal{A}(U)[\mathcal{E}(U)]} \mathcal{B}'(U))
$$
d'après le lemme \ref{lemma-differentials-sheafify}.
Par Algèbre, lemme \ref{algebra-lemma-NL-homotopy}, on en déduit que
$H^{-1}(\NL(\alpha))$ est le faisceau associé au préfaisceau
$U \mapsto H_1(L_{\mathcal{B}'(U)/\mathcal{A}(U)})$. Il reste donc à montrer
que les applications
$H_1(L_{\mathcal{B}'(U)/\mathcal{A}(U)}) \to
H_1(L_{\mathcal{B}(U)/\mathcal{A}(U)})$ induisent un isomorphisme
$\mathcal{H}'_1 \to \mathcal{H}_1$ entre faisceautisés.

\medskip\noindent
Injectivité de $\mathcal{H}'_1 \to \mathcal{H}_1$. Soit
$f \in H_1(L_{\mathcal{B}'(U)/\mathcal{A}(U)})$ d'image nulle dans
$\mathcal{H}_1(U)$. Il faut montrer que $f$ a une image nulle dans
$\mathcal{H}'_1(U)$ est nulle. L'hypothèse signifie qu'il existe un recouvrement
$\{U_i \to U\}$ tel que $f$ s'annule dans
$H_1(L_{\mathcal{B}(U_i)/\mathcal{A}(U_i)})$ pour tout $i$.
En remplaçant $U$ par $U_i$, on se ramène au cas où $f$ s'annule dans
$H_1(L_{\mathcal{B}(U)/\mathcal{A}(U)})$.
Par Algèbre, lemme \ref{algebra-lemma-colimits-NL}, on peut trouver une
sous-algèbre de type fini $\mathcal{B}'(U) \subset B \subset \mathcal{B}(U)$
telle que $f$ s'annule dans $H_1(L_{B/\mathcal{A}(U)})$.
Comme $\mathcal{B} = (\mathcal{B}')^\#$, il existe un recouvrement
$\{U_i \to U\}$ tel que $B \to \mathcal{B}(U_i)$ se factorise par
$\mathcal{B}'(U_i)$. Ainsi $f$ s'annule dans
$H_1(L_{\mathcal{B}'(U_i)/\mathcal{A}(U_i)})$, comme voulu.

\medskip\noindent
La surjectivité de $\mathcal{H}'_1 \to \mathcal{H}_1$ se démontre
exactement de la même manière.
\end{proof}

\begin{lemma}
\label{lemma-pullback-NL}
Soit $f : \Sh(\mathcal{C}) \to \Sh(\mathcal{D})$ un morphisme de topos.
Soit $\mathcal{A} \to \mathcal{B}$ un homomorphisme de faisceaux d'anneaux
sur $\mathcal{D}$. Alors $f^{-1}\NL_{\mathcal{B}/\mathcal{A}} =
\NL_{f^{-1}\mathcal{B}/f^{-1}\mathcal{A}}$.
\end{lemma}

\begin{proof}
La démonstration est omise. Indication : appliquer le
lemme \ref{lemma-pullback-differentials}.
\end{proof}

\noindent
Le complexe cotangent d'un morphisme de topos annelés se définit à l'aide du
complexe cotangent introduit ci-dessus.

\begin{definition}
\label{definition-cotangent-complex-morphism-ringed-topoi}
Soient $X = (\Sh(\mathcal{C}), \mathcal{O})$ et
$Y = (\Sh(\mathcal{C}'), \mathcal{O}')$ des topos annelés.
Soit $(f, f^\sharp) : X \to Y$ un morphisme de topos annelés.
Le {\it complexe cotangent naïf} $\NL_f = \NL_{X/Y}$ de ce morphisme de
topos annelés est $\NL_{\mathcal{O}/f^{-1}\mathcal{O}'}$.
Nous écrivons parfois $\NL_{X/Y} = \NL_{\mathcal{O}/\mathcal{O}'}$.
\end{definition}









\section{Germes des modules}
\label{section-stalks}

\noindent
Il faut être quelque peu prudent lorsqu'on prend les germes aux points, car la
colimite qui définit un germe (voir Sites, équation
\ref{sites-equation-stalk}) peut ne pas être filtrante\footnote{Bien entendu,
dans presque tous les cas naturels, la colimite est filtrante et une partie de
la discussion de cette section peut être simplifiée.}. En revanche, par
définition d'un point d'un site, le foncteur des germes est exact et commute aux
colimites arbitraires. Autrement dit, il se comporte exactement comme si la
colimite était filtrante.

\begin{lemma}
\label{lemma-stalk-exact}
Soit $\mathcal{C}$ un site.
Soit $p$ un point de $\mathcal{C}$.
\begin{enumerate}
\item On a $(\mathcal{F}^\#)_p = \mathcal{F}_p$
pour tout préfaisceau d'ensembles sur $\mathcal{C}$.
\item Le foncteur des germes
$\Sh(\mathcal{C}) \to \textit{Sets}$,
$\mathcal{F} \mapsto \mathcal{F}_p$ est exact (voir
Catégories, définition \ref{categories-definition-exact})
et commute aux colimites arbitraires.
\item Le foncteur des germes
$\textit{PSh}(\mathcal{C}) \to \textit{Sets}$,
$\mathcal{F} \mapsto \mathcal{F}_p$ est exact (voir
Catégories, définition \ref{categories-definition-exact})
et commute aux colimites arbitraires.
\end{enumerate}
\end{lemma}

\begin{proof}
Par Sites, lemme \ref{sites-lemma-point-pushforward-sheaf}, on a (1).
Par Sites, lemme \ref{sites-lemma-adjoint-point-push-stalk}, on voit que
$\textit{PSh}(\mathcal{C}) \to \textit{Sets}$,
$\mathcal{F} \mapsto \mathcal{F}_p$ est un adjoint à gauche, et, par
Sites, lemme \ref{sites-lemma-point-pushforward-sheaf}, il en va de même pour
$\Sh(\mathcal{C}) \to \textit{Sets}$,
$\mathcal{F} \mapsto \mathcal{F}_p$.
Le foncteur des germes commute donc aux colimites arbitraires (voir
Catégories, lemme \ref{categories-lemma-adjoint-exact}).
Il résulte de la définition d'un point d'un site, voir
Sites, définition \ref{sites-definition-point}, que
$\Sh(\mathcal{C}) \to \textit{Sets}$,
$\mathcal{F} \mapsto \mathcal{F}_p$
est exact. Comme la faisceautisation est exacte
(Sites, lemme \ref{sites-lemma-sheafification-exact}), il s'ensuit que
$\textit{PSh}(\mathcal{C}) \to \textit{Sets}$,
$\mathcal{F} \mapsto \mathcal{F}_p$
est exact.
\end{proof}

\noindent
En particulier, puisque le foncteur des germes
$\mathcal{F} \mapsto \mathcal{F}_p$ sur les préfaisceaux commute à toutes les
limites et colimites finies, on peut appliquer le raisonnement de la
démonstration de Sites, proposition
\ref{sites-proposition-functoriality-algebraic-structures-topoi}.
On obtient ainsi que, si $\mathcal{F}$ est un (pré)faisceau de structures
algébriques figurant dans Sites, proposition
\ref{sites-proposition-functoriality-algebraic-structures-topoi}, alors le germe
$\mathcal{F}_p$ porte naturellement une structure algébrique du même type.
Détaillons le cas où $\mathcal{F}$ est un préfaisceau abélien. L'addition
$+ : \mathcal{F} \times \mathcal{F} \to \mathcal{F}$ induit une application
$$
+ :
\mathcal{F}_p \times \mathcal{F}_p
=
(\mathcal{F} \times \mathcal{F})_p
\longrightarrow
\mathcal{F}_p
$$
où l'égalité utilise le fait que le foncteur des germes sur les préfaisceaux
d'ensembles commute aux limites finies. Ceci définit une structure de groupe
sur le germe $\mathcal{F}_p$. On obtient ainsi le foncteur des germes
$$
\textit{PAb}(\mathcal{C}) \longrightarrow \textit{Ab}, \quad
\mathcal{F} \longmapsto \mathcal{F}_p
$$
Par construction, l'ensemble sous-jacent à $\mathcal{F}_p$ est le germe du
préfaisceau d'ensembles sous-jacent. Cela définit aussi le foncteur des germes pour
les faisceaux de groupes abéliens, par précomposition avec l'inclusion
$\textit{Ab}(\mathcal{C}) \subset \textit{PAb}(\mathcal{C})$.

\begin{lemma}
\label{lemma-stalk-exact-abelian}
Soit $\mathcal{C}$ un site.
Soit $p$ un point de $\mathcal{C}$.
\begin{enumerate}
\item Le foncteur
$\textit{Ab}(\mathcal{C}) \to \textit{Ab}$,
$\mathcal{F} \mapsto \mathcal{F}_p$ est exact.
\item Le foncteur des germes
$\textit{PAb}(\mathcal{C}) \to \textit{Ab}$,
$\mathcal{F}  \mapsto  \mathcal{F}_p$
est exact.
\item Pour $\mathcal{F} \in \Ob(\textit{PAb}(\mathcal{C}))$, on a
$\mathcal{F}_p = \mathcal{F}^\#_p$.
\end{enumerate}
\end{lemma}

\begin{proof}
Cela résulte formellement du lemme \ref{lemma-stalk-exact} et de la
construction précédente du foncteur des germes.
\end{proof}

\noindent
Passons maintenant au cas des faisceaux de modules.
Soit $(\mathcal{C}, \mathcal{O})$ un site annelé.
(Pour cette discussion, il suffit que $\mathcal{O}$ soit un préfaisceau
d'anneaux.) Soit $\mathcal{F}$ un préfaisceau de $\mathcal{O}$-modules.
Soit $p$ un point de $\mathcal{C}$. On obtient alors une application
$$
\cdot :
\mathcal{O}_p \times \mathcal{O}_p
=
(\mathcal{O} \times \mathcal{O})_p
\longrightarrow
\mathcal{O}_p
$$
qui est le germe de l'application de multiplication, et
$$
\cdot :
\mathcal{O}_p \times \mathcal{F}_p
=
(\mathcal{O} \times \mathcal{F})_p
\longrightarrow
\mathcal{F}_p
$$
qui est le germe de l'application de multiplication. Nous omettons la
vérification que ceci définit une structure d'anneau sur $\mathcal{O}_p$ et une
structure de $\mathcal{O}_p$-module sur $\mathcal{F}_p$. On obtient ainsi un foncteur
$$
\textit{PMod}(\mathcal{O}) \longrightarrow \textit{Mod}(\mathcal{O}_p), \quad
\mathcal{F} \longmapsto \mathcal{F}_p
$$
Par construction, l'ensemble sous-jacent à $\mathcal{F}_p$ est le germe du
préfaisceau d'ensembles sous-jacent. Cela définit aussi le foncteur des germes pour
les faisceaux de $\mathcal{O}$-modules, par précomposition avec l'inclusion
$\textit{Mod}(\mathcal{O}) \subset \textit{PMod}(\mathcal{O})$.

\begin{lemma}
\label{lemma-stalk-exact-modules}
Soit $(\mathcal{C}, \mathcal{O})$ un site annelé.
Soit $p$ un point de $\mathcal{C}$.
\begin{enumerate}
\item Le foncteur
$\textit{Mod}(\mathcal{O}) \to \textit{Mod}(\mathcal{O}_p)$,
$\mathcal{F} \mapsto \mathcal{F}_p$ est exact.
\item Le foncteur des germes
$\textit{PMod}(\mathcal{O}) \to \textit{Mod}(\mathcal{O}_p)$,
$\mathcal{F} \mapsto \mathcal{F}_p$
est exact.
\item Pour $\mathcal{F} \in \Ob(\textit{PMod}(\mathcal{O}))$, on a
$\mathcal{F}_p = \mathcal{F}^\#_p$.
\end{enumerate}
\end{lemma}

\begin{proof}
Cela résulte formellement du lemme \ref{lemma-stalk-exact-abelian}, de la
construction précédente du foncteur des germes et du lemme \ref{lemma-abelian}.
\end{proof}

\begin{lemma}
\label{lemma-pullback-stalk}
Soit
$(f, f^\sharp) :
(\Sh(\mathcal{C}), \mathcal{O}_\mathcal{C})
\to
(\Sh(\mathcal{D}), \mathcal{O}_\mathcal{D})$
soit un morphisme de topos annelés ou de sites annelés.
Soit $p$ un point de $\mathcal{C}$ ou de $\Sh(\mathcal{C})$, et posons
$q = f \circ p$. Alors
$$
(f^*\mathcal{F})_p =
\mathcal{F}_q \otimes_{\mathcal{O}_{\mathcal{D}, q}}
\mathcal{O}_{\mathcal{C}, p}
$$
pour tout $\mathcal{O}_\mathcal{D}$-module $\mathcal{F}$.
\end{lemma}

\begin{proof}
On a
$$
f^*\mathcal{F} =
f^{-1}\mathcal{F} \otimes_{f^{-1}\mathcal{O}_\mathcal{D}}
\mathcal{O}_\mathcal{C}
$$
par définition. Comme la prise du germe en $p$ (c'est-à-dire l'application de
$p^{-1}$) commute à $\otimes$ par le
lemme \ref{lemma-tensor-product-pullback}, le résultat découle de la relation
entre le germe en $p$ d'une image inverse et les germes en $q$, exposée dans
Sites, lemme \ref{sites-lemma-point-morphism-sites}, ou
Sites, lemme \ref{sites-lemma-point-morphism-topoi}.
\end{proof}






\section{Faisceaux gratte-ciel}
\label{section-skyscraper}

\noindent
Soit $p$ un point d'un site $\mathcal{C}$ ou d'un topos
$\Sh(\mathcal{C})$. Dans cette section, nous étudions les propriétés
d'exactitude du foncteur qui associe à un groupe abélien $A$ le faisceau
gratte-ciel $p_*A$. Rappelons d'abord que
$p_* : \textit{Sets} \to \Sh(\mathcal{C})$ possède de nombreuses propriétés
d'exactitude; voir Sites, lemmes \ref{sites-lemma-stalk-skyscraper} et
\ref{sites-lemma-skyscraper-functor-exact}.

\begin{lemma}
\label{lemma-skyscraper-exact}
Soit $\mathcal{C}$ un site. Soit $p$ un point de $\mathcal{C}$ ou de son
topos associé.
\begin{enumerate}
\item Le foncteur $p_* : \textit{Ab} \to \textit{Ab}(\mathcal{C})$,
$A \mapsto p_*A$, est exact.
\item Il existe une décomposition fonctorielle en somme directe
$$
p^{-1}p_*A = A \oplus I(A)
$$
pour $A \in \Ob(\textit{Ab})$.
\end{enumerate}
\end{lemma}

\begin{proof}
Par Sites, lemme \ref{sites-lemma-stalk-skyscraper}, il existe des
applications fonctorielles $A \to p^{-1}p_*A \to A$ dont la composée est
$\text{id}_A$. On obtient donc une décomposition fonctorielle en somme directe
comme en (2), où $I(A)$ est le noyau de l'application d'adjonction
$p^{-1}p_*A \to A$. Le foncteur $p_*$ est exact à gauche par le
lemme \ref{lemma-exactness-pushforward-pullback}. Le foncteur $p_*$ transforme
les surjections en surjections par Sites,
lemme \ref{sites-lemma-skyscraper-functor-exact}. D'où (1).
\end{proof}

\noindent
Pour procéder de même avec les faisceaux de modules, soit $p$ un point d'un
topos annelé $(\Sh(\mathcal{C}), \mathcal{O})$. Rappelons que $p^{-1}$ est
précisément le foncteur des germes. On peut donc considérer $p$ comme un morphisme
de topos annelés
$$
(p, \text{id}_{\mathcal{O}_p}) :
(\Sh(pt), \mathcal{O}_p)
\longrightarrow
(\Sh(\mathcal{C}), \mathcal{O}).
$$
On obtient ainsi un foncteur image inverse
$p^* : \textit{Mod}(\mathcal{O}) \to \textit{Mod}(\mathcal{O}_p)$
qui coïncide avec le foncteur des germes étudié dans le
lemme \ref{lemma-stalk-exact-modules}. Dans cette section, nous considérons le foncteur
$p_* : \textit{Mod}(\mathcal{O}_p) \to \textit{Mod}(\mathcal{O})$.

\begin{lemma}
\label{lemma-skyscraper-modules-exact}
Soit $(\Sh(\mathcal{C}), \mathcal{O})$ un topos annelé.
Soit $p$ un point du topos $\Sh(\mathcal{C})$.
\begin{enumerate}
\item Le foncteur
$p_* : \textit{Mod}(\mathcal{O}_p) \to \textit{Mod}(\mathcal{O})$,
$M \mapsto p_*M$ est exact.
\item La surjection canonique $p^{-1}p_*M \to M$ est $\mathcal{O}_p$-linéaire.
\item La décomposition fonctorielle en somme directe
$p^{-1}p_*M = M \oplus I(M)$ du lemme \ref{lemma-skyscraper-exact}
n'est {\bf pas} $\mathcal{O}_p$-linéaire en général.
\end{enumerate}
\end{lemma}

\begin{proof}
La partie (1) et la surjectivité dans (2) résultent immédiatement du résultat
correspondant pour les faisceaux abéliens dans le
lemme \ref{lemma-skyscraper-exact}. Comme
$p^{-1}\mathcal{O} = \mathcal{O}_p$, on a $p^{-1} = p^*$; par conséquent,
$p^{-1}p_*M \to M$ est la co-unité $p^*p_*M \to M$ de l'adjonction pour les
modules, et est donc linéaire.

\medskip\noindent
Démonstration de (3). Soit $G$ un groupe. Considérons le topos
$G\textit{-Sets} = \Sh(\mathcal{T}_G)$
et le point $p : \textit{Sets} \to G\textit{-Sets}$. Voir Sites,
section \ref{sites-section-example-sheaf-G-sets}, et
exemple \ref{sites-example-point-G-sets}. Ici, $p^{-1}$ est le foncteur
d'oubli de l'action de $G$, et $p_*$ est son adjoint à droite, qui envoie
$M$ sur $\text{Map}(G, M)$. Les applications de la décomposition en somme
directe sont
$$
M \to \text{Map}(G, M) \to M
$$
où la première envoie $m \in M$ sur l'application constante de valeur $m$, et
la seconde est l'évaluation en l'élément neutre $1$ de $G$.
Supposons ensuite que $R$ soit un anneau muni d'une action de $G$. Celle-ci
détermine un faisceau d'anneaux $\mathcal{O}$ sur $\mathcal{T}_G$. La catégorie
des $\mathcal{O}$-modules est celle des $R$-modules $M$ munis d'une action de
$G$ compatible à celle sur $R$. La structure de $R$-module sur
$\text{Map}(G, M)$ est donnée par
$$
( r f ) (\sigma) = \sigma(r) f(\sigma)
$$
pour $r \in R$ et $f \in \text{Map}(G, M)$. C'est en effet l'unique structure
$G$-invariante de $R$-module compatible avec l'évaluation en $1$. Le lecteur
observera qu'en général l'image de $M \to \text{Map}(G, M)$ n'est pas un
sous-$R$-module (prendre par exemple $M = R$ et supposer l'action de $G$ non
triviale), ce qui achève la démonstration.
\end{proof}

\begin{example}
\label{example-ring-with-group-action}
Soit $G$ un groupe. Considérons le site $\mathcal{T}_G$ et son point $p$;
voir Sites, exemple \ref{sites-example-point-G-sets}. Soit $R$ un anneau muni
d'une action de $G$, correspondant à un faisceau d'anneaux $\mathcal{O}$ sur
$\mathcal{T}_G$. Alors $\mathcal{O}_p = R$, où l'on oublie l'action de $G$.
Dans ce cas, $p^{-1}p_*M = \text{Map}(G, M)$,
$I(M) = \{f : G \to M \mid f(1_G) = 0\}$, et
$M \to \text{Map}(G, M)$ associe à $m \in M$ la fonction constante de valeur $m$.
\end{example}




\section{Localisation et points}
\label{section-localize-points}

\begin{lemma}
\label{lemma-stalk-j-shriek}
Soit $(\mathcal{C}, \mathcal{O})$ un site annelé.
Soit $p$ un point de $\mathcal{C}$. Soit $U$ un objet de $\mathcal{C}$.
Pour $\mathcal{G}$ dans $\textit{Mod}(\mathcal{O}_U)$, on a
$$
(j_{U!}\mathcal{G})_p =
\bigoplus\nolimits_q \mathcal{G}_q
$$
où la somme directe porte sur les points $q$ de $\mathcal{C}/U$ situés
au-dessus de $p$; voir Sites, lemme \ref{sites-lemma-points-above-point}.
\end{lemma}

\begin{proof}
Nous utilisons la description de $j_{U!}\mathcal{G}$ comme faisceau associé au préfaisceau
$V \mapsto
\bigoplus\nolimits_{\varphi \in \Mor_\mathcal{C}(V, U)}
\mathcal{G}(V/_\varphi U)$
du lemme \ref{lemma-extension-by-zero}. Le germe de $j_{U!}\mathcal{G}$ en
$p$ est égale à celle de ce préfaisceau; voir le
lemme \ref{lemma-stalk-exact-modules}. Soit
$u : \mathcal{C} \to \textit{Sets}$ le foncteur correspondant à $p$
(voir Sites, section \ref{sites-section-points}). On voit alors que
$$
(j_{U!}\mathcal{G})_p = \colim_{(V, y)}
\bigoplus\nolimits_{\varphi : V \to U} \mathcal{G}(V/_\varphi U)
$$
où la colimite est prise dans la catégorie des groupes abéliens.
À un quadruplet $(V, y, \varphi, s)$ intervenant dans cette colimite, on peut
associer $x = u(\varphi)(y) \in u(U)$. On obtient ainsi
$$
(j_{U!}\mathcal{G})_p =
\bigoplus\nolimits_{x \in u(U)}
\colim_{(\varphi : V \to U, y), \ u(\varphi)(y) = x} \mathcal{G}(V/_\varphi U).
$$
C'est l'expression du lemme, d'après la description des points $q$ situés
au-dessus de $x$ donnée dans Sites,
lemme \ref{sites-lemma-points-above-point}.
\end{proof}

\begin{remark}
\label{remark-not-pushforward}
Attention : le résultat du lemme \ref{lemma-stalk-j-shriek} n'a pas d'analogue
pour $j_{U, *}$.
\end{remark}












\section{Images inverses de modules plats}
\label{section-pullback-flat}

\noindent
L'image inverse d'un module plat par un morphisme de topos annelés est plate.
La démonstration demande quelque soin.

\begin{lemma}
\label{lemma-pullback-flat}
\begin{reference}
\cite[Expos\'e V, Corollary 1.7.1]{SGA4}
\end{reference}
Soit
$(f, f^\sharp) :
(\Sh(\mathcal{C}), \mathcal{O}_\mathcal{C})
\to
(\Sh(\mathcal{D}), \mathcal{O}_\mathcal{D})$
soit un morphisme de topos annelés ou de sites annelés.
Alors $f^*\mathcal{F}$ est un $\mathcal{O}_\mathcal{C}$-module plat lorsque
$\mathcal{F}$ est un $\mathcal{O}_\mathcal{D}$-module plat.
\end{lemma}

\begin{proof}
Choisissons un diagramme comme dans le
lemme \ref{lemma-morphism-ringed-topoi-comes-from-morphism-ringed-sites}.
Rappelons que la propriété d'être un module plat est intrinsèque
(voir la section \ref{section-intrinsic} et la
définition \ref{definition-flat}). Il suffit donc de démontrer le lemme pour
le morphisme $(h, h^\sharp) :
(\Sh(\mathcal{C}'), \mathcal{O}_{\mathcal{C}'})
\to
(\Sh(\mathcal{D}'), \mathcal{O}_{\mathcal{D}'})$.
Autrement dit, on peut supposer que nos sites $\mathcal{C}$ et $\mathcal{D}$
admettent toutes les limites finies et que $f$ est un morphisme de sites induit
par un foncteur continu $u : \mathcal{D} \to \mathcal{C}$ qui commute aux
limites finies.

\medskip\noindent
Rappelons que $f^*\mathcal{F} =
\mathcal{O}_\mathcal{C} \otimes_{f^{-1}\mathcal{O}_\mathcal{D}}
f^{-1}\mathcal{F}$ (définition \ref{definition-pushforward}).
Par le lemme \ref{lemma-flat-change-of-rings}, il suffit de montrer que
$f^{-1}\mathcal{F}$ est un $f^{-1}\mathcal{O}_\mathcal{D}$-module plat.
Avec le paragraphe précédent, cela nous ramène à la situation du paragraphe suivant.

\medskip\noindent
Supposons que $\mathcal{C}$ et $\mathcal{D}$ soient des sites admettant toutes
les limites finies, et que $u : \mathcal{D} \to \mathcal{C}$ soit un foncteur
continu commutant aux limites finies. Soit $\mathcal{O}$ un faisceau d'anneaux
sur $\mathcal{D}$ et soit $\mathcal{F}$ un $\mathcal{O}$-module plat. Alors $u$
définit un morphisme de sites $f : \mathcal{C} \to \mathcal{D}$
(Sites, proposition \ref{sites-proposition-get-morphism}). Il faut montrer que
$f^{-1}\mathcal{F}$ est un $f^{-1}\mathcal{O}$-module plat. Soit $U$ un objet
de $\mathcal{C}$ et soit
$$
f^{-1}\mathcal{O}|_U \xrightarrow{(f_1, \ldots, f_n)}
f^{-1}\mathcal{O}|_U^{\oplus n} \xrightarrow{(s_1, \ldots, s_n)}
f^{-1}\mathcal{F}|_U
$$
un complexe de $f^{-1}\mathcal{O}|_U$-modules. Notre but est de construire une
factorisation de $(s_1, \ldots, s_n)$ sur les membres d'un recouvrement de $U$,
comme dans le lemme \ref{lemma-flat-eq}, partie (2).
Considérons les éléments $s_a \in f^{-1}\mathcal{F}(U)$ et
$f_a \in f^{-1}\mathcal{O}(U)$. Comme $f^{-1}\mathcal{F}$ et, de même,
$f^{-1}\mathcal{O}$ s'obtiennent par faisceautisation de $u_p\mathcal{F}$, on
peut, après avoir remplacé $U$ par les membres d'un recouvrement, supposer que
$s_a$ est l'image d'un élément
$s'_a \in u_p\mathcal{F}(U)$ et que $f_a$ est l'image d'un élément
$f'_a \in u_p\mathcal{O}(U)$. Après un nouveau remplacement de $U$ par les
membres d'un recouvrement, on peut supposer que $\sum f'_as'_a$ est nul dans
$u_p\mathcal{F}(U)$. Rappelons que la catégorie
$(\mathcal{I}_U^u)^{opp}$ est filtrante
(Sites, lemme \ref{sites-lemma-directed}) et que
$u_p\mathcal{F}(U) = \colim_{(\mathcal{I}_U^u)^{opp}} \mathcal{F}(V)$ et
$u_p\mathcal{O}(U) = \colim_{(\mathcal{I}_U^u)^{opp}} \mathcal{O}(V)$.
On peut donc supposer qu'il existe une paire
$(V, \phi) \in \Ob(\mathcal{I}_U^u)$, où $V$ est un objet de $\mathcal{D}$ et
$\phi$ est un morphisme $\phi : U \to u(V)$ de $\mathcal{D}$, ainsi que des éléments
$s''_a \in \mathcal{F}(V)$ et $f''_a \in \mathcal{O}(V)$ dont les images dans
$u_p\mathcal{F}(U)$ et $u_p\mathcal{O}(U)$ sont $s'_a$ et $f'_a$, et tels que
$\sum f''_a s''_a = 0$ dans $\mathcal{F}(V)$. On obtient alors un complexe
$$
\mathcal{O}|_V \xrightarrow{(f''_1, \ldots, f''_n)}
\mathcal{O}|_V^{\oplus n} \xrightarrow{(s''_1, \ldots, s''_n)}
\mathcal{F}|_V
$$
et l'autre implication du lemme \ref{lemma-flat-eq} montre qu'il existe un
recouvrement $\{V_i \to V\}$ dans $\mathcal{D}$ et, pour chaque $i$, une factorisation
$$
\mathcal{O}|_{V_i}^{\oplus n}
\xrightarrow{B''_i}
\mathcal{O}|_{V_i}^{\oplus l_i} \xrightarrow{(t''_{i1}, \ldots, t''_{il_i})}
\mathcal{F}|_{V_i}
$$
de $(s''_1, \ldots, s''_n)|_{V_i}$ telle que
$B''_i \circ (f''_1, \ldots, f''_n)|_{V_i} = 0$.
Posons $U_i = U \times_{\phi, u(V)} u(V_i)$, notons
$B_i \in \text{Mat}(l_i \times n, f^{-1}\mathcal{O}(U_i))$ l'image de $B''_i$
et $t_{ij} \in f^{-1}\mathcal{F}(U_i)$ l'image de $t''_{ij}$. On obtient une factorisation
$$
f^{-1}\mathcal{O}|_{U_i}^{\oplus n}
\xrightarrow{B_i}
f^{-1}\mathcal{O}|_{U_i}^{\oplus l_i}
\xrightarrow{(t_{i1}, \ldots, t_{il_i})}
\mathcal{F}|_{U_i}
$$
de $(s_1, \ldots, s_n)|_{U_i}$ telle que
$B_i \circ (f_1, \ldots, f_n)|_{U_i} = 0$.
Cela achève la démonstration.
\end{proof}

\begin{lemma}
\label{lemma-stalk-flat}
Soit $(\mathcal{C}, \mathcal{O})$ un site annelé et soit $p$ un point de
$\mathcal{C}$. Si $\mathcal{F}$ est un $\mathcal{O}$-module plat, alors
$\mathcal{F}_p$ est un $\mathcal{O}_p$-module plat.
\end{lemma}

\begin{proof}
Dans la section \ref{section-skyscraper}, nous avons vu que l'on peut
considérer $p$ comme un morphisme de topos annelés
$$
(p, \text{id}_{\mathcal{O}_p}) :
(\Sh(pt), \mathcal{O}_p)
\longrightarrow
(\Sh(\mathcal{C}), \mathcal{O}).
$$
tel que le foncteur image inverse
$p^* : \textit{Mod}(\mathcal{O}) \to \textit{Mod}(\mathcal{O}_p)$
coïncide avec le foncteur des germes. Le lemme résulte donc du
lemme \ref{lemma-pullback-flat}.
\end{proof}

\begin{lemma}
\label{lemma-check-flat-stalks}
Soit $(\mathcal{C}, \mathcal{O})$ un site annelé. Soit $\mathcal{F}$ un
faisceau de $\mathcal{O}$-modules. Soit $\{p_i\}_{i \in I}$ une famille
conservative de points de $\mathcal{C}$. Alors $\mathcal{F}$ est plat si et
seulement si $\mathcal{F}_{p_i}$ est un $\mathcal{O}_{p_i}$-module plat pour
tout $i \in I$.
\end{lemma}

\begin{proof}
Le lemme \ref{lemma-stalk-flat} donne une implication. Pour la réciproque,
utilisons l'égalité
$(\mathcal{F} \otimes_\mathcal{O} \mathcal{G})_p =
\mathcal{F}_p \otimes_{\mathcal{O}_p} \mathcal{G}_p$
fournie par le lemme \ref{lemma-tensor-product-pullback} (puisque le germe en
$p$ est donnée par $p^{-1}$), ainsi que le
lemme \ref{lemma-check-exactness-stalks}.
\end{proof}

\begin{lemma}
\label{lemma-pullback-ses}
Soit
$f : (\Sh(\mathcal{C}'), \mathcal{O}') \to (\Sh(\mathcal{C}), \mathcal{O})$
soit un morphisme de topos annelés. Soit
$0 \to \mathcal{F} \to \mathcal{G} \to \mathcal{H} \to 0$
une suite exacte courte de $\mathcal{O}$-modules, où $\mathcal{H}$ est un
$\mathcal{O}$-module plat. Alors la suite
$0 \to f^*\mathcal{F} \to f^*\mathcal{G} \to f^*\mathcal{H} \to 0$
est elle aussi exacte.
\end{lemma}

\begin{proof}
Comme $f^{-1}$ est exact, on a la suite exacte courte
$0 \to f^{-1}\mathcal{F} \to f^{-1}\mathcal{G} \to f^{-1}\mathcal{H} \to 0$
de $f^{-1}\mathcal{O}$-modules. Par le lemme \ref{lemma-pullback-flat}, le
$f^{-1}\mathcal{O}$-module $f^{-1}\mathcal{H}$ est plat. Le
lemme \ref{lemma-flat-tor-zero} montre donc que la suite reste exacte après
tensorisation sur $f^{-1}\mathcal{O}$ par $\mathcal{O}'$. Comme
$f^*\mathcal{F} = f^{-1}\mathcal{F} \otimes_{f^{-1}\mathcal{O}} \mathcal{O}'$
et de même pour $\mathcal{G}$ et $\mathcal{H}$, on conclut.
\end{proof}





\section{Topos localement annelés}
\label{section-locally-ringed}


\noindent
Une référence pour cette section est
\cite[Expos\'e IV, Exercice 13.9]{SGA4}.

\begin{lemma}
\label{lemma-locally-ringed}
Soit $(\mathcal{C}, \mathcal{O})$ un site annelé. Les conditions suivantes
sont équivalentes :
\begin{enumerate}
\item Pour tout objet $U$ de $\mathcal{C}$ et tout $f \in \mathcal{O}(U)$,
il existe un recouvrement $\{U_j \to U\}$ tel que, pour chaque $j$, soit
$f|_{U_j}$, soit $(1 - f)|_{U_j}$ soit inversible.
\item Pour $U \in \Ob(\mathcal{C})$, $n \geq 1$, et
$f_1, \ldots, f_n \in \mathcal{O}(U)$ engendrant l'idéal unité de
$\mathcal{O}(U)$, il existe un recouvrement $\{U_j \to U\}$ tel que, pour
chaque $j$, il existe $i$ pour lequel $f_i|_{U_j}$ est inversible.
\item L'application de faisceaux d'ensembles
$$
(\mathcal{O} \times \mathcal{O})
\amalg
(\mathcal{O} \times \mathcal{O})
\longrightarrow
\mathcal{O} \times \mathcal{O}
$$
qui envoie $(f, a)$ dans la première composante sur $(f, af)$ et $(f, b)$ dans
la seconde sur $(f, b(1 - f))$ est surjective.
\end{enumerate}
\end{lemma}

\begin{proof}
Il est clair que (2) implique (1). Pour montrer que (1) implique (2),
raisonnons par récurrence sur $n$. Le premier cas est $n = 2$ (le cas $n = 1$
étant trivial). On a alors $a_1f_1 + a_2f_2 = 1$ pour certains
$a_1, a_2 \in \mathcal{O}(U)$. Par hypothèse, il existe un recouvrement
$\{U_j \to U\}$ tel que, pour chaque $j$, soit $a_1f_1|_{U_j}$, soit
$a_2f_2|_{U_j}$ soit inversible. Ainsi $f_1|_{U_j}$ ou $f_2|_{U_j}$ est
inversible, comme voulu. Pour $n > 2$, on a
$a_1f_1 + \ldots + a_nf_n = 1$ pour certains
$a_1, \ldots, a_n \in \mathcal{O}(U)$.
Le cas $n = 2$ donne un recouvrement $\{U_j \to U\}_{j \in J}$ tel que, pour
chaque $j$, soit $f_n|_{U_j}$, soit
$a_1f_1 + \ldots + a_{n - 1}f_{n - 1}|_{U_j}$ soit inversible.
Disons que le premier cas se produit pour $j \in J_n$ et posons
$J' = J \setminus J_n$. Par l'hypothèse de récurrence, pour chaque
$j \in J'$, il existe un recouvrement
$\{U_{jk} \to U_j\}_{k \in K_j}$ tel que, pour chaque $k \in K_j$, il existe
un $i \in \{1, \ldots, n - 1\}$ tel que $f_i|_{U_{jk}}$ soit inversible. Par les
axiomes d'un site, la famille de morphismes
$\{U_j \to U\}_{j \in J_n} \cup \{U_{jk} \to U\}_{j \in J', k \in K_j}$
est un recouvrement possédant la propriété voulue.

\medskip\noindent
Supposons (1). Pour voir que l'application de (3) est surjective, soit
$(f, c)$ une section de $\mathcal{O} \times \mathcal{O}$ sur $U$. Par
hypothèse, il existe un recouvrement $\{U_j \to U\}$ tel que, pour chaque $j$,
$f$ ou $1 - f$ se restreigne en une section inversible. Dans le premier cas, on
prend $a = c|_{U_j} (f|_{U_j})^{-1}$; dans le second,
$b = c|_{U_j} (1 - f|_{U_j})^{-1}$. Ainsi $(f, c)$ appartient à l'image sur
chaque membre du recouvrement. Réciproquement, supposons (3). Pour tout $U$ et
$f \in \mathcal{O}(U)$, il existe un recouvrement $\{U_j \to U\}$ de $U$ tel que la
section $(f, 1)|_{U_j}$ appartienne à l'image de l'application de (3) sur
$U_j$. Cela signifie précisément que $f$ ou $1 - f$ se restreint en une section
inversible sur $U_j$, d'où (1).
\end{proof}

\begin{lemma}
\label{lemma-locally-ringed-stalk}
Soit $(\mathcal{C}, \mathcal{O})$ un site annelé. Considérons les conditions suivantes :
\begin{enumerate}
\item Pour tout objet $U$ de $\mathcal{C}$ et tout $f \in \mathcal{O}(U)$,
il existe un recouvrement $\{U_j \to U\}$ tel que, pour chaque $j$, soit
$f|_{U_j}$, soit $(1 - f)|_{U_j}$ soit inversible.
\item Pour tout point $p$ de $\mathcal{C}$, le germe $\mathcal{O}_p$ est
soit l'anneau nul, soit un anneau local.
\end{enumerate}
On a toujours (1) $\Rightarrow$ (2). Si $\mathcal{C}$ a assez de points,
alors (1) et (2) sont équivalentes.
\end{lemma}

\begin{proof}
Supposons (1). Soit $p$ un point de $\mathcal{C}$ donné par un foncteur
$u : \mathcal{C} \to \textit{Sets}$, et soit $f_p \in \mathcal{O}_p$.
Puisque $\mathcal{O}_p$ est calculé par Sites, équation
(\ref{sites-equation-stalk}), on peut représenter $f_p$ par un triplet
$(U, x, f)$, où $x \in U(U)$ et $f \in \mathcal{O}(U)$. Par hypothèse, il existe
un recouvrement $\{U_i \to U\}$ tel que, pour chaque $i$, $f$ ou $1 - f$ soit
inversible sur $U_i$. Comme $u$ définit un point du site, il existe, pour un
certain $i$, un $x_i \in u(U_i)$ qui s'envoie sur $x \in u(U)$. La discussion
autour de Sites, équation (\ref{sites-equation-stalk}), montre que $(U, x, f)$ et
$(U_i, x_i, f|_{U_i})$ définissent le même élément de $\mathcal{O}_p$. Ainsi
$f_p$ ou $1 - f_p$ est inversible. Par conséquent, dans l'anneau
$\mathcal{O}_p$, pour tout élément $a$, $a$ ou $1 - a$ est inversible. Cela
signifie que $\mathcal{O}_p$ est nul ou local; voir Algèbre,
lemme \ref{algebra-lemma-characterize-local-ring}.

\medskip\noindent
Supposons (2) et que $\mathcal{C}$ ait assez de points. Considérons
l'application de faisceaux d'ensembles
$$
\mathcal{O} \times \mathcal{O} \amalg \mathcal{O} \times \mathcal{O}
\longrightarrow
\mathcal{O} \times \mathcal{O}
$$
du lemme \ref{lemma-locally-ringed}, partie (3). Pour tout anneau local $R$,
l'application correspondante
$(R \times R) \amalg (R \times R) \to R \times R$
est surjective; voir par exemple Algèbre,
lemme \ref{algebra-lemma-characterize-local-ring}. Comme chaque
$\mathcal{O}_p$ est un anneau local ou nul, l'application est surjective sur
les germes. Puisque $\mathcal{C}$ a assez de points, elle est donc surjective,
ce qui conclut.
\end{proof}

\noindent
Dans Modules, section \ref{modules-section-pathology}, nous avons signalé
qu'un espace annelé $(X, \mathcal{O}_X)$ peut avoir un sous-espace ouvert sur
lequel le faisceau structural est nul. Pour l'exclure, on peut imposer que les
sections $1$ et $0$ aient des valeurs distinctes dans tout germe de $X$.
Pour les topos et sites annelés, la condition est que
\begin{equation}
\label{equation-one-is-never-zero}
\emptyset^\# \longrightarrow
\text{Equalizer}(0, 1 : * \longrightarrow \mathcal{O})
\end{equation}
soit un isomorphisme de faisceaux. Ici, $*$ est le faisceau singleton et
$\emptyset^\#$ le « faisceau vide », c'est-à-dire respectivement l'objet final
et l'objet initial de la catégorie des faisceaux; voir Sites,
exemple \ref{sites-example-singleton-sheaf}, respectivement
section \ref{sites-section-almost-cocontinuous}. Autrement dit, si
$U \in \Ob(\mathcal{C})$ n'est pas vide au sens faisceautique, alors
$1, 0 \in \mathcal{O}(U)$ ne sont pas égaux. Énonçons le lemme obligé.

\begin{lemma}
\label{lemma-ringed-stalk-not-zero}
Soit $(\mathcal{C}, \mathcal{O})$ un site annelé. Considérons les assertions :
\begin{enumerate}
\item (\ref{equation-one-is-never-zero}) est un isomorphisme, et
\item pour tout point $p$ de $\mathcal{C}$, le germe $\mathcal{O}_p$ n'est
pas l'anneau nul.
\end{enumerate}
On a toujours (1) $\Rightarrow$ (2), et, si $\mathcal{C}$ a assez de points,
alors (1) $\Leftrightarrow$ (2).
\end{lemma}

\begin{proof}
La démonstration est omise.
\end{proof}

\noindent
Les lemmes \ref{lemma-locally-ringed},
\ref{lemma-locally-ringed-stalk} et
\ref{lemma-ringed-stalk-not-zero} motivent la définition suivante.

\begin{definition}
\label{definition-locally-ringed}
Un site annelé $(\mathcal{C}, \mathcal{O})$ est dit
{\it localement annelé} si (\ref{equation-one-is-never-zero}) est un
isomorphisme et si les propriétés équivalentes du
lemme \ref{lemma-locally-ringed} sont satisfaites.
\end{definition}

\noindent
Dans \cite[Expos\'e IV, Exercice 13.9]{SGA4}, la condition que
(\ref{equation-one-is-never-zero}) soit un isomorphisme manque, ce qui conduit
à une notion légèrement différente de site et de topos localement annelés.
Comme nous sommes guidés par la notion d'espace localement annelé, nous avons
choisi d'ajouter cette condition (voir l'explication ci-dessus).

\begin{lemma}
\label{lemma-locally-ringed-intrinsic}
Être un site localement annelé est une propriété intrinsèque. Plus précisément :
\begin{enumerate}
\item si $f : \Sh(\mathcal{C}') \to \Sh(\mathcal{C})$ est un morphisme de
topos et si $(\mathcal{C}, \mathcal{O})$ est un site localement annelé, alors
$(\mathcal{C}', f^{-1}\mathcal{O})$ est un site localement annelé; et
\item si
$(f, f^\sharp) : (\Sh(\mathcal{C}'), \mathcal{O}')
\to (\Sh(\mathcal{C}), \mathcal{O})$
est une équivalence de topos annelés, alors $(\mathcal{C}, \mathcal{O})$ est
localement annelé si et seulement si $(\mathcal{C}', \mathcal{O}')$ l'est.
\end{enumerate}
\end{lemma}

\begin{proof}
Il est clair que (2) résulte de (1). Pour démontrer (1), remarquons que, comme
$f^{-1}$ est exact, on a $f^{-1}* = *$ et
$f^{-1}\emptyset^\# = \emptyset^\#$; de plus, $f^{-1}$ commute aux produits
et aux égalisateurs, et transforme isomorphismes et surjections en
isomorphismes et surjections. Ainsi, $f^{-1}$ transforme l'isomorphisme
(\ref{equation-one-is-never-zero}) en son analogue pour
$f^{-1}\mathcal{O}$, et la surjection du
lemme \ref{lemma-locally-ringed}, partie (3), en la surjection correspondante
pour $f^{-1}\mathcal{O}$.
\end{proof}

\noindent
En fait, le lemme \ref{lemma-locally-ringed-intrinsic}, partie (2), est
l'analogue de Schémas, lemme
\ref{schemes-lemma-isomorphism-locally-ringed}. Il garantit que la définition
suivante a un sens.

\begin{definition}
\label{definition-locally-ringed-topos}
Un topos annelé $(\Sh(\mathcal{C}), \mathcal{O})$ est dit
{\it localement annelé} si le site annelé sous-jacent
$(\mathcal{C}, \mathcal{O})$ est localement annelé.
\end{definition}

\noindent
Voici un exemple de conséquence de la propriété d'être localement annelé.

\begin{lemma}
\label{lemma-invertible-is-locally-free-rank-1}
Soit $(\Sh(\mathcal{C}), \mathcal{O})$ un topos annelé. Tout
$\mathcal{O}$-module localement libre de rang $1$ est inversible.
Si $(\mathcal{C}, \mathcal{O})$ est localement annelé, la réciproque est
également vraie (ce qui n'est pas le cas en général).
\end{lemma}

\begin{proof}
Supposons $\mathcal{L}$ localement libre de rang $1$ et considérons
l'application d'évaluation
$$
\mathcal{L} \otimes_\mathcal{O}
\SheafHom_\mathcal{O}(\mathcal{L}, \mathcal{O})
\longrightarrow \mathcal{O}
$$
Pour tout objet $U$ de $\mathcal{C}$, en se restreignant aux membres d'un
recouvrement qui trivialise $\mathcal{L}$, on voit que cette application est
un isomorphisme (détails omis). Le lemme \ref{lemma-invertible} montre donc que
$\mathcal{L}$ est inversible.

\medskip\noindent
Supposons $(\Sh(\mathcal{C}), \mathcal{O})$ localement annelé et soit $U$ un
objet de $\mathcal{C}$. Dans la démonstration du lemme \ref{lemma-invertible},
nous avons vu qu'il existe un recouvrement $\{U_i \to U\}$ tel que
$\mathcal{L}|_{\mathcal{C}/U_i}$ soit facteur direct d'un
$\mathcal{O}_{U_i}$-module libre fini. Après avoir remplacé $U$ par $U_i$,
soit
$p : \mathcal{O}_U^{\oplus r} \to \mathcal{O}_U^{\oplus r}$
un projecteur dont l'image est isomorphe à $\mathcal{L}|_{\mathcal{C}/U}$.
Alors $p$ correspond à une matrice
$$
P = (p_{ij}) \in \text{Mat}(r \times r, \mathcal{O}(U))
$$
qui est un projecteur : $P^2 = P$. Posons $A = \mathcal{O}(U)$, de sorte que
$P \in \text{Mat}(r \times r, A)$. Par Algèbre,
lemme \ref{algebra-lemma-finite-projective}, l'image de $P$ est un module
$M$ fini localement libre sur $A$. Il existe donc
$f_1, \ldots, f_t \in A$ engendrant l'idéal unité tels que $M_{f_i}$ soit libre
fini. Par le lemme \ref{lemma-locally-ringed}, après avoir remplacé $U$ par les
membres d'un recouvrement ouvert, on peut supposer $M$ libre. Cela signifie
que $\mathcal{L}|_U$ est libre (détails omis). Puisque $\mathcal{L}$ est
inversible, cela n'est possible que si le rang de $\mathcal{L}|_U$ vaut $1$,
ce qui achève la démonstration.
\end{proof}

\noindent
Nous voulons maintenant préciser ce que signifie être un morphisme d'espaces
localement annelés. Nous utiliserons le lemme suivant.

\begin{lemma}
\label{lemma-locally-ringed-morphism}
Soit
$(f, f^\sharp) : (\Sh(\mathcal{C}), \mathcal{O}_\mathcal{C})
\to (\Sh(\mathcal{D}), \mathcal{O}_\mathcal{D})$
soit un morphisme de topos annelés. Considérons les conditions suivantes :
\begin{enumerate}
\item Le diagramme de faisceaux
$$
\xymatrix{
f^{-1}(\mathcal{O}^*_\mathcal{D}) \ar[r]_-{f^\sharp} \ar[d] &
\mathcal{O}^*_\mathcal{C} \ar[d] \\
f^{-1}(\mathcal{O}_\mathcal{D}) \ar[r]^-{f^\sharp} &
\mathcal{O}_\mathcal{C}
}
$$
est cartésien.
\item Pour tout point $p$ de $\mathcal{C}$, en posant $q = f \circ p$, le diagramme
$$
\xymatrix{
\mathcal{O}^*_{\mathcal{D}, q} \ar[r] \ar[d] &
\mathcal{O}^*_{\mathcal{C}, p} \ar[d] \\
\mathcal{O}_{\mathcal{D}, q} \ar[r] &
\mathcal{O}_{\mathcal{C}, p}
}
$$
d'ensembles est cartésien.
\end{enumerate}
On a toujours (1) $\Rightarrow$ (2). Si $\mathcal{C}$ a assez de points,
alors (1) et (2) sont équivalentes. Si
$(\Sh(\mathcal{C}), \mathcal{O}_\mathcal{C})$
et
$(\Sh(\mathcal{D}), \mathcal{O}_\mathcal{D})$
sont des topos localement annelés, alors (2) équivaut à :
\begin{enumerate}
\item[(3)] Pour tout point $p$ de $\mathcal{C}$, en posant $q = f \circ p$,
l'homomorphisme d'anneaux
$\mathcal{O}_{\mathcal{D}, q} \to \mathcal{O}_{\mathcal{C}, p}$ est local.
\end{enumerate}
En fait, la propriété (2), ou la propriété (3) pour une famille conservative
de points, implique (1).
\end{lemma}

\begin{proof}
Le lemme se démontre de lui-même : il suffit de développer les définitions.
\end{proof}

\begin{definition}
\label{definition-morphism-locally-ringed-topoi}
Soit $(f, f^\sharp) : (\Sh(\mathcal{C}), \mathcal{O}_\mathcal{C})
\to (\Sh(\mathcal{D}), \mathcal{O}_\mathcal{D})$
soit un morphisme de topos annelés. Supposons que
$(\Sh(\mathcal{C}), \mathcal{O}_\mathcal{C})$
et
$(\Sh(\mathcal{D}), \mathcal{O}_\mathcal{D})$
soient des topos localement annelés. On dit que $(f, f^\sharp)$ est un
{\it morphisme de topos localement annelés} si et seulement si le diagramme de faisceaux
$$
\xymatrix{
f^{-1}(\mathcal{O}^*_\mathcal{D}) \ar[r]_-{f^\sharp} \ar[d] &
\mathcal{O}^*_\mathcal{C} \ar[d] \\
f^{-1}(\mathcal{O}_\mathcal{D}) \ar[r]^-{f^\sharp} &
\mathcal{O}_\mathcal{C}
}
$$
(voir le lemme \ref{lemma-locally-ringed-morphism}) est cartésien. Si
$(f, f^\sharp)$ est un morphisme de sites annelés, on dit que c'est un
{\it morphisme de sites localement annelés} si le morphisme de topos annelés
associé est un morphisme de topos localement annelés.
\end{definition}

\noindent
Il est clair qu'un isomorphisme de topos annelés entre topos localement
annelés est automatiquement un isomorphisme de topos localement annelés.

\begin{lemma}
\label{lemma-composition-morphisms-locally-ringed-topoi}
Soient
$(f, f^\sharp) :
(\Sh(\mathcal{C}_1), \mathcal{O}_1)
\to (\Sh(\mathcal{C}_2), \mathcal{O}_2)$ et
$(g, g^\sharp) :
(\Sh(\mathcal{C}_2), \mathcal{O}_2) \to
(\Sh(\mathcal{C}_3), \mathcal{O}_3)$
soient des morphismes de topos localement annelés. Alors leur composée
$(g, g^\sharp) \circ (f, f^\sharp)$ (voir la
définition \ref{definition-ringed-topos}) est encore un morphisme de topos
localement annelés.
\end{lemma}

\begin{proof}
La démonstration est omise.
\end{proof}

\begin{lemma}
\label{lemma-locally-ringed-intrinsic-morphism}
Soit $f : \Sh(\mathcal{C}') \to \Sh(\mathcal{C})$ un morphisme de topos.
Si $\mathcal{O}$ est un faisceau d'anneaux sur $\mathcal{C}$, alors
$$
f^{-1}(\mathcal{O}^*) = (f^{-1}\mathcal{O})^*.
$$
En particulier, si $\mathcal{O}$ fait de $\mathcal{C}$ un site localement
annelé, alors, en posant $f^\sharp = \text{id}$, le morphisme de topos annelés
$$
(f, f^\sharp) :
(\Sh(\mathcal{C}'), f^{-1}\mathcal{O})
\to
(\Sh(\mathcal{C}), \mathcal{O})
$$
est un morphisme de topos localement annelés.
\end{lemma}

\begin{proof}
Remarquons que le diagramme
$$
\xymatrix{
\mathcal{O}^* \ar[rr] \ar[d]_{u \mapsto (u, u^{-1})} & &
{*} \ar[d]^{1} \\
\mathcal{O} \times \mathcal{O} \ar[rr]^-{(a, b) \mapsto ab} & &
\mathcal{O}
}
$$
est cartésien. Comme $f^{-1}$ est exact, on en déduit que
$$
\xymatrix{
f^{-1}(\mathcal{O}^*)
\ar[d]_{u \mapsto (u, u^{-1})} \ar[rr] & &
{*} \ar[d]^{1} \\
f^{-1}\mathcal{O} \times f^{-1}\mathcal{O} \ar[rr]^-{(a, b) \mapsto ab} & &
f^{-1}\mathcal{O}
}
$$
est cartésien, ce qui prouve la première assertion. Pour la seconde,
remarquons que $(\mathcal{C}', f^{-1}\mathcal{O})$ est un site localement annelé
par le lemme \ref{lemma-locally-ringed-intrinsic}; l'assertion a donc un sens.
La première partie montre alors que le morphisme est un morphisme de topos
localement annelés.
\end{proof}

\begin{lemma}
\label{lemma-localize-morphism-locally-ringed-topoi}
Localisation des sites et topos localement annelés.
\begin{enumerate}
\item Soit $(\mathcal{C}, \mathcal{O})$ un site localement annelé et soit $U$
un objet de $\mathcal{C}$. Alors la localisation
$(\mathcal{C}/U, \mathcal{O}_U)$ est un site localement annelé, et le morphisme
de localisation
$$
(j_U, j_U^\sharp) :
(\Sh(\mathcal{C}/U), \mathcal{O}_U)
\to
(\Sh(\mathcal{C}), \mathcal{O})
$$
est un morphisme de topos localement annelés.
\item Soit $(\mathcal{C}, \mathcal{O})$ un site localement annelé et soit
$f : V \to U$ un morphisme de $\mathcal{C}$. Alors le morphisme
$$
(j, j^\sharp) :
(\Sh(\mathcal{C}/V), \mathcal{O}_V)
\to
(\Sh(\mathcal{C}/U), \mathcal{O}_U)
$$
du lemme \ref{lemma-relocalize} est un morphisme de topos localement annelés.
\item Soit
$(f, f^\sharp) :
(\mathcal{C}, \mathcal{O})
\longrightarrow
(\mathcal{D}, \mathcal{O}')$
un morphisme de sites localement annelés, où $f$ est donné par le foncteur
continu $u : \mathcal{D} \to \mathcal{C}$. Soit $V$ un objet de $\mathcal{D}$
et posons $U = u(V)$. Alors le morphisme
$$
(f', (f')^\sharp) :
(\Sh(\mathcal{C}/U), \mathcal{O}_U)
\to
(\Sh(\mathcal{D}/V), \mathcal{O}'_V)
$$
du lemme \ref{lemma-localize-morphism-ringed-sites} est un morphisme de sites
localement annelés.
\item Soit
$(f, f^\sharp) :
(\mathcal{C}, \mathcal{O})
\longrightarrow
(\mathcal{D}, \mathcal{O}')$
un morphisme de sites localement annelés, où $f$ est donné par le foncteur
continu $u : \mathcal{D} \to \mathcal{C}$. Soient
$V \in \Ob(\mathcal{D})$, $U \in \Ob(\mathcal{C})$ et $c : U \to u(V)$.
Alors le morphisme
$$
(f_c, (f_c)^\sharp) :
(\Sh(\mathcal{C}/U), \mathcal{O}_U)
\to
(\Sh(\mathcal{D}/V), \mathcal{O}'_V)
$$
du lemme \ref{lemma-relocalize-morphism-ringed-sites} est un morphisme de topos
localement annelés.
\item Soit $(\Sh(\mathcal{C}), \mathcal{O})$ un topos localement annelé et
soit $\mathcal{F}$ un faisceau sur $\mathcal{C}$. Alors la localisation
$(\Sh(\mathcal{C})/\mathcal{F}, \mathcal{O}_\mathcal{F})$
est un topos localement annelé, et le morphisme de localisation
$$
(j_\mathcal{F}, j_\mathcal{F}^\sharp) :
(\Sh(\mathcal{C})/\mathcal{F}, \mathcal{O}_\mathcal{F})
\to
(\Sh(\mathcal{C}), \mathcal{O})
$$
est un morphisme de topos localement annelés.
\item Soit $(\Sh(\mathcal{C}), \mathcal{O})$ un topos localement annelé.
Soit $s : \mathcal{G} \to \mathcal{F}$ une application de faisceaux sur
$\mathcal{C}$. Alors le morphisme
$$
(j, j^\sharp) :
(\Sh(\mathcal{C})/\mathcal{G}, \mathcal{O}_\mathcal{G})
\longrightarrow
(\Sh(\mathcal{C})/\mathcal{F}, \mathcal{O}_\mathcal{F})
$$
du lemme \ref{lemma-relocalize-ringed-topos} est un morphisme de topos
localement annelés.
\item Soit
$f :
(\Sh(\mathcal{C}), \mathcal{O})
\longrightarrow
(\Sh(\mathcal{D}), \mathcal{O}')$
un morphisme de topos localement annelés. Soit $\mathcal{G}$ un faisceau sur
$\mathcal{D}$ et posons $\mathcal{F} = f^{-1}\mathcal{G}$. Alors le morphisme
$$
(f', (f')^\sharp) :
(\Sh(\mathcal{C})/\mathcal{F}, \mathcal{O}_\mathcal{F})
\longrightarrow
(\Sh(\mathcal{D})/\mathcal{G}, \mathcal{O}'_\mathcal{G})
$$
du lemme \ref{lemma-localize-morphism-ringed-topoi} est un morphisme de topos
localement annelés.
\item Soit
$f :
(\Sh(\mathcal{C}), \mathcal{O})
\longrightarrow
(\Sh(\mathcal{D}), \mathcal{O}')$
un morphisme de topos localement annelés. Soient $\mathcal{G}$ un faisceau sur
$\mathcal{D}$ et $\mathcal{F}$ un faisceau sur $\mathcal{C}$, et soit
$s : \mathcal{F} \to f^{-1}\mathcal{G}$ un morphisme de faisceaux. Alors le morphisme
$$
(f_s, (f_s)^\sharp) :
(\Sh(\mathcal{C})/\mathcal{F}, \mathcal{O}_\mathcal{F})
\longrightarrow
(\Sh(\mathcal{D})/\mathcal{G}, \mathcal{O}'_\mathcal{G})
$$
du lemme \ref{lemma-relocalize-morphism-ringed-topoi} est un morphisme de topos
localement annelés.
\end{enumerate}
\end{lemma}

\begin{proof}
La partie (1) est claire, puisque $\mathcal{O}_U$ n'est autre que la
restriction de $\mathcal{O}$; les lemmes
\ref{lemma-locally-ringed-intrinsic} et
\ref{lemma-locally-ringed-intrinsic-morphism} s'appliquent. La partie (2) est
claire, car $(j, j^\sharp)$ est une localisation d'un site localement annelé,
donc (1) s'applique. La partie (3) est également claire, puisque
$(f')^\sharp$ est la restriction de $f^\sharp$ au topos
$\Sh(\mathcal{C}/U)$; voir la démonstration du
lemme \ref{lemma-localize-morphism-ringed-topoi}. Le diagramme de la
définition \ref{definition-morphism-locally-ringed-topoi} pour $f'$ est donc
la restriction du diagramme correspondant pour $f$, et la restriction est un
foncteur exact. La partie (4) résulte formellement de (2) et (3). Les parties
(5), (6), (7) et (8) résultent respectivement de (1), (2), (3) et (4), en
agrandissant les sites comme dans le
lemme \ref{lemma-morphism-ringed-topoi-comes-from-morphism-ringed-sites} et en
traduisant tout en termes de sites et de morphismes de sites au moyen des comparaisons des lemmes
\ref{lemma-localize-compare},
\ref{lemma-relocalize-compare},
\ref{lemma-localize-morphism-compare} et
\ref{lemma-relocalize-morphism-compare}.
(On pourrait aussi appliquer directement aux parties (5), (6), (7) et (8) les
mêmes arguments que dans les démonstrations de (1), (2), (3) et (4).)
\end{proof}






\section{Image directe exceptionnelle pour les modules}
\label{section-lower-shriek-modules}

\noindent
Dans cette section, nous étendons au cas des faisceaux de modules la
construction de $g_!$ étudiée dans la section
\ref{section-exactness-lower-shriek}.

\begin{lemma}
\label{lemma-lower-shriek-modules}
Soit $u : \mathcal{C} \to \mathcal{D}$ un foncteur continu et cocontinu entre
sites. Notons $g : \Sh(\mathcal{C}) \to \Sh(\mathcal{D})$ le morphisme de
topos associé. Soit $\mathcal{O}_\mathcal{D}$ un faisceau d'anneaux sur
$\mathcal{D}$ et posons
$\mathcal{O}_\mathcal{C} = g^{-1}\mathcal{O}_\mathcal{D}$. Ainsi $g$ devient
un morphisme de topos annelés tel que $g^* = g^{-1}$. Il existe alors un foncteur
$$
g_! :
\textit{Mod}(\mathcal{O}_\mathcal{C})
\longrightarrow
\textit{Mod}(\mathcal{O}_\mathcal{D})
$$
adjoint à gauche de $g^*$.
\end{lemma}

\begin{proof}
Soit $U$ un objet de $\mathcal{C}$. Pour tout
$\mathcal{O}_\mathcal{D}$-module $\mathcal{G}$, on a
\begin{align*}
\Hom_{\mathcal{O}_\mathcal{C}}(j_{U!}\mathcal{O}_U, g^{-1}\mathcal{G})
& =
g^{-1}\mathcal{G}(U) \\
& =
\mathcal{G}(u(U)) \\
& =
\Hom_{\mathcal{O}_\mathcal{D}}(j_{u(U)!}\mathcal{O}_{u(U)}, \mathcal{G})
\end{align*}
car $g^{-1}$ se décrit par restriction; voir Sites,
lemme \ref{sites-lemma-when-shriek}. Une formule analogue vaut évidemment
pour une somme directe de modules de la forme $j_{U!}\mathcal{O}_U$. Par
Homologie, lemme \ref{homology-lemma-partially-defined-adjoint}, et par le
lemme \ref{lemma-module-quotient-flat}, on voit que $g_!$ existe.
\end{proof}

\begin{remark}
\label{remark-when-shriek-equal}
Attention ! Soient $u : \mathcal{C} \to \mathcal{D}$, $g$,
$\mathcal{O}_\mathcal{D}$ et $\mathcal{O}_\mathcal{C}$ comme dans le
lemme \ref{lemma-lower-shriek-modules}. En général, le diagramme
$$
\xymatrix{
\textit{Mod}(\mathcal{O}_\mathcal{C}) \ar[r]_{g_!} \ar[d]_{forget} &
\textit{Mod}(\mathcal{O}_\mathcal{D}) \ar[d]^{forget} \\
\textit{Ab}(\mathcal{C}) \ar[r]^{g^{Ab}_!} &
\textit{Ab}(\mathcal{D})
}
$$
ne commute {\bf pas} (ici, $g^{Ab}_!$ est celui du
lemme \ref{lemma-g-shriek-adjoint}). Il existe une transformation de foncteurs
$$
g_!^{Ab} \circ forget \longrightarrow forget \circ g_!
$$
La démonstration du lemme \ref{lemma-lower-shriek-modules} montre que c'est un
isomorphisme si et seulement si
$g^{Ab}_!j_{U!}\mathcal{O}_U \to g_!j_{U!}\mathcal{O}_U$
est un isomorphisme pour tout objet $U$ de $\mathcal{C}$. Comme
$g_!j_{U!}\mathcal{O}_U = j_{u(U)!}\mathcal{O}_{u(U)}$, cette condition équivaut à ce que
$$
g^{Ab}_!j_{U!}\mathcal{O}_U \longrightarrow j_{u(U)!}\mathcal{O}_{u(U)}
$$
soit un isomorphisme pour tout objet $U$ de $\mathcal{C}$. Pour un tel $U$,
on obtient un diagramme commutatif
$$
\xymatrix{
\mathcal{C}/U \ar[r]_-{u'} \ar[d]_{j_U} & \mathcal{D}/u(U) \ar[d]^{j_{u(U)}} \\
\mathcal{C} \ar[r]^u & \mathcal{D}
}
$$
de foncteurs cocontinus de sites; voir Sites,
lemme \ref{sites-lemma-localize-cocontinuous}. Par conséquent,
$g^{Ab}_!j_{U!} = j_{u(U)!}(g')^{Ab}_!$, où
$g' : \Sh(\mathcal{C}/U) \to \Sh(\mathcal{D}/u(U))$ est le morphisme de
topos induit par le foncteur cocontinu $u'$. On voit donc que
$g_! = g^{Ab}_!$ si l'application canonique
\begin{equation}
\label{equation-compare-on-localizations}
(g')^{Ab}_!\mathcal{O}_U \longrightarrow \mathcal{O}_{u(U)}
\end{equation}
est un isomorphisme pour tout objet $U$ de $\mathcal{C}$.
\end{remark}

\noindent
Les deux résultats suivants sont de nature légèrement différente.

\begin{lemma}
\label{lemma-special-square-cocontinuous}
Supposons donné un diagramme commutatif
$$
\xymatrix{
(\Sh(\mathcal{C}'), \mathcal{O}_{\mathcal{C}'})
\ar[r]_{(g', (g')^\sharp)} \ar[d]_{(f', (f')^\sharp)} &
(\Sh(\mathcal{C}), \mathcal{O}_\mathcal{C}) \ar[d]^{(f, f^\sharp)} \\
(\Sh(\mathcal{D}'), \mathcal{O}_{\mathcal{D}'}) \ar[r]^{(g, g^\sharp)} &
(\Sh(\mathcal{D}), \mathcal{O}_\mathcal{D})
}
$$
de topos annelés. Supposons que :
\begin{enumerate}
\item $f$, $f'$, $g$ et $g'$ correspondent aux foncteurs cocontinus
$u$, $u'$, $v$ et $v'$ comme dans Sites,
lemme \ref{sites-lemma-cocontinuous-morphism-topoi};
\item $v \circ u' = u \circ v'$,
\item $v$ et $v'$ sont à la fois continus et cocontinus;
\item pour tout objet $V'$ de $\mathcal{D}'$, le foncteur
${}^{u'}_{V'}\mathcal{I} \to {}^{\ \ \ u}_{v(V')}\mathcal{I}$
induit par $v$ est cofinal; et
\item $g^{-1}\mathcal{O}_{\mathcal{D}} = \mathcal{O}_{\mathcal{D}'}$
et $(g')^{-1}\mathcal{O}_{\mathcal{C}} = \mathcal{O}_{\mathcal{C}'}$.
\end{enumerate}
Alors, sur les modules, on a $f'_* \circ (g')^* = g^* \circ f_*$ et
$g'_! \circ (f')^{-1} = f^{-1} \circ g_!$.
\end{lemma}

\begin{proof}
On a $(g')^*\mathcal{F} = (g')^{-1}\mathcal{F}$ et
$g^*\mathcal{G} = g^{-1}\mathcal{G}$ en vertu de (5). La première égalité
résulte donc immédiatement de l'égalité correspondante dans Sites,
lemme \ref{sites-lemma-special-square-cocontinuous}. Comme les adjoints à
gauche $g_!$ et $g'_!$ de $g^*$ et $(g')^*$ existent par le
lemme \ref{lemma-lower-shriek-modules}, la seconde égalité résulte de l'unicité
des foncteurs adjoints.
\end{proof}

\begin{lemma}
\label{lemma-special-square-continuous}
Considérons un diagramme commutatif
$$
\xymatrix{
(\Sh(\mathcal{C}'), \mathcal{O}_{\mathcal{C}'})
\ar[r]_{(g', (g')^\sharp)} \ar[d]_{(f', (f')^\sharp)} &
(\Sh(\mathcal{C}), \mathcal{O}_\mathcal{C}) \ar[d]^{(f, f^\sharp)} \\
(\Sh(\mathcal{D}'), \mathcal{O}_{\mathcal{D}'}) \ar[r]^{(g, g^\sharp)} &
(\Sh(\mathcal{D}), \mathcal{O}_\mathcal{D})
}
$$
de topos annelés et supposons donnés des foncteurs
$$
\xymatrix{
\mathcal{C}' \ar[r]_{v'} &
\mathcal{C} \\
\mathcal{D}' \ar[r]^v \ar[u]^{u'} &
\mathcal{D} \ar[u]_u
}
$$
tels que (avec les notations de Sites, sections
\ref{sites-section-morphism-sites} et
\ref{sites-section-cocontinuous-morphism-topoi}) :
\begin{enumerate}
\item $u$ et $u'$ sont continus et donnent les morphismes $f$ et $f'$;
\item $v$ et $v'$ sont cocontinus et donnent les morphismes $g$ et $g'$;
\item $u \circ v = v' \circ u'$,
\item $v$ et $v'$ sont à la fois continus et cocontinus; et
\item $g^{-1}\mathcal{O}_{\mathcal{D}} = \mathcal{O}_{\mathcal{D}'}$
et $(g')^{-1}\mathcal{O}_{\mathcal{C}} = \mathcal{O}_{\mathcal{C}'}$.
\end{enumerate}
Alors, sur les modules, $f'_* \circ (g')^* = g^* \circ f_*$ et
$g'_! \circ (f')^{-1} = f^{-1} \circ g_!$.
\end{lemma}

\begin{proof}
On a $(g')^*\mathcal{F} = (g')^{-1}\mathcal{F}$ et
$g^*\mathcal{G} = g^{-1}\mathcal{G}$ en vertu de (5). La première égalité
résulte immédiatement de l'égalité correspondante dans Sites,
lemme \ref{sites-lemma-special-square-continuous}. Comme les adjoints à gauche
$g_!$ et $g'_!$ de $g^*$ et $(g')^*$ existent par le
lemme \ref{lemma-lower-shriek-modules}, la seconde égalité résulte de l'unicité
des foncteurs adjoints.
\end{proof}







\section{Faisceaux constants}
\label{section-constant}

\noindent
Soit $E$ un ensemble et soit $\mathcal{C}$ un site. Le faisceau constant de
valeur $E$ est noté $\underline{E}$; voir Sites,
exemple \ref{sites-example-constant-sheaf}. Si $E$ est un groupe abélien, un
anneau, un module, etc., alors $\underline{E}$ est un faisceau de groupes
abéliens, d'anneaux, de modules, etc.

\begin{lemma}
\label{lemma-constant-exact}
Soit $\mathcal{C}$ un site. Si $0 \to A \to B \to C \to 0$ est une suite
exacte courte de groupes abéliens, alors
$0 \to \underline{A} \to \underline{B} \to \underline{C} \to 0$
est une suite exacte de faisceaux abéliens; elle est même exacte comme suite
de préfaisceaux abéliens.
\end{lemma}

\begin{proof}
Comme la faisceautisation est exacte, il est clair que
$0 \to \underline{A} \to \underline{B} \to \underline{C} \to 0$
est une suite exacte de faisceaux abéliens. Ainsi
$0 \to \underline{A} \to \underline{B} \to \underline{C}$
est une suite exacte de préfaisceaux abéliens. Pour voir que
$\underline{B} \to \underline{C}$ est surjective, choisissons une section
ensembliste $s : C \to B$. Elle induit une section
$\underline{s} : \underline{C} \to \underline{B}$ de faisceaux d'ensembles,
inverse à droite de la surjection $\underline{B} \to \underline{C}$.
\end{proof}

\begin{lemma}
\label{lemma-tensor-with-finitely-presented}
Soit $\mathcal{C}$ un site. Soit $\Lambda$ un anneau et soient $M$ et $Q$ des
$\Lambda$-modules. Si $Q$ est un $\Lambda$-module de présentation finie, alors
$\underline{M \otimes_\Lambda Q}(U) = \underline{M}(U) \otimes_\Lambda Q$
pour tout $U \in \Ob(\mathcal{C})$.
\end{lemma}

\begin{proof}
Choisissons une présentation
$\Lambda^{\oplus m} \to \Lambda^{\oplus n} \to Q \to 0$. Elle donne une
suite exacte $M^{\oplus m} \to M^{\oplus n} \to M \otimes Q \to 0$.
Par le lemme \ref{lemma-constant-exact}, on obtient la suite exacte
$$
\underline{M}(U)^{\oplus m} \to
\underline{M}(U)^{\oplus n} \to
\underline{M \otimes Q}(U) \to 0
$$
ce qui démontre le lemme. (Remarquons que la prise de sections sur $U$ commute
toujours aux sommes directes finies, mais non aux sommes directes arbitraires.)
\end{proof}

\begin{lemma}
\label{lemma-flat-sections}
Soit $\mathcal{C}$ un site. Soit $\Lambda$ un anneau cohérent et soit $M$ un
$\Lambda$-module plat. Pour $U \in \Ob(\mathcal{C})$, le module
$\underline{M}(U)$ est un $\Lambda$-module plat.
\end{lemma}

\begin{proof}
Soit $I \subset \Lambda$ un idéal de type fini. Par Algèbre,
lemme \ref{algebra-lemma-flat}, il suffit de montrer que
$\underline{M}(U) \otimes_\Lambda I \to \underline{M}(U)$
est injective. Comme $\Lambda$ est cohérent, $I$ est un $\Lambda$-module de
présentation finie. Par le lemme \ref{lemma-tensor-with-finitely-presented},
$\underline{M}(U) \otimes I = \underline{M \otimes I}$. Puisque $M$ est plat,
l'application $M \otimes I \to M$ est injective; il en va donc de même de
$\underline{M \otimes I} \to \underline{M}$.
\end{proof}

\begin{lemma}
\label{lemma-completion-flat}
Soit $\mathcal{C}$ un site. Soit $\Lambda$ un anneau noethérien et soit
$I \subset \Lambda$ un idéal. Le faisceau
$\underline{\Lambda}^\wedge = \lim \underline{\Lambda/I^n}$
est une $\underline{\Lambda}$-algèbre plate. De plus, on a les
identifications canoniques
$$
\underline{\Lambda}/I\underline{\Lambda} =
\underline{\Lambda}/\underline{I} =
\underline{\Lambda}^\wedge/I\underline{\Lambda}^\wedge =
\underline{\Lambda}^\wedge/\underline{I} \cdot \underline{\Lambda}^\wedge =
\underline{\Lambda}^\wedge/\underline{I}^\wedge =
\underline{\Lambda/I}
$$
où $\underline{I}^\wedge = \lim \underline{I/I^n}$.
\end{lemma}

\begin{proof}
Pour montrer que $\underline{\Lambda}^\wedge$ est plat, il suffit de vérifier
que $\underline{\Lambda}^\wedge(U)$ est un $\Lambda$-module plat pour tout
$U \in \Ob(\mathcal{C})$; voir les lemmes
\ref{lemma-flatness-presheaves} et \ref{lemma-flatness-sheafification}.
Par le lemme \ref{lemma-flat-sections},
$$
\underline{\Lambda}^\wedge(U) = \lim \underline{\Lambda/I^n}(U)
$$
est la limite d'un système de $\Lambda/I^n$-modules plats. Le
lemme \ref{lemma-constant-exact} montre que les applications de transition
sont surjectives. On conclut par Compléments d'algèbre,
lemme \ref{more-algebra-lemma-limit-flat}.

\medskip\noindent
Pour établir les égalités, remarquons que
$\underline{\Lambda}(U)/I\underline{\Lambda}(U) = \underline{\Lambda/I}(U)$
par le lemme \ref{lemma-tensor-with-finitely-presented}. Il s'ensuit que
$\underline{\Lambda}/I\underline{\Lambda} =
\underline{\Lambda}/\underline{I} = \underline{\Lambda/I}$. Le système de
suites exactes courtes
$$
0 \to \underline{I/I^n}(U) \to \underline{\Lambda/I^n}(U) \to
\underline{\Lambda/I}(U) \to 0
$$
a des applications de transition surjectives et donne donc une suite exacte courte
$$
0 \to \lim \underline{I/I^n}(U) \to \lim \underline{\Lambda/I^n}(U) \to
\lim \underline{\Lambda/I}(U) \to 0
$$
voir Homologie, lemme \ref{homology-lemma-Mittag-Leffler}. Ainsi,
$\underline{\Lambda}^\wedge/\underline{I}^\wedge =
\underline{\Lambda/I}$. Comme
$$
I \underline{\Lambda}^\wedge \subset
\underline{I} \cdot \underline{\Lambda}^\wedge \subset
\underline{I}^\wedge
$$
il suffit de montrer que
$I \underline{\Lambda}^\wedge(U) = \underline{I}^\wedge(U)$
pour tout $U$. Choisissons des générateurs $I = (f_1, \ldots, f_r)$.
Pour tout $n$, on obtient une suite exacte courte
$$
0 \to K_n/(I^n)^{\oplus r} \to (\Lambda/I^n)^{\oplus r}
\xrightarrow{(f_1, \ldots, f_r)}
I/I^{n + 1} \to 0
$$
où $K_n =
\{(x_1, \ldots, x_r) \in \Lambda^{\oplus r} \mid \sum x_i f_i \in I^{n + 1}\}$.
On obtient des suites exactes courtes
$$
0 \to \underline{K_n/(I^n)^{\oplus r}}(U) \to
\underline{(\Lambda/I^n)^{\oplus r}}(U) \to
\underline{I/I^{n + 1}}(U) \to 0
$$
Un calcul donne $K_n = K + (I^n)^{\oplus r}$; les applications de transition
$K_{n + 1}/(I^{n + 1})^{\oplus r} \to K_n/(I^n)^{\oplus r}$ sont donc
surjectives. Le système de modules de gauche a ainsi des applications de
transition surjectives et vérifie a fortiori la condition de Mittag--Leffler.
On voit donc que
$(f_1, \ldots, f_r) :
(\underline{\Lambda}^\wedge)^{\oplus r}(U) \to \underline{I}^\wedge(U)$
est surjective par Homologie, lemme \ref{homology-lemma-Mittag-Leffler}, ce
qu'il fallait démontrer.
\end{proof}

\begin{lemma}
\label{lemma-locally-constant-finite-type}
Soit $\mathcal{C}$ un site. Soit $\Lambda$ un anneau et soit $M$ un
$\Lambda$-module. Supposons que $\Sh(\mathcal{C})$ ne soit pas le topos vide. Alors
\begin{enumerate}
\item $\underline{M}$ est un faisceau de $\underline{\Lambda}$-modules de type
fini si et seulement si $M$ est un $\Lambda$-module fini; et
\item $\underline{M}$ est un faisceau de $\underline{\Lambda}$-modules de
présentation finie si et seulement si $M$ est un $\Lambda$-module de
présentation finie.
\end{enumerate}
\end{lemma}

\begin{proof}
Démonstration de (1). Si $M$ est engendré par $x_1, \ldots, x_r$, alors
$x_1, \ldots, x_r$ définissent des sections globales de $\underline{M}$ qui
l'engendrent; ainsi
$\underline{M}$ est de type fini. Réciproquement, supposons
$\underline{M}$ de type fini. Soit $U \in \mathcal{C}$ un objet qui n'est pas
vide au sens faisceautique (Sites, définition \ref{sites-definition-empty}).
Un tel objet existe puisque $\Sh(\mathcal{C})$ n'est pas le topos vide. Il
existe alors un recouvrement $\{U_i \to U\}$ et un nombre fini de sections
$s_{ij} \in \underline{M}(U_i)$ engendrant $\underline{M}|_{U_i}$. Après
raffinement du recouvrement, on peut supposer que les $s_{ij}$ proviennent
d'éléments $x_{ij}$ de $M$. Les $x_{ij}$ définissent des sections globales de
$\underline{M}$ dont les restrictions à $U$ engendrent $\underline{M}$.

\medskip\noindent
Supposons qu'il existe des éléments $x_1, \ldots, x_r$ de $M$ définissant des
sections globales de $\underline{M}$ qui engendrent $\underline{M}$ comme
faisceau de $\underline{\Lambda}$-modules. Montrons que $x_1, \ldots, x_r$
engendrent $M$ comme $\Lambda$-module. Soit $x \in M$. Il existe un recouvrement
$\{U_i \to U\}_{i \in I}$ et des $f_{i, j} \in \underline{\Lambda}(U_i)$ tels
que $x|_{U_i} = \sum f_{i, j} x_j|_{U_i}$. Après raffinement, on peut supposer
$f_{i, j} \in \Lambda$. Puisque $U$ n'est pas vide au sens faisceautique, il existe
au moins un $i \in I$ tel que $U_i$ ne soit pas non plus vide au sens
faisceautique. L'application
$M \to \underline{M}(U_i)$ est alors injective (détails omis). On conclut que
$x = \sum f_{i, j}x_j$ dans $M$, comme voulu.

\medskip\noindent
Démonstration de (2). Supposons $\underline{M}$ de présentation finie comme
$\underline{\Lambda}$-module. Par (1), $M$ est de type fini. Choisissons des
générateurs $x_1, \ldots, x_r$ de $M$ comme $\Lambda$-module. On obtient une
suite exacte courte $0 \to K \to \Lambda^{\oplus r} \to M \to 0$, qui devient
$$
0 \to \underline{K} \to \underline{\Lambda}^{\oplus r} \to \underline{M} \to 0
$$
par le lemme \ref{lemma-constant-exact}. Le
lemme \ref{lemma-kernel-surjection-finite-onto-finite-presentation} montre que
$\underline{K}$ est de type fini. Par (1), $K$ est donc un
$\Lambda$-module fini, et $M$ est donc un $\Lambda$-module de présentation finie.
\end{proof}




\section{Faisceaux localement constants}
\label{section-locally-constant}

\noindent
Voici la définition générale.

\begin{definition}
\label{definition-locally-constant}
Soit $\mathcal{C}$ un site. Soit $\mathcal{F}$ un faisceau d'ensembles, de
groupes, de groupes abéliens, d'anneaux, de modules sur un anneau fixé
$\Lambda$, etc.
\begin{enumerate}
\item On dit que $\mathcal{F}$ est un {\it faisceau constant} d'ensembles, de
groupes, de groupes abéliens, d'anneaux, de modules sur $\Lambda$, etc., s'il
est isomorphe, comme faisceau d'ensembles, de groupes, de groupes abéliens,
d'anneaux ou de modules sur $\Lambda$, à un faisceau
constant $\underline{E}$ comme dans la section \ref{section-constant}.
\item On dit que $\mathcal{F}$ est {\it localement constant} si, pour tout
objet $U$ de $\mathcal{C}$, il existe un recouvrement $\{U_i \to U\}$ tel que
$\mathcal{F}|_{U_i}$ soit un faisceau constant.
\item Si $\mathcal{F}$ est un faisceau d'ensembles ou de groupes, on dit que $\mathcal{F}$ est
{\it fini localement constant} si les valeurs constantes sont des ensembles
finis ou des groupes finis.
\end{enumerate}
\end{definition}

\begin{lemma}
\label{lemma-pullback-locally-constant}
Soit $f : \Sh(\mathcal{C}) \to \Sh(\mathcal{D})$ un morphisme de topos.
Si $\mathcal{G}$ est un faisceau localement constant d'ensembles, de groupes,
de groupes abéliens, d'anneaux, de modules sur un anneau fixé $\Lambda$, etc.,
sur $\mathcal{D}$, alors $f^{-1}\mathcal{G}$ possède la même propriété sur
$\mathcal{C}$.
\end{lemma}

\begin{proof}
La démonstration est omise.
\end{proof}

\begin{lemma}
\label{lemma-morphism-locally-constant}
Soit $\mathcal{C}$ un site possédant un objet final $X$.
\begin{enumerate}
\item Soit $\varphi : \mathcal{F} \to \mathcal{G}$ une application de
faisceaux d'ensembles localement constants sur $\mathcal{C}$. Si
$\mathcal{F}$ est fini localement constant, il existe un recouvrement
$\{U_i \to X\}$ tel que $\varphi|_{U_i}$ soit l'application de faisceaux
constants associée à une application d'ensembles.
\item Soit $\varphi : \mathcal{F} \to \mathcal{G}$ une application de
faisceaux de groupes abéliens localement constants sur $\mathcal{C}$. Si
$\mathcal{F}$ est fini localement constant, il existe un recouvrement
$\{U_i \to X\}$ tel que $\varphi|_{U_i}$ soit l'application de faisceaux
abéliens constants associée à un homomorphisme de groupes abéliens.
\item Soit $\Lambda$ un anneau et soit
$\varphi : \mathcal{F} \to \mathcal{G}$ une application de faisceaux de
$\Lambda$-modules localement constants sur $\mathcal{C}$. Si $\mathcal{F}$
est de type fini, il existe un recouvrement $\{U_i \to X\}$ tel que
$\varphi|_{U_i}$ soit l'application de faisceaux constants de
$\Lambda$-modules associée à une application de $\Lambda$-modules.
\end{enumerate}
\end{lemma}

\begin{proof}
La démonstration est omise.
\end{proof}

\begin{lemma}
\label{lemma-locally-constant}
Soit $\mathcal{C}$ un site. Soit $\Lambda$ un anneau. Soient $M$, $N$ des
$\Lambda$-modules et $\mathcal{F}, \mathcal{G}$ des faisceaux localement
constants de $\Lambda$-modules.
\begin{enumerate}
\item Si $M$ est de présentation finie, alors
$$
\underline{\Hom_\Lambda(M, N)} =
\SheafHom_{\underline{\Lambda}}(\underline{M}, \underline{N})
$$
\item Si $M$ et $N$ sont tous deux de présentation finie, alors
$$
\underline{\text{Isom}_\Lambda(M, N)} =
\mathit{Isom}_{\underline{\Lambda}}(\underline{M}, \underline{N})
$$
\item Si $\mathcal{F}$ est de présentation finie, alors
$\SheafHom_{\underline{\Lambda}}(\mathcal{F}, \mathcal{G})$
est un faisceau localement constant de $\Lambda$-modules.
\item Si $\mathcal{F}$ et $\mathcal{G}$ sont tous deux de présentation finie, alors
$\mathit{Isom}_{\underline{\Lambda}}(\mathcal{F}, \mathcal{G})$
est un faisceau d'ensembles localement constant.
\end{enumerate}
\end{lemma}

\begin{proof}
Démonstration de (1). Posons $E = \Hom_\Lambda(M, N)$. Nous voulons montrer que
l'application canonique
$$
\underline{E}
\longrightarrow
\SheafHom_{\underline{\Lambda}}(\underline{M}, \underline{N})
$$
est un isomorphisme. Le module $M$ admet une présentation
$\Lambda^{\oplus s} \to \Lambda^{\oplus t} \to M \to 0$.
Alors $E$ s'insère dans une suite exacte
$$
0 \to E \to \Hom_\Lambda(\Lambda^{\oplus t}, N) \to
\Hom_\Lambda(\Lambda^{\oplus s}, N)
$$
et l'on a de même
$$
0 \to
\SheafHom_{\underline{\Lambda}}(\underline{M}, \underline{N}) \to
\SheafHom_{\underline{\Lambda}}(\underline{\Lambda^{\oplus t}}, \underline{N})
\to
\SheafHom_{\underline{\Lambda}}(\underline{\Lambda^{\oplus s}}, \underline{N})
$$
Cela ramène la question au cas où $M$ est libre fini, dans lequel le résultat
est clair.

\medskip\noindent
Démonstration de (3). La question est locale sur $\mathcal{C}$; on peut donc
supposer $\mathcal{F} = \underline{M}$ et $\mathcal{G} = \underline{N}$ pour des
$\Lambda$-modules $M$ et $N$. Par le
lemme \ref{lemma-locally-constant-finite-type}, $M$ est de présentation finie.
Le résultat découle alors de (1).

\medskip\noindent
Les parties (2) et (4) résultent de (1) et (3), ainsi que du fait que
$\mathit{Isom}$ peut être vu comme le sous-faisceau des sections de
$\SheafHom_{\underline{\Lambda}}(\mathcal{F}, \mathcal{G})$
qui admettent un inverse dans
$\SheafHom_{\underline{\Lambda}}(\mathcal{G}, \mathcal{F})$.
\end{proof}

\begin{lemma}
\label{lemma-kernel-finite-locally-constant}
Soit $\mathcal{C}$ un site.
\begin{enumerate}
\item La catégorie des faisceaux d'ensembles finis localement constants est
stable par limites et colimites finies dans $\Sh(\mathcal{C})$.
\item La catégorie des faisceaux abéliens finis localement constants est une
sous-catégorie de Serre faible de $\textit{Ab}(\mathcal{C})$.
\item Soit $\Lambda$ un anneau noethérien. La catégorie des faisceaux de
$\Lambda$-modules localement constants et de type fini sur $\mathcal{C}$ est
une sous-catégorie de Serre faible de $\textit{Mod}(\mathcal{C}, \Lambda)$.
\end{enumerate}
\end{lemma}

\begin{proof}
Démonstration de (1). On peut travailler localement sur $\mathcal{C}$. Par le
lemme \ref{lemma-morphism-locally-constant}, on peut donc supposer donné un
diagramme fini d'ensembles finis tel que notre diagramme de faisceaux soit le
diagramme associé de faisceaux constants. Il suffit alors de prendre la limite
ou la colimite dans la catégorie des ensembles, puis le faisceau constant
associé. Certains détails sont omis.

\medskip\noindent
Pour démontrer (2) et (3), appliquons le critère d'Homologie,
lemme \ref{homology-lemma-characterize-weak-serre-subcategory}. L'existence
des noyaux et conoyaux s'établit comme ci-dessus. L'hypothèse noethérienne dans
(3) garantit que le noyau d'une application de $\Lambda$-modules finis est un
$\Lambda$-module fini. Pour voir que la catégorie est stable par extensions
(dans le cas des faisceaux de $\Lambda$-modules), considérons une extension de
faisceaux de $\Lambda$-modules
$$
0 \to \mathcal{F} \to \mathcal{E} \to \mathcal{G} \to 0
$$
sur $\mathcal{C}$, où $\mathcal{F}$ et $\mathcal{G}$ sont de type fini et
localement constants. Après localisation sur $\mathcal{C}$, on peut supposer
$\mathcal{F}$ et $\mathcal{G}$ constants, c'est-à-dire que l'on obtient
$$
0 \to \underline{M} \to \mathcal{E} \to \underline{N} \to 0
$$
pour certains $\Lambda$-modules $M, N$. Choisissons des générateurs
$y_1, \ldots, y_m$ de $N$, d'où une suite exacte courte
$0 \to K \to \Lambda^{\oplus m} \to N \to 0$ de $\Lambda$-modules. Après une nouvelle
localisation, on peut supposer que $y_j$ se relève en une section $s_j$ de
$\mathcal{E}$. Ainsi $\mathcal{E}$ est un poussé comme dans le diagramme suivant
$$
\xymatrix{
0 \ar[r] &
\underline{K} \ar[d] \ar[r] &
\underline{\Lambda^{\oplus m}} \ar[d] \ar[r] &
\underline{N} \ar[d] \ar[r] & 0 \\
0 \ar[r] &
\underline{M} \ar[r] &
\mathcal{E} \ar[r] &
\underline{N} \ar[r] & 0
}
$$
Par une nouvelle application du lemme \ref{lemma-morphism-locally-constant}
(et parce que $K$ est un $\Lambda$-module fini puisque $\Lambda$ est
noethérien), on voit que l'application
$\underline{K} \to \underline{M}$ est localement constante, ce qui permet de conclure.
\end{proof}

\begin{lemma}
\label{lemma-tensor-product-locally-constant}
Soit $\mathcal{C}$ un site et soit $\Lambda$ un anneau. Le produit tensoriel de
deux faisceaux localement constants de $\Lambda$-modules sur $\mathcal{C}$ est
un faisceau localement constant de $\Lambda$-modules.
\end{lemma}

\begin{proof}
La démonstration est omise.
\end{proof}











\section{Localisation des faisceaux d'anneaux}
\label{section-localizing-sheaves-rings}

\noindent
Soit $(\mathcal{C}, \mathcal{O})$ un site annelé. Soit
$\mathcal{S} \subset \mathcal{O}$ un sous-préfaisceau d'ensembles tel que,
pour tout $U \in \Ob(\mathcal{C})$, l'ensemble
$\mathcal{S}(U) \subset \mathcal{O}(U)$ soit une partie multiplicative; voir
Algèbre, définition \ref{algebra-definition-multiplicative-subset}. On peut
alors considérer le préfaisceau d'anneaux
$$
\mathcal{S}^{-1}\mathcal{O} :
U \longmapsto \mathcal{S}(U)^{-1}\mathcal{O}(U).
$$
L'application de restriction envoie la section $f/s$, avec
$f \in \mathcal{O}(U)$ et $s \in \mathcal{S}(U)$, sur
$(f|_V)/(s|_V)$ pour $V \to U$ dans $\mathcal{C}$.

\begin{lemma}
\label{lemma-simple-invert}
Dans la situation ci-dessus, l'application vers le faisceautisé
$$
\mathcal{O} \longrightarrow (\mathcal{S}^{-1}\mathcal{O})^\#
$$
est un homomorphisme de faisceaux d'anneaux ayant la propriété universelle
suivante : pour tout homomorphisme de faisceaux d'anneaux
$\mathcal{O} \to \mathcal{A}$ qui envoie chaque section locale de
$\mathcal{S}$ sur une section inversible de $\mathcal{A}$, il existe une
unique factorisation $(\mathcal{S}^{-1}\mathcal{O})^\# \to \mathcal{A}$.
\end{lemma}

\begin{proof}
La démonstration est omise.
\end{proof}

\noindent
Soit $(\mathcal{C}, \mathcal{O})$ un site annelé. Soit
$\mathcal{S} \subset \mathcal{O}$ un sous-préfaisceau d'ensembles tel que,
pour tout $U \in \mathcal{C}$, l'ensemble
$\mathcal{S}(U) \subset \mathcal{O}(U)$ soit une partie multiplicative. Soit
$\mathcal{F}$ un faisceau de $\mathcal{O}$-modules. On peut alors considérer
le préfaisceau de $\mathcal{S}^{-1}\mathcal{O}$-modules
$$
\mathcal{S}^{-1}\mathcal{F} :
U \longmapsto \mathcal{S}(U)^{-1}\mathcal{F}(U).
$$
L'application de restriction envoie la section $t/s$, avec
$t \in \mathcal{F}(U)$ et $s \in \mathcal{S}(U)$, sur
$(t|_V)/(s|_V)$ si $V \to U$ est un morphisme de $\mathcal{C}$. Ainsi,
$\mathcal{S}^{-1}\mathcal{F}$ est un préfaisceau de
$\mathcal{S}^{-1}\mathcal{O}$-modules.

\begin{lemma}
\label{lemma-simple-invert-module}
Dans la situation ci-dessus, l'application vers le faisceautisé
$$
\mathcal{F} \longrightarrow (\mathcal{S}^{-1}\mathcal{F})^\#
$$
possède la propriété universelle suivante : pour tout homomorphisme de
$\mathcal{O}$-modules $\mathcal{F} \to \mathcal{G}$ tel que chaque section
locale de $\mathcal{S}$ agisse de façon inversible sur $\mathcal{G}$, il
existe une unique factorisation
$(\mathcal{S}^{-1}\mathcal{F})^\# \to \mathcal{G}$. De plus, on a
$$
(\mathcal{S}^{-1}\mathcal{F})^\#
=
(\mathcal{S}^{-1}\mathcal{O})^\# \otimes_\mathcal{O} \mathcal{F}
$$
comme faisceaux de $(\mathcal{S}^{-1}\mathcal{O})^\#$-modules.
\end{lemma}

\begin{proof}
La démonstration est omise.
\end{proof}









\section{Faisceaux d'ensembles pointés}
\label{section-pointed}

\noindent
Dans cette section, nous rassemblons quelques faits sur les faisceaux
d'ensembles pointés, que nous n'avions jusqu'ici mentionnés que pour les
faisceaux abéliens.

\medskip\noindent
Un ensemble pointé est une paire $(S, 0)$, où $S$ est un ensemble et
$0 \in S$ est un élément de $S$. Un morphisme $(S, 0) \to (S', 0')$ d'ensembles pointés est simplement
une application $S \to S'$ envoyant $0$ sur $0'$. Par abus de notation, nous
dirons « soit $S$ un ensemble pointé » pour signifier que $S$ est muni d'un
élément distingué $0 \in S$. Un faisceau d'ensembles pointés est la même chose
qu'un faisceau d'ensembles $\mathcal{F}$ muni d'un « pointage »
$0 : * \to \mathcal{F}$, où $*$ est le faisceau final
(Sites, exemple \ref{sites-example-singleton-sheaf}).

\medskip\noindent
Pour un morphisme de sites ou de topos, il existe des foncteurs image directe
et image inverse sur les catégories de faisceaux d'ensembles pointés; voir
Sites, section \ref{sites-section-sheaves-algebraic-structures}. Ils se
construisent en prenant l'image directe, respectivement l'image inverse, du
faisceau d'ensembles sous-jacent et en le pointant convenablement (en utilisant
que l'image inverse du faisceau final est le faisceau final).

\medskip\noindent
Soit $u : \mathcal{C} \to \mathcal{D}$ un foncteur continu et cocontinu entre
sites. Soit $g : \Sh(\mathcal{C}) \to \Sh(\mathcal{D})$ le morphisme de topos
associé à $u$; voir Sites,
lemme \ref{sites-lemma-cocontinuous-morphism-topoi}. Sur les faisceaux
d'ensembles pointés, $g^{-1}$ admet également un adjoint à gauche $g_!$. Sa
construction est entièrement analogue à celle de $g_!$ sur les faisceaux
abéliens dans la section \ref{section-exactness-lower-shriek}.

\medskip\noindent
De même, si $j : \mathcal{C}/U \to \mathcal{C}$ est comme dans la
section \ref{section-localize}, alors il existe un adjoint à gauche $j_!$ du
foncteur $j^{-1}$ sur les faisceaux d'ensembles pointés.

\medskip\noindent
Si nous avons un jour besoin de ces faits et constructions, nous énoncerons et
démontrerons ici précisément les lemmes correspondants.






\begin{multicols}{2}[\section{Autres chapitres}]
\noindent
Préliminaires
\begin{enumerate}
\item \hyperref[introduction-section-phantom]{Introduction}
\item \hyperref[conventions-section-phantom]{Conventions}
\item \hyperref[sets-section-phantom]{Théorie des ensembles}
\item \hyperref[categories-section-phantom]{Catégories}
\item \hyperref[topology-section-phantom]{Topologie}
\item \hyperref[sheaves-section-phantom]{Faisceaux sur les espaces}
\item \hyperref[sites-section-phantom]{Sites et faisceaux}
\item \hyperref[stacks-section-phantom]{Champs}
\item \hyperref[fields-section-phantom]{Corps}
\item \hyperref[algebra-section-phantom]{Algèbre commutative}
\item \hyperref[brauer-section-phantom]{Groupes de Brauer}
\item \hyperref[homology-section-phantom]{Algèbre homologique}
\item \hyperref[derived-section-phantom]{Catégories dérivées}
\item \hyperref[simplicial-section-phantom]{Méthodes simpliciales}
\item \hyperref[more-algebra-section-phantom]{Compléments d'algèbre}
\item \hyperref[smoothing-section-phantom]{Lissage des morphismes d'anneaux}
\item \hyperref[modules-section-phantom]{Faisceaux de modules}
\item \hyperref[sites-modules-section-phantom]{Modules sur les sites}
\item \hyperref[injectives-section-phantom]{Objets injectifs}
\item \hyperref[cohomology-section-phantom]{Cohomologie des faisceaux}
\item \hyperref[sites-cohomology-section-phantom]{Cohomologie sur les sites}
\item \hyperref[dga-section-phantom]{Algèbre différentielle graduée}
\item \hyperref[dpa-section-phantom]{Algèbre à puissances divisées}
\item \hyperref[sdga-section-phantom]{Faisceaux différentiels gradués}
\item \hyperref[hypercovering-section-phantom]{Hyperrecouvrements}
\end{enumerate}
Schémas
\begin{enumerate}
\setcounter{enumi}{25}
\item \hyperref[schemes-section-phantom]{Schémas}
\item \hyperref[constructions-section-phantom]{Constructions de schémas}
\item \hyperref[properties-section-phantom]{Propriétés des schémas}
\item \hyperref[morphisms-section-phantom]{Morphismes de schémas}
\item \hyperref[coherent-section-phantom]{Cohomologie des schémas}
\item \hyperref[divisors-section-phantom]{Diviseurs}
\item \hyperref[limits-section-phantom]{Limites de schémas}
\item \hyperref[varieties-section-phantom]{Variétés}
\item \hyperref[topologies-section-phantom]{Topologies sur les schémas}
\item \hyperref[descent-section-phantom]{Descente}
\item \hyperref[perfect-section-phantom]{Catégories dérivées de schémas}
\item \hyperref[more-morphisms-section-phantom]{Compléments sur les morphismes}
\item \hyperref[flat-section-phantom]{Compléments sur la platitude}
\item \hyperref[groupoids-section-phantom]{Schémas en groupoïdes}
\item \hyperref[more-groupoids-section-phantom]{Compléments sur les schémas en groupoïdes}
\item \hyperref[etale-section-phantom]{Morphismes étales de schémas}
\end{enumerate}
Thèmes de théorie des schémas
\begin{enumerate}
\setcounter{enumi}{41}
\item \hyperref[chow-section-phantom]{Homologie de Chow}
\item \hyperref[intersection-section-phantom]{Théorie de l'intersection}
\item \hyperref[pic-section-phantom]{Schémas de Picard des courbes}
\item \hyperref[weil-section-phantom]{Théories cohomologiques de Weil}
\item \hyperref[adequate-section-phantom]{Modules adéquats}
\item \hyperref[dualizing-section-phantom]{Complexes dualisants}
\item \hyperref[duality-section-phantom]{Dualité pour les schémas}
\item \hyperref[discriminant-section-phantom]{Discriminants et différentes}
\item \hyperref[derham-section-phantom]{Cohomologie de de Rham}
\item \hyperref[local-cohomology-section-phantom]{Cohomologie locale}
\item \hyperref[algebraization-section-phantom]{Géométrie algébrique et formelle}
\item \hyperref[curves-section-phantom]{Courbes algébriques}
\item \hyperref[resolve-section-phantom]{Résolution des surfaces}
\item \hyperref[models-section-phantom]{Réduction semi-stable}
\item \hyperref[functors-section-phantom]{Foncteurs et morphismes}
\item \hyperref[equiv-section-phantom]{Catégories dérivées de variétés}
\item \hyperref[pione-section-phantom]{Groupes fondamentaux de schémas}
\item \hyperref[etale-cohomology-section-phantom]{Cohomologie étale}
\item \hyperref[crystalline-section-phantom]{Cohomologie cristalline}
\item \hyperref[proetale-section-phantom]{Cohomologie pro-étale}
\item \hyperref[relative-cycles-section-phantom]{Cycles relatifs}
\item \hyperref[more-etale-section-phantom]{Compléments de cohomologie étale}
\item \hyperref[trace-section-phantom]{La formule des traces}
\end{enumerate}
Espaces algébriques
\begin{enumerate}
\setcounter{enumi}{64}
\item \hyperref[spaces-section-phantom]{Espaces algébriques}
\item \hyperref[spaces-properties-section-phantom]{Propriétés des espaces algébriques}
\item \hyperref[spaces-morphisms-section-phantom]{Morphismes d'espaces algébriques}
\item \hyperref[decent-spaces-section-phantom]{Espaces algébriques décents}
\item \hyperref[spaces-cohomology-section-phantom]{Cohomologie des espaces algébriques}
\item \hyperref[spaces-limits-section-phantom]{Limites d'espaces algébriques}
\item \hyperref[spaces-divisors-section-phantom]{Diviseurs sur les espaces algébriques}
\item \hyperref[spaces-over-fields-section-phantom]{Espaces algébriques sur des corps}
\item \hyperref[spaces-topologies-section-phantom]{Topologies sur les espaces algébriques}
\item \hyperref[spaces-descent-section-phantom]{Descente et espaces algébriques}
\item \hyperref[spaces-perfect-section-phantom]{Catégories dérivées d'espaces}
\item \hyperref[spaces-more-morphisms-section-phantom]{Compléments sur les morphismes d'espaces}
\item \hyperref[spaces-flat-section-phantom]{Platitude sur les espaces algébriques}
\item \hyperref[spaces-groupoids-section-phantom]{Groupoïdes dans les espaces algébriques}
\item \hyperref[spaces-more-groupoids-section-phantom]{Compléments sur les groupoïdes dans les espaces}
\item \hyperref[bootstrap-section-phantom]{Amorçage}
\item \hyperref[spaces-pushouts-section-phantom]{Sommes amalgamées d'espaces algébriques}
\end{enumerate}
Thèmes de géométrie
\begin{enumerate}
\setcounter{enumi}{81}
\item \hyperref[spaces-chow-section-phantom]{Groupes de Chow des espaces}
\item \hyperref[groupoids-quotients-section-phantom]{Quotients de groupoïdes}
\item \hyperref[spaces-more-cohomology-section-phantom]{Compléments sur la cohomologie des espaces}
\item \hyperref[spaces-simplicial-section-phantom]{Espaces simpliciaux}
\item \hyperref[spaces-duality-section-phantom]{Dualité pour les espaces}
\item \hyperref[formal-spaces-section-phantom]{Espaces algébriques formels}
\item \hyperref[restricted-section-phantom]{Algébrisation des espaces formels}
\item \hyperref[spaces-resolve-section-phantom]{Résolution des surfaces revisitée}
\end{enumerate}
Théorie des déformations
\begin{enumerate}
\setcounter{enumi}{89}
\item \hyperref[formal-defos-section-phantom]{Théorie formelle des déformations}
\item \hyperref[defos-section-phantom]{Théorie des déformations}
\item \hyperref[cotangent-section-phantom]{Le complexe cotangent}
\item \hyperref[examples-defos-section-phantom]{Problèmes de déformation}
\end{enumerate}
Champs algébriques
\begin{enumerate}
\setcounter{enumi}{93}
\item \hyperref[algebraic-section-phantom]{Champs algébriques}
\item \hyperref[examples-stacks-section-phantom]{Exemples de champs}
\item \hyperref[stacks-sheaves-section-phantom]{Faisceaux sur les champs algébriques}
\item \hyperref[criteria-section-phantom]{Critères de représentabilité}
\item \hyperref[artin-section-phantom]{Axiomes d'Artin}
\item \hyperref[quot-section-phantom]{Espaces de Quot et de Hilbert}
\item \hyperref[stacks-properties-section-phantom]{Propriétés des champs algébriques}
\item \hyperref[stacks-morphisms-section-phantom]{Morphismes de champs algébriques}
\item \hyperref[stacks-limits-section-phantom]{Limites de champs algébriques}
\item \hyperref[stacks-cohomology-section-phantom]{Cohomologie des champs algébriques}
\item \hyperref[stacks-perfect-section-phantom]{Catégories dérivées de champs}
\item \hyperref[stacks-introduction-section-phantom]{Introduction aux champs algébriques}
\item \hyperref[stacks-more-morphisms-section-phantom]{Compléments sur les morphismes de champs}
\item \hyperref[stacks-geometry-section-phantom]{La géométrie des champs}
\end{enumerate}
Thèmes de théorie des espaces de modules
\begin{enumerate}
\setcounter{enumi}{107}
\item \hyperref[moduli-section-phantom]{Champs de modules}
\item \hyperref[moduli-curves-section-phantom]{Espaces de modules de courbes}
\end{enumerate}
Divers
\begin{enumerate}
\setcounter{enumi}{109}
\item \hyperref[examples-section-phantom]{Exemples}
\item \hyperref[exercises-section-phantom]{Exercices}
\item \hyperref[guide-section-phantom]{Guide bibliographique}
\item \hyperref[desirables-section-phantom]{Desiderata}
\item \hyperref[coding-section-phantom]{Style de programmation}
\item \hyperref[obsolete-section-phantom]{Obsolète}
\item \hyperref[fdl-section-phantom]{Licence de documentation libre GNU}
\item \hyperref[index-section-phantom]{Index généré automatiquement}
\end{enumerate}
\end{multicols}


\bibliography{my}
\bibliographystyle{amsalpha}

\end{document}
